% !TEX encoding = UTF-8 Unicode
%------------------------------------------------------------------------------
% Beginning of rh-GFERF.tex
%------------------------------------------------------------------------------
%
 

% AMS-LaTeX version 2 sample file for journals, based on amsart.cls.
%
\documentclass{amsart}
\usepackage{mathrsfs}
\usepackage[all]{xy}
\usepackage{mathabx} 

\usepackage{wasysym}
\usepackage{amssymb}

\usepackage{mathtools}
\usepackage{mathbbol}

\usepackage{ulem} 
%\textwidth=145mm \textheight=225mm
%%%%%%%%%%%%%%%%%%%%%%%%%%%%%%%%%%%%%%%%%%%%%%%%%%%%%%%%%%%%%%%%%%%%%%%%%%%%
%% Margin comments, obtained from Dani Wise
%%%%%%%%%%%%%%%%%%%%%%%%%%%%%%%%%%%% %%%%%%%%%%%%%%%%%%%%%%%%%%%%%%%%%%%%%%%%

%\usepackage{ifthen}
 
\newcommand{\showcomments}{yes}
%\newcommand{\showcomments}{no}

\newsavebox{\commentbox}
\newenvironment{comment}%
% begin comment
{\ifthenelse{\equal{\showcomments}{yes}}%
% then begin comment in margin
{\footnotemark
    \begin{lrbox}{\commentbox}
    \begin{minipage}[t]{1.25in}\raggedright\sffamily\upshape\tiny
    \footnotemark[\arabic{footnote}]}%
% else eat contents of the environment
{\begin{lrbox}{\commentbox}}}%
% end comment
{\ifthenelse{\equal{\showcomments}{yes}}%
% then end comment
{\end{minipage}\end{lrbox}\marginpar{\usebox{\commentbox}}}%
% else finish eating
{\end{lrbox}}}

\usepackage{ifpdf}
\ifpdf %if using pdfLaTeX in PDF mode
  \usepackage[pdftex]{graphicx}
  \DeclareGraphicsExtensions{.pdf,.png,.jpg,.jpeg,.mps}
  \usepackage{pgf}
  \usepackage{tikz}
\else %if using LaTeX or pdfLaTeX in DVI mode
  \usepackage{graphicx}
  \DeclareGraphicsExtensions{.eps,.bmp}
  \DeclareGraphicsRule{.emf}{bmp}{}{}% declare EMF filename extension
  \DeclareGraphicsRule{.png}{bmp}{}{}% declare PNG filename extension
  \usepackage{pgf}
  %\usepackage{tikz}
  \usepackage{pstricks}
\fi
\usepackage{epic,bez123}
%\usepackage{floatflt}% package for floatingfigure environment
\usepackage{wrapfig}% package for wrapfigure environment

\DeclareMathAlphabet{\mathpzc}{OT1}{pzc}{m}{it}


\usepackage{ifpdf}
\ifpdf %if using pdfLaTeX in PDF mode
  \usepackage[pdftex]{graphicx}
  \DeclareGraphicsExtensions{.pdf,.png,.jpg,.jpeg,.mps}
  \usepackage{pgf}
  %\usepackage{tikz}
\else %if using LaTeX or pdfLaTeX in DVI mode
  \usepackage{graphicx}
  \DeclareGraphicsExtensions{.eps,.bmp}
  \DeclareGraphicsRule{.emf}{bmp}{}{}% declare EMF filename extension
  \DeclareGraphicsRule{.png}{bmp}{}{}% declare PNG filename extension
  \usepackage{pgf}
  %\usepackage{tikz}
  \usepackage{pstricks}
\fi
\usepackage{epic,bez123}
%\usepackage{floatflt}% package for floatingfigure environment
\usepackage{wrapfig}% package for wrapfigure environment

\usepackage{geometry} % to change the page dimensions
%\geometry{a4paper} % or letterpaper (US) or a5paper or....
\geometry{margin=1in} % for example, change the margins to 2 inches all round

%\newtheorem*{mainthm}{Main Theorem}
\newtheorem*{thmA}{Theorem A}
\newtheorem*{thmB}{Theorem B}
\newtheorem*{maincor}{Corollary}
\newtheorem{thm}{Theorem}[section]
\newtheorem{prop}[thm]{Proposition}
\newtheorem{cor}[thm]{Corollary}
\newtheorem{lem}[thm]{Lemma}
\newtheorem{conj}[thm]{Conjecture}
\newtheorem*{claim}{Claim}
 

\newtheorem{assump}{}
\renewcommand{\theassump}{Assumption \Alph{assump}}

\newtheorem{conv}[thm]{Convention} 

\theoremstyle{definition}
\newtheorem{defn}[thm]{Definition}

\theoremstyle{remark}
\newtheorem{quest}[thm]{Question}
\newtheorem{rem}[thm]{Remark}
\newtheorem{example}[thm]{Example}
\newtheorem{examples}[thm]{Examples}
\newtheorem{prob}[thm]{Problem}
\newtheorem*{ack}{Acknowledgment}

%\numberwithin{equation}{section}

%Cayley graph
\newcommand{\Ga }{\mathscr{G} (G, \mathcal A)}
%Relative Cayley graph
\newcommand{\G }{\mathscr{G} (G, S\cup \mathcal P)}
%Cayley graph
\newcommand{\Gx }{\mathrm{Cay}(\Gamma, S)}
% GP pair
\newcommand{\GP }{(G, \mathcal P)}
%Relative word distance
\newcommand{\dxp }{d_{S\cup{\mathbb P}}}
%Word distance
\newcommand{\dx }{d_S}

\DeclarePairedDelimiter\ceil{\lceil}{\rceil}
\DeclarePairedDelimiter\floor{\lfloor}{\rfloor}

% Floyd Distance with basepont O
\newcommand{\fdo }{\rho_{\lambda, o}}
% Floyd Length with basepoint O
\newcommand{\flo }{\mathfrak l_{\lambda, o}}
% Shorcut Floyd distance with basepoint O 
\newcommand{\sfdo }{\bar{\rho}_{\lambda, o}}

% Special Floyd distance/length without \lambda
\newcommand{\fds }{\rho}
\newcommand{\fls }{\mathfrak l}

\newcommand{\ax}{\mathrm{Ax}}


% action 
\newcommand{\act}{\curvearrowright}

%Floyd completion with the group G
\newcommand{\Gf}{\overline{G}_\lambda}
\newcommand{\Gb}{\bar{G}}
 
 
\newcommand{\path}[1]{\mathbf{\ushort{#1}}}
  

% Projection complex
\newcommand{\QT}{\mathcal C_L(\mathbb S)}
\newcommand{\PC}{\mathcal P_K(\mathbb S)}

\newcommand{\CS}{\mathbb S}
 
 
 
%Ends space of a tree T in Section 4
\newcommand{\pT}{\partial{T}}
%Gromov boundary of a hyperbolic X
\newcommand{\pX}{\partial{X}}
%Uniformly Conical points in Limit set of G
 
 
\newcommand{\mG}{\partial_{r\mu} G}
\newcommand{\mGX}{\partial_{r\mu}^\infty (G,X)}
\newcommand{\mGm}{\partial_{r\mu}^{min} G}
\newcommand{\mGr}{\partial_{r\mu} G}
\newcommand{\pG}{\partial{\Gamma}}
 
%Poincare Series

\newcommand{\PS}[1]{\mathcal P_{#1}(s)}
\newcommand{\Ps}[1]{\mathcal P_{#1}(s,1)}
\newcommand{\PSS}[1]{\mathcal P'_{#1}(s,o)}

%Dirac measures
\newcommand{\dirac}[1]{{\mbox{Dirac}}{(#1)}}
% Support of measures...
\newcommand{\supp}{\mbox{supp}}
\newcommand{\isom}{\mbox{Isom}} 

%Growth rate
\newcommand{\g}[1]{\delta_{#1, S}}

\newcommand{\e}[1]{\omega(#1)}

\newcommand{\Rmu}{{R_\mu}}

%Hausdorff dimension
\newcommand{\Hdim}{{\rm Hdim}}
 
\newcommand{\pmu}{\bar\mu}

 
%Diameter of subset
\newcommand{\diam }[1]{{\textbf{diam}\big(#1\big)}}
%Projection
\newcommand{\proj }{\textbf{Proj}}
% Label of path
\newcommand{\lab }{\textbf{Lab}}
% Length of a path
\newcommand{\len }{\ell}
% cardinality
\newcommand{\card}[1]{\lvert#1\rvert}
%  Absolute value notation
\newcommand{\abs}[1]{\lvert#1\rvert}
%  Word norm value notation
\newcommand{\norm}[2]{\lvert #1 \rvert_{#2}}
%  Parallel notation
\newcommand{\para}[1]{\|#1\|}
% Label of path
 
\newcommand{\traj}{\mathbf{Traj}}
 
\renewcommand{\vec}[1]{\boldsymbol{#1}}


\def\red{\textcolor{red}}
\def\blue{\textcolor{blue}}
\def\green{\textcolor{green}}
\def\str{\sout}
\newcommand{\lw}[1]{{\color{cyan}{#1}}}
\newcommand{\ywy}[1]{{\color{red}{#1}}}

\usepackage[bookmarks=true, pdfauthor={YANG Wenyuan}]{hyperref}
\usepackage{etoolbox}
\apptocmd{\sloppy}{\hbadness 10000\relax}{}{}
\apptocmd{\sloppy}{\vbadness 10000\relax}{}{}
\usepackage{enumerate}
\begin{document}

\title{Branching Random Walks on relatively hyperbolic groups}

 
%    Information for first author
\author{Matthieu Dussaule}
\address{Matthieu Dussaule : D\'epartement de math\'ematiques, Facult\'e des Sciences Batiment I,
2 Boulevard Lavoisier
49045 Angers cedex 01}
\email{matthieu.dussaule@hotmail.fr}

\author{Longmin Wang}
\address{Longmin Wang : School of Statistics and Data Science, LPMC, Nankai University, 94 Weijin Road, Nankai District, Tianjin, China }
\email{wanglm@nankai.edu.cn}

\author{Wenyuan Yang}
\address{Wenyuan Yang : Beijing International Center for Mathematical Research (BICMR), Beijing University, No. 5 Yiheyuan Road, Haidian District, Beijing, China}
\email{yabziz@gmail.com}
 

%    \thanks will become a 1st page footnote.
\thanks{M.D.   has received funding from the European Research Council  (ERC) under the European Union's Horizon 2020 research and innovation  program under the Grant Agreement No 759702. L.W. is supported by National Natural Science Foundation of China (12171252).  W.Y. is supported by National Key R \& D Program of China (SQ2020YFA070059).}

%    Information for second author
% \author{Author Two}
% \address{Mathematical Research Section, School of Mathematical Sciences,
% Australian National University, Canberra ACT 2601, Australia}
% \email{two@maths.univ.edu.au}
% \thanks{Support information for the second author.}

%    General info
\subjclass[2010]{Primary 20F65, 20P05. Secondary 05C81, 20F67, 28A78}

%\date{\today}

\dedicatory{}

\keywords{Growth rate, Hausdorff dimension, Bowditch and Floyd boundaries, Parabolic gap}

\begin{abstract}
 Let $\Gamma$ be a non-elementary relatively hyperbolic group with a finite generating set.  Consider a finitely supported admissible and symmetric probability measure $\mu$ on $\Gamma$ and a probability measure $\nu$ on $\mathbb{N}$ with mean $r$.  Let $\mathrm{BRW}(\Gamma,\nu,\mu)$ be the branching random walk on $\Gamma$ with offspring distribution $\nu$ and base motion given by the random walk with step distribution $\mu$.  It is known that for $1 < r \leq R$ with $R$ the radius of convergence for the Green function of the random walk, the population of $\mathrm{BRW}(\Gamma,\nu,\mu)$ survives forever, but eventually vacates every finite subset of $\Gamma$.  We prove that in this regime, the growth rate of the trace of the branching random walk is equal to the growth rate $\omega_\Gamma(r)$ of the Green function of the underlying random walk. We also prove that the Hausdorff dimension of the limit set $\Lambda(r)$, which is the random subset of the Bowditch boundary consisting of all accumulation points of the trace of $\mathrm{BRW}(\Gamma,\nu,\mu)$, is equal to a constant times $\omega_\Gamma(r)$.
\end{abstract}
%\begin{comment}
%    I have made the text wider to shorten the paper. In my opinion,  a very long paper would scare the readers and referee as well. But I'm OK with the original formatting.
%    {\color{blue} I put it back because we couldn't read the comments in the margin anymore. But I'm okay to put it as you wish once we have finished discussing.}
%\end{comment}
\maketitle



%\section{Growth of Green functions}

%\subsection{A critical gap criterion}
%A proper action of $G$ on a metric space $(\mathrm Y, d)$ naturally induces a proper left-invariant pseudo-metric $d_G$ on $G$ as follows: 
%$$
%\forall g, h\in G: \;d_G(g, h)=d(go, ho)$$
%for a basepoint $o\in \mathrm Y$. Using the pseudo-metric $d_G$, we can define the ball set $N(o, n)$, the critical exponent $\e \Gamma$ of a subset $\Gamma$ correspondingly. To emphasize the distance, we use the notation $\omega_{d_G}(\Gamma)$ only in this subsection. In other places, the distance should be clear in context.

%To present the criterion,  one needs to introduce the  Poincar\'e series 
%$$\PS{\Gamma, d_G} = \sum\limits_{g \in \Gamma} \exp(-s\cdot d_G(1, g)), \; s \ge 0,$$
%which clearly  diverges for $s<\omega_{d_G}(\Gamma)$ and converges for $s>  \omega_{d_G}(\Gamma)$. 


%Dalbo et al. \cite{DPPS} presented a very useful criterion to differentiate the critical exponents of two Poincar\'e series, which is the key tool to establish growth-tightness. We formulate it with   purpose to exploit the extension map. 



%Since it plays an important role in proving Theorem
%\ref{tightfloyd} and its proof is rather short, we include their
%proof for the completeness.

 
%Let $A, B$ be two sets in $G$. Denote by $\mathbb W(A, B)$ the set of all words over the alphabet set $A\sqcup B$ with letters alternating in $A$ and $B$. We consider a left-invariant pseudo-metric $d_G$ on $G$ (which might not be coming from the pullback via a proper action). 
%\begin{lem}\label{degrowth}
%Assume that there exists an injective  map $\iota:  \mathbb W(A)\to  \mathbb W(A, B)$ such that the evaluation map $\Phi: \mathbb W(A, B) \to G$ is injective on the subset $\iota(\mathbb W(A))$ as well. Denote $X:=\Phi(\iota(\mathbb W(A)))\subset G$. If $B$ is finite, then $\PS{A, d_G}$ converges at $s=\omega_{d_G}(X)$. In particular,   we have the critical gap:
%  $$\omega_{d_G}(X)
%> \omega_{d_G}(A)$$ provided that $\PS{A, d_G}$ diverges at $s=\omega_{d_G}(A)$. 
%\end{lem}



%\begin{proof}
%Since $\Phi: \mathbb W(A, B) \to G$ is injective, each element in the image $X$ has a unique  alternating product form over $A\sqcup B$. For a word $W=a_1b_1a_2\cdots a_nb_n \in \mathbb W(A, B)$, we have $$
%d_G(1, a_1b_1\cdots a_nb_n) \le \sum_{1\le i\le n} \big (d_G(1, a_i)+D\big),
%$$ 
%where $D:=\max \{d_G(1, b): {b\in B}\}<\infty$ since $B$ is a finite set.
%As a consequence, we estimate the Poincar\'e series  $\PS{X, d_G}$  of $X$ as follows:
%\begin{align*}
%&\sum\limits_{g \in X} \exp(-s\cdot d_G(1, g)) \\
%\ge&  \sum\limits_{n =1}^{\infty} \left(\sum\limits_{a_1, \ldots, a_n \in A} \exp(-s\cdot d_G(1, a_1 b_1 \cdots a_n  b_n)) \right)\\
%\ge &   \sum\limits_{n =1}^{\infty} \left(\sum\limits_{a \in A} \exp(-s\cdot d_G (1, a))\right)^n\cdot   \exp(-sD)^n.
%\end{align*}
 

%By way of contradiction, assume that $\PS{A, d_G}$ diverges at $s=\omega_{d_G}(X)$, so   there exists some $s_0 >
%\omega_{d_G}(X)$ such that $\sum\limits_{a \in A} \exp(-s_0\cdot d_G (1,
%a))> 1/D$. Consequently, $\PS{X, d_G}$ diverges at $s=s_0$, so  implies that  $\omega_{d_G}(X) > s_0$. This  contradiction concludes the proof of lemma. 
%\end{proof}
%\begin{rem}
%From now on, we shall always consider the distance $d_G$ as the pullback of the distance $d$ on $\mathrm Y$ via the proper action. Hence, the subindex $d_G$ is omitted for simplicity. 
%\end{rem}

%\subsection{Convergence of Poincare series on parabolic subgroups}
%Let $d_r(x,y)=-\log F_r(x,y)=-\log G_r(x,y)+\log G_r(1,1)$ be the Green distance on $G$.

%\begin{lem}
%$\PS{P, d_r} (1) = \sum\limits_{g \in P} \exp(-s\cdot d_r(1, g))= \sum_{g\in P} F_r(x,y) <\infty$.
%\end{lem}
%\begin{proof}
%By the critical criterion, it suffices to prove that the critical exponent of $G$ with respect to $d_r$ equals 1:
%$$
%\sum_{g\in G} (F_r(1,g))^s =   \sum_{g\in G} \exp(-sd_r(1,g)) 
%$$
%diverges at $s<1$ and converges at $s>1$. By Gouezel-Lalley, we have that $\sum_{g\in G} F_r(1,g)=\infty$ for any $r\ge 1$. Thus, the critical exponent of $d_r$ is strictly greater than 1.
%\end{proof}

%\section{}
%\newpage


%Consistency of notations.

%It seems we are now using
%\begin{itemize}
%    \item $r$ for the parameter of the Green function,
%    \item $\rho$ for the spectral radius and $R$ for its inverse,
%    (Although it might seem unnecessary, I think it is good to keep both letters. Because when introducing induced transition kernels on parabolic subgroups, we already have lot of indices, so having a power -1 would add confusion. For example, the radius of convergence on an $\eta$ neighborhood of a parabolic subgroup $P$ with parameter $r$ would be $\rho_{\eta,P}(r)^{-1}$. Using the notation $R_{\eta,P}(r)$ is simpler.
%    But if you want to keep the letter $R$ for something else, it is okay.)
%    \item $\lambda$ for the parameter of the Floyd function,
%    \item $\Lambda(r)$ for the limit set of the BRW,
%    \item $\delta$ for the Floyd function itself,
%    \item $\omega$ for the critical exponent of the Poincaré series with weight $G_r(e,x)$,
%    \item $H_n$ for the sum of the Green function on spheres,
%    \item $\mathbb{P}$ for the set of parabolic cosets, $\mathbb P_0\subset \mathbb P$ for the parabolic subgroup representatives.
%    \item $\mathcal P$ for the trace of BRW
%    \item $\mathbf{P}$ for probability,
%    \item $\mathbf{E}$ for expectation,
%    \item $\Omega$ for cones and $\Pi$ for shadows
%    \item $\kappa$, $K$, $D$ or $C$ according to the context for bounding constants and radii of balls,
%\end{itemize}

\tableofcontents
%{\color{blue}About the distinction between metric and distance that Wenyuan earlier indicated.
%I've been reading some scientific forums on the internet.
%The word metric seems to refer to a more general notion than the word distance.
%Roughly speaking, it refers to a way of "measuring" things, while the word distance refers to the usual notion of distance between points.
%The distinction seems to be relevant in Physics, especially in the field of general relativity (eg. the Minkowsky metric is not a distance) and I suspect that mathematicians working on (pseudo)-Riemannian and Finsler geometry encounter this distinction as well.
%This might be the reason why Paulin mentioned it. In pure mathematical language, both words are synonyms, but I understand this is somehow important and I agree we should choose the word distance everywhere.}
\section{Introduction}

\subsection{Background and motivation}%\begin{comment}
    %It is great you start with the introduction of BRC. Maybe we should provide a general reference for this? As our results are somehow indiscipline, I felt we should give a more friendly introduction to the geometers and probabilist...  
    %{\color{blue}I didn't find a satisfying introduction to BMC to refer to. Maybe one of ou has an idea ?} Probably you see a recent paper by Kaimanovich-Woess : https://arxiv.org/abs/2205.13615 about BRW. Maybe we can give some effort to be more consistant with their terms (of course, they work with BRW very differently, but they are leaders in this area at least from more geometric perspective) 
%\end{comment}
Let $(V,E)$ be a locally finite connected infinite graph. %\begin{comment}
%How about we state it for a countable state space $\mathcal S$ or $V$? Then we endow the graph structure afterwards when introduce relative hyperbolicity.
%{\color{blue}I find this a bit artificial and anyway, any state space of a Markov chain can be considered as a graph by adding edges between $x$ and $y$ if $p(x,y)>0$.}
%\end{comment}
A Branching Markov chain (BMC) on $(V,E)$ is defined as follows.
One starts with a single particle at a fixed vertex $v_0\in V$.
For $n\geq 1$, each particle still alive at time $n$ dies and gives birth to an independent random number of offspring particles, according to a probability measure $\nu$ on $\mathbb N = \{1, 2, 3, \ldots\}$, each of them independently moving on $(V,E)$ according to an underlying Markov chain on $(V,E)$ driven by a transition kernel $p(x,y),x,y\in V$.
Sometimes in literature, the measure $\nu$ is distributed on $\mathbb Z_{\geq 0}=\{0,1,2,\ldots \}$ but
we will always assume that there is at least one offspring particle to avoid the extinction of the system.
This is not serious restriction, as conditioning on non-extinction, one can assume that $\nu$ is distributed on $\mathbb N$, see \cite[Chapter~1]{AN04}.
We will also assume that the underlying Markov chain is irreducible, i.e.\ every vertex of the graph can be visited by the Markov chain with positive probability and that it is symmetric, i.e.\ $p(x,y)=p(y,x)$ for all $x,y\in V$.
There is a large body of work dedicated to branching Markov chains on the real line and we refer to \cite{Shi} and references therein for more details.
We also refer to \cite{AN04} for a general discussion and historical perspective on branching processes.



A branching Markov chain is called recurrent if
%\str{it eventually fills up the whole graph, i.e.\ }
% \begin{comment}
%   In the non-symmetric case, a transient BRW may visit all vertices in the graph.
%   %{\color{blue} This is super interesting. Can you provide a reference ?}
%   \lw{For examples, consider the $p$-biased BRW on regular tree $\mathbb{T}^{d + 1}$, that is, the base motion is given by a $p$-biased RW with transition probability $p(x, y) = \frac{1}{d + p}$ for $d(o, x) < d(o, y)$ and $p(x, y) = \frac{p}{d + p}$ for $d(o, x) > d(o, y)$, where $x\sim y$ and $o$ is the root. Then the dimension of the limit set is $h(r) = \log \left[ \frac{d + p - \sqrt{(d + p)^2 - 4 d p r^2}}{2 p r} \right]$. When $p > 1/d$, $h(r)$ has a gap at $R = \frac{d+p}{2 \sqrt{d p}}$; If $p < 1/d$, then for $\frac{d+p}{d p + 1} < r < R$, $h(r) = \log d$ and the BRW visits all vertices of the tree. This example is from a working project of a student on BRW on trees with finitely many cone types. Note that this is not a BRW on group. For non-symmetric BRW on groups, I am not sure if the dimension $h(r)$ can attain the dimension of the whole boundary before $R$. }
%   %{\color{blue}Thank you, this is very interesting !}
% \end{comment}
with positive probability (and thus with probability 1), some (and thus every) vertex of the graph is visited by infinitely many particles of the BMC.
It is called transient otherwise.
Recurrence or transience of the BMC is governed by the expectation of the offspring distribution
$$\mathbf E[\nu]=\sum_{k\geq 1}k\nu(k)$$
and by the spectral radius $\rho$ of the Markov chain defined by
$$\rho=\limsup_np_n(x,y)^{1/n},$$
which is independent of $x$ and $y$, provided the underlying Markov chain is irreducible.
Here $p_n$ is the $n$-th convolution power of $p$ defined by
$$p_n(x,y)=\sum_{z_1,\ldots,z_{n-1}\in V}p(x,z_1)p(z_1,z_2)\cdots p(z_{n-1},y).$$
More precisely, if $\mathbf E[\nu]\leq \rho^{-1}$, then the branching Markov chain is transient, otherwise, it is recurrent, see \cite{BP94, GM06} and references therein.


Let us now restrict our attention to the following context.
Consider a finitely generated group $\Gamma$ endowed with a finite generating set and a probability measure $\mu$ on $\Gamma$.
For a given probability measure $\nu$ on $\mathbb N$, the \textit{branching random walk} $\mathrm{BRW}(\Gamma,\nu,\mu)$ is the branching Markov chain on the Cayley graph of $\Gamma$ driven by $\nu$ and by the $\mu$-random walk, which is the Markov chain whose transition probability is given by $p(x,y):=\mu(x^{-1}y)$.
Since $p$ is assumed to be symmetric, $\mu$ is symmetric in the sense that $\mu(x)=\mu(x^{-1})$ for every $x\in \Gamma$.
In this case, irreducibility of the random walk means that the support of $\mu$ generates $\Gamma$ as a group.
%Assuming further that $\mu$ is symmetric, i.e.\ $\mu(x)=\mu(x^{-1})$, this is equivalent to the support of $\mu$ generating $\Gamma$ as a group.
We also say that the random walk driven by $\mu$ is \textit{admissible}.

In the transient case, we introduce the trace $\mathcal{P}$ of the branching random walk, which is the set of vertices that are ever visited by $\mathrm{BRW}(\Gamma,\nu,\mu)$.
It is a random subset of $\Gamma$ and it has been a fruitful line of research to study the geometric properties of $\mathcal{P}$.
When the group $\Gamma$ is endowed with a geometric boundary $\partial \Gamma$, one can define the limit set $\Lambda$ of $\mathcal{P}$ as the closure of $\mathcal{P}$ in $\partial \Gamma$, i.e.\ $\Lambda=\mathbf{cl}(\mathcal{P})\cap \partial\Gamma$.
In hyperbolic context, the growth rate of $\mathcal{P}$ and the Hausdorff dimension of $\Lambda$ for suited distance on the boundary $\partial \Gamma$ has been related to asymptotic quantities involving $\Gamma$ and $\mu$ as we explain below.
Let us first mention that Benjamini and M\"uller \cite{BM12} studied qualitative properties of $\mathcal{P}$ and listed a certain number of conjectures. They proved in particular that $\mathcal{P}$ has exponential growth, while their method does not give quantitative results on the growth rate.
Hutchcroft proved in \cite{Hutchcroft20} that on any non-amenable group, two independent branching random walks almost surely intersect at most finitely often, which imply that $\mathcal{P}$ has infinitely many ends almost surely.
This answers some of the questions in \cite{BM12}.
Let us also mention that in a very recent work \cite{KaimanovichWoess}, Kaimanovich and Woess studied limit behaviour of branching random walks in terms of geometric features of $\Gamma$ with a very new angle.

\subsection{Earlier results on hyperbolic groups}
We introduce the Green function defined as
$$G_r(x,y)=\sum_{n\geq 0}p_n(x^{-1}y)r^n,x,y\in \Gamma.$$
Its radius of convergence $R$ is independent of $x,y$ provided that the random walk driven by $\mu$ is admissible. We have $R=\rho^{-1}$, where $\rho$ is the spectral radius introduced above.
The groups under consideration in this paper will always be non-amenable, so by a landmark result of Kesten, $R>1$, see \cite[Corollary~12.5]{Woessbook}.
For every $r\leq R$, we set
$$H_r(n)=\sum_{x\in S_n}G_r(e,x)$$
where $S_n$ is the sphere of radius $n$ centered at $e$ for the word distance
and
$$\omega_\Gamma(r)=\limsup_{n\to \infty}\frac{1}{n}\log H_r(n).$$
We call $\omega_\Gamma(r)$ the growth rate of the Green function.
It depends both on $\mu$ and on the chosen finite generating set of $\Gamma$.
We also set $\mathcal{P}_n=\mathcal{P}\cap S_n$ and $M_n=\sharp \mathcal{P}_n$.
We define the growth rate of $\mathrm{BRW}(\Gamma,\nu,\mu)$ as $\limsup \frac{1}{n}\log M_n$.

Assume that $\Gamma$ is a finitely generated free group endowed with its standard generating set and consider a branching random walk on $\Gamma$ with $\mathbf E[\nu]=r\in [1,\rho^{-1}]$.
Let $\lambda<1$ and for $x,y\in \Gamma$, set $d_\lambda(x,y)=\lambda^{n}$, where $n$ is the biggest integer such that the prefixes of length $n$ of $x$ and $y$ coincide.
Then, $d_\lambda$ extends to a distance on $\Gamma \cup \partial \Gamma$, where $\partial \Gamma$ is the set of ends of $\Gamma$.
Liggett~\cite{Liggett96} and Hueter-Lalley \cite{HL00} proved that whenever the underlying random walk is % a (possibly anisotropic)
a symmetric, possibly anisotropic, nearest neighbor random walk on $\Gamma$, then the limit
$$\theta(r)=\lim_{n \to \infty} \left( M_n \right)^{1/n}$$
is well defined and is almost surely a constant.
Moreover, the Hausdorff dimension of the limit set $\Lambda$ in $\partial \Gamma$ is equal to $- \log_{\lambda} \theta(r)$.
Finally, the function $\theta(r)$ is continuous and strictly increasing on $[1,\rho^{-1}]$ and has critical exponent $1/2$ at $\rho^{-1}$, which means that
there exists $C>0$ such that
$$\theta(\rho^{-1})-\theta(r)\underset{r\to \rho^{-1}}{\sim}C\sqrt{\rho^{-1}-r}$$
and
$$\log \theta(\rho^{-1})-\log \theta(r)\underset{r\to \rho^{-1}}{\sim}\frac{C}{\theta(\rho^{-1})}\sqrt{\rho^{-1}-r}.$$
Furthermore, $\log \theta(r)\leq \frac{1}{2}v$, where $v=\log (2q-1)$ is the volume growth rate of the free group $\Gamma$ with $q$ generators. %Note that  $v=\log (2q-1)$ if $\Gamma$ is the free group with $q$ generators.

In \cite{SWX}, the authors extended these results to all hyperbolic groups and expressed $\log \theta$ as the growth rate of the Green function, i.e.
$$\theta(r)=\mathrm{e}^{\omega_\Gamma(r)}.$$
Here, $\partial \Gamma$ is the Gromov boundary of $\Gamma$ endowed with a visual distance $d_\lambda$ satisfying
$d_\lambda(\xi,\zeta)\asymp \lambda^{(\xi|\zeta)}$, where $(\cdot|\cdot)$ is the Gromov product,  see \cite[Section~2.1]{SWX} for more details.
In particular,   the Hausdorff dimension of the limit set was proven there to be   $\omega_\Gamma(r)$ for $r<\rho^{-1}$, but the critical case $r=\rho^{-1}$ remained open.
As a particular case of our work, we will show that this is still true at $r=\rho^{-1}$, see Corollary~\ref{corohausdorffdimensionhyperbolic} below.




\subsection{Main results}
%The class of relatively hyperbolic groups  admits many equivalent formulations. We first give a geometric definition modelled on the class of geometrically finite Kleinian groups.

%Let $\Gamma$ be a finitely generated  group acting properly and isometrically on a proper geodesic hyperbolic   space $X$.  We say that the action $\Gamma$ on $X$  is a \textit{cusp-uniform action}  if there exists a $\Gamma$-invariant family of sufficiently separated horoballs centered at   points in Gromov boundary of $X$   such that $\Gamma$ acts co-compactly on their complement. There are finitely many conjugacy classes of horoball stabilizers called \textit{maximal parabolic subgroups}. For convenience  we    choose    a complete (finite) set denoted by $\mathbb P_0$  of conjugacy representatives, and say that the pair $(\Gamma, \mathbb P_0)$ is   \textit{relatively hyperbolic} or $\Gamma$ is \textit{hyperbolic relative to $\mathbb P_0$}. The Bowditch boundary $\partial_{\mathbf B} (\Gamma, \mathbb P_0)$ is defined to be the Gromov boundary of $X$. 
%Bowditch showed that any two cusp-uniform actions of $\Gamma$ sharing  the same collection of maximal parabolic subgroups have the homeomorphic Gromov boundaries and thus  $\partial_{\mathbf B} (\Gamma, \mathbb P_0)$ is well-defined. We emphasize that a different set $\mathbb P_0$ of conjugacy representatives leads to the same Bowditch boundary, and for this reason, we usually omit $\mathbb P_0$ and write $\partial_{\mathbf B} \Gamma.$  

%Recall that the Gromov boundary of $X$ consists of   geodesic rays up to a finite Hausdorff distance. By the very definition of cusp-uniform action, 
%the Bowditch boundary  consists of conical points and bounded parabolic points. A point $\xi$ in $\partial_{\mathbf B} \Gamma$ is called conical if there exists geodesic ray terminating at $\xi$ which exits every horoballs it enters.  A parabolic point then corresponds to the geodesic ray eventually trapped in a horoball so the horoball is centered at this parabolic point. 

%By the fundamental lemma of Geometric Group Theory, since $\Gamma$ acts co-compactly on the complement, the Cayley graph of $G$ is quasi-isometric to the complement with induced distance. This leads to the   equivalent formulation of relatively hyperbolic group in terms of Cayley graphs without directly involving hyperbolic geometry.   



 
%Fix a generating set  $S$ on $\Gamma$. Let      $\Gx$ be the Cayley graph with respect to $S$, equipped with combinatorial distance $d$ called word distance. Let $S_n=\{x\in \Gamma: d(e, x)=n\}$ be the $n$-sphere centered at the identity $e$. We also frequently write $|x|:=d(e,x)$. The Bowditch boundary turns out to compactify  $\Gx$ in a nice way so that the the induced action on the compactification is a  convergence group action. Moreover, it  is a geometrically finite action as the limit set contains exactly conical and bounded parabolic points (dynamically defined in contrast to the above geometrically defined). This actually   serves as an equivalent and dynamic formulation of relatively hyperbolic groups.   

%Generalizing a theorem of Floyd in Kleinian groups, Gerasimov proved that the Bowditch boundary could be parametrized as quotient of Floyd boundary of the Cayley graph $\Gx$. This construction, due to Floyd, works for any     graph with a preferred re-scaling function. For the purpose of this paper, we restrict to an exponential function with the Cayley graph.    

%Fix $\lambda\in (0,1)$ and a basepoint $o\in \Gamma$. We   define a Floyd length distance $\delta_o$ on $\Gx$ which is the induced length distance  by declaring each edge $e$ with length $\lambda^{|x|}$. If $\Gamma$ is infinite, then the distance $\delta_o(x,y)$ is incomplete and its Cauchy completion is called Floyd completion. The remainder $\partial_{\mathbf F} \Gamma$ is the Floyd boundary of $\Gx$.


%\subsection{Main results}

%With Bowditch and Floyd completion of the state space $\Gamma$, we are interested in the limit behaviour of the trace of branching random walks $(\Gamma,\nu,\mu)$  on the corresponding boundaries. 
Our goal in this paper is to generalize some of the aforementioned results to the class of   relatively hyperbolic groups, whose precise definition is recalled   in Section~\ref{sectionrelativelyhyperbolicgroups}.
Briefly, a finitely generated group $\Gamma$ is called relatively hyperbolic if it admits a geometrically finite action on a proper geodesic hyperbolic space $X$.
The Bowditch boundary of $\Gamma$ is then the limit set of the orbit of a fixed base point $x_0$ in the Gromov boundary of $X$.
It is unique up to $\Gamma$-equivariant homeomorphism.
We say that $\Gamma$ is non-elementary if its Bowditch boundary is infinite.
Since their introduction by Gromov,  these groups have been extensively studied by many authors and from different point of views. Besides the class of Gromov hyperbolic groups, the following   groups of geometric and algebraic interests are the main archetypal examples: 
\begin{enumerate}

    \item
    Fundamental groups of finite volume hyperbolic manifolds, and of more general finite volume Riemaniann manifolds with negatively pinched sectional curvature. 
    \item
    Infinitely ended groups, equivalently by Stalling's theorem, all non-trivial amalgamated free products and HNN extension over finite groups. 

    
\end{enumerate}
In particular, free products are  the simplest combinatorial examples of relatively hyperbolic groups, on which branching random walks  were studied by Candellero, Gilch and M\"uller \cite{CandelleroGilchMuller}. 
We recover some of their results in this paper.

A finitely generated group $\Gamma$ endowed with a finite generating set $S$ is equipped with the \textit{word distance}.
Its volume growth rate is the growth rate of the balls for the word distance.
We again refer to Section~\ref{sectionrelativelyhyperbolicgroups} for more details.

\begin{thm}\label{Thmvolumegrowthconst}
Let $\Gamma$ be a non-elementary relatively hyperbolic group endowed with a finite generating set.
Consider a finitely supported admissible and symmetric probability measure $\mu$ on $\Gamma$ and a probability measure $\nu$ on $\mathbb{N}$ with mean $r\in [1,R]$.
Consider a branching random walk $(\Gamma,\nu,\mu)$ starting at $e$.
Then, almost surely,
$$\limsup_{n\to\infty}\frac{1}{n}\log M_n=\omega_\Gamma(r).$$
Moreover, the function $r\mapsto \omega_\Gamma(r)$ is increasing and continuous and satisfies that $\omega_\Gamma(r)\leq v/2$, where $v$ is the volume growth rate of the group for the word distance.
\end{thm}

We then investigate the limit behaviours of the branching random walk at infinity and we compute the Hausdorff dimension of the limit set of the trace in various compactifications of relatively hyperbolic groups. Introduced by Floyd \cite{Floyd}, the Floyd boundary can be constructed as a compactification for  any locally finite graph, equipped with a rescaling of the graph distance called the Floyd distance. 
Generalizing Floyd's theorem  \cite{Floyd} for geometrically finite Kleinian groups,   Gerasimov \cite{Gerasimov} proved that for any relatively hyperbolic groups, the Bowditch boundary can be realized as a quotient of  the Floyd boundary of the Cayley graph.  The Floyd distance depending on a parameter $\lambda\in (0,1)$ is then transferred to the Bowditch boundary  as soon as $\lambda\geq \lambda_0$, where $\lambda_0$ is given by \cite[Map Theorem]{Gerasimov}.
The corresponding distance on the Bowditch boundary is called the \textit{shortcut distance} and will be described in Section~\ref{SectionFloydboundary}. %{\color{blue}I think we should be explicit about the constant $\lambda_0$ in our main theorems.} Yes!

\begin{thm}\label{ThmHdimconst}
%Let $\Gamma$ be a non-elementary relatively hyperbolic group endowed with a finite generating set.
%Consider a finitely supported admissible and symmetric probability measure $\mu$ on $\Gamma$ and a probability measure $\nu$ on $\mathbb{N}$ with mean $r\in [1,R]$. 
%Consider a branching random walk $(\Gamma,\nu,\mu)$ starting at $e$. 
Under the assumption of Theorem \ref{Thmvolumegrowthconst}, denote by $\Lambda(r)$ the limit set of the branching random walk in the Bowditch boundary, endowed with the shortcut distance of parameter $\lambda\in [\lambda_0,1)$.
Then, almost surely,
$$\Hdim\big(\Lambda(r)\big)=\frac{\omega_\Gamma(r)}{-\log \lambda}.$$
\end{thm}
\begin{rem}
The lower bound on the Hausdorff dimension actually holds for the whole limit set of the branching random walk in the Floyd boundary endowed with the Floyd distance, see Proposition~\ref{lowerboundHdim}.
 We can only prove the upper bound for a subset of the Floyd boundary, see Proposition~\ref{upperboundHausdorff}. In many interesting cases, the Floyd boundary is homeomorphic to the Bowditch boundary (eg. if  parabolic subgroups are amenable). 
Moreover, under the technical condition that  the volume growth rate $v_S(P)$ for every parabolic subgroup $P$ is smaller than $\omega_{\Gamma}(r)$, the above conclusion is still true for the  Floyd boundary, see Corollary~\ref{corowholeFloyd}.
In general, the identification of the Hausdorff dimension  remains open for the full limit set in the Floyd boundary.
%\begin{comment}
%{\color{blue} We removed the section where we proved this last statement about $v_S(P)$ and besides I'm not sure we completely proved it. Maybe we should remove this statement.}\ywy{This is farly easy to prove, as the critical exponent of a quasi-convex subgroup is a upper bound on the Hausdorff dimension of its limit set: we just cover the whole parabolic limit set by the shadows of points in the subgroup. Then it is done. }
%\end{comment}
\end{rem}

%Regarding the infinitely-ended groups,   the following corollary  recovers partially the main result \cite[Theorem 3.5]{CandelleroGilchMuller} on free product of groups with standard generating sets in a greater generality. An infinitely-ended group is  called accessible if any two ends can be separated by removing a fixed number of edges in its Cayley graph. Fixing   $\lambda\in (0,1)$, the visual distance between two ends is defined to be $\lambda^n$ where $n$ is the minimal   radius of a ball centered at a basepoint $o$ separating the two ends.
%Alternatively, such groups admit a splitting as a graph of groups over finite edge groups so that the vertex groups are either finite or one-ended.
%It is a famous result of Dunwoody \cite{Dun85} that finitely presented group are accessible.


%\begin{comment}
%{\color{blue} This is already stated in \cite{CandelleroGilchMuller}. Besides, the proof for free products is way simpler.} Yes, you are right. After a careful check, I think they have a complete proof for the Hausdorff dimension for the whole boundary. 
%{\color{blue}The corollary for accessible groups is interesting, but we need to prove the lower bound, because we only proved it for the shortcut distance on the Bowditch boundary.} Yes, you are right. I will add the proof (2-3 sentences). We only need to show in this case every bottleneck for thin ends is close to a transition point. 
 
%\end{comment}
As above-mentioned, groups with infinitely many ends form a special class of relatively hyperbolic groups.
Such groups can be compactified with the ends boundary, introduced by Freudenthal \cite{Freudenthal}. A natural family of visual distances depending on a parameter $\lambda\in(0,1)$ on the ends boundary      will be described in Section~\ref{sectionends}.
Along the way, we prove the following Theorem, which both extends  \cite[Theorem~3.5]{CandelleroGilchMuller} to all groups with infinitely many ends and fix a gap in their proof, as we will explain in Section~\ref{sectionends}.
\begin{thm}\label{coroaccessiblegroups}
Under the assumption of Theorem \ref{Thmvolumegrowthconst}, if $\Gamma$ is  a group with infinitely many ends  and $\Lambda(r)$ is the   limit set of the branching random walk in the ends boundary, endowed with the visual distance of parameter $\lambda\in (0,1)$, then almost surely,
$$\Hdim\big(\Lambda(r)\big)=\frac{\omega_\Gamma(r)}{-\log \lambda}.$$  
\end{thm}


Finally, for hyperbolic groups, the Bowditch boundary and the Gromov boundary coincide and the shortcut distance is bi-Lipschitz to the visual distance.
We thus deduce from Theorem~\ref{ThmHdimconst} the following, which  resolves \cite[Conjecture 1.4]{SWX}.
\begin{cor}\label{corohausdorffdimensionhyperbolic}
Under the assumption of Theorem \ref{Thmvolumegrowthconst}, if   $\Gamma$ is a non-elementary hyperbolic group and if $\Lambda(R)$ is the limit set of the branching random walk with mean offspring $R$ in the Gromov boundary of $\Gamma$ endowed with a visual distance of parameter $\lambda\in [\lambda_0,1)$,
then almost surely,
$$\Hdim\big(\Lambda(R)\big)=\frac{\omega_\Gamma(R)}{-\log \lambda}.$$
\end{cor}

In the course of this investigation, we obtain much more precise information on the trace $\mathcal{P}$ of the branching random walk.
We prove that $\mathcal{P}$ almost surely tracks the logarithm neighborhood of transition points along geodesics ending at conical limit points in the limit set.
Recall that a point on a geodesic is called a transition point if it is not deep in a parabolic coset.
We refer to Section~\ref{sectionrelativelyhyperbolicgroups} for a precise definition.
%Let $\mathbb P=\{gP: g\in G, P\in\mathbb P_0\}$ be the collection of all parabolic cosets.

\begin{thm}\label{thmlogarithmictracking}
Under the assumption of Theorem \ref{Thmvolumegrowthconst}, {there exists $0 < \kappa < \infty$ such that} 
almost surely,  for every conical point $\xi\in \Lambda(r)$, we have
$$
\limsup_{|x|\to\infty} \frac{d(x,\mathcal P)}{\log |x|}\le \kappa, %<\infty
$$
where $x$ is taken over  the set of transition points on $[e,\xi]$.
\end{thm}

This should be compared with analogous results for random walks.
Under finite first moment condition, a random walk on a hyperbolic group almost surely sub-linearly tracks geodesics from the basepoint $e$ to the limit point of the random walk in the Gromov boundary.
We refer to  Kaimanovich \cite[Theorem~7.3]{Kaimanovich} for a proof, see also \cite{Ledrappier} for the case of trees.
Sub-linear tracking of geodesics is one of the most important result in the study of random walks on groups with hyperbolic properties and is related to a celebrated multiplicative ergodic theorem of Furstenberg and Kesten \cite{FurstenbergKesten}.
It was first coined by Kaimanovich in the context of symmetric spaces \cite{Kaimanovichsymmetricspaces} and was then extended to groups with non-positive curvature by Karlsson and Margulis \cite{KarlssonMargulis}
and to more general classes of group by Tiozzo \cite{Tiozzo}, including mapping class groups.
If the random walk has finite support, then the tracking is in fact logarithmic and this is true for all weakly hyperbolic groups, i.e.\ groups admitting a non-elementary action by isometries on a Gromov hyperbolic space, see \cite[Theorem~1.3]{MaherTiozzo}.
If the group is relatively hyperbolic, then the random walk actually sub-linearly tracks transition points on the word geodesic in the Cayley graphs, see \cite[Proposition~3.2]{DussauleYang}.
Theorem~\ref{thmlogarithmictracking} can thus be thought as a generalization of these results to branching random walks on relatively hyperbolic groups and is new, even for hyperbolic groups.



%However most of their results still hold and some of them are recovered here in the more general setting of relatively hyperbolic groups. 
%\begin{comment}
%I felt too long for the reader to reach the main result of the paper.   So I would suggest move the next paragraphs to the end of introduction, with a more discussion on the possibility of extending the phase transition ....  
%\end{comment}
\medskip
Let us finally say that we will not investigate the problem of the critical exponent of $\omega_\Gamma(r)$ in this paper.
In \cite{CandelleroGilchMuller}, the authors show that for free products this critical exponent is not always $1/2$, depending on $\mu$ and more precisely depending on whether the first derivative of the Green function is finite or infinite at its radius of convergence, see \cite[Theorem~3.10]{CandelleroGilchMuller}.
It might be possible to prove that it is in fact $1/2$ whenever the underlying random walk is spectrally non-degenerate, combining techniques of \cite{SWX}, \cite{DussauleLLT2} and of the present paper.
We refer to Definition~\ref{defnspectraldegeneracy} for the definition of spectral degeneracy of a random walk.

\subsection{Parabolic gap and purely exponential growth of Green functions}
Among the results of \cite{CandelleroGilchMuller}, the authors claimed that the Hausdorff dimension of the limit set  intersected with the set of ends  of each free factor is strictly less than that of  the  whole limit set of the branching random walk (see their Corollary~3.7).
However, their proof is incorrect on assuming that the quantity $H_r(n)$ as defined above is sub-multiplicative; we refer to Remark \ref{CGMGap}   for more details.
In our study, the sub-multiplicativity of $H_r(n)$   turns out to be a  subtle property and we propose a sufficient criterion called a parabolic gap condition, which is inspired by the work of \cite{DOP} and that we now introduce.
Note that assuming that the parabolic gap condition holds, we recover their result \cite[Corollary~3.7]{CandelleroGilchMuller}, see     Remark \ref{CGMCor}.


Let $\Gamma$ be a non-elementary relatively hyperbolic group.
Let $P$ be a parabolic subgroup.
We set
$$H_{P,r}(n)=\sum_{x\in S_n\cap P}G_r(e,x)$$
and
$$\omega_P(r)=\limsup_{n\to \infty}\frac{1}{n}\log H_{P,r}(n).$$
Thus, $\omega_\Gamma(r)$ is the growth rate of the sum of the Green functions along spheres, while $\omega_P(r)$ is the similarly defined growth rate, but restricted to the parabolic subgroup $P$.
%We will see below that $\omega_\Gamma$ and $\omega_P$ can be expressed as the critical exponent of a twisted Poincaré series.
It follows from the definition that $\omega_P\leq \omega_\Gamma$.

\begin{defn}
If $\omega_P(r)<\omega_\Gamma(r)$, we say that $\Gamma$ has a \textit{parabolic gap} along $P$ for the Green function at $r$.
If for every $P$, for every $r\in (1,R]$, $\omega_P(r)<\omega_\Gamma(r)$, then we say that $\Gamma$ has a \textit{parabolic gap} for the Green function.
\end{defn}

Note that this notion only depends on the underlying random walk driven by $\mu$, not on the offspring distribution of the branching random walk.

\begin{thm}\label{mainthmparabolicgap}
Let $\Gamma$ be a non-elementary relatively hyperbolic group.
Consider a finitely supported admissible and symmetric probability measure $\mu$ on $\Gamma$.
If $\Gamma$ has a parabolic gap for the Green function,
then the sum of Green functions along spheres is roughly multiplicative and has purely exponential growth, in the sense that for all $1\leq r\leq R$, there exist $C=C(r)\geq 1$ and $C'=C'(r)\geq 1$
%\begin{comment}
%{\color{blue}In Step 2 of the proof of Lemma~\ref{LowerBDLem}, the constant $c_1$ depends on $r$, right ? This raises a difficulty in proving that $\omega_\Gamma$ is continuous.} \ywy{Yes, it is true from the current argument.} \lw{I have a proof for the continuity. Please check Lemma \ref{FactsHrLem}.}
%{\color{blue}This is very good !}
%\end{comment}
such that for all $n$,
$$\frac{1}{C}H_{r} (n+m)\leq  H_r(n)H_r(m)\le C H_{r} (n+m)
$$
and
$$\frac{1}{C'}\mathrm{e}^{n\omega_\Gamma(r)}\leq H_r(n)\leq C' \mathrm{e}^{n\omega_\Gamma(r)}.$$
%Moreover, the function $r\mapsto \omega_\Gamma(r)$ is continuous.
\end{thm}

\begin{rem}
Actually, $C$ and $C'$ can be chosen independently of $r$ in the two upper bounds, which do not require a parabolic gap to hold by Lemma~\ref{UpperBDLem}.
\end{rem}


We will prove two criteria for having a parabolic gap, see Corollary~\ref{parabolicdivergentgrowthgap} and Proposition~\ref{secondcriterionparabolicgap}.
In particular, we   prove that whenever parabolic subgroups are amenable, $\omega_P(r)<\omega_\Gamma(r)$ for all $r<R$.
Moreover, if parabolic subgroups have sub-exponential growth, then we also have $\omega_P(R)<\omega_\Gamma(R)$.


\medskip
Let us compare our notion of parabolic gap   to similar notions in different settings.
Consider a finitely generated group $\Gamma$ acting via isometries on a metric space $(X,d)$.
Then, one can endow $\Gamma$ with a left-invariant pseudo-distance by declaring
$$d(g,h)=d(g\cdot x_0,h\cdot x_0),$$
where $x_0$ is a fixed point in $X$.
 We define the volume growth rate for any subgroup $P<\Gamma$ as follows:
%$$v_X(\Gamma)=\limsup_{n\to \infty} \frac{1}{n}\log \sharp B(e,n)=\limsup_{n\to \infty}\frac{1}{n}\log \sharp \left \{g\in \Gamma,d(x_0,g\cdot x_0)\leq n\right \}$$
%and for $P$ a parabolic subgroup,
$$v_X(P)=\limsup_{n\to \infty} \frac{1}{n}\log \big( \sharp B(e,n)\cap P\big)=\limsup_{n\to \infty}\frac{1}{n}\log \big (\sharp \left \{g\in P: d(x_0,g\cdot x_0)\leq n\right \}\big).$$
Choosing $X$ to be the Cayley graph associated with a finite generating set $S$, endowed with the graph distance, we recover the word distance on $\Gamma$.
Then the volume growth rate of a subgroup $P$, also denoted by $v_S(P)$, is the standard terminology.


If $\Gamma$ is relatively hyperbolic, then it acts by isometries on a proper geodesic hyperbolic space $(X,d)$.
%We fix a base point $x_0$ and consider the induced distance on $\Gamma$.
We say in this context that $\Gamma$ has a parabolic gap (also referred as critical gap in literature) if for every parabolic subgroup $P$, $v_X(P)<v_X(\Gamma)$.
This definition makes sense in larger contexts than relatively hyperbolic groups and
this property was studied a lot in literature, see for instance \cite{DOP}, \cite{DPPS}, \cite{PS18}, \cite{PTV}, \cite{Rob}, \cite{ST21}, \cite{Vidotto} and references therein.
For typical cases, we have $v_X(P)<v_X(\Gamma)$.
On the other hand, exotic examples   of  geometrically finitely groups acting on a negatively-pinched Cartan-Hadamard manifold with the critical gap property, i.e.\ $v_X(\Gamma)=v_X(P)$ for some $P$, were first constructed in \cite{DOP}, see \cite{Peigne} for other examples.



On the contrary, when endowing $\Gamma$ with a word distance given by a finite generating set $S$, this cannot happen,
%More precisely, we set a finitely generating set $S$ and let $v_S(\Gamma)$ be the growth rate of $\Gamma$ and $v_S(P)$ be the growth rate of $P$ for the induced distance on $P$, i.e.
%$$v_S(\Gamma)=\limsup_{n\to \infty} \frac{1}{n}\log \sharp B(e,n)$$
%and
%$$v_S(P)=\limsup_{n\to \infty} \frac{1}{n}\log \sharp B(e,n)\cap P.$$
since in this context, it was proved in \cite{DFW} that one always has $v_S(P)<v_S(\Gamma)$.

\bigskip
Before going further, recall that the critical exponents $v_X(\Gamma)$ and $v_X(P)$ coincide with the exponential radius of convergence of a suited Poincaré series.
Namely, define
$$\Theta_s(\Gamma)=\sum_{g\in \Gamma}\mathrm{e}^{-sd(x_0,g\cdot x_0)}=\sum_{n\geq 0}\sum_{n\leq d(x_0,g\cdot x_0)<n+1}\mathrm{e}^{-sn}.$$
Then, for $s<v_X(\Gamma)$, $\Theta_s(\Gamma)$ diverges and for $s>v_X(\Gamma)$, it converges.
Similarly, replacing $\Gamma$ by any subgroup $P$ in the above formula defines the corresponding $\Theta_s(P)$ and $v_X(P)$.
%$$\Theta_s(P)=\sum_{g\in P}\mathrm{e}^{-sd(x_0,g\cdot x_0)}=\sum_{n\geq 0}\sum_{g\in P,n\leq d(x_0,g\cdot x_0)<n+1}\mathrm{e}^{-sn}.$$
%Again, for $s<v_X(P)$, $\Theta_s(P)$ diverges and for $s>v_X(P)$, it converges.

Now, consider a probability measure $\mu$ on a relatively hyperbolic group $\Gamma$ and for $r>0$, set
$$I(r)=\sum_{g\in \Gamma}G_r(e,g)G_r(g,e).$$
The following result is well-known, see for instance \cite[Proposition~1.9]{GouezelLalley}.
For every $r$, we have
$$I(r)=\frac{d}{dr}\big (rG_r(e,e)\big ).$$
Thus, for $r<R$, this quantity converges and for $r>R$, it diverges.
Consider the $r$-symmetrized Green distance $d_{G,r}$ defined as
$$d_{G,r}(g,h)=-\log \frac{G_r(g,h)}{G_r(e,e)}-\log \frac{G_r(h,g)}{G_r(e,e)}$$
which is a generalization of the Green distance introduced by Blach\`ere and Brofferio \cite{BlachereBrofferio}.
Then, the quantity $I(r)$ is exactly the Poincaré series associated with the distance $d_{G,r}$.
The only difference with the previous setting is that the parameter $r$ is part of the definition of the distance.
The quantity $-\log R$ is analogous to the critical exponent $v_X(\Gamma)$.
It is more complicated to define a notion of parabolic gap, because we cannot interpret 
$-\log R$ as the radius of convergence of the Poincaré series, which is not a power series in $r$.
However, the analogous notion which was coined in \cite{DG21} is called spectral degeneracy.
We will properly introduce this notion below and
we refer to \cite[Section~3.3]{DussauleLLT1} and \cite[Section~1.3]{DPT} for more explanations on this analogy.
Anyway, by results of Cartwright \cite{Cartwright1}, \cite{Cartwright2} and of Candellero and Gilch \cite{CandelleroGilch}, it is possible to construct both a spectrally degenerate random walk and a spectrally non-degenerate one on a relatively hyperbolic group, although every known example is in the class of free products.


\medskip
%Notice that in the geometric context, the parabolic gap property only depends on the distance $d$ and in the random walk context, it only depends on the measure $\mu$.
Back to our critical gap condition, we will see that the critical exponent $\omega_\Gamma(r)$ is the radius of convergence of a twisted Poincaré series $s\mapsto \Theta_{r,s}(\Gamma)$ defined by~(\ref{defPoincareseries}), involving both the Green function $G_r(e,x)$ and the word distance.
We saw above that there are sufficient conditions to have a parabolic gap, but it is difficult to tell if one can construct an example where $\omega_P(r)=\omega_\Gamma(r)$.
Answering this question would require new material.
However, the last example of Section~\ref{sectionexamples} below  suggests that it might happen (see the famous examples in the geometric context in \cite{DOP} and \cite{Peigne}).
We prove there that if there exists a finitely generated group $P$ endowed with a finitely supported admissible random walk with convergent twisted Poincaré series $\Theta_{r,s}$ at some $r\leq R$, then for a suited random walk on the free product $\Gamma=P*\mathbb Z^d$, we have $\omega_\Gamma(r')=\omega_P(r')$ for some $r'$.
See the comments at the end of Section~\ref{sectionexamples} for further details.

\subsection{Organization of the paper}
We now outline the contents of the paper and explain the overall strategy of our proofs.
In Section~\ref{sectionrelativelyhyperbolicgroups}, we recall the definition of relatively hyperbolic groups, of the Floyd distance and of the Floyd boundary.
The Floyd distance is then related with geometric properties of such groups via a number of preliminary results  used throughout the paper.
Finally, we recall the relative Ancona inequalities that will be a crucial tool.

In Section~\ref{SectiongrowthrateGreenfunction}, we study the growth rate of the Green function $\omega_\Gamma(r)$.
We prove in particular that it is increasing and continuous and bounded by $v/2$, see Corollary~\ref{corolowersemicontinuous}, Corollary~\ref{coroboundv/2}.
We also prove Theorem~\ref{mainthmparabolicgap}, i.e.\ purely exponential growth, assuming there is a parabolic gap, see Lemma~\ref{LowerBDLem}.
This is done by using classical methods for Poincar\'e series, inspired by \cite{YangSCC}.
Finally, we discuss the notion of parabolic gap through examples in the last part of the section.

Section~\ref{SectiongrowthrateBRW} is dedicated to the growth rate of the branching random walk and we prove that it is almost surely equal to the growth rate $\omega_\Gamma(r)$ of the Green function, see Proposition~\ref{P:vgrBRW}.
In particular, this ends the proof of Theorem~\ref{Thmvolumegrowthconst}.
The upper-bound $\limsup \frac{1}{n}\log M_n\leq \omega_\Gamma(r)$ is very general and does not involve relative hyperbolicity.
The lower-bound is more difficult to obtain and we have to restrict our attention to points $x\in \Gamma$ such that a geodesic from $e$ to $x$ spends a uniformly bounded amount of time $L$ in every parabolic coset.
Such geodesics are Morse in the sense of \cite{Cordes}, see also \cite{Tran} for a study a Morse geodesics in relatively hyperbolic groups.
It turns out Morse geodesic rays form a proper subset of the whole Bowditch boundary, which is typically too small to serve as a model for the Poisson boundary and thus is too small to give much information about asymptotic properties of finitely supported random walks, see in particular the comments in the introduction of \cite{QRT}, where a bigger boundary called the sub-linearly Morse boundary is introduced.
We can nevertheless prove that the growth rate of the Green function restricted to Morse directions converges to the growth rate of the whole Green function, as $L$ tends to infinity, see Lemma~\ref{omegaL}.
This is enough to adapt the arguments of \cite{SWX} for hyperbolic groups, which allows us to conclude the proof.

Finally, in the two last sections, we study the Hausdorff dimension of the limit set and we prove Theorem~\ref{ThmHdimconst}, Theorem~\ref{coroaccessiblegroups} and Corollary~\ref{corohausdorffdimensionhyperbolic}.
We start with the lower bound in Section~\ref{SectionlowerboundHausdorff}.
Following \cite{SWX}, we use the Frostman lemma and show that for every $h<\omega_\Gamma(r)/-\log \lambda$, there exists with positive probability a finite measure $\chi$ on the limit set $\Lambda$ such that the integral
$$\int \int \overline{\delta}_e(x,y)^{-h} d\chi(x)d\chi(y)$$ is finite, where $\overline{\delta}_e$ is the shortcut distance on the Bowditch boundary.
However, the proof in \cite{SWX} has a gap and we need to find a new strategy to construct the measure $\chi$, which will be defined as a random Patterson-Sullivan type measure on the limit set.
The construction is performed by using convergence results for random finite measures, whose proofs are postponed to the Appendix.
Our strategy also applies to groups with infinitely many ends and we prove Theorem~\ref{coroaccessiblegroups}.
The upper-bound for the Hausdorff dimension is proved in Section~\ref{SectionupperboundHausdorff}.
We first prove Theorem~\ref{thmlogarithmictracking}, i.e.\ logarithmic tracking by the trace of transition points on geodesic rays joining the limit set, see Lemma~\ref{LogTrackingLem}.
This allows us to cover the limit set with suited shadows and to conclude as in \cite{SWX}.
Note that the covering by shadows in \cite{SWX} only works for $r<\rho^{-1}$ and Theorem~\ref{thmlogarithmictracking} is one of the needed step to cover the case $r=\rho^{-1}$.





\section{Relatively hyperbolic groups and random walks}\label{sectionrelativelyhyperbolicgroups}
\subsection{Relative hyperbolic groups} %\begin{comment}
%{\color{blue}Shouldn't we move this definition to the introduction ?} I was wondering the main   results obtained so far should hold for BRW on hyperbolic spaces. This may have some interests to people. I wrote this definition to allow  such extension later on. Let us decide how to do such extension. I would suggest to complete the current argument for Cayley graphs and then adapt the arguments later on for the hyperbolic space (I expect not much difference between them, only with technical complexities.    
%{\color{blue}I don't understand what you mean here.
%What do you think can be generalized and how do you define $\omega_\Gamma(r)$ in this context ? I don't think there is a correct way to define the Green function associated to an action of $\Gamma$ on a space $X$.}
%I meant we do the sum of Green functions over the sphere now with respect to the hyperbolic distance. Tecnically, this is the same as the we endow the graph with a different proper distance (by pulling back from the group action on a hyperbolic space).{\color{blue}Okay I see. But then, Ancona inequalities are not known to hold and probably do not, so we need other tools. This is the main difficulty in the recent work of Boulanger and al. Also, I think this is a different (though quite interesting) topic and the proofs would be closer to \cite{SWX}. I'm not sure we should include this in this paper since it is already very long.} Okay, let us leave it out. 
%\end{comment}
%Let $\Gamma$ be a finitely generated  group acting properly and isometrically on a proper geodesic hyperbolic   space $X$.  We say that the action $\Gamma$ on $X$  is a \textit{cusp-uniform action} if there exists a $\Gamma$-invariant family of sufficiently separated horoballs centered at   points in Gromov boundary of $X$   such that $\Gamma$ acts co-compactly on its complement. There are finitely many conjugacy classes of horoball stabilizers called \textit{maximal parabolic subgroups}, and    choose    a complete (finite) set denoted by $\mathbb P_0 $  of conjugacy representatives. We also say the pair $(\Gamma, \mathbb P_0)$ is \textit{relatively hyperbolic} or $\Gamma$ is \textit{hyperbolic relative to $\mathbb P_0$}. The Bowditch boundary $\partial_{\mathbf B} (\Gamma, \mathbb P_0)$ is defined to be the Gromov boundary of $X$. Bowditch showed that any two cusp-uniform actions of $\Gamma$ sharing  the same collection of maximal parabolic subgroups have the homeomorphic Gromov boundary and thus  $\partial_{\mathbf B} (\Gamma, \mathbb P_0)$ is well-defined.


%{\color{blue}I think that defining a relatively hyperbolic group as a group acting geometrically finitely on a proper geodesic hyperbolic space would be better, because within the definition, we introduce the two types of limit points : parabolic or conical. Since the definition in terms of cusp-uniform actions is also clear and gives a geometric perspective, we can keep it as a complementary definition. I suggest the following.}

Consider a finitely generated group $\Gamma$ acting properly via isometries on a proper Gromov hyperbolic space $X$.
Define the \textit{limit set} $\Lambda_\Gamma$ as the closure of $\Gamma$ in the Gromov boundary $\partial X$ of $X$, that is, fixing a base point $x_0$ in $X$, $\Lambda_\Gamma$ is the set of all possible limits of sequences $g_n\cdot x_0$ in $\partial X$, $g_n\in \Gamma$.
A point $\xi\in \Lambda_\Gamma$ is called conical if there is a sequence $g_{n}$ of $\Gamma$ and distinct points $\xi_1,\xi_2$ in $\Lambda_\Gamma$ such that
$g_{n}\xi$ converges to $\xi_1$ and $g_{n}\zeta$ converges to $\xi_2$ for all $\zeta\neq \xi$ in $\Lambda_\Gamma$.
A point $\xi\in \Lambda_\Gamma$ is called \textit{parabolic} if its stabilizer in $\Gamma$ is infinite, fixes exactly $\xi$ in $\Lambda_\Gamma$ and contains no loxodromic element.
A parabolic limit point $\xi$ in $\Lambda_\Gamma$ is called \textit{bounded parabolic} if its stabilizer in $\Gamma$ is infinite and acts co-compactly on $\Lambda_\Gamma \setminus \{\xi\}$.
Say that the action of $\Gamma$ on $X$ is \textit{geometrically finite} if the limit set only consists of conical limit points and bounded parabolic limit points.

There are in literature several equivalent definitions of relatively hyperbolic groups.
Following Bowditch \cite{Bowditch}, we say that a finitely generated group $\Gamma$ is \textit{relatively hyperbolic} with respect to a collection of subgroups $\mathbb P_0$ if it acts via a geometrically finite action on a proper geodesic Gromov hyperbolic space $X$, such that the stabilizers of parabolic limit points for this action are exactly the conjugates of the groups in $\mathbb P_0$,
which are called maximal parabolic subgroups or simply parabolic subgroups if there is no ambiguity.
By \cite[Proposition~6.15]{Bowditch}, there is a finite number of conjugacy classes of parabolic subgroups, so in other words, $\mathbb P_0$ needs to be finite.

The limit set $\Lambda_\Gamma$ is called the \textit{Bowditch boundary} of $\Gamma$.
It is unique up to equivariant homeomorphism and we will denote it by $\partial_\mathcal{B}\Gamma$ in the sequel.
A relatively hyperbolic group is called \textit{non-elementary} if its Bowditch boundary is infinite; equivalently, if some parabolic subgroup $P\in \mathbb P_0$ is of infinite index in $\Gamma$.

Relatively hyperbolic groups are modelled on finite co-volume Kleinian groups.
In this case, the group acts via a geometrically finite action on the hyperbolic space $\mathbb H^n$ and there is a collection of separated horoballs such that the action on the complement of these horoballs is co-compact.
The parabolic subgroups are exactly the stabilizers of the horoballs.
Moreover, the Bowditch boundary is the ideal boundary $\mathbb S^{n-1}$ of $\mathbb H^n$ and parabolic limit points are the centers of the horoballs.

In \cite{Bowditch}, Bowditch gives another definition of relative hyperbolicity, mimicking the above geometric description of Kleinian groups.
Given a hyperbolic space $X$, one can define a coarse notion of horoballs.
A finitely generated group $\Gamma$ acting properly via isometries on a proper geodesic hyperbolic space $X$ is \textit{relatively hyperbolic} if only if there exists a $\Gamma$-invariant family of sufficiently separated horoballs centered at points in the Gromov boundary of $X$ such that $\Gamma$ acts co-compactly on the complement of these horoballs.
The parabolic subgroups are then exactly the stabilizers of these horoballs.
We also refer to \cite{Osin}, \cite{Farb}, \cite{DrutuSapir} and references therein for alternative definitions of relatively hyperbolic groups.

We set
$$\mathbb P=\{gP: g\in \Gamma, P\in\mathbb P_0\}$$ 
and we call $\mathbb P$ the collection of all parabolic cosets.

 Let us also fix some notations.
Given a finite generating $S$ on $\Gamma$, let $\Gx$ be the Cayley graph with respect to $S$.
The graph combinatorial distance is called the \textit{word distance.}
We denote the $n$-sphere centered at the identity $e$ by
$S_n=\{x\in \Gamma: d(e, x)=n\}$.
We will frequently write $|x|:=d(e,x)$. 
Finally, we will denote by $v$ the \textit{volume growth rate} of $\Gamma$ with respect to $S$, which is defined by
$$v=\limsup \frac{1}{n}\log \sharp B(e,n).$$
The word distance on relatively hyperbolic groups has purely exponential growth in the following sense.

\begin{lem}\label{purelyexponentialgrowth}\cite[Theorem 1.9]{YangPS}
There exists $c>0$ such that for every $n\geq0$, we have
$$\frac{1}{c}\mathrm{e}^{vn}\leq \sharp S_n\leq c \mathrm{e}^{vn}.$$
\end{lem}

\subsection{Floyd boundary}\label{SectionFloydboundary} 
 
We first recall the definition of   the Floyd boundary and their relation with the Bowditch boundary.
This boundary was introduced by Floyd in \cite{Floyd} and we also refer to \cite{Karlssonfreesubgroups} for more details.

Let $\Gamma$ be a finitely generated group and let $\Gx$ denote its Cayley graph associated with a finite generating set $S$.
Let $f:\mathbb{N}\to \mathbb{R}_{\ge 0}$ be a function satisfying that
$$\sum_{n\geq 0}f(n)<\infty$$ and that there exists $\lambda\in (0,1)$ such that
$$1\geq f({n+1})/f(n)\geq\lambda$$ for all $n \in \mathbb{N}$.
The function $f$ is then called the {\it rescaling function} or the \textit{Floyd function}.
In the following, we will always choose an exponential Floyd function, that is $f(n)=\lambda^n$ for some $\lambda\in (0,1)$.
Fix a basepoint $o\in \Gamma$ and rescale $\Gx$ by declaring the length of an edge $\sigma$ to be $f(d(o,\sigma))$.
The induced length distance on $\Gx$ is called the {\it Floyd distance} with respect to the basepoint $o$ and Floyd function $f$ and is denoted by $\delta_{f,o}(.,.)$.
Whenever $f$ is of the form $f(n)=\lambda^n$, we will write $\delta_{\lambda,o}=\delta_{f,o}$ and if $\lambda$ is fixed, $\delta_o=\delta_{\lambda,o}$.

The Floyd compactification $\overline{\Gamma}_{\mathcal{F}}$  is the Cauchy completion of $\Gx$ endowed with the Floyd distance.
The Floyd boundary is then defined as $\partial_{\mathcal{F}}\Gamma=\overline{\Gamma}_{\mathcal{F}}\setminus \Gx$. Different choices of base-points yield bi-Lipschitz Floyd distances because
\begin{equation}\label{changeofbasepointFloyd}
\forall x,y\in \Gamma,\; \delta_{f,o}(x,y) \le \lambda^{-d(o,o')} \delta_{f,o'}(x,y) 
\end{equation}
so the corresponding Floyd compactifications are bi-Lipschitz. Note  that the topology  may  depend on the choice of the rescaling function and the generating set. 

The following fact proved in   \cite{Karlssonfreesubgroups} plays a crucial role in understanding the Floyd geometry.
\begin{lem}\label{Karlssonlem}\cite[Lemma~1]{Karlssonfreesubgroups}
For any $\delta>0$, there exists a function $\kappa=\kappa(\delta)$ with the following property.
If  $x,y,z\in \Gamma$ are three points so that $\delta_x(y,z)\ge \delta$ then $d(x, [y,z])\le \kappa$.    
\end{lem}

If $\sharp \partial_{\mathcal{F}}\Gamma\ge 3$, Karlsson proved in \cite{Karlssonfreesubgroups}  that $\Gamma$ acts  by homeomorphism on $\partial_{\mathcal{F}}\Gamma$ as a convergence group action. By the general theory of convergence groups,   the elements in $\Gamma$ are subdivided into the classes of elliptic, parabolic and loxodromic elements. The latter two   being   infinite order  elements have exactly one and two fixed points in $\partial_{\mathcal{F}}\Gamma$ accordingly.
Moreover,  in this case the Floyd boundary contains uncountable many points and so the cardinality of $\partial_{\mathcal{F}}\Gamma$ is either 0, 1, 2 or uncountably infinite. By   \cite[Proposition~7]{Karlssonfreesubgroups}, $\sharp \partial_{\mathcal{F}}\Gamma= 2$ exactly when the group $\Gamma$ is virtually infinite cyclic.
Following Karlsson, we say that the Floyd boundary is \textit{trivial} if it is finite. The non-triviality of Floyd boundary does not depend on the choice of generating sets \cite[Lemma~7.1]{Yanggrowthtightness}.
We will only have to deal with groups with non-trivial Floyd boundary.


\medskip
We now assume that $\Gamma$ is non-elementary relatively hyperbolic.
We denote by $\partial_{\mathcal{B}}\Gamma$ its Bowditch boundary.
The following is due to Gerasimov.
\begin{thm}\label{mapGerasimov}\cite[Map Theorem]{Gerasimov}
There exists $\lambda_0\in (0,1)$ such that for every $\lambda\in[\lambda_0,1)$, the identity on $\Gamma$ extends to a continuous and equivariant surjection $\phi$ from the Floyd compactification to the Bowditch compactification of $\Gamma$.
\end{thm}

Actually, Gerasimov only stated the existence of the map $\phi$ for one Floyd function $f_0=\lambda_0^n$, but then
Gerasimov and Potyagailo proved that the same result holds for any Floyd function $f\geq f_0$, see \cite[Corollary~2.8]{GePoJEMS}.
They also proved that the preimage of a conical limit point is reduced to a single point and described the preimage of a parabolic limit point in terms of the action of $\Gamma$ on $\partial_{\mathcal{F}}\Gamma$, see precisely \cite[Theorem~A]{GePoJEMS}.
From now on, the parameter $\lambda$ will always be assumed to be contained in $[\lambda_0,1)$.


The Floyd distance can be transferred to a distance on the Bowditch boundary using the map $\phi$.
The resulting distance is called the \textit{ shortcut distance} and we denote it by $\overline{\delta}_{e,\lambda}$ or $\overline{\delta}_e$ if $\lambda$ is fixed.
It is the largest distance on the Bowditch boundary satisfying that for every $\xi ,\zeta\in \partial_\mathcal{F}\Gamma$,
\begin{equation}\label{shortcutsmallerFloyd}
    \bar{\delta}_{e,\lambda}(\phi(\xi),\phi(\zeta))\leq \delta_{e,\lambda}(\xi,\zeta).
\end{equation}
We refer to \cite[Section~4]{GePoGGD} for more details on its construction.

If $\Gamma$ is hyperbolic, then the Gromov, Bowditch and Floyd boundary all coincide.
Thus, the shortcut distance and the Floyd distance are the same.
Furthermore, by \cite[Proposition~6.1]{PY}, the visual distance and the Floyd distance are bi-Lipschitz.

The next couple of lemmas will be used later on.





\begin{lem}\label{ExtensionLem}
Suppose that $\Gamma$ admits a non-trivial Floyd boundary. Then there exist a finite set $F$ of elements and constants $c\ge 1, \delta>0$ with the following property:

for any two elements $g, h\in \Gamma$, there exists $f\in F$ such that $gfh$ labels a $c$-quasi-geodesic and 
$$
\max\{d(g, [e, gfh]), \; d(gf, [e, gfh])\}\le \epsilon
$$
and $\delta_g(e, gfh)\ge \delta$.
\end{lem}
\begin{proof}
Note that if the Floyd boundary of $\Gamma$ is nontrivial, then $\Gamma$ is not virtually cyclic, and every hyperbolic element is strongly contracting \cite{Yanggrowthtightness}.
Thus, the extension lemma in \cite{YangSCC} applies in this setting.
Namely, let $F$ any set of three independent hyperbolic elements. Set $F^n=\{f^n: f\in F\}$ for given $n\ge 1$. Then for any sufficiently large $n_0$, and for any $g,h\in \Gamma$ there exists $f\in F^{n_0}$ such that $gfh$ labels a $c$-quasi-geodesic for a uniform constant $c$. 

It remains to prove that $\delta_g(e, gf^{2n_0}h)$ has a uniformly lower bound $\delta$ when $n_0$ is large.   Indeed, since  every $ f\in F$ is a hyperbolic element with two distinct fixed points, there exists $\delta=\delta(F)>0$ such that $\delta_e(f^{-n_0}, f^{n_0})\ge \delta$ for any $f\in F$ and $n\ge 1$.
Since $gf^{2n_0}h$ labels a $c-$quasi-geodesic, we see that $d(e, f^{n_0}[e,h])$ and $d(e, f^{-n_0}[e,g^{-1}])$  increase to $\infty$  as $n\to\infty$.
By Lemma \ref{Karlssonlem}, we have  for $n\ge n_0$ $\delta_e(f^n,f^nh)<\delta/4$  and $\delta_e(f^{-n_0},f^{-n_0}g^{-1})<\delta/4$.
Thus, $\delta_e(f^{-n_0}g^{-1}, f^{n_0}h)>\delta/2$, and then $\delta_{gf^{n_0}}(e, gf^{2n_0}h)\ge \delta/2$. Consequently, there exists $\delta'=\delta'(\delta, n_0)$ such that $\delta_{g}(e, gf^{2n_0}h)\ge \delta'$.
\end{proof}

%\begin{lem}\label{FloydRayLem}\begin{comment}
%{\color{blue}We do not use this lemma anymore. Should we remove it ?}
%\end{comment}
%Let $\delta_{e}(x,\xi) \le  \lambda^{n}$ for two points $x\in\Gamma, \xi\in \partial_{\mathcal{F}}\Gamma$. Then there exists $D>0$ such that $d(e, \gamma)>  n -D$ for any Floyd geodesic $\gamma$ between $x$ and $\xi$. 
%\end{lem}
%\begin{proof} 
%Choose $y\in\gamma$ such that $d(e,y)=d(e, \gamma)$. By definition, the Floyd length of any path from $y$ to $\xi$ must be bigger than that of the radial geodesic ray from $y$ extending the direction $[e,y]$. Thus, if $m:=d(e,y)\le n -D$, we have   
%$$
%\delta_{e}(y,\xi) \ge \sum_{i\ge m} \lambda^i\ge \frac{\lambda^{n-D+1}}{1-\lambda}
%$$
%Since $y$ is a point on a Floyd geodesic, we have      $\delta_{e}(y,\xi) \le \delta_{e}(x,\xi) \le  \lambda^{n}$, we obtain  $$
%(1-\lambda)\ge \lambda^{-D+1}.
%$$
%This  
%would be impossible for a sufficiently large $D$, so the proof is complete.
%\end{proof}

Floyd and Bowditch boundaries are visual: any two distinct points can be connected by a geodesic.  This enables us to define 
the notion of  shadows on both of them, that will be used in our arguments. 
Given $K>0$ and $x\in \Gamma$, let $\Pi_K(x)$ be the set of boundary points $\xi$ for which some geodesic between $e$ and $\xi$ intersects $B(x,K)$.
We call $\Pi_K(x)$ the big shadow at $x$ of width $K$. 
Balls and shadows are related by \cite[Lemmas~4.13, 4.14, 4.15]{PY}.
We prove here a slight generalization of these results.

\begin{lem}\label{diametershadow}
There exists $C$ such that the diameter of the big shadow $\Pi_K(g)$ is bounded by $CK\lambda^{|x|-K}$ for either the Floyd distance on the Floyd boundary or the shortcut distance on the Bowditch boundary.
\end{lem}

\begin{proof}
By~(\ref{shortcutsmallerFloyd}) we only need to give the proof for the Floyd boundary.
Let $\xi,\zeta$ in $\Pi_K(g)$ and let $[e,\xi]$ and $[e,\zeta]$ be two geodesics intersecting $B(x,K)$ at $y$ and $z$ respectively.
Then, following back $[e,\xi]$ from $\xi$ to $y$, then following a path from $y$ to $z$ that stays inside $B(x,K)$ and finally following the geodesic $[e,\zeta]$ from $z$ to $\zeta$ yields a path from $\xi$ to $\zeta$ of Floyd length at most
$$\sum_{k\geq |x|-K}\lambda^k+2K\lambda^{|x|-K}+\sum_{k\geq |x|-K}\lambda^k\leq CK\lambda^{|x|-K}.$$
This concludes the proof.
\end{proof}

\subsection{Transition points and Floyd geometry}\label{sectiontransitionpoints}
In contrast with hyperbolic groups, the Cayley graph of relatively hyperbolic groups is not Gromov hyperbolic anymore, so the thin triangle property and the Morse property do not hold in general. However, a certain kind of ``relative" Morse property persists and is manifested in a notion of transition points introduced in \cite{Hruska} (see also \cite{DrutuSapir, GePoGGD}). 

Recall that $\mathbb P=\{gP: g\in \Gamma, P\in\mathbb P_0\}$ is the collection of all parabolic cosets. 
\begin{defn}
Let $P\in\mathbb P$ be a parabolic coset and $\eta, L>0$ be fixed constants. A point $p$ on a geodesic $\alpha$ is called \textit{$(\eta, L)$-deep} in $P$ if $B(p, 2L)\cap \alpha\subseteq N_\eta(P)$. It is called \textit{$(\eta, L)$-transitional} if it is not \textit{$(\eta, L)$-deep} in any parabolic coset $P\in \mathbb P$.  
\end{defn}

According to the definition, it is clear that an $(\eta, L_1)$-transition point is an $(\eta, L_2)$-transition point for $L_1\le L_2.$ 
The parameters $\eta, L>0$ are usually chosen via the \textit{bounded intersection} property of the collection $\mathbb P$ (see \cite[Lemma~4.7]{DrutuSapir}): for any $\eta>0$ there exists $L=L(\eta)>0$ such that for any two $P\ne P'\in\mathbb P$ we have $\diam{N_\eta(P)\cap N_\eta(P')}\le L.$


\begin{lem}\label{EntryTransLem}\cite[Proposition~5.6]{GePoGGD}, \cite[Theorem~4.1]{DrutuSapir}.
For large enough $\eta$, there exists $L=L(\eta)$ such that any point of a geodesic $\alpha$ can  be $(\eta, L)$-deep in at most one  $P\in \mathbb P$. 
Moreover, if it is $(\eta, L)$-deep in $P$, the  entry and exit points of $\alpha$ at $N_\eta(P)$ are $(\eta, L)$-transitional. 
\end{lem}


 

 

The following result which refines Lemma~\ref{Karlssonlem} explains the application of the Floyd geometry in relatively hyperbolic groups. 
\begin{lem}\label{TransLargeLem}\cite[Corollary~5.10]{GePoGGD} 
For every large enough $\eta$, there exist $L=L(\eta)$ and $\delta=\delta(\eta)>0$ such that if $x$ is an $(\eta, L)$-transition point  on a geodesic $\alpha$ with endpoints $\alpha_-$ and $\alpha_+$, then $\delta_{x}(\alpha_-,\alpha_+)\ge \delta$. 
%Conversely, for any $\delta,\eta>0$, there exists $L=L(\delta,\eta)>0$ such that if $x$ is a point on a geodesic $\alpha$ with $\delta_{x}(\alpha_-,\alpha_+)\ge \delta$, then $x$ is an $(\eta, L)$-transition point.
\end{lem}



%The Cayley graph of a relatively hyperbolic group enjoys  the following version of a relative thin triangle property.
%\begin{lem}\label{RelThinTriangleLem'}
%There exist universal constants $\tau, R>0$ with the following property.
%For any geodesic triangle $\Delta(x_0,x_1,x_2)$ in $\Gamma$, we have 
%\begin{enumerate}
%    \item 
%    either there exists  a point $w\in \Gamma$ such that $d(w, [x_i,x_{i+1}])\le R$ for each $i \mod 3$.
%    \item
%    or there exists a unique parabolic coset $Y\in \mathbb Y$ such that $diam(N_\tau(Y)\cap [x_i,x_{i+1}])\ge R$ for each $i \mod 3$.
%\end{enumerate} 

%\end{lem}

%\begin{lem}\label{RelThinTriangleLem}
%There exists a universal constant $\delta>0$ with the following property.
%For any triple of elements $(x_0,x_1,x_2)$ in $\Gamma$, there exists a triple of elements $(w_0,w_1,w_2)$ in $\Gamma$  such that the following holds for $i \mod 3$:
%\begin{enumerate}
%    \item 
%$
%d(x_i,x_{i+1}) \sim d(x_i,w_i)+d(w_i,w_{i+1})+d(w_{i+1},x_{i+1}).
%$
%\item 
%$\delta_{w_i}(x_i,x_{i+1})\ge \delta$ and $\delta_{w_{i+1}}(x_i,x_{i+1})\ge \delta$.
%\item 
%If $diam(\{w_0, w_1,w_2\})>M$, then there exists a unique parabolic coset $Y\in \mathbb Y$ such that $d(w_i, Y)\le M$.
%\end{enumerate}
%\end{lem}

%The (ordered) triple $(w_0,w_1,w_2)$ in the statement is called a \textit{center} of the triangle $\Delta(x_0,x_1,x_2)$.




The following is called the relative thin triangle property for transition points.
It is folklore and was proved at several places, using different terminology, see for instance \cite[Proposition~7.1.1]{GePoGGD}, \cite[Proposition~4.6]{Sistometricrelhyp}.
It can also be derived from \cite[Theorem~1.12]{DrutuSapir}, \cite[Section~8]{Hruska}, \cite[Proposition~3.15]{Osin}.
Usually, it is stated for points $x,y,z$ in the $\Gamma$.
The following version for points in the boundary is proved in \cite[Lemma~2.4]{DG20}.

%\begin{lem}\label{ThinTransitional}
%For large enough $\eta$, there exists $L(\eta)$ such that for every $L\geq L(\eta)$, the following holds.
%There exists a  constant $C = C(\eta, L)>0$ such that for given geodesic triangle, any $(\eta, L)$-transitional point on a side is $C$-close to a $(\eta, L)$-transitional point on one of the two other sides.
%\end{lem}
%\begin{comment}
%    I have no objection to your new formulation. The things bothering me are as follows: we need to make a consistent choice of notations. For me, I would like to reserve $\epsilon$ as arbitrary numbers in the arguments of computing HDim. How about $(\eta, R)$ or other choices? Also $\mathbb Y$ might be changed to $\mathbb P$? These things are tiny, but annoying.

%    Lastly, we must   give proper credit  to  Lemma 2.15. As this kind of result is well-known now, different authors may have different (almost equivalent) formulations. We should cite not just Sisto, but also Hruska, Drutu-Sapir and also Potyagailo-Gerasimov, to avoid complaints ...  
    
%\end{comment}
%{\color{blue} I absolutely agree with your comment. I will clean this and add the proper references.
%About notations, the problem is that $\mathbb P$ is already used for probabilities and $r$ and $R$ are used for the parameter of the Green function. In one paper with Ilya, we had the same problem, and we chose $(\eta_1,\eta_2)$-transition points. But to be consistent with what follows, we could maybe choose $(L_1,L_2)$-transition points and then $L$-transition points means $(L_1,L_2)$-transition points for suited $L_1,L_2)$, where $l_2$ depends on $L_1$.
%Then, the letter $\eta$ can be used for radii of balls to avoid the letters $r$ and $R$ there.}
%{\color{red}
%I think we can give simpler arguments.
%Moreover, the exposition is a bit messy, because when choosing $w$ thanks to Lemma~2.9 below, this choice is not unique at all. If one only require $w$ to be within bounded distance of transition points on $[e,x]$ and $[e,y]$, then any points close to $e$ will match.
%Moreover, the thin triangle property along transition points was already used at several places and was proved in our paper with Ilya entropy and drift.
%Finally, I don't find the references to these technical lemmae clear enough in the proof of Lemma~2.16.
%Also, Lemma~2.9 is not proved, but we can refer to Lemma~4.16 in my paper Local limit theorem I, which state a simpler form of Lemma~2.9 with relative geodesics instead of transition points on actual geodesics.

%If you agree, I would replace this small discussion by the following and put it at the beginning of the section and add the details in red in the proof of Lemma~2.16.

%Let $\mathbb Y=\{gP: g\in G, P\in\mathbb P\}$ be the collection of all parabolic cosets.
%\begin{defn}
%Let $Y$ be a parabolic coset and $L_1, L_2>0$ be fixed constants. A point $p$ on a geodesic $\alpha$ is called \textit{$(L_1, L_2)$-deep} in $Y$ if $B(p, 2L_2)\cap \alpha\subseteq N_{L_1}(Y)$. It is called \textit{$(L_1, L_2)$-transitional} if it is not \textit{$(L_1, L_2)$-deep} in any parabolic coset $Y\in \mathbb Y$.  
%\end{defn}




\begin{lem}\label{ThinTransitional2}
For large enough $\eta$, there exist $L(\eta)$ such that for every $L\geq L(\eta)$, there exists $C=C(\eta,L)$ such that for every triple of points $(x,y,z)$ that are either conical limit points or elements of $\Gamma$, any $(\eta,L)$-transition point on one of the side of the geodesic triangle with vertices $x,y,z$ is within $C$ of an $(\eta,L)$-transition point on one of the two other sides.
\end{lem}
In what follows, we fix $\eta$ and $L(\eta)$ so that any pair $(\eta,L)$ satisfies the above lemmas for $L\geq L(\eta)$.










\subsection{Random walks on relatively hyperbolic groups}
Let $\Gamma$ be a relatively hyperbolic group and let $\mu$ be a finitely supported symmetric admissible probability measure on $\Gamma$.
Let $\rho$ be the spectral radius of the random walk and $R$ its inverse.
We collect here some results that will be used all along the paper.
Recall that $\eta$ and $L(\eta)$ are fixed such that for every $L\geq L(\eta)$, any $(\eta,L)$-transition point satisfies the results in Section~\ref{sectiontransitionpoints}.


A very useful set of inequalities relating the Green functions along geodesics were first proved   by Ancona \cite{Ancona} in hyperbolic groups and used to  identify  the Martin boundary   with the Gromov boundary.
These inequalities were recently extended up to the spectral radius by Gou\"ezel-Lalley in co-compact Fuchisan groups \cite{GouezelLalley} and by Gou\"ezel in general hyperbolic groups \cite{Gouezel}.
They state that there exists $C$, depending only on the hyperbolicity parameters of the group, such that for every $x,y,z\in \Gamma$ such that $y$ is on a geodesic from $x$ to $z$, for every $r\leq R$,
$$\frac{1}{C}G_r(x,y)G_r(y,z)\leq G_r(x,z)\leq C G_r(x,y)G_r(y,z).$$
The only non-trivial part is the upper-bound.
In relatively hyperbolic groups, a relative version of Ancona inequalities in terms of Floyd distance was obtained in \cite{GGPY} to establish a surjective map from the Martin boundary to the Floyd boundary.

Ancona inequalities are one of the main ingredient in \cite{SWX} for studying branching random walks on hyperbolic groups.
In the present paper, we will make very crucial use of the relative Ancona inequalities extended up to the spectral radius in \cite{DG21}.

\begin{prop}[Relative Ancona inequalities]\cite[Theorem~3.6]{DG21}
For every $L\geq L(\eta)$ and $K\geq0$, there exists $C=C(\eta, K)$ such that the following holds.
Let $x,y,z\in \Gamma$ and assume that $y$ is within $K$ of an $(\eta,L)$-transition point on $[x,z]$.
Then for every $r\leq R$, we have
$$\frac{1}{C}G_r(x,y)G_r(y,z)\leq G_r(x,z)\leq C G_r(x,y)G_r(y,z).$$
\end{prop}

Note that there also exists a strong form of relative Ancona inequalities in \cite{DG21}, although we will not need them in this paper.
The following result is one of the step into proving relative Ancona inequalities in \cite{DG21}.
It will be useful in this paper.
If $A\subset \Gamma$ and if $x,y\in \Gamma$, for every $r\leq R$ we denote by $G_r(x,y;A)$ the Green function from $x$ to $y$ restricted to trajectories staying in $A$, except maybe at the end points, i.e.
$$G_r(x,y;A)=\sum_{n\geq 0}\sum_{z_1, \ldots ,z_{n-1}\in A}\mu(x^{-1}z_1)\mu(z_1^{-1}z_2) \cdots \mu(z_{n-1}^{-1}y)r^n.$$

\begin{lem}\label{SupExpDecayLem}\cite[Proposition 3.5]{DG21}
For every $L\geq L(\eta)$, there exist  $\delta>1$ and $K_0>0$   such that the following holds.  
For every $x,y,z\in \Gamma$ such that  $y$ is an $(\eta, L)$-transition point on $[x,z]$ and for every $K\geq K_0$, we have
$$G_R(x,z; B(y,K)^c)\le e^{- e^{\delta K}}.$$
\end{lem}





Finally, we will also use the following result, proved in \cite{DG21}.
If $P\in \mathbb P$ is a parabolic coset, $\eta\geq 0$ and $x\in \Gamma$, we denote by  $$\pi_{N_\eta(P)}(x) :=\{y\in N_\eta(P): d(x,y)=d(x, N_\eta(P))\}$$ the set of its shortest projections on the $\eta$-neighborhood $N_\eta(P)$ of $P$. 
For $x,y\in \Gamma$, we set $$d_{N_\eta(P)}(x,y):=\diam{\pi_{N_\eta(P)}(x)\cup\pi_{N_\eta(P)}(y)}.$$
It follows from \cite[Corollary~8.2]{Hruska} that the shortest projection is coarsely Lipschitz: 
$$d_{N_\eta(P)}(x,y)\le kd(x,y)+k$$ for a fixed $k\ge 1$ depending only on $\eta$. Thus, $\pi_{N_\eta(P)}(x)$ has  bounded diameter.

\begin{lem}\label{lemmaexponentialmoments}
Let $P\in \mathbb P$ be a parabolic coset.
For every $M\geq 0$, there exists $\eta_0$ such that for $\eta\geq \eta_0$, we have
$$G_R(e,x;N_\eta(P)^c)\leq \mathrm{e}^{-Md_{N_\eta(P)}(e,x)}.$$
%  \lw{We should use $d_{N_{\eta}(P)}(e, x)$ instead of $d(x, \pi_{N_{\eta}(P)}(e))$, or add some assumption on $x$, right?  Since it may happen that $x$ is on the geodesic from $e$ to $N_{\eta}(P)$ and so $G_R(e, x; N_{\eta}(P)^c) \approx G_R(e, x)$ is not exponentially decay in $d(x, \pi_{N_{\eta}(P)}(e))$.   This does not affect our application since we only use this lemma to $x$ with $d(x, \pi_{N_{\eta}(P)}(e)) \leq d_{N_{\eta}(P)}(e, x)$.{\color{blue}Yes of course, sorry. In the application of \cite{DG21} in the proof below, the number $l$ is in fact $\asymp$ to $d_{N_\eta(P)}(x_0,x)$, not $d(\pi_{N_\eta(P)}(x_0),x)$.}}
\end{lem}

\begin{proof}
In \cite[Lemma~4.6]{DG21} this result is stated for the first return kernel to $N_\eta(P)$, i.e. the quantity
$G_R(x_0,x;N_\eta(P)^c)$ where $x_0\in N_\eta(P)$, but the proof can be applied here.
Indeed it is shown without assuming that $x_0\in N_\eta(P)$ that the whole contribution of trajectories from $x_0$ to $x$ staying outside $N_\eta(P)$ is bounded by
$\mathrm{e}^{-l h(\eta)}G_R(x_0,x)$, where $\frac{1}{c}l\leq d_{N_\eta(P)}(x_0,x)\leq c l$ for some $c>0$ 
and $h(\eta)$ is a function of $\eta$ going to infinity as $\eta$ goes to infinity, see the before last equation of the proof of \cite[Lemma~4.6]{DG21}.
In particular, applying this to $x_0=e$, we get that for large enough $\eta$,
$$G_R(e,x;N_\eta(P)^c)\leq \mathrm{e}^{-M d_{N_\eta(P)}(e,x)}.$$
This concludes the proof.
\end{proof}
%\subsection{Hausdorff dimension}
%Let $X$ be a subset of $\Gamma$.  
%We say that $\widetilde{X}$ is a \textit{tight augmentation} of $X$ if $X\subset \widetilde{X}$ and the remainder $\widetilde{X}\setminus X$ has growth rate strictly smaller than that of $X$: $v(X)>v({\widetilde{X}\setminus X})$. 

%Recall that the Poincaré series of $X$ as follows 
%$$
%\Theta_s(X) = \sum_{x\in X} e^{-s|x|} = \sum_{n\ge 0} |S_n\cap X| e^{-sn}
%$$
%diverges at $s<v(X)$ and converges $s>v(X)$. Thus, the Poincaré series of $\widetilde{X}$ has the same convergence property as that of $X$. 


%\begin{lem}\label{ComputeHdim}
%Let $X\subset \Gamma$.  
%Suppose a subset $\Lambda\subset \partial \Gamma$ is contained in the limit super of the shadows $\{\Pi(x,\epsilon |x|): x\in X\}$ for any given $\epsilon>0$. Then $$\Hdim(\Lambda, \delta_e)\le -v(X)/\log \lambda.$$     
%\end{lem}
%\begin{proof}
%For any $C>0$, there exists $c=c(C)$ such that  $$\Pi(x, C) \subset B_{\delta_e}(\xi, c\lambda^{|x|})$$ for any $\xi\in \Pi(x,C)$ and $x\in \Gamma$. Thus, for any $s>0$, $$[\diam{\Pi(x,\epsilon |x|)}]^s\prec \lambda^{s|x|(1+\epsilon)} \prec e^{-s|x| (1+\epsilon)\log \lambda }$$ 
%If $s>-v(X)/\log \lambda$, there exists $\epsilon>0$ such that $-s  (1+\epsilon)\log \lambda>v(X)$. Thus, $$\mathcal H^s(\Lambda) \prec \sum_{x\in X} e^{-s|x| (1+\epsilon)\log \lambda } \le \sum_{n\ge 1} |X_n| e^{-s n (1+\epsilon)\log \lambda} <\infty $$ The conclusion follows.
%\end{proof}

\section{The growth rate of the Green function}\label{SectiongrowthrateGreenfunction}


\subsection{Preliminary results}
Recall the following definitions from the introduction.
Let $\Gamma$ be a finitely generated group endowed with a finite generating set.
Let $\mu$ be a symmetric
probability measure whose finite support generates $\Gamma$. Denote by $\rho$ the spectral
radius of $\mu$ and by $R$ its inverse.
%When there is no ambiguous context, we will simply write $\rho=\rho_\mu$ and $R=R_\mu$.
Set $p_n(x,y)=\mu^{\star n}(x^{-1}y)$ for $n\ge 1$ and $p_0(x,y)=\delta_x(y)$. Then, $R$ is the radius of convergence of the
Green function
$$
G_r(x, y) :=\sum_{n=0}^\infty r^np_n(x,y).
$$
 


For $1\le r\le R$, we consider   the sum of the $r$-Green function over the $n$-sphere
   $$H_r(n) := \sum_{x\in S_n} G_r(e,x)$$    
and define the \textit{growth rate  of the Green function} as follows 
$$\omega_\Gamma(r):= \limsup_{n\to \infty} \frac{\log H_r(n)}{n}$$
We first record a few simple facts about $H_r(n)$.
\begin{lem}\label{FactsHrLem}
The following statements are true:
\begin{enumerate}
    \item 
      % For any $1\le r\le R_\mu$, there exists $C=C(r)>1$ such that for $n\ge 1$
      There exists $C > 1$ such that for any $1 \leq  r \leq R$ and $n \geq 1$
    $$C^{-1} H_{r}(n+1)\le H_r(n) \le CH_{r}(n+1).$$
    \item
    There exists a constant $C>0$ such that $\frac{1}{C}\leq H_1(n)\le C$ for any $n\ge 1$. In particular, $\omega_{\Gamma}(1) = 0$. 
    \item
      % For every $1<r\le R$, $H_r(n)\ge 1$ grows exponentially, i.e.\ $\omega_\Gamma(r)>0$.
      The function $r \mapsto \omega_{\Gamma}(r)$ is strictly increasing on $[1,R]$ and continuous on $[1,  R)$. Consequently, $\omega_{\Gamma}(r) > 0$ for $r > 1$.
\end{enumerate}

\end{lem}
\begin{proof}
Since the random walk is irreducible and $\Gamma$-invariant, there exist % a uniform number $p_0>0$ such that $p(x,y)>p_0$ for any $x,y\in \Gamma$ with $d(x,y)=1$. Thus, if $x\ne 1\ne y$, we have $p_n(e,x)\ge p_0\cdot  p_{n-1}(e,y)$ for $n\ge 1$. This implies the existence of a constant $C=C(r)>0$ such that $G_r(e, x)\ge C G_r(e,y)$. Thus, we have $H_{r}(n)\le CH_{r}(n+1)$. For the other inequality, note that every  $y\in S_{n+1}$ is adjacent   to at most $N=\sharp S_1$ elements  $x\in S_{n}$.  
% We then obtain $H_{r}(n+1)\le \frac{N}{C} H_{r}(n)$.
$l \ge 1$ and a uniform number $p_0 > 0$ such that $p_l(x, y) > p_0$ for any $x,y\in \Gamma$ with $d(x,y)=1$. Thus, if $x\ne 1\ne y$, we have $p_n(e,x)\ge p_0\cdot  p_{n-l}(e,y)$ for $n\ge l$. This implies $G_r(e, x)\ge p_0 G_r(e,y)$. Thus, we have $H_{r}(n+1)\ge p_0 H_{r}(n)$.
For the other inequality, note that every  $y\in S_{n+1}$ is adjacent   to at most $N=\sharp S_1$ elements  $x\in S_{n}$.  
We then obtain $H_{r}(n+1)\le \frac{N}{p_0} H_{r}(n)$.
This proves the first statement.


We introduce the partial shadow $\widetilde{\Pi}_K(x)$ of width $K$ at $x$ as the set of limit points such that some geodesic $[e,\xi]$ intersects the ball $B(x,K)$ at a transition point. By the Shadow Lemma \cite[Proposition~4.4]{DG20} for harmonic measures, there exists $K>0$ such that $G_1(e,x)\asymp \nu(\widetilde{\Pi}_K(x))$ for any $x\in \Gamma$, which means that the ratio of these two quantities is bounded away from 0 and infinity.
Note that for each $n\ge 1$, any conical limit point can be covered in a uniform number of shadows at $x\in S_n$. We thus obtain that the sum $\sum_{x\in S_n} \nu(\widetilde{\Pi}_K(x))$ coarsely gives the measure of the whole set of conical limit points, so is uniformly bounded from above and below. The second statement follows.


Finally let us prove the third statement.
Let $1\le s<r$.
Since the random walk is finitely supported, there exists $c_1>0$ such that for every $x$, $p_m(e,x)=0$ for every $m\leq c_1|x|$.
Thus, we have
$$G_s(e,x)=\sum_{m\geq c_1 |x|}s^m p_m(e,x)
\leq \left(\frac{s}{r}\right)^{c_1 |x|}\sum_{m \geq c_1 |x|}r^m p_m(e,x)
=\left(\frac{s}{r}\right)^{c_1 |x|}G_r(e,x)$$ and so
\begin{equation}\label{exponentialdifferenceGreen}
G_s(e,x)\leq \left(\frac{s}{r}\right)^{c_1 |x|}G_r(e,x).
\end{equation}
% Hence, there is $\alpha>0$ such that $H_s(n)\leq \mathrm{e}^{-\alpha n}H_r(n)$.
% By \cite[Note~1.7]{GouezelLalley}, $H_s(r)$ diverges when $n$ goes to infinity, so $\omega_\Gamma(r)\geq \alpha>0$.
Thus, $H_s(n) \leq \left( \frac{s}{r} \right)^{c_1 n} H_r(n)$ and $\omega_{\Gamma}(r) -  \omega_{\Gamma}(s) \ge c_1 \left( \log r  - \log s\right)$.  Therefore $\omega_{\Gamma}(r)$ is strictly increasing on $[1, R]$.

  For $\delta > 0$ we choose $c_2$ so that $v \le c_2 \left( \log R - \log (R - \delta) \right)$.
  Note that since the underlying random walk is symmetric, for every $x$ and every $m$, we have
  $p_m(e,x)p_m(e,x)\leq p_{2m}(e,e)$ and by \cite[Lemma~1.9]{Woessbook}, $p_{2m}(e,e)\leq R^{-2m}$.
  Thus,
  $p_m(e, x) \leq R^{-m}$ for every $x \in \Gamma$ and $n \ge 0$.
  Consequently, by Lemma~\ref{purelyexponentialgrowth}, we have for $1 \le r \le R - \delta$,
\[
\sum_{|x| = n} \sum_{m > c_2 n} r^m p_m(e, x) \le c\mathrm{e}^{v n} \sum_{m > c_2 n} \left( \frac{r}{R} \right)^m \le c_3 
\]
for some constant $c_3 > 0$. Now for $1 \le s < r \le R - \delta$,
\begin{align*}
      H_s(n) = \sum_{|x| = n} \sum_{m = 0}^{\infty} s^m p_m(e, x)
  &\ge \left( \frac{s}{r} \right)^{c_2 n} \sum_{|x| = n} \sum_{m = 0}^{c_2 n} r^m p_m(e, x)\\
  &\ge \left( \frac{s}{r} \right)^{c_2 n} \left( H_r(n) - c_3 \right). 
\end{align*}
It follows that $\omega_{\Gamma}(r) - \omega_\Gamma(s) \le c_2 \left( \log r - \log s \right)$. Since $\delta > 0$ is arbitrary, we prove that $\omega_{\Gamma}(r)$ is continuous in $1 \le r < R$.
\end{proof}


Our goal in the remainder of this section is to compare $H_r(n)$ with $\mathrm{e}^{n \omega_\Gamma(r)}$.

\subsection{Upper bound} 
We start with the following lemma.

\begin{lem}\label{UpperBDLem}
There exists  constants $C,C'>1$ such that for any $1\le r\le R$ and any  integer $n, m\ge 1$, we have
$$H_r(m) H_r(n) \le C H_{r}(m+n) ,$$
 and $H_r(n)\le C' e^{n \omega_\Gamma(r)}$. 
In particular, the following limit exists
$$\omega_\Gamma(r)= \lim_{n\to \infty} \frac{\log H_r(n)}{n}=\sup_n \frac{\log C^{-1}H_r(n)}{n}.$$
\end{lem}



Set $l=\max\{d(e,f): f\in F\}+4\epsilon$, where $F$ is a finite set given by Lemma \ref{ExtensionLem}.
and let
$A(n, l)=\bigcup\limits_{-l\le i\le l}S_{n+i}$ be the annulus of width $l$ and radius $n$ and
define $$\Phi: S_m\times S_n\to A(n+m, l)$$ by setting $\Phi(g,h)=gfh$, where $f\in F$ is provided by Lemma~\ref{ExtensionLem}.
\begin{lem}\label{BoundedToOne}
The map $\Phi$ is uniformly bounded to one.
\end{lem}
\begin{proof}
Indeed, assume that $\Phi(g,h)=gf_1h$ and $\Phi(x,y)=xf_2y$ for $f_1,f_2\in F$. If $\Phi(g,h)=\Phi(x,y)$, we obtain $d(g,x), d(h^{-1},y^{-1})\le 4\epsilon$ from Lemma \ref{ExtensionLem}. Thus, there are at most $3 (\sharp B(e,4\epsilon))^2$ pairs of  elements $(x,y)$ such that  $\Phi(g,h)=\Phi(x,y)$.
\end{proof}
\begin{proof}[Proof of Lemma \ref{UpperBDLem}]
For any $g\in S_n, h\in S_m$, we have 
$$
G_r(e,g) G_r(e,h) \le c_1 G_r(e,g) G_r(e,f) G_r(e,h)\le c_2 G_r(e, gfh) 
$$
Thus, 
$$
H_r(n) H_r(m) \le c_2\sum_{z\in A(n+m,l)} G_r(e,z).
$$
Note that $c_3^{-1} H_{r}(n+1)\le H_r(n) \le c_3H_{r}(n+1)$ by Lemma \ref{FactsHrLem}.
Thus, we have
$$H_r(n) H_r(m) \le c_4 H_r(n+m).$$
This proves the first part of the lemma.
The second part follows from the Fekete sub-additive lemma.
\end{proof}

\begin{cor}\label{lowersemicontinuity}\label{corolowersemicontinuous}
The function $r\in [1,R]\mapsto \omega_\Gamma(r)$ is increasing and continuous. % lower semi-continuous.
\end{cor}

\begin{proof}
  % By~(\ref{exponentialdifferenceGreen}), $\omega_\Gamma$ is increasing.
We proved in Lemma \ref{FactsHrLem} that $\omega_{\Gamma}$ is increasing and continuous on $[1, R)$.
By Lemma~\ref{UpperBDLem}, it can be expressed as a supremum of continuous functions, hence it is lower semi-continuous by \cite[Lemma~2.41]{AliprantisBorder}.
We deduce that it is left continuous at $R$.
\end{proof}

\begin{cor}\label{coroboundv/2}
For every $r\leq R$, $\omega_\Gamma(r)\leq v/2$.
\end{cor}

\begin{proof}
By the Cauchy-Schwarz inequality,
$$\left(\sum_{x\in S_n}G_r(e,x)\right)^2\leq |S_n|\sum_{x\in S_n}G_r(e,x)^2.$$
For any $r<R$, by \cite[Proposition~1.9]{GouezelLalley},
$$c(r)=\sum_{x\in \Gamma}G_r(e,x)^2<\infty,$$
so
$$\omega_\Gamma(r)\leq \limsup \frac{1}{2n}\left(\log c(r)+\log |S_n|\right)=\frac{1}{2}v.$$
This proves the desired inequality for $r<R$.
By Corollary~\ref{lowersemicontinuity}, $r\mapsto \omega_{\Gamma}(r)$ is %lower semi-continuous,
continuous, so the inequality also holds at $R$.
\end{proof}


\subsection{Lower bound via parabolic gap} 

For $1\le r\le R$ and $s\ge 0$, we consider the following Poincar\'e series:

\begin{equation}\label{defPoincareseries}
\Theta_{r, s}(\Gamma) := \sum_{h\in \Gamma} G_r(e, h) e^{-sd(e, h)}.
\end{equation}
 We can rearrange the terms in $\Theta_{r, s}(\Gamma)$ as follows

$$\Theta_{r, s}(\Gamma) = \sum_{n\ge 0} H_r(n) e^{-sn}$$
so that  $H_r(n)$ appears in the place of $\sharp S_n$ in the usual Poincare series. Thus, for each $r$ fixed, $\omega_\Gamma(r)$ is the exponential radius of convergence of the series $\Theta_{r, s}(\Gamma)$ in $s$.



Similarly, for any subgroup $P\subset \Gamma$, we can consider the associated Poincar\'e series
$$\Theta_{r, s}(P)=\sum_{n\geq 0}\sum_{x\in P\cap S_n}G_r(e,x)\mathrm{e}^{-sd(e,x)}$$ and its growth rate
$$ \omega_P(r)=\limsup \frac{1}{n}\log \sum_{x\in P\cap S_n}G_r(e,x),$$
which is the exponential radius of convergence of the series $s\mapsto \Theta_{r,s}(P)$.

\begin{defn}
We say that $\Gamma$ has \textit{parabolic gap}   for the Green functions if for every parabolic subgroup $P\in \mathbb{P}_0$, $\omega_P(r)<\omega_\Gamma(r)$ for every $1< r\le R$.
\end{defn}


\begin{lem}\label{LowerBDLem}
Suppose that $\Gamma$ has a parabolic gap for the Green function.
For every $1<r\leq R$, there exists $C=C(r)>1$ such that 
$$H_{r} (n+m) \le C H_r(n) H_r(m)
$$
and there exists a constant $C'=C'(r)>1$ such that for any $1<r\le R_\mu$, we have 
$$
\frac{1}{C'} e^{n \omega_\Gamma(r)}\le H_r(n)\le C' e^{n \omega_\Gamma(r)}
$$
for every $n\ge 1$.
\end{lem}


\begin{proof}
Let $\mathbb P_0$ be the finite set of maximal parabolic subgroup up to conjugacy.
Set $$K_r(n) := \max_{P\in\mathbb P_0}\sum_{x\in S_n\cap P} G_r(e,x).$$ 



\textbf{Step 1.} First, the following holds  
\begin{equation}\label{threeunion}
H_{r}(n+m) \le \sum_{0\le k \le n} \sum_{0\le j \le m}  c_0 H_{r}(k) \cdot H_{r}(j) \cdot   K_{r}(n+m-k-j)  
\end{equation}
for any $n, m\ge 0$. 

Indeed, for given $x\in S_{n+m}$, consider a  geodesic $\gamma=[e,x]$ and the point $y\in \gamma$ such that $d(e,y)=n$. If $y$ is a transition point, then by the relative Ancona inequalities, we have 
$$
G_r(e,x) \le C G_r(e, y) G_r(e, y^{-1}x).
$$
This corresponds to the case $k=j=0$.

If $y$ is deep in some $P$-coset $X$, let $u,v$ be the entry and exit points of $\gamma$ in $N_\eta(X)$ (possibly $u=e$ or $v=x$). Then $u,v$ are transition points so again the relative Ancona inequalities show   
$$
G_r(e,x) \le C G_r(e, u) G_r(e, u^{-1}v)  G_r(e, v^{-1}x).
$$

Summing up $G_r(e,x)$ over all $x\in S_n$ according to $k=d(u,y)$ and $l=d(y,v)$, we obtain (\ref{threeunion}). 

 
\textbf{Step 2.} 
By assumption,  $\omega_P(r) < \omega_\Gamma(r)$ for every $P\in \mathbb P_0$. Then for any given $\omega_\Gamma(r)  > \omega> \omega_P(r)$, there exists $c_1=c_1(r)>0$ such that $K_r(i)\le c_1 e^{i \cdot \omega }$ for any $i\ge 1$.
For $\omega>0$, we define $$a^\omega(n)= e^{-\omega n} \cdot H_{r}(n).$$ 
Then a re-arrangement of (\ref{threeunion}) gives rise to the form as follows:  
\begin{equation}\label{SubMulEQ}
a^\omega(n+m) \le c_1 \left( \sum_{1\le k\le n} a^\omega(k) \right) \cdot \left(\sum_{1\le j \le m} a^\omega(j) \right),
\end{equation}
for any $n, m \ge 0$.
We conclude as in \cite[Theorem 5.3]{YangSCC}.
\end{proof}

Theorem~\ref{mainthmparabolicgap} is now a consequence of Lemma~\ref{FactsHrLem} for the lower bound at $r=1$, Lemma~\ref{LowerBDLem} for the upper bound and Lemma~\ref{UpperBDLem} for the lower bound at $r>1$.
%\begin{cor}\label{continuityomega}
%If $\Gamma$ has a parabolic gap for the Green function, then $\omega_\Gamma(r)$ is a continuous function of $r$.
%\end{cor}

%\begin{proof}

%\end{proof}


\subsection{Criteria for Green parabolic gap}
A possible way to get this parabolic gap is the following divergence criterion based on \cite[Lemma~2.23]{YangSCC}. 


Let $A, B$ be two subsets of $G$.
Denote by $\mathbb W(A, B)$ the set of all words over the alphabet set $A\sqcup B$ with letters alternating in $A$ and $B$. 

\begin{lem}\label{degrowth}
Assume that there exists an injective  map $\iota:  \mathbb W(A)\to  \mathbb W(A, B)$ such that the evaluation map $\Phi: \mathbb W(A, B) \to G$ is injective on the subset $\iota(\mathbb W(A))$ as well. Set $X:=\Phi(\iota(\mathbb W(A,B))$  and assume that   $B$ is finite. Then   $\Theta_{r,s}(A)$ converges  at $s=\omega_{X}(r)$. In particular, if $\Theta_{r,s}(A)$ diverges at $s=\omega_{A}(r)$ 
then $\omega_{X}(r)
> \omega_{A}(r)$.
\end{lem}


\begin{proof}
Since $\Phi: \mathbb W(A, B) \to G$ is injective, each element in the image $X$ has a unique  alternating product form over $A\sqcup B$. Set  $C_1=\max_{b\in B} \{d(e, b)\}<\infty$, $C_2=\min_{b\in B} \Big\{\frac{G_r(1, b)}{G_r(1,1)}\Big\}<\infty$ where  $B$ is a finite set by assumption. For a word $W=a_1b_1a_2\cdots a_nb_n \in \mathbb W(A, B)$, we have $$
d(e, a_1b_1\cdots a_nb_n) \le \sum_{1\le i\le n} \big (d(e, a_i)+C_1\big),
$$and 
$$
G_r(e, a_1b_1\cdots a_nb_n) \ge C_2^n \prod_{1\le i\le n} \big (G_r(e, a_i) \big).
$$
As a consequence, we estimate the Poincar\'e series of $X$ as follows:
$$
\begin{array}{rl}
&\sum\limits_{g \in X} G_r(e,g)  e^{-s\cdot d(e, g)} \\
 
\ge &   \sum\limits_{n =1}^{\infty} \left(\sum\limits_{a \in A} G_r(e,a) e^{-s\cdot d (e, a)}\right)^n\cdot   (C_2e^{-sC_1}) ^n.
\end{array}
$$


By contradiction, assume that $\sum\limits_{a \in A} G_r(e,a) e^{-\omega_X(r) \cdot d (e, a)} = \infty$. Then there exists some $s >
\omega_X(r)$ such that $\sum\limits_{a \in A} G_r(e,a)  e^{-s\cdot d (e,
a)}\cdot C_2e^{-sC_1}
> 1$. By the above estimates, this implies the  series $\Theta_X(r,s)$ diverges at $s$, so $s\le \omega_X(r)$, which is a contradiction.
\end{proof}

\begin{cor}\label{parabolicdivergentgrowthgap}
Let $P\in \mathbb{P}_0$. 
If the series $$\Theta_{r,s}(P)=\sum_{p\in P} G_r(e,p) e^{-s d(e,p)} $$ diverges at $s=\omega_P(r)$ then $\omega_P(r)<\omega_\Gamma(r)$.
\end{cor}
\begin{proof} 
For any hyperbolic element $h$ and sufficiently large  integer $n\ge 1$, the subgroup generated by $P$ and $h^n$ is a free product $P* \langle h^n\rangle$. 
Thus, the evaluation map $\Phi: \mathbb W(A, B) \to G$ is injective on $\mathbb W(A)$ seen as a subset of $\mathbb W(A,B)$, where $A=P$ and $B=\{h^n,h^{-n}\}$.
\end{proof}

%\begin{rem}
%For usual Poincar\'e series, $\sharp S_{n+m} \le \sharp S_n\sharp S_M$ so the divergence follows automatically. However, this is hard to verify for Green functions, since this inequality needs to use  the Anacona inequality for $P$.  That said, in your suggested example $\Gamma=\Gamma_0\star \Gamma_1$ hyperbolic relative to free groups $\Gamma_0=\mathbb F_2$, $\mathbb F_2$ can be shown divergent.
%\end{rem}

Here is a second criterion for the parabolic gap.

\begin{prop}\label{secondcriterionparabolicgap}
Let $P\in \mathbb P_0$.
For every $1<r\leq R$, if $\omega_P(r)\leq 0$, then $\omega_P(r)<\omega_\Gamma(r)$.
\end{prop}

\begin{proof}
By Lemma~\ref{FactsHrLem}, for every $1<r$, $\omega_\Gamma(r)>0$.
\end{proof}

\subsection{Examples}\label{sectionexamples}
We now give various examples of different possible situations.
In Example~A, we give a criterion for having $\omega_P(r)\leq 0$, which automatically implies that $\omega_P(r)<\omega_\Gamma(r)$.
In Example~B, we construct an example where $\omega_P(r)>0$, but we still have $\omega_P(r)<\omega_\Gamma(r)$.
Example~C is devoted to construct an example where $\omega_P(r)=\omega_\Gamma(r)$, assuming that $P$ satisfies some properties.

We first recall some terminology from \cite{DG21}.
Let $r\leq R$ and let $P$ be a parabolic subgroup.
Denote by $p^{r,P}$ the first return transition kernel to $P$ associated with $r\mu$, i.e.\ for every $x,y\in P$,
\begin{equation}\label{deffirstreturnkernel}
p^{r,P}(x,y)=\sum_{n\geq 1}\hspace{.1cm}\sum_{z_1,...,z_n\notin P}r^n\mu(x^{-1}z_1)\mu(z_1^{-1}z_2)...\mu(z_n^{-1}y).
\end{equation}
Also denote by $G_t^{r,P}$ its associated Green function at $t$, i.e.
$$G_t^{r,P}(x,y)=\sum_{n\geq 0}t^np_n^{r,P}(x,y),$$
where $p^{r,P}_n$ is the $n$th convolution power of $p^{r,P}$.
Finally, denote by $R_P(r)$ the inverse of the spectral radius of $p^{r,P}$.
Then by \cite[Lemma~4.4]{DG21}, $G_1^{r,P}(e,x)=G_r(e,x)$ which is finite, so in particular $R_P(r)\geq 1$.

\begin{defn}\label{defnspectraldegeneracy}
We say that the random walk is \textit{spectrally degenerate} along $P$ if $R_P(R)=1$.
It is called \textit{spectrally non-degenerate} if it is not spectrally degenerate along any parabolic subgroup.
\end{defn}


\subsubsection{Example~A}
\begin{prop}
Let $P$ be a parabolic subgroup.
Assume that $P$ is amenable.
Then for every $1<r< R$, $\omega_P(r)< 0$.
In particular, for every $1<r< R$, $\omega_P(r)<\omega_\Gamma(r)$.
Moreover, if the random walk is not spectrally degenerate along $P$, then $\omega_P(R)\leq 0<\omega_\Gamma(R)$.
\end{prop}

\begin{proof}
Fix $r<R$ and choose $s$ such that $r<s<R$.
Denote by $p^{s,P}$ the first return transition kernel to $P$ associated with $s\mu$ and by $G^{s,P}$ its associated Green function as above.
By \cite[Lemma~4.15]{DG21}, the spectral radius of $p^{s,P}$ is strictly less than 1, because $s<R$.
Since the underlying random walk driven by $\mu$ is symmetric, $p^{s,P}$ is a symmetric $P$-invariant transition kernel on $P$.
By amenability, this transition kernel is necessarily sub-Markov, see \cite[Corollary~12.5]{Woessbook}.
Let $t<1$ be such that $\tilde{p}^s=\frac{1}{t}p^{s,P}$ is Markov.
Then, letting $\widetilde{G}^s$ be the Green function associated with $\tilde{p}^s$, for any $x\in P$ we have $G_1^{s,P}(e,x)=\widetilde{G}_t^s(e,x)$.
We now prove that
\begin{equation}\label{equationGreenparabolicfinite}
\sum_{x\in P}\widetilde{G}_t^s(e,x)<\infty.
\end{equation}
Consider the convolution operator
$$\widetilde{P}_s:f\mapsto \left ( x\mapsto \sum_{y}\tilde{p}^s(x,y)f(y)\right )$$
acting on the space of $\ell_1$ functions, i.e.\ summable functions on $P$.
The $\ell_1$-norm of $\widetilde{P}_s$ is 1, so its $\ell_1$-spectral radius is bounded by 1.
In particular, $1/t$ is bigger than the $\ell_1$-spectral radius, so by definition, $I-t\widetilde{P}_s$ is invertible in the space of summable functions.
Moreover, the inverse is of the form 
$$\widetilde{Q}_s=\sum_{n\geq 0} t^n\widetilde{P}_s^n.$$
Consider the function $f$ defined by $f(x)=1$ if $x=e$ and 0 otherwise.
Then, $\widetilde{Q}_sf(x)=\widetilde{G}_t^s(e,x^{-1})$.
Since the function $f$ is summable, $\widetilde{Q}_sf$ also is summable, and so the Green function $\widetilde{G}^s$ at $t$ is summable.
This proves~(\ref{equationGreenparabolicfinite}).

Since for $x\in P$, $G_s(e,x)=G_1^{s,P}(e,x)=\widetilde{G}_t^s(e,x)$,
we deduce that $\sum_{x\in P}G_s(e,x)<\infty$ and so
$\omega_P(s)\leq 0$.
Moreover, by~(\ref{exponentialdifferenceGreen}), we have
$$G_r(e,x)\leq \left(\frac{r}{s}\right)^{c|x|}G_s(e,x)$$ for every $x\in \Gamma$.
Summing over $P\cap S_n$, we see that $\omega_P(r)<\omega_P(s)$ and so $\omega_P(r)<0$.

Finally, if the random walk is not spectrally degenerate along $P$, then by definition the spectral radius of $p^{R,P}$ is strictly less than 1.
The same proof shows that $\omega_P(R)\leq 0$.
\end{proof}
 
%Thus, if every parabolic subgroup is amenable, the function $\omega_\Gamma(r)$ is continuous on $(1,R)$.
%Since it is non-decreasing and lower semi-continuous, it is also continuous at $R$.
 
 When the parabolic subgroup $P$ has sub-exponential growth, we do not need spectral non-degeneracy to get that $\omega_P(R)\leq 0$.
 \begin{prop}\label{subexpgrowthparabolicgap}
 Let $P$ be a maximal parabolic subgroup.
Assume that $P$ has sub-exponential growth.
Then for every $1\leq r\leq R$, $\omega_P(r)\leq 0$.
In particular, for every $1<r\leq R$, $\omega_P(r)<\omega_\Gamma(r)$.
 \end{prop}
 
 \begin{proof}
 Let $s>0$.
There exists $C\geq 0$ such that for every $x\in \Gamma$, $G_R(e,x)\leq C$.
 In particular, for every $r\leq R$,
 $$\Theta_{r,s}(P)\leq C\sum_{n\geq 0}\sharp (P\cap S_n)\mathrm{e}^{-sn}<+\infty.$$
 Thus, $\Theta_{r,s}(P)$ is finite for every positive $s$ and so $\omega_P(r)\leq 0$.
 \end{proof}
 
 \subsubsection{Adapted random walks on free products}
Before giving other examples, let us briefly recall some terminology and basic properties of random walks on free products.
Consider a free product $\Gamma=\Gamma_0*\Gamma_1$.
Let $\mu_0$ be an admissible probability measure on $\Gamma_0$, $\mu_1$ an admissible probability measure on $\Gamma_1$ and define $$\mu_\alpha=\alpha \mu_1+(1-\alpha)\mu_0.$$
Then $\mu_\alpha$ is an admissible probability measure on $\Gamma$.
Moreover, if both $\mu_0$ and $\mu_1$ are symmetric, respectively finitely supported, then so is $\mu_\alpha$.

Such a probability measure is called adapted to the free product structure and it can only move inside one of the free factors $\Gamma_0$ or $\Gamma_1$ at each step.
Adapted probability measures on free products have been considered by many authors, see \cite{Cartwright1}, \cite{Cartwright2}, \cite{CandelleroGilch}, \cite{CandelleroGilchMuller}, \cite{WoessMartin} and \cite{WoessLLT} for instance.
For convenience we will assume that the random walk driven by $\mu_i$ on $\Gamma_i$ is transient at the spectral radius, i.e.\ the Green function is finite at its radius of convergence.
This is not very restrictive, since by the Varapoulos Theorem, only groups with quadratic growth, i.e.\ groups that are virtually $\mathbb{Z}$ or $\mathbb Z^2$, can carry an admissible random walk which is not transient at the spectral radius.
Actually, Varapoulos \cite{Varapoulos} proved that only groups with quadratic growth can carry a non-transient (at $r=1$) random walk, but a standard argument due to Guivarc'h \cite{Guivarch} allows one to reduce non-transience at $r=R$ to non-transience at $r=1$ by the use of a suited $h$-process, see \cite{Woessbook} for more details and in particular \cite[Theorem~7.8]{Woessbook} for a complete proof.

Denote by $G^{\mu_i}$ the Green function on $\Gamma_i$ associated with $\mu_i$, $i=0,1$ and by $R_{\mu_i}$ the radius of convergence of $G^{\mu_i}$, i.e.\ the inverse of the spectral radius of $\mu_i$.
Also, as in the previous example, denote by $p^{s,\Gamma_0}$ the first return kernel to $\Gamma_0$ associated with $s\mu_\alpha$ and by $G_t^{s,\Gamma_0}$ the associated Green function.
We first relate $G^{\mu_0}$ and $G^{s,\Gamma_0}$ which are two Green functions associated with different transition kernels on the same group $\Gamma_0$.
Because $\mu_\alpha$ is adapted to the free product structure, the first return kernel $p^{s,\Gamma_0}$ can be written as
$$p^{s,\Gamma_0}(e,x)=(1-\alpha)s\mu_0+w_{s,\alpha}\delta_{e,x},$$
where $w_{s,\alpha}$ is the weight of the first return to $e$, starting with a step in $\Gamma_1$.
Thus, \cite[Lemma~9.2]{Woessbook} shows that for any $x,y\in \Gamma_0$,
\begin{equation}\label{superWoess}
G_t^{s,\Gamma_0}(x,y)=\frac{1}{1-w_{s,\alpha}t}G_\frac{(1-\alpha)st}{1-w_{s,\alpha}t}^{\mu_0} (x,y).
\end{equation}
Define
$$\zeta_0(s)=\frac{(1-\alpha)s}{1-w_{s,\alpha}}.$$
As in the previous example, by \cite[Lemma~4.4]{DG21}, $G_1^{s,\Gamma_0}(e,x)=G_s(e,x)$, where as usual $G_s$ denotes the Green function associated with the initial random walk driven by $\mu_\alpha$ at $s$.
So in particular, applying~(\ref{superWoess}) at $t=1$,
\begin{equation}\label{superWoess2}
G_s(x,y)=G_1^{s,\Gamma_0}(x,y)=\frac{1}{1-w_{s,\alpha}}G_{\zeta_0(s)}^{\mu_0}(x,y).
\end{equation}


We also set
$$\zeta_1(s)=\frac{\alpha s}{1-w'_{s,\alpha}},$$
where $w'_{s,\alpha}$ is the weight of the first return to $e$, starting with a step in $\Gamma_0$. 
Then, by symmetry, we also have for any $x,y\in \Gamma_1$,
\begin{equation}\label{superWoess2'}
G_s(x,y)=\frac{1}{1-w'_{s,\alpha}}G_{\zeta_1(s)}^{\mu_1}(x,y).
\end{equation}

Let $R_\alpha$ be the inverse of the spectral radius of $\mu_\alpha$ and let $R_{\mu_0}$ be the inverse of the spectral radius of $\mu_0$ on $\Gamma_0$.
\begin{lem}\label{lemalphato0}
As $\alpha$ tends to 0, $R_\alpha$ and $\zeta_0(R_\alpha)$ both converge to $R_{\mu_0}$.
Moreover, $w_{r,\alpha}$ converges to 0 and the convergence is uniform in $r\leq R_\alpha$.
\end{lem}

\begin{proof}
We first show that $R_\alpha$ is a continuous function of $\alpha$.
Let $P_\alpha$ be the convolution operator
$$P_\alpha:f\mapsto \left (x\mapsto \sum_{y}\mu_\alpha(x^{-1}y)f(y)\right).$$

\begin{claim}
In the $\ell^2$ operator norm, $P_\alpha$ is a continuous function of $\alpha$.
\end{claim}

\begin{proof}[Proof of the claim]
We first show that $P_\alpha$ is continuous in the $\ell_1$ and $\ell_{\infty}$ operator norms.
Fix $\alpha_0$.
Note that $\mu_\alpha$ has a finite support included in a fixed finite set $\Sigma$, so $\mu_\alpha(x)$ uniformly converges to $\mu_{\alpha_0}(x)$ as $\alpha$ tends to $\alpha_0$.
Let $f$ be an $\ell_\infty$ function.
Then, for every $x$,
$$\left |\bigg (P_\alpha-P_{\alpha_0}\bigg)f(x)\right |\leq \sum_{y\in \Gamma}\left |\mu_{\alpha}(x^{-1}y)-\mu_{\alpha_0}(x^{-1}y)\right|\|f\|_{\infty}.$$
The only possible $y$ such that we do not have $\mu_{\alpha_0}(x^{-1}y)=0$ and $\mu_\alpha(x^{-1}y)=0$ are in $x\Sigma$.
This yields
$$\sum_{y\in \Gamma}\left |\mu_{\alpha}(x^{-1}y)-\mu_{\alpha_0}(x^{-1}y)\right|=\sum_{y'\in \Sigma}\left |\mu_{\alpha}(y')-\mu_{\alpha_0}(y')\right |.$$
Therefore, for every $\epsilon>0$, if $\alpha$ is close enough to $\alpha_0$,
$$\left |\bigg (P_\alpha-P_{\alpha_0}\bigg)f(x)\right |\leq\epsilon \|f\|_\infty.$$
This proves continuity for the $\ell_\infty$ operator norm.
Now, let $f$ be an $\ell_1$ function. Then,
$$\left \|\bigg(P_\alpha-P_{\alpha_0}\bigg)f\right \|_1\leq \sum_{x,y\in \Gamma}\big|\mu_\alpha(y^{-1}x)-\mu_{\alpha_0}(y^{-1}x)\big||f(y)|.$$
Fix $y$.
Then, as above,
$$\sum_{x\in \Gamma}\big|\mu_\alpha(y^{-1}x)-\mu_{\alpha_0}(y^{-1}x)\big|=\sum_{x'\in \Sigma}\big|\mu_\alpha(x')-\mu_{\alpha_0}(x')\big|.$$
Let $\epsilon>0$.
Then if $\alpha$ is close enough to $\alpha_0$ this last sum is bounded by $\epsilon$, independently of $y$.
Consequently,
$$\left \|\bigg(P_\alpha-P_{\alpha_0}\bigg)f\right \|_1\leq \epsilon  \sum_{y\in \Gamma}|f(y)|=\epsilon \|f\|_1.$$
This proves continuity for the $\ell_1$ operator norm.
The Riesz-Thorin interpolation theorem \cite[Theorem~(6.27)]{Folland} shows that $P_\alpha$ converges to $P_{\alpha_0}$ for the $\ell_p$ operator norm, for every $1\leq p\leq +\infty$.
This proves the claim, taking $p=2$.
\end{proof}


Since $R_\alpha$ is the inverse of the spectral radius of $P_\alpha$, continuity of $R_\alpha$ follows.
To conclude it thus suffices to prove that $w_{r,\alpha}$ uniformly converges to 0.
By \cite[Proposition~9.18]{Woessbook},
$$w_{r,\alpha}=\sum_{n\geq 1}\mathbf P_\alpha\left (X_n=e,X_k\neq e,1\leq k<n, \text{ first step chosen using }\alpha\mu_1\right )r^n$$
and by the Markov property,
$$w_{r,\alpha}=r\alpha\mu_1(e)+r\sum_{x\neq e\in \Gamma_1}\mathbf P_\alpha(X_1=x)\sum_{n\geq 1}\mathbf P_\alpha\left (X_n=x^{-1},X_k\neq x^{-1},k<n\right )r^n.$$
Set
\begin{equation}\label{defF}
F_r(e,x)=\sum_{n\geq 1}\mathbf P_\alpha\left (X_n=x,X_k\neq x,k<n\right )r^n
\end{equation}
and for $i=0,1$, $x\in \Gamma_i$,
\begin{equation}\label{defFi}
F_r^i(e,x)=\sum_{n\geq 1}\mathbf P_{\mu_i}\left (X_n=x,X_k\neq x,k<n\right )r^n.
\end{equation}
Then, by \cite[Lemma~1.13~(b)]{Woessbook}, for $x\in \Gamma$,
$$G_r(e,x)=F_r(e,x)G_r(e,e)$$
and for $x\in \Gamma_i$,
$$G_r^i(e,x)=F_r^i(e,x)G_r^i(e,e).$$
Thus, by~(\ref{superWoess2}), we have
$$F_r(e,x)=F_{\zeta_i(r)}^i(e,x)$$
and so we recover \cite[Proposition~9.18]{Woessbook}.
In particular,
$$w_{r,\alpha}\leq \alpha R_\alpha \left (\mu_1(e)+\sum_{x\in \Gamma_1}\mu_1(x)F_{R_{\mu_1}}^1(e,x^{-1})\right ).$$
Finally, $F_{R_{\mu_1}}^1(e,x^{-1})$ is finite since we assume that the random walk on $\Gamma_1$ is transient at the spectral radius.
Thus, there exists a constant $C$ such that
$$w_{r,\alpha}\leq \alpha C$$
and so $w_{r,\alpha}$ uniformly converges to 0.
\end{proof}




\begin{lem}\label{lemalphato1}
There exists $\alpha_0$ such that for $\alpha\in [0,\alpha_0]$, the quantity $w'_{r,\alpha}$ stays bounded away from 1.
In particular, as $\alpha$ tends to 0, $\zeta_1(r)$ converges to 0 and
the convergence is uniform in $1\leq r\leq R_\alpha$.
\end{lem}

\begin{proof}
By \cite[Proposition~9.18]{Woessbook},
$$w'_{r,\alpha}\leq U(r)=\sum_{n\geq 1}\mathbf P_\alpha (X_n=e,X_k\neq e,1\leq k<n)r^n.$$
Also, by \cite[Lemma~1.13]{Woessbook},
$$G_r(e,e)=\frac{1}{1-U(r)}.$$
Note that $G_r(e,e)$ depends on $\alpha$, but according to Lemma~\ref{lemalphato0} and~(\ref{superWoess2}), $G_r(e,e)$ is uniformly bounded for $\alpha$ in a fixed neighborhood of 0.
Consequently, $U(r)$ stays bounded away from 1.
\end{proof}

\subsubsection{Example~B}

Here is now an example where $\omega_P(r)>0$ and $\omega_P(r)<\omega_\Gamma(r)$.
Consider the group
$\Gamma=\mathbb F_2*\mathbb{Z}^d$, where $\mathbb F_2$ is the free group with two generators.
In other words, with the notations above, we set $\Gamma_0=\mathbb F_2$ and $\Gamma_1=\mathbb{Z}^d$.
Then $\Gamma$ is hyperbolic relative to $\Gamma_0$ and $\Gamma_1$.

We choose $d\geq 3$ so that every finitely supported admissible random walk on $\Gamma_1$ is transient at the spectral radius.
Choose an adapted measure $\mu_\alpha=\alpha\mu_1+(1-\alpha)\mu_0$ as above.
Then, $\mu_0$ is a probability measure on the non-amenable group $\mathbb F_2$ whose finite support generates $\mathbb F_2$,
hence, $R_{\mu_0}>1$.
According to Lemma~\ref{lemalphato0}, we can thus fix $\alpha$ so that $\zeta_0(r)>1$ for every $r$ in a neighborhood of $R_\alpha$, say $r_0< r\leq R_\alpha$.
Now that $\alpha$ is fixed, we omit it in the notations.
By~(\ref{superWoess2}),
$$\sum_{x\in \mathbb F_2\cap S_n}G_r(e,x)=\frac{1}{1-w_r}\sum_{x\in \mathbb F_2\cap S_n}G_{\zeta_0(r)}^{\mu_0}(e,x).$$
Since $\zeta_0(r)>1$, \cite[Note~1.7]{GouezelLalley} shows that
$\sum_{x\in \mathbb F_2\cap S_n}G_{\zeta_0(r)}^{\mu_0}(e,x)$ diverges as $n$ tends to infinity.
In particular, we see that $\omega_{\mathbb F_2}(r)\geq 0$.
On the other hand, by~(\ref{exponentialdifferenceGreen}), for $s>r$, $\omega_{\mathbb F_2}(s)>\omega_{\mathbb F_2}(r)$, so for large enough $r$, $\omega_{\mathbb F_2}(r)>0$.

We also deduce from~(\ref{superWoess2}) that $\omega_{\mathbb F_2}(r)=\omega_{\mu_0}(\zeta_0(r))$, where $\omega_{\mu_0}$ is the growth rate of the Green function associated with $\mu_0$ on $\mathbb F_2$.
Since $\mathbb F_2$ is hyperbolic, by \cite[Theorem~3.1]{SWX}, the Green function has purely exponential growth, i.e.
$$\sum_{x\in \mathbb F_2\cap S_n}G_r^{\mu_0}(e,x)\asymp \mathrm{e}^{n\omega_{\mu_0}(r)}.$$
Consequently, the Poincaré series
$$\Theta_{r,s}(\mathbb F_2)=\sum_{x\in \mathbb F_2}G_r(e,x)\mathrm{e}^{-sd(e,x)}$$
diverges at $s=\omega_{\mathbb F_2}(r)$.
By Corollary~\ref{parabolicdivergentgrowthgap}, $\omega_{\mathbb F_2}(r)<\omega_\Gamma(r)$.



\subsubsection{Example~C}
Finally, here is a last example.
We assume that there exists a finitely generated group $\Gamma_0$ endowed with an admissible finitely supported probability measure $\mu_0$ such that for some $r_0< R_{\mu_0}$, the Poincaré series
$$\Theta_{r_0,s}(\Gamma_0)=\sum_{x\in \Gamma_0}G_{r_0}^{\mu_0}(e,x)\mathrm{e}^{-sd(e,x)}$$
converges at $s=\omega_{\mu_0}(r_0)>0$.


We consider the free product $\Gamma=\Gamma_0*\Gamma_1$ where $\Gamma_1=\mathbb{Z}^d$, $d\geq 3$.
As above, $\Gamma$ is hyperbolic relative to $\Gamma_0$ and $\Gamma_1$.
We consider the adapted measure $\mu_\alpha=\alpha \mu_1+(1-\alpha)\mu_0$.
We will prove that for some $r$, $\omega_{\Gamma_0}(r)=\omega_\Gamma(r)$.
By Lemma~\ref{lemalphato0}, $\zeta_0(R_\alpha)$ converges to $R_{\mu_0}$ as $\alpha$ converges to 0.
Thus, for small enough $\alpha$, there exists $r_\alpha$ such that $\zeta_0(r_\alpha)=r_0$.
Also, by Lemma~\ref{lemalphato1}, $\zeta_1(r_\alpha)$ converges to 0 and $w'_{r,\alpha}$ stays bounded away from 1 as $\alpha$ tends to 0.
By~(\ref{superWoess2'}), for every $\epsilon>0$, there exists $\alpha$ such that for every $x\neq e\in \Gamma_1$,
\begin{equation}\label{smallGs}
G_r(e,x)\leq \epsilon.
\end{equation}

Every element $x\in \Gamma$ can be written as $x=a_1b_1...a_kb_k$, where $a_i\in \Gamma_0$, $b_i\in \Gamma_1$ and $a_i\neq e$ except maybe $a_1$ and $b_i\neq e$ except maybe $b_k$.
Moreover, since the random walk is adapted to the free product structure, it has to pass through $a_1b_1...a_k$ before reaching $a_1b_1...a_kb_k$.
Consequently,
$$\frac{G_r(e,x)}{G_r(e,e)}=\frac{G_r(e,a_1b_1...a_k)}{G_r(e,e)}\frac{G_r(e,b_k)}{G_r(e,e)},$$
see also \cite[(3.3)]{WoessMartin}.
Note that this is an exact version of the relative Ancona inequalities in the specific case of adapted random walks on free products.
We thus get

\begin{equation}\label{boundPoincare}
\begin{split}
&\Theta_{r,s}(\Gamma)=\sum_{x\in \Gamma}G_r(e,x)\mathrm{e}^{-sd(e,x)}\\
&\leq G_r(e,e)\sum_{k\geq 0}\left (\sum_{a\in \Gamma_0\setminus\{e\}}\frac{G_r(e,a)}{G_r(e,e)}\mathrm{e}^{-sd(e,a)}\right )^k\left (\sum_{b\in \Gamma_1\setminus\{e\}}\frac{G_r(e,b)}{G_r(e,e)}\mathrm{e}^{-sd(e,b)}\right )^k.
\end{split}
\end{equation}

By~(\ref{superWoess2}), $\omega_{\Gamma_0}(r_\alpha)=\omega_{\mu_0}(r_0)$, so $\omega_\Gamma(r_\alpha)\geq \omega_{\mu_0}(r_0)>0$.
Therefore,
\begin{align*}
\sum_{x\in \Gamma_0}\frac{G_{r_\alpha}(e,x)}{G_{r_\alpha}(e,e)}\mathrm{e}^{-\omega_\Gamma(r_\alpha)d(e,x)}\leq &\sum_{x\in \Gamma_0}\frac{G_{r_0}^{\mu_0}(e,x)}{G_{r_0}^{\mu_0}(e,e)}\mathrm{e}^{-\omega_{\mu_0}(r_0)d(e,x)}\\
&=\frac{1}{G_{r_0}^{\mu_0}(e,e)}\Theta_{r_0,\omega_{\mu_0}(r_0)}(\Gamma_0)<+\infty.
\end{align*}
If $x\in \Gamma_1$, then as explained above, we have
$$\frac{G_r(e,x)}{G_r(e,e)}=F_r(e,x)=F_{\zeta_1(r)}^1(e,x)$$ where $F$ and $F^1$ are defined by~(\ref{defF}) and~(\ref{defFi}).
Since the random walk on $\Gamma_1=\mathbb Z^d$ is transient at the spectral radius, there exists $C_1$ such that
$$F_{\zeta_1(r)}^1(e,x)\leq F_{R_{\mu_1}}^1(e,x)\leq C_1.$$
Also, since $\Gamma_1=\mathbb{Z}^d$ has polynomial growth, there exists a constant $C_2$ such that
$$\sum_{x\in \Gamma_1}\frac{G_r(e,x)}{G_r(e,e)}\mathrm{e}^{-\omega_{\Gamma}(r_\alpha)d(e,x)}\leq C_1 \sum_{x\in \Gamma_1}\mathrm{e}^{-\omega_{\mu_0}(r_0)d(e,x)}\leq C_2.$$
We now choose $\epsilon>0$ such that
$$\epsilon C_2 \frac{\Theta_{r_0,\omega_{\mu_0}(r_0)}(\Gamma_0)}{G_{r_0}^{\mu_0}(e,e)}=C_3<1$$
and we fix $\alpha$ such that
by~(\ref{smallGs}),
$$\sum_{b\in \Gamma_1\setminus\{e\}}\frac{G_r(e,b)}{G_r(e,e)}\mathrm{e}^{-sd(e,b)}\leq \epsilon \sum_{x\in \Gamma_1}\mathrm{e}^{-\omega_{\Gamma}(r_\alpha)d(e,x)}\leq \epsilon C_2.$$
Consequently, by~(\ref{boundPoincare}),
$$\Theta_{r_\alpha,\omega_\Gamma(r_\alpha)}(\Gamma)\leq \sum_{k\geq 0}C_3^k.$$
This proves that the Poincaré series $\Theta_{r_\alpha,\omega_\Gamma(r_\alpha)}(\Gamma)$ is convergent.
According to Lemma~\ref{LowerBDLem}, we deduce that $\omega_{P}(r_\alpha)=\omega_\Gamma(r_\alpha)$ for some parabolic group $P$.
Since $\Gamma_1=\mathbb Z^d$ has polynomial growth, by Proposition~\ref{subexpgrowthparabolicgap}, we necessarily have $P=\Gamma_0$.

\begin{rem}
We assumed that the Poincaré series of $\Gamma_0$ was convergent at its critical exponent for some $r_0<R_{\mu_0}$.
However, this was only for convenience.
If this Poincaré series is convergent for $r=R_{\mu_0}$, then we can make the same construction and choose instead $r_\alpha=R_\alpha$.
We only need to know that $\zeta_0(R_\alpha)=R_{\mu_0}$ for small enough $\alpha$.
The fact that $\zeta_0(R_\alpha)=R_{\mu_0}$ is equivalent to the fact that $\mu_\alpha$ is spectrally degenerate along $\Gamma_0$, so we just need to ensure that for small enough $\alpha$, $\mu_\alpha$ is spectrally degenerate along $\Gamma_0$.
Now, if the Poincaré series is convergent for $r=R_{\mu_0}$, then $\sum_{x\in S_n\cap \Gamma_0}G^{\mu_0}_{R_{\mu_0}}(e,x)\leq C\mathrm{e}^{-\omega_{\mu_0}(R_{\mu_0})}$ and so
$$\sum_{x\in \Gamma_0}G^{\mu_0}_{R_{\mu_0}}(e,x)G^{\mu_0}_{R_{\mu_0}}(x,e)\leq C\Theta_{r_0,\omega_{\mu_0}}(r_0)<\infty.$$
By \cite[Proposition~1.9]{GouezelLalley}, this implies that the first derivative of the Green function $t\mapsto G^{\mu_0}_t(e,e)$ is finite at $R_{\mu_0}$.
Using the work of \cite[Section~7]{CandelleroGilch}, we can then construct such a spectrally degenerate probability $\mu_\alpha$ along $\Gamma_0$ for small enough $\alpha$.
\end{rem}

\begin{rem}\label{CGMGap}
In \cite[Lemma~4.7]{CandelleroGilchMuller}, the authors prove that $\omega_P(r)<\omega_\Gamma(r)$ always holds.
However, their proof is incorrect.
Indeed, they use the following inequality
$$\sum_{x\in S_{n+m}\cap P}G_r(e,x)\leq C\sum_{x\in S_n\cap P}G_r(e,x)\sum_{x\in S_m\cap P}G_r(e,x).$$
Proving such an inequality would require that for any $x\in S_{n+m}$ and any $y\in S_n$ on a geodesic from $e$ to $x$, we have
$$G_r(e,x)\leq C G_r(e,y)G_r(y,x).$$
This in turn would require Ancona inequalities for the group $P$.
It is an open question whether the fact that Ancona inequalities hold for any geodesic implies that the group is hyperbolic, but it is easy to prove that they do not hold in a parabolic subgroup $P$ if $P$ is virtually abelian, see for instance \cite[Section~5.2]{DG20}.
\end{rem}

%This example, combined with Corollary~\ref{parabolicdivergentgrowthgap}, suggests that the following conjecture is true.
%\begin{conj}
%We have the following dichotomy. \begin{comment}
%I have proofread this part that is very nice math. This conjecture is also interesting. But I think the two statements are mathematically exclusive: once one statement is true, then the other would be vacuous, right? (personally I tend to believe examples in case 2 exists). If yes, stating it in this case would be confusing.
%\end{comment}

%\begin{itemize}
%    \item If for every finitely generated group $\Gamma_0$, for every admissible finitely supported probability measure $\mu_0$ on $\Gamma_0$ and for every $r\leq R_{\mu_0}$, the Poincaré series $\Theta_{r,s}(\Gamma_0)$ diverges at $s=\omega_{\Gamma_0}(r)$, then for every relatively hyperbolic group $\Gamma$, every admissible finitely supported probability measure $\mu$ and every parabolic subgroup $P$, $\omega_P(r)<\omega_\Gamma(r)$.
%    \item On the contrary, if there exists $\Gamma_0$, $\mu_0$ and $r_0$ as above such that this Poincaré series converges, then there is a relatively hyperbolic group $\Gamma$ and a parabolic subgroup $P$ for which $\omega_P(r)=\omega_\Gamma(r)$ at some $r$.
%\end{itemize}
%\end{conj}
%By the above construction, the second statement in the conjecture is true.
%However, we cannot deduce the first one from Corollary~\ref{parabolicdivergentgrowthgap} since the first return kernel to a parabolic subgroup is typically not finitely supported. \begin{comment}
%I also do not understand what you mean here. Perhaps this deserves to be further clarified.    
%\end{comment}

%{\color{blue}
%Okay, thank you for your remark. I suggest to remove the conjecture above and to write things as follows.

Let us give some final remarks to conclude this discussion.
The last example raises the following question.
\begin{quest}\label{question}
Does there exist a finitely generated group $\Gamma_0$ endowed with an admissible finitely supported probability measure $\mu_0$ such that for some $r\leq R_{\mu_0}$, the Poincaré series $\Theta_{r,s}(\Gamma_0)$ is convergent at the radius of convergence $s=\omega_{\Gamma_0}(r)>0$ ?
\end{quest}

If there is a positive answer to this question, then as we saw, there exists a relatively hyperbolic group $\Gamma$ endowed with an admissible finitely supported probability measure $\mu$ and a parabolic subgroup $P$ for which $\omega_P(r)=\omega_\Gamma(r)$ at some $r$.
Moreover, if the measure $\mu_0$ is symmetric, then we can choose the measure $\mu$ to be symmetric as well.

On the contrary, if this question has a negative answer, then Corollary~\ref{parabolicdivergentgrowthgap} suggests that we always have $\omega_P(r)<\omega_\Gamma(r)$.
However, this corollary requires the Poincar\'e series
$$\Theta_{r,s}(P)=\sum_{p\in P} G_r(e,p) e^{-s d(e,p)} $$
to be divergent, where $G_s(e,p)$ is the Green function associated with the measure $\mu$ on $\Gamma$.
Using \cite[Lemma~4.4]{DG21} as above, we can rewrite this Poincar\'e series as
$$\Theta_{r,s}(P)=\sum_{p\in P} G_1^{r,P}(e,p) e^{-s d(e,p)},$$
where $G_1^{r,P}$ is the Green function at 1 associated with the first return kernel $p^{r,P}$ defined in~(\ref{deffirstreturnkernel}).
Unfortunately, this first return kernel is not in general finitely supported, so even if Question~\ref{question} has a negative answer, we cannot deduce that this Poincar\'e series is divergent, hence we cannot deduce that $\omega_P(r)<\omega_\Gamma(r)$.
%}
%{\color{blue}
%We nevertheless make the following conjecture.
%\begin{comment}
%{\color{red}Personnaly I do not think we have enough evidence to support this conjecture. ALso if the first return kernel was known in some examples to be finitely generated, then your conjecture would be actually a proposition. In my opinion, it is not essential as the previous question. Or whether there exists examples of first return kerl is finitely generated. Anyways, if you insist on it, I'm also Okay.}
%{\color{blue}No I don't hold on to this conjecture and you're right we do not have clear evidence.  Note however that in the very particular case of adapted random walks on free products (which is the setting of \cite{CandelleroGilchMuller}), the first return kernel is indeed finitely supported.}
%\ywy{I sugegst at least to remove the first sentence "Assume that Question~\ref{question} has a negative answer for symmetric measures,". It is bizzare to me that a conjecture assumes a negative answer to another question. So what should be expected: the question to be negative or positive?? If the authors are uncertain, then a further conjecture would be the square of uncertainty. }
%{\color{blue}Okay I see your point. The conjecture in my mind was more : if the question has negative answer, then it also has a negative eanswer for transition kernels which are not necessarily finitely supported but have exponential moments and then say that the first return kernel has exponential moments to conclude. But I understand this is not satisfying and I deleted the conjecture.} I agree what you conjectured is interesting. Maybe we should work on it ourselies, instead of releasing it to public, as it is bit technical for the readers. 
%\end{comment}
%\begin{conj}
%Assume that Question~\ref{question} has a negative answer for symmetric measures, i.e.\ for every finitely generated group $\Gamma_0$, for every admissible symmetric finitely supported probability measure $\mu_0$ on $\Gamma_0$ and for every $r\leq R_{\mu_0}$, the Poincaré series $\Theta_{r,s}(\Gamma_0)$ diverges at $s=\omega_{\Gamma_0}(r)$.
%Then, for every relatively hyperbolic group $\Gamma$, every admissible symmetric finitely supported probability measure $\mu$ on $\Gamma$ and every parabolic subgroup $P$, $\omega_P(r)<\omega_\Gamma(r)$ for every $1< r\leq R_\mu$.
%\end{conj}
%}
 
%\section{Purely exponential growth of Green functions}
%\begin{prop}
%Suppose that $\Gamma$ is hyperbolic relative to peripheral subgroups of subexponental growth. Then,
%$$H_r(n+m)\asymp H_r(n)H_r(m)$$
%and
%there exists a constant $C>1$ such that for any $1\le r\le R_\mu$, we have 
%$$
%C^{-1} e^{n \omega_\Gamma(r)}\le H_r(n)\le C e^{n \omega_\Gamma(r)}
%$$
%for any $n\ge 1$.
%\end{prop}
%\begin{proof}
%The upper bound follows from Lemma \ref{UpperBDLem}. We shall prove that $\Gamma$ has parabolic gap for Green function, so the lower bound follows by Lemma \ref{UpperBDLem}. 

%Since a group with sub-exponential growth is amenable, every  maximal parabolic subgroup $H$ is amenable. Let $p_{H, r}(1, x)$ be the first return $r$-kernel of $\mu$-random walk to $x\in P$. We have  two cases to consider according to whether the spectral radius of $p_{H, r}$ equal to 1: \textit{spectrally degenerate} or \textit{spectrally non-degenerate}. 

%\textbf{Case 1.} The parabolic subgroup $H$ is spectrally non-degenerate: $\rho_H(r)<1$. By your proposition 2.1,  $\sum_{x\in H} G_{r}(1,x)<\infty$, so $\omega_H(r){\color{red}\leq}0$. 

%\textbf{Case 2.} The parabolic subgroup $H$ is spectrally degenerate: $\rho_H(r)<1$. Then $\rho_H$ has finite exponential moment by \cite[Lemma 4.5]{DG21}. The Kesten's spectral characterization of  amenability implies that $p_{H,r}$ is a probability measure. It is well-known that under the finite second moment, $H$ is not virtually $\mathbb Z^2$ if and only if some (or any) symmetric random walk is transient. Since the induced Green function $G_r(1,x)=G_{H,r}(1,x)$ on $H$ under $p_{H,r}$  is finite, the parabolic subgroup cannot be virtually $\mathbb Z^2$. In other words, $
%\mathbb Z^2$ must be spectrally non-degenerate, treated in Case 1.

%Thus, the induced $p_{H,r}$-random walk on $P$ is transient, so we have $G_{H,r}(1,x)\le G_{H,r}(1,1)<\infty$ for any $x\in H$. Since $\sharp (S_n\cap H)$ is of sub-exponential function, we conclude again that  $\omega_H(r){\color{red}\leq}0$. 

%Since $\omega_\Gamma(r)>0$ for $1<r\le R_\mu$, the gap property for Green function follows.
%{\color{red} One can greatly simplify the proof here, as follows.\begin{comment}I agree with this simplified proof, but I think this inequality is non-necessary: $G_{H,r}(1,x)\le G_{H,R}(1,1)$.\end{comment}
%We have $G_{H,r}(1,x)\le G_{H,R}(1,1)$. By \cite[Lemma~4.']{DG21}, $G_{H,R}(1,1)=G_R(1,1)$ and by a famous result of Guivarc'h \cite[p.~85, remark~b)]{Guivarch}, this last quantity is finite. Since $\sharp (S_n\cap H)$ is of sub-exponential growth, we conclude again that  $\omega_H(r)\leq 0$. 
%}
%\end{proof}

%When the random walk is not spectrally degenerate, we can loosen a bit our assumptions on parabolic subgroups.

%\begin{prop}
%Assume that the parabolic subgroup $P$ is amenable and that $\mu$ is not spectrally degenerate along $P$.
%Then, for any $r\leq R_\mu$, the sum
%$$\sum_{x\in P}G_r(e,x)$$
%is finite.
%\end{prop}




%\begin{proof}
%We only need to prove the result at $r=R_\mu$ for the Green function is non-increasing in $r$.
%The sum in the proposition can be rewritten as
%$$\sum_{x\in P}G_{P,R_\mu}(e,x|1).$$
%For simplicity, we will drop the indices $P,R_\mu$.
%Since the random walk is not spectrally degenerate along $P$ and $P$ is amenable, the transition kernel $p$ is strictly sub-Markov.
%Let $t<1$ be such that $\tilde{p}=\frac{1}{t}p$ is Markov.
%Then, letting $\widetilde{G}$ be the Green function associated with $\tilde{p}$, for any $x\in \mathcal{H}$ we have $G(e,x|1)=\widetilde{G}(e,x|t)$.
%Hence, the sum above can be written as
%$$\sum_{x\in P}\widetilde{G}(e,x|t).$$
%Consider the convolution operator
%$$\widetilde{P}:f\mapsto \left ( x\mapsto \sum_{y}\tilde{p}(x,y)f(y)\right )$$
%acting on the space of $\ell_1$ functions, i.e.\ summable functions on $P$.
%The $\ell_1$-norm of $\widetilde{P}$ is 1, so its $\ell_1$-spectral radius is bounded by 1.
%In particular, $1/t$ is bigger than the $\ell_1$-spectral radius, so by definition, $I-t\widetilde{P}$ is invertible in the space of summable functions.
%Moreover, the inverse is of the form 
%$$\widetilde{Q}=\sum_{n\geq 0} t^n\widetilde{P}^n.$$
%Consider the function $f$ defined by $f(x)=1$ if $x=e$ and 0 otherwise.
%Then, $\widetilde{Q}f(x)=\widetilde{G}(e,x^{-1}|t)$.
%Since the function $f$ is summable, $\widetilde{Q}f$ also is summable, and so the Green function at $t$ is summable.
%This concludes the proof.
%\end{proof}

%\begin{cor}
%Let $P$ be a maximal parabolic subgroup and assume that $P$ is amenable.
%If the random walk is not spectrally degenerate along $P$, then
%$\omega_P(r)<\omega_\Gamma(r)$.
%In particular, if every parabolic subgroup is amenable and if the random walk is non-spectrally degenerate, then
%$$H_r(n+m)\asymp H_r(n)H_r(m)$$
%and
%there exists a constant $C>1$ such that for any $1\le r\le R_\mu$, we have 
%$$
%C^{-1} e^{n \omega_\Gamma(r)}\le H_r(n)\le C e^{n \omega_\Gamma(r)}
%$$
%for any $n\ge 1$.
%\end{cor}

%\begin{proof}
%Since the sum
%$$\sum_{x\in P}G_r(e,x)$$
%is finite, we have $\omega_P(r)\leq 0$.
%Hence, $\omega_P(r)<\omega_\Gamma(r)$.
%\end{proof}



%Finally, here is a third situation where $\omega_P<\omega_\Gamma$ and $H_r(n)$ has pure exponential growth.
%\begin{defn}
%Say that the random walk is isotropic if the driving measure $\mu$ is the uniform measure on the set of generators $S$.
%\end{defn}


%\begin{prop}
%Assume that the random walk is isotropic.
%Then,
%$$H_r(n+m)\asymp H_r(n)H_r(m)$$
%and
%there exists a constant $C>1$ such that for any $1\le r\le R_\mu$, we have 
%$$
%C^{-1} e^{n \omega_\Gamma(r)}\le H_r(n)\le C e^{n \omega_\Gamma(r)}
%$$
%for any $n\ge 1$.
%\end{prop}

%\begin{proof}
%Since the random walk is isotropic, the Green function $G_r(e,x)$ is constant along spheres.
%Fixing $r$ and $n$, set $w_n(r)=G_r(e,x)$, where $x\in S_n$.
%Then, for any $s$,
%$$\sum_{x\in \Gamma}G_r(e,x)\mathrm{e}^{-sd(e,x)}=\sum_{n\geq 0}w_n(r)\sharp (S_n)\mathrm{e}^{-sn}$$
%and for every parabolic subgroup $P$,
%$$\sum_{x\in P}G_r(e,x)\mathrm{e}^{-sd(e,x)}=\sum_{n\geq 0}w_n(r)\sharp (S_n\cap P)\mathrm{e}^{-sn}.$$
%By pure exponential growth (cf. \cite[Theorem 1.9]{YangPS}),
%$$\sharp (S_n)\asymp \mathrm{e}^{vn}.$$
%\begin{lem}
%There exists $C\geq 0$ and $v_P<v$ such that
%$$\sharp (S_n\cap P)\leq C\mathrm{e}^{v_Pn}.$$
%\end{lem}

%\begin{proof}
%{\color{red} Note that a parabolic subgroup $P$ is relatively quasiconvex of infinite index. By \cite[Theorem 2]{DFW}, the conclusion follows from the growth tightness of those subgroups. (The same result holds for more general class of subgroups called weakly convex subgroups by \cite[Theorem 4.8]{YangSCC}).}
%\end{proof}


%Hence, $\omega_\Gamma(r)=v+\limsup \frac{w_n(r)}{n}$
%and $\omega_P(r)\leq v_P+\limsup \frac{w_n(r)}{n}$.
%Consequently, $\omega_P(r)<\omega_\Gamma(r)$.
%This being true for every parabolic subgroup, we can conclude as above.
%\end{proof}





\section{The growth rate of the trace of the branching random walk}\label{SectiongrowthrateBRW}
\label{s:vgrbrw}


Recall that $\mathcal{P}_n$ is the set of points in $S_n$ that are eventually visited by some particle of the branching random walk and that $M_n=\sharp \mathcal{P}_n$.
In this section, we compare the growth rate of the Green function $\omega_\Gamma(r)$ with the growth rate of the branching random walk $\log \limsup M_n^{1/n}$.
Our goal is to prove the following proposition.
\begin{prop}\label{P:vgrBRW}
For $r\in [1, \rho^{-1}]$, $\omega_{\Gamma}(r) = \limsup_{n \to \infty} \frac{1}{n} \log M_n$ almost surely.
\end{prop}

Theorem~\ref{Thmvolumegrowthconst} is then a consequence of Proposition~\ref{P:vgrBRW}, Corollary~\ref{corolowersemicontinuous} and Corollary~\ref{coroboundv/2}.

\subsection{Upper bound}
We first prove the following.



\begin{prop}\label{upperboundvgrBRW}
Almost surely, we have
$$\limsup_{n \to \infty} \frac{1}{n} \log M_n \leq \omega_\Gamma(r).$$
\end{prop}


The proof of \cite{SWX}, which relies on the Borel-Cantelli lemma uses the purely exponential growth of  the Green functions over spheres.
However, only a small adaptation is needed to apply it here.
We rewrite it for convenience.


\begin{proof}
For $x\in \Gamma$, we denote by $Z_x$ the number of particles of the branching random walk that ever visit $x$.
The many-to-one formula states that
$$\mathbf E[Z_x]=G_r(e,x).$$
Then,
$M_n=\sum_{x\in S_n}\mathbb{1}_{Z_x\geq 1}$, so by the Markov inequality,
$$\mathbf E[M_n]=\sum_{x\in S_n}\mathbf P (Z_x\geq 1)\leq \sum_{x\in S_n}\mathbf E[Z_x]=H_r(n).$$

Let $\epsilon>0$.
By the Markov inequality,
$$\mathbf{P}\left (M_n^{1/n}\geq \mathrm{e}^{\omega_\Gamma(r)}+\epsilon\right)\leq \frac{\mathbf{E}[M_n]}{(\mathrm{e}^{\omega_\Gamma(r)}+\epsilon)^n}\leq \frac{H_r(n)}{(\mathrm{e}^{\omega_\Gamma(r)}+\epsilon)^n}.$$
By definition,
$\omega_\Gamma(r)=\log \limsup H_r(n)^{1/n}$,
so there are at most finitely many $n$ such that
$$H_r(n)^{1/n}\geq \mathrm{e}^{\omega_\Gamma(r)}+\epsilon/2.$$
Therefore,
$$\sum_n \frac{H_r(n)}{(\mathrm{e}^{\omega_\Gamma(r)}+\epsilon)^n}<\infty.$$
The statement of the lemma is thus a consequence of the Borel-Cantelli lemma.
\end{proof}

\subsection{Lower bound}
Before proving the lower bound, we first recall some geometric lemmas about relatively hyperbolic groups.
Let $\mathbb P=\{gP: g\in \Gamma, P\in\mathbb P_0\}$ be the collection of all parabolic cosets.
 Recall that $\eta$ and $L(\eta)$ are fixed so that for $L\geq L(\eta)$, any $(\eta,L)$-transition point satisfies the results of Section~\ref{sectiontransitionpoints}.



\begin{defn}\label{deftransitional}
For $L\geq L(\eta)$, an $(\eta,L)$-transition point on a geodesic is called an $L$-transition point.
A geodesic $\alpha$ is called \textit{$L$-transitional} if every point on $\alpha$ is an $L$-transition point.
\end{defn}



\begin{lem}\label{RelThinTriangleLem2}
Let $L\geq L(\eta)$.
There exists $K$ such that the following holds.
For every $x,y,z\in \Gamma$,
if both $[x,z]$ and $[y,z]$ are $L$-transitional, then there exists a point $w$ within $K$ of an $L$-transition point on $[x,z]$, an $L$-transition point on $[y,z]$ and an $L$-transition point on $[x,y]$.
%Moreover, any two points $w,w'$ satisfying this property are within $K$ of each other.
\end{lem}

\begin{proof}
Let $x,y,z$ satisfy the statement of the lemma.
Applying Lemma~\ref{ThinTransitional2}, consider the last point $w_0$ on $[z,x]$ which is within $C$ of $[z,y]$ and let $w$ be the next point on $[z,x]$.
Since $w$ also is $L$-transitional, by definition of $w_0$, $w$ is within $C$ of an $L$-transition point on $[x,y]$.
Moreover, $w$ is within $C+1$ of an $L$-transitional point on $[y,z]$.
\end{proof}

\begin{lem}\label{existenceofcenter}
There exists a constant $C$ such that the following holds.
Let $x_0,x_1,x_2$ be three points in $\Gamma$.
Then, there exist $w_0,w_1,w_2$ such that for $i \mod 3$,
$$d(x_i,x_{i+1}) \geq d(x_i,w_i)+d(w_i,w_{i+1})+d(w_{i+1},x_{i+1})-C.$$
Moreover, $w_i$ and $w_{i+1}$ are within $C$ of an $L$-transition point on $[x_i,x_{i+1}]$.
Finally, if $[x_i,x_{i+1}]$ is $L$-transitional, then $d(w_i,w_{i+1})\leq C'$, where $C'$ only depends on $L$. 
\end{lem}

It will be convenient to rely on similar results proved in \cite{DussauleLLT1}.
However, the terminology is a bit different and \cite{DussauleLLT1} uses the notion of relative geodesics.
Let us briefly introduce this notion.
Let $\mathbb P_0$ be the chosen set of representatives of conjugacy classes of parabolic subgroups and let $S$ be a finite generating set of $\Gamma$.
Following Osin \cite{Osin}, the relative graph is the Cayley graph of $\Gamma$ endowed with the (possibly infinite) generating set $S\bigcup \bigcup_{P\in \mathbb P_0}P$.
It is quasi-isometric to the coned-off graph introduced by Farb \cite{Farb} who gave one of the first definitions of relatively hyperbolic groups.
The relative graph is hyperbolic.
A relative geodesic is a geodesic in the relative graph.
By \cite[Proposition~8.13]{Hruska}, if $x,y\in \Gamma$, then any point on a relative geodesic from $x$ to $y$ is within a uniformly bounded distance of a transition point on a geodesic from $x$ to $y$ (in the Cayley graph of $\Gamma$).

\begin{proof}[Proof of Lemma~\ref{existenceofcenter}]
Consider the projection of $x_0$ on a relative geodesic from $x_1$ to $x_2$ in the relative Cayley graph.
If several projections exist, choose the closest possible to $x_1$.
Denote this projection by $w_1$ and
let $w_2$ be the point on this relative geodesic just after $w_1$.

By \cite[Lemma~4.16]{DussauleLLT1}, any relative geodesic from $x_0$ to $x_1$ passes at a point $v$ within a bounded distance of $w_1$.
We prove by contradiction that $v$ is within a bounded distance of the projection of $x_2$ on a relative geodesic from $x_0$ to $x_1$ the closest to $x_1$.
Denote by $v'$ such a projection.
Then applying again \cite[Lemma~4.16]{DussauleLLT1}, the relative geodesic from $x_2$ to $x_1$ we chose would pass at a point $w_1'$ within a bounded distance of $v'$.
If $d(v,v')$ is large, then $d(w_1,w_1')$ is also large.
Now, if $w_1$ is before $w_1'$ on the relative geodesic, this contradicts the definition of $v'$ and if $w_1'$ is before $w_1$, this contradicts the definition of $w_1$.

Finally, denote by $w_0$ the point just before $v$ on the relative geodesic from $x_0$ to $x_1$.
Then, applying \cite[Lemma~4.16]{DussauleLLT1} one last time, a relative geodesic from $x_i$ to $x_{i+1}$ passes within a bounded distance of $w_i$ and $w_{i+1}$.
Since points on a relative geodesics are within a bounded distance of transition points (see \cite[Proposition~8.13]{Hruska}), this proves the two first properties of the points $w_i$.

Notice that the points $w_i,w_{i+1}$ are chosen within a bounded distance of successive points on a relative geodesic from $x_i$ to $x_{i+1}$.
If the geodesic $[x_i,x_{i+1}]$ is $L$-transitional, then the corresponding relative geodesic has bounded jumps in parabolic subgroups, hence the distance between $w_i$ and $w_{i+1}$ is bounded.
This concludes the last part of the lemma.
\end{proof}


We define
\[
  S_{n, L} = \left\{ x \in S_n \colon [e, x] \text{ is } L \text{-transitional} \right\}.  
\]
Set $M_{n, L} = \sharp \mathcal{P}_n \cap S_{n, L}$. 
We first consider the lower bound for the quantity $\lim_{n \to \infty} \frac{1}{n} \log M_{n, L}$. 
To this end, mimicking the strategy of \cite{SWX}, we need first and second moments estimates for $M_{n, L}$.  Set
\[
  H_{r, L}(n) = \sum_{x \in S_{n, L}} G_r(e, x), 
\]
and 
\[
  \omega_{\Gamma, L}(r) = \limsup_{n \to \infty} \frac{1}{n} \log H_{r, L}(n). 
\]



\begin{prop}\label{puregrowthtransitional}
  \label{GreenL}
  For every $L \geq L(\eta)$ and $r \in [1, \rho^{-1}]$, there is a constant $c_L > 0$ such that
  \[
    c_L^{-1} \mathrm{e}^{\omega_{\Gamma, L}(r)} \leq H_{r, L}(n) \leq c_L \mathrm{e}^{\omega_{\Gamma, L}(r)}, \quad n \geq 0. 
  \]
  
\end{prop}

\begin{proof}
  For $x \in S_{n + m, L}$, the point $y = [e, x] \cap S_n$ is in $S_{n, L}$ and $y^{-1}x \in S_{m, L}$.  By the relatively Ancona inequalities, there is a constant $c_1 > 0$ such that
\[
  G_r(e, x) \leq c_1 G_r(e, y) G_r(e, y^{-1} x).
\]
Thus we have that
\begin{equation}
  \label{e:GreenLU}
  H_{r, L}(n + m) \leq c_1 \sum_{y \in S_{n, L}} G_r(e, y) \sum_{z \in S_{m, L}} G_r(e, z) = c_1 H_{r, L}(n) H_{r, L}(m). 
\end{equation}


Let $F$ be the finite set given by Lemma~\ref{ExtensionLem}.  Set
\[
  l = \max \left\{ d(e, f) \colon f \in F \right\} + 4 \epsilon. 
\]
For $x \in S_{n, L}$ and $y \in S_{m, L}$ there is $f \in F$ such that
\[
  d(x, [e, x f y]) \leq \epsilon, \quad d(x f, [e, x f y]) \leq \epsilon. 
\]
Also, there are positive constants $c_2$ and $c_3$ such that
\[
  G_r(e, x) G_r(e, y) \leq c_2 G_r(e, x) G_r(x, x f) G_r(x f, x f y) \leq c_3 G_r(e, x f y).  
\]
Note that $x f y \in \bigcup_{-l \leq i \leq l} S_{n + i, L}$.  Thus
\begin{equation}
  \label{e:GreenLD}
  H_{r, L}(n) H_{r, L}(m) \leq c_4 H_{r, L}(n+m)
\end{equation}
for some $c_4 > 0$.  This proposition follows by the Fekete Subadditive Lemma and~\eqref{e:GreenLU}, \eqref{e:GreenLD}. 
\end{proof}


\begin{lem}\label{omegaL}
  For $r \in [1, \rho^{-1}]$, $\lim_{L \to \infty} \omega_{\Gamma, L}(r) = \omega_{\Gamma}(r)$. 
\end{lem}


\begin{proof}
  For $n_0 \in \mathbb{N}$ and $x_i \in S_{n_0}$, $1 \leq i \leq m$, we can choose $f_i \in F$ so that the conditions in Lemma~\ref{ExtensionLem} hold with $g$, $h$ replaced by $x_1f_1x_2\cdots x_{i-1} f_{i-1} x_i$, $x_{i+1}$.  In particular, 
  \[
    x = x_1 f_1 x_2 \cdots f_{m-1} x_m \in \bigcup_{k = (n_0 - l) m}^{(n_0 + l) m} S_{k, \, L}
  \]
  for $L$ sufficiently large,  where $l = \max \left\{ d(e, f) \colon f \in F \right\} + 4 \epsilon$.  Since $F$ is a fixed finite set, we have that 
  \[
    G_r(e, x) \geq c_1^m G_r(e, x_1) \cdots G_r(e, x_m) 
  \]
  with $c_1 = \min_{f \in F} \frac{G_r(e, \, f)}{G_r(e, \, e)}$.  By Lemma~\ref{BoundedToOne}, there is $c_2 > 0$ such that each $x$ has at most $c_2^m$ possible representations in the form of $x = x_1 f_1 x_2 \cdots f_{m-1} x_m$.  Therefore 
  \[
    \sum_{k=(n_0 - l) m}^{(n_0 + l) m} H_{r, L}(r) \geq c_3^m \sum_{x_1, \ldots, x_m \in S_{n_0}} G_r(e, x_1) \cdots G_r(e, x_m) = c_3^m \left[ H_r(n_0) \right]^m  
  \]
  with $c_3 = c_1 c_2^{-1}$.    This and Proposition~\ref{GreenL} imply that
  \[
    2 c_L l m \mathrm{e}^{(n_0 + l) m \omega_{\Gamma, L}(r)} \geq c_3^m \left[ H_r(n_0) \right]^m.   
  \]
  Letting first $m \to \infty$ and then $L \to \infty$, we have for every $n_0 \in \mathbb{N}$, 
  \[
    \liminf_{L \to \infty} \omega_{\Gamma, L}(r) \geq \frac{1}{n_0 + l} \left( \log c_3 + \log H_r(n_0) \right), 
  \]
  which completes the proof of this lemma. 
\end{proof}


%}
Now we are ready to estimate the second moment of $M_{n, L}$, which will help us find a lower bound. 
Let $x$ and $y$ be in $S_{n, L}$.  By Lemma~\ref{RelThinTriangleLem2}, there exists $w = w(x, y) \in \Gamma$ such that $w$ is within a bounded distance of transition points of $[e, x]$ and $[e, y]$ respectively.
\begin{lem}\label{Correlationxy}
  For $1 \leq r < \rho^{-1}$ and $x$, $y \in S_{n, L}$, there exists a positive constant $c > 0$ such that
  \[
    \sum_{z \in \Gamma} G_r(e, z) G_r(z, x) G_r(z, y) \leq c G_r(e, w) G_r(w, x) G_r(w, y).  
  \]
  
\end{lem}

\begin{proof}
  Let $\kappa =C+K$ where $C$ and $K$ are given by Lemma~\ref{ThinTransitional2} and Lemma~\ref{RelThinTriangleLem2}.  Define
  \[
    \Omega_1 = \left\{ z \in \Gamma \colon d(w, u) \leq \kappa \text{ for some transition point } u \in [e, z] \right\}, 
  \]
  and
  \[
    \Omega_2 = \left\{ z \in \Gamma \colon d(w, u_i) \leq \kappa \text{ for transition points } u_1 \text{ on } [z, x] \text{ and } u_2 \text{ on } [z, y] \right\}.  
  \]
By Lemma~\ref{ThinTransitional2} and Lemma~\ref{RelThinTriangleLem2}, $\Gamma = \Omega_1 \cup \Omega_2$.  

 
Assume $z \in \Omega_1$. Applying Lemma~\ref{existenceofcenter} to $(x,y,z)$, we get the existence of transition points $v,v'$ on $[x,y]$ such that $v$ is within a bounded distance of a transition point on $[z,x]$ and $v'$ is within a bounded distance of a transition point on $[z,y]$.
Moreover, since $[e,x]$ and $[e,y]$ are $L$-transitional, $d(v,v')$ is bounded.
Combining all this, we see that $v$ is within a bounded distance of transition points on $[x,y]$, $[x,z]$ and $[y,z]$.

Now, applying Lemma~\ref{RelThinTriangleLem2} to $(w,x,y)$, we have that $v$ is within a bounded distance of either a transition point on $[w,x]$ or on $[w,y]$. We assume without loss of generality that the latter holds.

  
  \begin{center}
\begin{tikzpicture}[scale=2.2]
\draw (-1.8,.5)--(-2,-.4);
\draw (2.2,.5)--(2,-.4);
\draw (-2,-.4)--(2,-.4);
\draw (-.2,-1.3)--(-.2,-.4);
\draw[dotted] (-.2,-.4)--(-.2,0);
\draw (-.2,0)--(.4,-.2);

\draw (-.2,0)--(.2,.4);

\draw[dotted] (.4,-.2)--(.4,-.4);
\draw (.4,-.4)--(.4,-1);
\draw (.4,-.2)--(.2,.4);
\draw (.2,.4)--(.2,1);
\draw[dotted] (.2,1)--(.2,1.4);
\draw (-2.3,1.9)--(-2.5,1);
\draw (.7,1.9)--(.5,1);
\draw (-2.5,1)--(.5,1);
\draw (.2,1.4)--(-1.6,1.7);
\draw (-1.6,1.7)--(-.5,1.3);
\draw (-.5,1.3)--(.2,1.4);
\draw (-1.6,1.7)--(-1.6,2.5);
\draw (-.5,1.3)--(-.5,2.3);
%\draw (-1.4,0)node{$w_1$};
\draw (-.5,-1.3) node{$e$};
%\draw (1.1,-.15) node{$w_3$};
%\draw (-.1,.5) node{$w_2$};
\draw (.1,.05) node{$w$};
\draw (.6,-1.1) node{$x$};
%\draw (.4,1.4) node{$v_1$};
%\draw (-.7,1.8) node{$v_2$};
\draw (-1.7,2.5) node{$z$};
%\draw (-2.2,1.3) node{$v_3$};
\draw (.3,1.4) node{$v$};
\draw (-.45,1.2) node{$v'$};
\draw (-.4,2.3) node{$y$};
\end{tikzpicture}
\end{center}





\begin{claim}
For every given $K$, there exists $K'$ such that if $w$ is within $K$ of a transition point on $[v,y]$, then $d(v,w)\leq K'$.
\end{claim}

\begin{proof}[Proof of the claim]
Since $v$ is within a bounded distance of a transition point on $[w,y]$ by assumption, we have
$$d(w,y)\geq d(w,v)+d(v,y)-C.$$
Thus, if $d(w,[v,y])\leq K$, then
$$d(v,y)\geq d(d(w,v)+d(w,y)-2K$$
and so
$$2d(v,w)\leq C+2K,$$
which proves the claim.
\end{proof}

\begin{claim}
For every given $K$, there exists $K'$ such that if $v$ is within $K$ of a transition point on $[w,x]$, then $d(v,w)\leq K'$.
\end{claim}

\begin{proof}[Proof of the claim]
By the previous claim, we can assume that $w$ is far from a transition point on $[v,y]$.
But then, since $w$ is within a bounded distance of a transition point on $[x,y]$, applying Lemma~\ref{RelThinTriangleLem2} to $(v,x,y)$ shows that $w$ is within a bounded distance of a transition point on $[v,x]$.
Therefore,
$$d(v,x)\geq d(v,w)+d(w,x)-C.$$
Thus, if $d(v,[w,x])\leq K$, then
$$d(w,x)\geq d(v,w)+d(v,x)-2K$$
and so
$$2d(v,w)\leq c+2K,$$
which proves the claim.
\end{proof}

These two claims show that either $d(v,w)$ is bounded or
$w$ is far from a transition point on $[v,y]$ and $v$ is far from a transition point on $[w,x]$.
Applying Lemma~\ref{RelThinTriangleLem2} to $(v,x,y)$ and then to $(z,w,x)$, in every case we get that
$w$ is within a bounded distance of a transition point on $[v,x]$ and $v$ is within a bounded distance of a transition point on $[w,z]$.

By the relatively Ancona inequalities, we thus have
  \begin{align*}
    G_r(e, z) &\leq c_1 G_r(e, w) G_r(w, v) G_r(v, z), \\
    G_r(z, x) &\leq c_1 G_r(z, v) G_r(v, w) G_r(w, x), \\
    G_r(z, y) &\leq c_1 G_r(z, v) G_r(v, y). 
  \end{align*}
Consequently,
  \begin{align*}
&\sum_{z \in \Omega_1} G_r(e, z) G_r(z, x) G_r(z, y) \\
    &\hspace{.2cm}\leq
    c_1^3 G_r(e, w) G_r(w, x) \sum_{z \in \Omega_1} G_r(v, y) G_r(w, v)G_r(v,w) G_r(v, z)G_r(z,v)^2. 
  \end{align*}
    
Note that $G_r(w, v) G_r(v, y) \leq c_2 G_r(w, y)$.  Therefore,
  \begin{align*}
    \sum_{z \in \Omega_1} G_r(e, z) G_r(z, x) G_r(z, y)
    \leq&
    c_3 G_r(e, w) G_r(w, x) G_r(w, y) \\
    &\hspace{.6cm}\sum_{v \in [w, y]} G_r(w, v) \sum_{z \in \Gamma} G_r(v,z)G_r(z, v)^2\\
    &
    \leq c_4 G_r(e, w) G_r(w, x) G_r(w, y). 
  \end{align*}
Here we used the facts that for $r < \rho^{-1}$, $G_r(w, v)$ is decaying exponentially in $d(w, v)$, which is a direct consequence of~(\ref{exponentialdifferenceGreen}) and that
$$\sum_{z \in \Gamma} G_r(v, z)G_r(z,v)^2\leq C\sum_{z\in \Gamma}G_r(v,z)G_r(z,v) < \infty$$
by \cite[Proposition~1.9]{GouezelLalley}.

\medskip
Now we consider the case $z \in \Omega_2$.  Using Lemma~\ref{existenceofcenter} again, there exists $v \in \Gamma$ such that
$v$ is within a bounded distance of a transition point on $[e,z]$, a transition point on $[z,w]$ and a transition point on $[e,w]$. 

  \begin{center}
\begin{tikzpicture}[scale=2.2]
\draw (-1.8,.5)--(-2,-.4);
\draw (2.2,.5)--(2,-.4);
\draw (-2,-.4)--(2,-.4);
\draw (-1.2,-1.3)--(-1.2,-.4);
\draw[dotted] (-1.2,-.4)--(-1.2,0);
\draw (-1.2,0)--(.8,-.1);
\draw (-1.2,0)--(-1,.4);
\draw[dotted] (.8,-.1)--(.8,-.4);
\draw (.8,-.4)--(.8,-1);
\draw (.8,-.1)--(-1,.4);
\draw (-1,.4)--(-1,1);
\draw[dotted] (-1,1)--(-1,1.4);
\draw (-2.3,1.9)--(-2.5,1);
\draw (.7,1.9)--(.5,1);
\draw (-2.5,1)--(.5,1);
\draw (-1,1.4)--(-1.2,1.7);
\draw (-1.2,1.7)--(-1.9,1.3);
\draw (-1.9,1.3)--(-1,1.4);
\draw (-1.2,1.7)--(-1.2,2.5);
\draw (-1.9,1.3)--(-1.9,2.3);
%\draw (-1.4,0)node{$u_1$};
\draw (-1.5,-1.3) node{$e$};
%\draw (1.1,-.15) node{$u_2$};
%\draw (-.1,.5) node{$u_3$};
\draw (-.9,.2) node{$v$};
\draw (1,-1) node{$z$};
%\draw (.4,1.4) node{$w_1$};
%\draw (-.7,1.8) node{$w_3$};
\draw (-1.1,2.5) node{$y$};
%\draw (-2.2,1.3) node{$w_2$};
\draw (-2.1,2.3) node{$x$};
\draw (-1.3,1.5) node{$w$};
\end{tikzpicture}
\end{center}
  
By the same argument as in the case $z \in \Omega_1$, we have that
  \[
    \sum_{z \in \Omega_2} G_r(e, z) G_r(z, x) G_r(z, y)
    \leq
    c_5 G_r(e, w) G_r(w, x) G_r(w, y)
  \]
This completes the proof of the lemma.  
\end{proof}



This lemma will help us estimate $\mathbf E \left [M_{n,L}^2\right ]$.
We first recall the following result from \cite{SWX} whose proof has nothing to do with hyperbolicity and holds for any finitely generated group $\Gamma$.
As above, for every $x\in \Gamma$, we denote by $Z_x$ the number of particles that ever visit $x$.

\begin{lem}\label{4.3SWX}
Assume that $\nu$ has finite second moment.
Then, there exists $C$ such that for every $x,y\in \Gamma$,
$$\mathbf E[Z_xZ_y]\leq  C \sum_{z\in \Gamma}G_r(e,z)G_r(z,x)G_r(z,y).$$
\end{lem}



We deduce the following result.

\begin{prop}\label{Coro4.4L}
Assume that $\nu$ has finite second moment and that $r<\rho^{-1}$.
Then there exists $C_L$ such that
$$\mathbf{E}[M_{n,L}]\geq C_L \mathrm{e}^{n\omega_{\Gamma,L}(r)}.$$
\end{prop}

\begin{proof}
As in the proof of \cite[Lemma~4.4]{SWX}, we deduce from Lemma~\ref{4.3SWX} that
$$\mathbf P(Z_x\geq 1)\geq c G_r(e,x).$$
Since $\mathbf{E} [M_{n,L}]=\sum_{x\in S_{n,L}}\mathbf P(Z_x\geq 1)$, the result follows from Proposition~\ref{puregrowthtransitional}.
\end{proof}


\begin{cor}\label{GreenLmoment}
Assume that $\nu$ has finite second moment.  For $r \in (1, \rho^{-1})$ and $L$ sufficiently large, there exists a positive constant $c$ such that
$$\mathbf{E} \left[ M_{n, L}^2 \right] \leq c \left( \mathbf{E}[M_{n, L}] \right)^2.$$ 
\end{cor}

\begin{proof}
By Lemma~\ref{omegaL} and the fact that $\omega_{\Gamma}(r) > 1$, we have $\omega_{\Gamma, L}(r) > 1$ for sufficiently large $L$.  Applying Proposition~\ref{GreenL} and Lemma~\ref{Correlationxy},
  \begin{align*}
    \mathbf{E} \left[ M_{n, L}^2 \right]
    \leq&
          c_1 \sum_{x, y \in S_{n, L}} \sum_{z \in \Gamma} G_r(e, z) G_r(z, x) G_r(z, y) \\
    \leq&
          c_2 \sum_{k=0}^n \sum_{w \in S_{k, L}} \sum_{x, y \in S_{n, L}} G_r(e, w) G_r(w, x) G_r(w, y) \\
    \leq&
          c_3 \sum_{k=0}^n \mathrm{e}^{(2n - k) \omega_{\Gamma, L}(r)} \\
    \leq&
          c_4 \left( \mathbf{E}[M_{n, L}] \right)^2.   
  \end{align*}
This yields the desired bound.
\end{proof}

We can now prove the lower bound and finish the proof of Proposition~\ref{P:vgrBRW}.
\begin{proof}[Proof of Proposition~\ref{P:vgrBRW}]
We first fix $r<\rho^{-1}$ and assume that $\nu$ has finite second moment.
By Proposition~\ref{Coro4.4L},
$$\mathbf{P}\left (M_{n,L}\geq \frac{c_1}{2}\mathrm{e}^{\omega_{\Gamma,L}(r)n}\right)\geq \mathbf{P}\left (M_{n,L}\geq \frac{1}{2}\mathbf{E}[M_{n,L}]\right)$$
and so, by the Paley-Zygmund inequality and Corollary~\ref{GreenLmoment}, for some $c_2>0$,
$$\mathbf{P}\left (M_{n,L}\geq \frac{c_1}{2}\mathrm{e}^{\omega_{\Gamma,L}(r)n}\right)\geq \frac{\big(\mathbf{E}[M_{n,L}]\big)^2}{\mathbf{E}\left[M_n^2\right]}\geq c_2.$$
Thus, with positive probability,
the events $\left\{M_{n,L}^{1/n}\geq \left(\frac{c_1}{2}\right)^{1/n}\mathrm{e}^{\omega_{\Gamma,L}(r)}\right\}$
occur for infinitely many $n$ and so, with positive probability, $\limsup \frac{1}{n}\log M_{n,L}\geq \omega_{\Gamma,L}(r)$.
By definition, $M_n\geq M_{n,L}$,
hence for every large enough $L$, with positive probability (a priori depending on $L$), we have
$\limsup \frac{1}{n}\log M_n\geq \omega_{\Gamma,L}(r)$.
By \cite[Lemma~4.7]{SWX}, $\limsup M_n^{1/n}$ is almost surely a constant.
Thus, for every $L$, almost surely we have
$\limsup \frac{1}{n}\log M_n\geq \omega_{\Gamma,L}(r)$.
Letting $L$ tend to infinity along a sequence, it follows from Lemma~\ref{omegaL} that
$\limsup \frac{1}{n}\log M_n\geq \omega_{\Gamma}(r)$.
Thus by Proposition~\ref{upperboundvgrBRW},
$$\limsup \frac{1}{n} \log M_n=\omega_\Gamma(r).$$


We conclude as in \cite{SWX}.
For every $\epsilon>0$, we can construct a probability measure $\nu'$ with mean $r-\epsilon$ and with finite second moment so that $\nu$ stochastically dominates $\nu'$.
Denoting by $M_n'$ the number of vertices in $S_n$ ever visited by a branching random walk driven by $\mu$ and $\nu'$, we see that $M_n$ stochastically dominates $M_n'$ and so
$$\mathbf{P}\left(\limsup \frac{1}{n}\log M_n\geq \omega_\Gamma(r-\epsilon)\right)\geq \mathbf{P}\left(\limsup \frac{1}{n}\log M_n'\geq \omega_\Gamma(r-\epsilon)\right)=1.$$
Since $\omega_P(r)$ is continuous by Corollary~\ref{corolowersemicontinuous}, we deduce that
$$\limsup \frac{1}{n} \log M_n\geq \omega_\Gamma(r)$$ almost surely.
\end{proof}








 

%\section{The volume growth rate of the branching random walk}
%Recall that $\mathcal{P}_n$ is the set of points in $S_n$ that are eventually visited by some particle of the branching random walk and that $M_n=\sharp \mathcal{P}_n$.
%In this section, we compare the growth rate of the Green function $\omega_\Gamma(r)$ with the growth rate of the branching random walk $\log \limsup M_n^{1/n}$.
%Our goal is to prove the following proposition.
%\begin{prop}\label{P:vgrBRW}
%Assume that for every parabolic subgroup $P$, $\omega_P\leq 0$.
%Then, for $\lambda\in [1,\rho^{-1}]$, $H(\lambda)=\limsup M_n^{1/n}$ almost surely.
%\end{prop}
%\subsection{Upper bound}
%For the following proposition, we do not need to assume that the Green function has pure exponential growth and so we do not need to assume that $\omega_P<\omega_\Gamma$.
%\begin{prop}\label{upperboundvgrBRW}
%Almost surely, we have
%$$\limsup M_n^{1/n}\leq \mathrm{e}^{\omega_\Gamma(r)}.$$
%\end{prop}

%The proof of \cite{SWX}, which relies on the Borel-Cantelli lemma, does not use the pure exponential growth and can be applied here.
%We rewrite it for convenience.

%\begin{proof}
%Let $\epsilon>0$.
%By the Markov inequality and the many-to-one formula,
%$$\mathbb{P}\left (M_n^{1/n}\geq \mathrm{e}^{\omega_\Gamma(r)}+\epsilon\right)\leq \frac{\mathbb{E}(M_n)}{(\mathrm{e}^{\omega_\Gamma(r)}+\epsilon)^n}=\frac{H_r(n)}{(\mathrm{e}^{\omega_\Gamma(r)}+\epsilon)^n}.$$
%By definition,
%$\omega_\Gamma(r)=\log \limsup H_r(n)^{1/n}$,
%so there are at most finitely many $n$ such that
%$H_r(n)^{1/n}\geq \mathrm{e}^{\omega_\Gamma(r)}+\epsilon/2$.
%Hence,
%$$\sum_n \frac{H_r(n)}{(\mathrm{e}^{\omega_\Gamma(r)}+\epsilon)^n}<\infty.$$
%The statement of the lemma is thus a consequence of the Borel-Cantelli lemma.
%\end{proof}

%\subsection{Lower bound}
%To prove the lower bound, we need to assume that $\omega_P\leq 0$.
%We first prove the following result.



%\begin{lem}\label{lemmaparabolicsumfinitesubcriticalcase}
%Assume that $\omega_P\leq 0$.
%For every $r<R$, there exists $M=M(r)$ such that
%$$\sum_{x\in P}G_r(e,x)\leq M.$$
%\end{lem}

%\begin{proof}
%Since the underlying random walk is finitely supported, there exists $c\geq 0$ such that for every $x\in\Gamma$, we have that $p_n(e,x)=0$ as long as $n\leq c|x|$.
%Hence,
%$$G_r(e,x)=\sum_{n\geq c|x|}r^np_n(e,x)=\sum_{n\geq c|x|}\left (\frac{r}{R}\right)^nR^np_n(e,x)\leq \left (\frac{r}{R}\right)^{c |x|}G_R(e,x).$$
%Thus, there exists $\alpha>0$ such that
%$$\sum_{x\in P}G_r(e,x)\leq \sum_{x\in P}G_R(e,x)\mathrm{e}^{-\alpha|x|}.$$
%This last sum is finite, since $\omega_P\leq 0$.
%\end{proof}


%The following result replaces \cite[Lemma~4.5]{SWX}.

%\begin{lem}
%For every $1<\lambda<\rho^{-1}$, there exists $c=c(\lambda)$ such that for every $0\leq k\leq 2n$,
%$$\sum_{\underset{d(x,y)=k}{x,y\in S_n}}\sum_{z\in \Gamma}G_\lambda(e,z)G_\lambda(z,x)G_\lambda(z,y)\leq cH(\lambda)^{n+k/2}.$$
%\begin{comment}
%Do you have the little $k$ before $H(\lambda)^{n+k/2}$ on the right side? I think this can be got rid off in the computation: the choice $w_2, w_1$ are fixed once $x,y$ are chosen, so there is no need to take $\sum_{j\leq k}$, right?

%Also in my computation, I'm reduced to the finiteness of the following series: $\sum_{z\in P}G_\lambda(e,z)G_\lambda(z,x)G_\lambda(z,y)$ (where $z$ is summed over a parabolic subgroup $P$ and $x,y$ lie in $P$). This is finite if $\sum_{z\in P}G_\lambda(e,z)$ is finite. I do not know whether this is a more general condition with interesting examples. )
%{\color{red} You're right, but since we first sum over $u_1,u_2,u_3,w_1,w_2,w_3$, we need the converse : once $u_1,u_2,u_3,w_1,w_2,w_3$ are fixed, then the size of $|w_2^{-1}x|$ and $|w_3^{-1}y|$ are fixed and satisfy that their sum is bounded by $k$ (up to an additive error). This is also true.} Yes, you are right. There should be additional $k$ here. Fortunately, this does not effect the proof in [SWX].
%\end{comment}
%\end{lem}

%\begin{proof}
%There are three situations to deal with, as suggested by the following pictures.
%The first one corresponds to $z\in \Omega_1(x,y)$ in \cite{SWX}.
%The second one corresponds to $z\in \Omega_2(x,y)$ or $z\in \Omega_3(x,y)$.
%The third one is new.

%\begin{center}
%\begin{tikzpicture}[scale=2.2]
%\draw (-1.8,.5)--(-2,-.4);
%\draw (2.2,.5)--(2,-.4);
%\draw (-2,-.4)--(2,-.4);
%\draw (-1.2,-1.3)--(-1.2,-.4);
%\draw[dotted] (-1.2,-.4)--(-1.2,0);
%\draw (-1.2,0)--(.8,-.1);
%\draw (-1.2,0)--(.2,.4);
%\draw[dotted] (.8,-.1)--(.8,-.4);
%\draw (.8,-.4)--(.8,-1);
%\draw (.8,-.1)--(.2,.4);
%\draw (.2,.4)--(.2,1);
%\draw[dotted] (.2,1)--(.2,1.4);
%\draw (-2.3,1.9)--(-2.5,1);
%\draw (.7,1.9)--(.5,1);
%\draw (-2.5,1)--(.5,1);
%\draw (.2,1.4)--(-1,1.7);
%\draw (-1,1.7)--(-1.9,1.3);
%\draw (-1.9,1.3)--(.2,1.4);
%\draw (-1,1.7)--(-1,2.5);
%\draw (-1.9,1.3)--(-1.9,2.3);
%\draw (-1.4,0)node{$u_1$};
%\draw (-1.5,-1.3) node{$e$};
%\draw (1.1,-.15) node{$u_2$};
%\draw (-.1,.5) node{$u_3$};
%\draw (1,-1) node{$z$};
%0000000000000\draw (.4,1.4) node{$w_1$};
%\draw (-.7,1.8) node{$w_3$};
%\draw (-.8,2.5) node{$y$};
%\draw (-2.2,1.3) node{$w_2$};
%\draw (-2.1,2.3) node{$x$};
%\end{tikzpicture}
%\end{center}


%Then, $d(e,w_1)\leq n-k/2$ and
%$d(w_2,x) +d(w_3,y)\leq k$.
%Moreover, when $w_1$ $w_2$, $w_3$ are fixed, then $d(w_2,x)=i$ and $d(w_3,y)=j$ are fixed.
%We can bound the sum by
%\begin{align*}
%&\sum_{|u_1|\leq n-k/2}G_\lambda(e,u_1)\sum_{u_2}G_\lambda(u_1,u_2)\sum_{z}G_\lambda(u_2,z)^3\sum_{u_3}G_\lambda(u_2,u_3)^2\\
%&\hspace{1cm}\sum_{w_1}G_\lambda(u_3,w_1)^2\sum_{w_2}G_\lambda(w_1,w_2)\sum_{w_3}G_\lambda(w_1,w_3)\\
%&\hspace{1cm}\sum_{|w_2^{-1}x|=i}G_\lambda(w_2,x)\sum_{|w_3^{-1}y|=j}G_\lambda(w_3,y).
%\end{align*}
%\begin{comment}
%Here, $u_1$ is also varied on $[e, w_1]$ after $w_1$ is chosen. This is the main difference with the proof below. The other difference is that I try to make clear the order of different sums: we are actually having two triples to vary. We must be careful which variables are fixed when the bounds are plugged in. 
%\end{comment} 

%By what precedes,
%$\sum G_\lambda(u_1,u_2)$, $\sum G_\lambda(u_2,u_3)^2$, $\sum G_\lambda(w_1,w_2)$ and $\sum G_\lambda(w_1,w_3)$ are all bounded by $M$. 
%We then bound $\sum G_\lambda(u_2,z)^3$ and $\sum G_\lambda (u_3,w)^2$ by some constant like in \cite{SWX}. 
%Finally, $\sum G_\lambda(e,u_1)$ is bounded by $H(\lambda)^{n-k/2}$ while $\sum G_\lambda(w_2,x)$ is bounded by $H(\lambda)^{i}$ and $\sum G_\lambda(w_3,y)$ by $H(\lambda)^{j}$, with $i+j\leq k$.
%This concludes the proof in the first situation.



%\begin{center}
%\begin{tikzpicture}[scale=2.2]
%\draw (-1.8,.5)--(-2,-.4);
%\draw (2.2,.5)--(2,-.4);
%\draw (-2,-.4)--(2,-.4);
%\draw (-1.2,-1.3)--(-1.2,-.4);
%\draw[dotted] (-1.2,-.4)--(-1.2,0);
%\draw (-1.2,0)--(.8,-.1);
%\draw (-1.2,0)--(.2,.4);
%\draw[dotted] (.8,-.1)--(.8,-.4);
%\draw (.8,-.4)--(.8,-1);
%\draw (.8,-.1)--(.2,.4);
%\draw (.2,.4)--(.2,1);
%\draw[dotted] (.2,1)--(.2,1.4);
%\draw (-2.3,1.9)--(-2.5,1);
%\draw (.7,1.9)--(.5,1);
%\draw (-2.5,1)--(.5,1);
%\draw (.2,1.4)--(-1,1.7);
%\draw (-1,1.7)--(-1.9,1.3);
%\draw (-1.9,1.3)--(.2,1.4);
%\draw (-1,1.7)--(-1,2.5);
%\draw (-1.9,1.3)--(-1.9,2.3);
%\draw (-1.4,0)node{$w_1$};
%\draw (-1.5,-1.3) node{$e$};
%\draw (1.1,-.15) node{$w_3$};
%\draw (-.1,.5) node{$w_2$};
%\draw (1,-1) node{$y$};
%\draw (.4,1.4) node{$v_1$};
%\draw (-.7,1.8) node{$v_2$};
%\draw (-.8,2.5) node{$z$};
%\draw (-2.2,1.3) node{$v_3$};
%\draw (-2.1,2.3) node{$x$};
%\end{tikzpicture}
%\end{center}

%Again, $d(e,w_1)\leq n-k/2$ and $d(w_2,x)+d(w_3,y)\leq k$. Note that we also have $d(v_3,x)+d(w_3,y)\leq k$. We bound the sum by
%\begin{align*}
%&\sum_{|w_1|\leq n-k/2}G_\lambda(e,w_1)\sum_{w_2}G_\lambda(w_1,w_2)\sum_{v_1}G_\lambda(w_2,v_1)^2\sum_{v_2}G_\lambda(v_1,v_2)^2\\
%&\hspace{1cm}\sum_{z}G_\lambda(v_2,z)^3\sum_{v_3}G_\lambda(v_2,v_3)\sum_{w_3}G_\lambda(w_2,w_3)\\
%&\hspace{1cm}\sum_{|v_3^{-1}x|=i}G_\lambda(v_3,x)\sum_{|w_3^{-1}y|=j}G_\lambda(w_3,y).
%\end{align*}


%Then, $\sum G_\lambda(v_1,v_2)^2$, $\sum G_\lambda(v_2,v_3)$, $\sum G_\lambda(w_1,w_2)$ and $\sum G_\lambda(w_2,w_3)$ are bounded by $M$,
%while $\sum G_\lambda(w_2,v_1)^2$ and $\sum G_\lambda (v_2,z)^3$ are bounded by some constant.
%Finally, $G_\lambda(e,w_1)$ is bounded by $H(\lambda)^{n-k/2}$, $\sum G_\lambda(v_3,x)$ by $H(\lambda)^i$ and $\sum G_\lambda(w_3,y)$ by $H(\lambda)^{j}$.
%This again concludes the proof.

%\begin{center}
%\begin{tikzpicture}[scale=2.2]
%\draw (-1.8,.5)--(-2,-.4);
%\draw (2.2,.5)--(2,-.4);
%\draw (-2,-.4)--(2,-.4);
%\draw (-1.2,-1.3)--(-1.2,-.4);
%\draw[dotted] (-1.2,-.4)--(-1.2,0);
%\draw (-1.2,0)--(.8,-.1);
%\draw (-1.2,0)--(.2,.4);
%\draw (-1.2,0)--(1,.3);
%\draw[dotted] (.8,-.1)--(.8,-.4);
%\draw (.8,-.4)--(.8,-1);
%\draw (1,.3)--(1,1.1);
%\draw (.8,-.1)--(1,.3);
%\draw (1,.3)--(.2,.4);
%\draw (.2,.4)--(.2,1.6);
%\draw (-1.4,0)node{$w_1$};
%\draw (-1.5,-1.3) node{$e$};
%\draw (1.1,-.15) node{$w_3$};
%\draw (1.2,.3) node{$u$};
%\draw (-.1,.5) node{$w_2$};
%\draw (1,-1) node{$y$};
%\draw (1.2,1.1) node{$z$};
%\draw (0,1.6) node{$x$};
%\end{tikzpicture}
%\end{center}

%In this last situation, $d(e,w_1)\leq n-k/2$ and $d(w_2,x)+d(w_3,y)\leq k$ as above.
%We bound the sum by
%\begin{align*}
%&\sum_{|w_1|\leq n-k/2}G_\lambda(e,w_1)\sum_{u}G_\lambda(w_1,u)\sum_{z}G_\lambda(u,z)^3\sum_{w_2}G_\lambda(u,w_2)\sum_{w_3}G_\lambda(u,w_3)\\
%&\hspace{1cm}\sum_{|w_2^{-1}x|=i}G_\lambda(w_2,x)\sum_{|w_3^{-1}y|=j}G_\lambda(w_3,y).
%\end{align*}

%Here, $\sum G_\lambda (w_1,u)$, $\sum G_\lambda(u,w_2)$ and $\sum G_\lambda(u,w_3)$ are bounded by $M$, while
%$\sum G_\lambda(u,z)^3$ is bounded by a constant.
%Finally,
%$G_\lambda(e,w_1)$ is bounded by $H(\lambda)^{n-k/2}$, $\sum G_\lambda(w_2,x)$ by $H(\lambda)^i$ and $\sum G_\lambda(w_3,y)$ by $H(\lambda)^{j}$.
%This concludes the proof in the last situation.
%\end{proof}


%{\color{blue}What follows is a more elaborate version of the proof of Lemma 5.3 after discussion with Longmin.}


%Here are a few facts about the relative thin triangle property (I will try to give a precise reference..). 
%\begin{lem}
%There exists $\delta>0$ and $C\geq 0$ with the following property.
%For any triple of elements $(x_0,x_1,x_2)$ in $\Gamma$, there exists a triple of elements $(w_0,w_1,w_2)$ in $\Gamma$  such that the following holds for $i \mod 3$
%$$
%\left |d(x_i,x_{i+1}) - \bigg(d(x_i,w_i)+d(w_i,w_{i+1})+d(w_{i+1},x_{i+1})\bigg)\right |\leq C.
%$$
%and $\delta_{w_i}(x_i,x_{i+1})\ge \delta$ and $\delta_{w_{i+1}}(x_i,x_{i+1})\ge \delta$.
%\end{lem}

%The (ordered) triple $(w_0,w_1,w_2)$ is called a \textit{center} of the triangle $\Delta(x_0,x_1,x_2)$. A center $(w_0,w_1,w_2)$ is called \textit{of type $(l, i, j)$} if $|d(x_0, w_0)- l|\leq C$, $|d(x_1,w_1)-i|\leq C$ and $|d(x_2, w_2)- j|\leq C$.
%\begin{lem}
%There exists a uniform constant $M>0$ such that any two centers $(w_0,w_1,w_2)$ and $(w_0',w_1',w_2')$ of $(x_0,x_1,x_2)$ are $M$-bounded: $d(w_i,w_i')\le M$ for $i=0,1,2.$
%Moreover, if $diam(w_0, w_1,w_2)>M$ then there exists a unique peripheral coset $P$ such that $d(w_i, P)\le M$.
%\end{lem}
 

 


%The following result replaces \cite[Lemma~4.5]{SWX}.

%\begin{lem}\label{lemma4.5SWX}
%For every $1<\lambda<\rho^{-1}$, there exists $c=c(\lambda)$ such that for every $0\leq k\leq 2n$,
%$$\sum_{\underset{d(x,y)=k}{x,y\in S_n}}\sum_{z\in \Gamma}G_\lambda(e,z)G_\lambda(z,x)G_\lambda(z,y)\leq cH(\lambda)^{n+k/2}.$$
%\end{lem}



%\begin{proof} 
%Let $(l, i, j)$  be  a triple of three non-negative integers satisfying   $0\le l\le n-k$ and $i+j\le 2k$. 

%Let $\Omega(l, i, j)$ be the set of triples $(x, y, z)\in S_n\times S_n\times \Gamma$ with $d(x,y)=2k$ such that the triangle $\Delta(1,x,y)$ has a $(l, i, j)$-center $(w_1,w_2,w_3)$. We then  decompose $\Omega(l, i, j)$ into the union of the following non-exclusive three sets.

%Roughly speaking (to be more precise later..), let $\Omega_1(l, i, j)$ be the set of triples $(x, y, z)\in \Omega(l, i, j)$  such that  the triangle $\Delta(1,x,z)$ has a center  $(u_1,u_2,u_3)$ such that $d(e, u_1)\le d(e, w_1)$. 

%The set $\Omega_2(l, i, j)$ consists of $(x, y, z)\in\Omega(l, i, j)$ such that the triangle $\Delta(1,x,z)$ or $\Delta(1,y,z)$ has a center $(v_1,v_2,v_3)$ such that $d(e, v_2)\ge d(e, w_2)$ or $d(e, v_3)\ge d(e, w_3)$. 

%The set $\Omega_3(l, i, j)$ consists of $(x, y, z)\in\Omega(l, i, j)$  such that the center  of $(1,x,z)$ or $(1,y,z)$ is    the center  of $(1,x,y)$. 
 
%\begin{center}
%\begin{tikzpicture}[scale=2.2]
%\draw (-1.8,.5)--(-2,-.4);
%\draw (2.2,.5)--(2,-.4);
%\draw (-2,-.4)--(2,-.4);
%\draw (-1.2,-1.3)--(-1.2,-.4);
%\draw[dotted] (-1.2,-.4)--(-1.2,0);
%\draw (-1.2,0)--(.8,-.1);
%\draw (-1.2,0)--(.2,.4);
%\draw[dotted] (.8,-.1)--(.8,-.4);
%\draw (.8,-.4)--(.8,-1);
%\draw (.8,-.1)--(.2,.4);
%\draw (.2,.4)--(.2,1);
%\draw[dotted] (.2,1)--(.2,1.4);
%\draw (-2.3,1.9)--(-2.5,1);
%\draw (.7,1.9)--(.5,1);
%\draw (-2.5,1)--(.5,1);
%\draw (.2,1.4)--(-1,1.7);
%\draw (-1,1.7)--(-1.9,1.3);
%\draw (-1.9,1.3)--(.2,1.4);
%\draw (-1,1.7)--(-1,2.5);
%\draw (-1.9,1.3)--(-1.9,2.3);
%\draw (-1.4,0)node{$u_1$};
%\draw (-1.5,-1.3) node{$e$};
%\draw (1.1,-.15) node{$u_2$};
%\draw (-.1,.5) node{$u_3$};
%\draw (1,-1) node{$z$};
%\draw (.4,1.4) node{$w_1$};
%\draw (-.7,1.8) node{$w_3$};
%\draw (-.8,2.5) node{$y$};
%\draw (-2.2,1.3) node{$w_2$};
%\draw (-2.1,2.3) node{$x$};
%\end{tikzpicture}
%\end{center}



%\textbf{Case 1.}  
%We  estimate the sum over the triples in $\Omega_1(l, i,j)$. The  configurations are  determined by choosing $\{u_1\in\Gamma: |u_1|\le l\}$ and  $u_2, u_3\in P$ with $|u_1^{-1}u_3|\le l$, and then by choosing $w_1\in\Gamma$ with $|u_3^{-1} w_1|\le l$, and $w_2,w_3\in P$ with $|w_2^{-1}x|=i,|w_2^{-1}y|=j$. Thus, we have
%\begin{align*}
%&\sum_{(x,y,z)\in\Omega_1(l,i,j)}   G_\lambda(e,z)G_\lambda(z,x)G_\lambda(z,y) \\ 
%\le&  \sum_{|u_1|\le l} G_\lambda(e,u_1) \Bigg[\left( \sum_{u_2,u_3\in u_1 P} \sum_{z\in\Gamma} G_\lambda(u_1,z)G_\lambda(z,u_2)^2 G_\lambda(u_2,u_3)^2 \right)  \\
%& \Bigg( \sum_{w_1\in\Gamma} G_\lambda(u_3,w_1)^2 \sum_{w_2,w_3\in w_1P} G_\lambda(w_1,w_2)G_\lambda(w_1,w_3)  \\
%&    \sum_{|w_2^{-1}x|=i}   G_\lambda(w_2,x) \sum_{|w_3^{-1}y|=j}  G_\lambda(w_3,y)\Bigg)\Bigg].
%\end{align*}

%Note that there exists a finite constant $M$ such that $\sum_{w_1\in\Gamma} G_\lambda^2(u_3,w_1)   \le M$  and $$ \sum_{z\in\Gamma} G_\lambda(u_1,z) G_\lambda(z,u_2)^2 \le  \sum_{z\in\Gamma} G_\lambda(e,e) G_\lambda(z,u_2)^2 \le M$$
%and 
%$$
%\sum_{u_2,u_3\in u_1P} G_\lambda(u_2,u_3)^2 \le \left(\sum_{u\in P} G_\lambda(e,u)\right)^2<M
%$$
%and for given $w_1\in\Gamma$,
%$$
%\sum_{w_2,w_3\in w_1P} G_\lambda(w_1,w_2)G_\lambda(w_1,w_3) \le  \left(\sum_{u\in P} G_\lambda(e,u)\right)^2<M.
%$$
 
%We vary $0\le i, j\le k$ and then $0\le l\le n-k$ get
%\begin{align*}
%&\sum_{l, i, j}\sum_{\Omega_1(l,i,j)}   G_\lambda(e,z)G_\lambda(z,x)G_\lambda(z,y)\\
%\prec  
%&\sum_{0\le l \le n-k}  \left(\sum_{|u_1|\leq l}G_\lambda(e,u_1)\sum_{|w_2^{-1}x|\le k}G_\lambda(w_2,x)\sum_{|w_3^{-1}y|\le k}G_\lambda(w_3,y)\right)\\
%\prec & H(\lambda)^{n-k} H(\lambda)^{2k}\prec  H(\lambda)^{n+k}.
%\end{align*}
 
 
 
 
%\textbf{Case 2.} The configurations of $\Omega_{2}(l, i, j)$  are given by choosing first $w_1\in S_l,w_2,w_3\in w_1P$ and then $v_1\in \Gamma,v_2,v_3\in v_1P$. Similarly,
%we can bound the sum  
%\begin{align*}
%&\sum_{(x,y,z)\in\Omega_2(l,i,j)}   G_\lambda(e,z)G_\lambda(z,x)G_\lambda(z,y) \\ 
%\le &\sum_{|w_1|=l}G_\lambda(e,w_1) \sum_{w_2, w_3\in w_1P} \Bigg[\sum_{v_1\in\Gamma}\Bigg ( \sum_{v_2,v_3\in v_1P}  \\
%&\left( \sum_{z\in\Gamma}  G_\lambda(z, v_2)^2 \right) G_\lambda(v_1,v_2)^2 G_\lambda(w_1,v_1) G_\lambda(v_2,v_3) G_\lambda(v_1,w_3)  \\
%&\sum_{|v_3^{-1}x|=i'} G_\lambda(v_3,x) \sum_{|w_3^{-1}y|= j}  G_\lambda(w_3,y)\Bigg)\Bigg].
%\end{align*}
 

%\begin{center}
%\begin{tikzpicture}[scale=2.2]
%\draw (-1.8,.5)--(-2,-.4);
%\draw (2.2,.5)--(2,-.4);
%\draw (-2,-.4)--(2,-.4);
%\draw (-1.2,-1.3)--(-1.2,-.4);
%\draw[dotted] (-1.2,-.4)--(-1.2,0);
%\draw (-1.2,0)--(.8,-.1);
%\draw (-1.2,0)--(.2,.4);
%\draw[dotted] (.8,-.1)--(.8,-.4);
%\draw (.8,-.4)--(.8,-1);
%\draw (.8,-.1)--(.2,.4);
%\draw (.2,.4)--(.2,1);
%\draw[dotted] (.2,1)--(.2,1.4);
%\draw (-2.3,1.9)--(-2.5,1);
%\draw (.7,1.9)--(.5,1);
%\draw (-2.5,1)--(.5,1);
%\draw (.2,1.4)--(-1,1.7);
%\draw (-1,1.7)--(-1.9,1.3);
%\draw (-1.9,1.3)--(.2,1.4);
%\draw (-1,1.7)--(-1,2.5);
%\draw (-1.9,1.3)--(-1.9,2.3);
%\draw (-1.4,0)node{$w_1$};
%\draw (-1.5,-1.3) node{$e$};
%\draw (1.1,-.15) node{$w_3$};
%\draw (-.1,.5) node{$w_2$};
%\draw (1,-1) node{$y$};
%\draw (.4,1.4) node{$v_1$};
%\draw (-.7,1.8) node{$v_2$};
%\draw (-.8,2.5) node{$z$};
%\draw (-2.2,1.3) node{$v_3$};
%\draw (-2.1,2.3) node{$x$};
%\end{tikzpicture}
%end{center}
 
%As above, we have $\sum_{z\in \Gamma} G_\lambda(z,v_2)^2\le M$, and for any $v_1, w_1\in\Gamma$, we have 
%$$\sum_{v_2, v_3\in v_1P} G_\lambda(w_1,v_1) G_\lambda(v_2,v_3)G_\lambda(v_2,v_1)\le G_\lambda(e,e) (\sum_{v\in P} G_\lambda(e,v))^2\le M$$ 
%So the above sum is bounded by
%\begin{align*}
%\prec &\sum_{|w_1|=l}G_\lambda(e,w_1)  \sum_{w_2, w_3\in w_1P} \Bigg(\sum_{v_1\in\Gamma}  G_\lambda(v_1, v_2)^2    \sum_{|v_3^{-1}x|=i'} G_\lambda(v_3,x) \sum_{|w_3^{-1}y|= j}  G_\lambda(w_3,y)\Bigg).
%\end{align*}
%Noting that $\sum_{v_1\in\Gamma}  G_\lambda(v_1, v_2)^2 \le M$ and varying $0\le l\le n-k, 0\le i, j\le k$, we have
%\begin{align*}
%&\sum_{l,i,j}\sum_{\Omega_2(l,i,j)}   G_\lambda(e,z)G_\lambda(z,x)G_\lambda(z,y)\\
%\prec &\sum_{|w_1|\leq n-k}G_\lambda(e,w_1) \sum_{|v_3^{-1}x|\le k}G_\lambda(v_3,x)\sum_{|w_3^{-1}y|\le k}G_\lambda(w_3,y)\\
%\prec  & H(\lambda)^{n-k}H(\lambda)^{k}H(\lambda)^{k}\prec  H(\lambda)^{n+k}.
%\end{align*}



 
%\begin{center}
%\begin{tikzpicture}[scale=2.2]
%\draw (-1.8,.5)--(-2,-.4);
%\draw (2.2,.5)--(2,-.4);
%\draw (-2,-.4)--(2,-.4);
%\draw (-1.2,-1.3)--(-1.2,-.4);
%\draw[dotted] (-1.2,-.4)--(-1.2,0);
%\draw (-1.2,0)--(.8,-.1);
%\draw (-1.2,0)--(.2,.4);
%\draw (-1.2,0)--(1,.3);
%\draw[dotted] (.8,-.1)--(.8,-.4);
%\draw (.8,-.4)--(.8,-1);
%\draw (1,.3)--(1,1.1);
%\draw (.8,-.1)--(1,.3);
%\draw (1,.3)--(.2,.4);
%\draw (.2,.4)--(.2,1.6);
%\draw (-1.4,0)node{$w_1$};
%\draw (-1.5,-1.3) node{$e$};
%\draw (1.1,-.15) node{$w_3$};
%\draw (1.2,.3) node{$u$};
%\draw (-.1,.5) node{$w_2$};
%\draw (1,-1) node{$y$};
%\draw (1.2,1.1) node{$z$};
%\draw (0,1.6) node{$x$};
%\end{tikzpicture}
%\end{center}

%\textbf{Case 3}.    
%The configurations of $\Omega_{3}(l, i, j)$  are obtained by choosing first $w_1\in S_l,w_2,w_3\in w_1P$ and then $u\in w_1P$. We derive similarly, 
%\begin{align*}
%&\sum_{(x,y,z)\in\Omega_3(l,i,j)} G_\lambda(e,z)G_\lambda(z,x)G_\lambda(z,y) \\ 
%\le  &\sum_{|w_1|=l}G_\lambda(e,w_1)\sum_{w_2,w_3, u\in w_1 P}\Bigg(\\
%&\left(\sum_{z\in\Gamma}G_\lambda(w_1,u) G_\lambda(u,z)^3\right) G_\lambda(u,w_2) G_\lambda(u,w_3)\\
%& \sum_{|w_2^{-1}x|=i}G_\lambda(w_2,x)\sum_{|w_3^{-1}y|=j}G_\lambda(w_3,y)\Bigg). 
%\end{align*}

%Here, $\sum_{z\in\Gamma} G_\lambda (u,z)^{3}$, and $\sum_{w_2,w_3\in w_1P} G_\lambda(u,w_2),G_\lambda(u,w_3)$  are bounded by $M$ for any $u\in\Gamma$. Varying $0\le l\le n-k, 0\le i, j\le k$, we have
%\begin{align*}
%&\sum_{ l,i,j}\sum_{ \Omega_3(l,i,j)} G_\lambda(e,z)G_\lambda(z,x)G_\lambda(z,y) \\
%\prec& \sum_{|w_1|\leq n-k}G_\lambda(e,w_1) \sum_{|w_2^{-1}x|\le k}G_\lambda(w_2,x)\sum_{|w_3^{-1}y|\le k}G_\lambda(w_3,y)\\
%\prec  & H(\lambda)^{n-k}H(\lambda)^{k}H(\lambda)^{k} \prec   H(\lambda)^{n+k}.
%\end{align*}
%Summarizing all the three cases, we have $$\sum_{l,i,j} \sum_{ \Omega(l,i,j)} G_\lambda(e,z)G_\lambda(z,x)G_\lambda(z,y)  \le  H(\lambda)^{n+k}$$
%concluding the proof.
%\end{proof}

%{\color{red}I agree with your proof.
%I'm trying to refine it so that we don't need the assumption $\omega_P\leq 0$ but only $\omega_P<\omega_\Gamma$.

%Here is a refined version where we do not need to use that $\sum_{u\in P}G_\lambda(e,u)$ is bounded and where we only need that for all $P$, $\omega_P(\lambda)<\omega_\Gamma(\lambda)$.
%We fix $\epsilon$ such that $\sup_P\omega_P(\lambda)+\epsilon<\omega_\Gamma(\lambda)$.

%Also, there is a differences of notations. In the statement of the lemma, $d(x,y)=k$, but inside the proof, $d(x,y)=2k$.
%I choose $d(x,y)=k$.

%Case 1 : $(u_1,u_2,u_3)$ is a center of $(e,w_1,z)$.

%We first fix $u_1$, then $u_2,u_3\in u_1P$ and then $w_1$ such that $d(e,u_1)+d(u_1,u_3)+d(u_3,w_1)\leq l+C$.
%We then fix $w_2$ and $w_3$ in $w_1Q$ (possibly some other parabolic subgroup $Q$ than $P$) such that $d(x,w_2)+d(w_2,w_3)+d(w_3,y)\leq k+C$.
%We use the notation $i=d(w_2,x)$, $j=d(w_3,y)$, $m=d(w_2,w_3)$ with $i+j+m\leq k+C$ and $l_1=d(e,u_1)$, $l_2=d(u_1,u_3)$, $l_3=d(u_3,w_1)$ with $l_1+l_2+l_3\leq l+C$. We will actually not need $l_3$ in the computations.
%We also set $m_2=d(w_1,w_2)$ and $m_3=d(w_1,w_3)$.
%We have $m+i+j\leq k+C$, $2l+m_1+m_2\leq 2n+C$ in this situation.
%We get the bound
%\begin{align*}
%&\sum_{(x,y,z)\in\Omega_1}   G_\lambda(e,z)G_\lambda(z,x)G_\lambda(z,y) \\ 
%\le& \sum_{l_1+l_2\leq n-k/2+C} \hspace{.2cm}\sum_{|u_1|\le l_1} G_\lambda(e,u_1) \sum_{\underset{d(u_1,u_3)\leq l_2)}{u_2,u_3\in u_1 P}}G_\lambda(u_1,u_2) G_\lambda(u_2,u_3)^2\\
%&\sum_{z\in\Gamma} G_\lambda(u_2,z)G_\lambda(z,u_2)^2   \sum_{w_1\in\Gamma} G_\lambda(u_3,w_1)^2 \\
%&\sum_{i+j+m\leq k+C}\hspace{.2cm}\sum_{\underset{d(w_2,w_3)=m}{w_2,w_3\in w_1P}} G_\lambda(w_1,w_2)G_\lambda(w_1,w_3)  \\
%&    \sum_{d(w_2,x)=i}   G_\lambda(w_2,x) \sum_{d(w_3,y)=j}  G_\lambda(w_3,y).
%\end{align*}

%}

%We are now ready to complete the proof of Proposition~\ref{P:vgrBRW}.
%\begin{proof}
%We first assume that $\lambda<R$ and that the measure $\nu$ has finite second moment.
%Then, by \cite[Lemma~4.3, Corollary~4.4]{SWX}, letting $Z_x$ be the number of particles of the branching random walk that ever visit $x$, we have
%\begin{equation}\label{equationlemma4.3}
%\mathbb{E}\left[Z_xZ_y\right]\lesssim \sum_{z\in \Gamma}G_\lambda(e,z)G_\lambda(z,x)G_\lambda(z,y).
%\end{equation}
%and
%\begin{equation}\label{equationcorollary4.4}
%\mathbb{E}\left[M_n\right]\gtrsim \mathrm{e}^{\omega_\Gamma(\lambda) n}.
%\end{equation}
%Thus,
%$$\mathbb{E}\left[M_n^2\right]\lesssim \sum_{x,y\in S_n}\mathbb{E}\left[Z_xZ_y\right]\lesssim \sum_{k\leq 2n}\sum_{\underset{d(x,y)=k}{x,y\in S_n}}\sum_{z\in \Gamma}G_\lambda(e,z)G_\lambda(z,x)G_\lambda(z,y)$$
%and so by Lemma~\ref{lemma4.5SWX}, we have
%$$\mathbb{E}\left[M_n^2\right]\lesssim \sum_{k\leq 2n}\mathrm{e}^{\omega_\Gamma(\lambda) (n+k/2)}\lesssim \mathrm{e}^{\omega_\Gamma(\lambda) 2n}.$$
%Therefore, we deduce from~(\ref{equationcorollary4.4}) that
%\begin{equation}\label{equationcorollary4.6}
%\mathbb{E}\left[M_n^2\right]\lesssim \left(\mathbb{E}[M_n]\right)^2.
%\end{equation}
%It follows that for some $c_1> 0$,
%$$\mathbb{P}\left (M_n\geq \frac{c_1}{2}\mathrm{e}^{\omega_\Gamma(\lambda)n}\right)\geq \mathbb{P}\left (M_n\geq \frac{1}{2}\mathbb{E}[M_n]\right)$$
%and so, by the Paley-Zygmund inequality, for some $c_2>0$,
%$$\mathbb{P}\left (M_n\geq \frac{c_1}{2}\mathrm{e}^{\omega_\Gamma(\lambda)n}\right)\geq \frac{\left(\mathbb{E}[M_n]\right)^2}{\mathbb{E}\left[M_n^2\right]}\geq c_2.$$
%Thus, with positive probability,
%the events $\left\{M_n^{1/n}\geq \left(\frac{c_1}{2}\right)^{1/n}\mathrm{e}^{\omega_\Gamma(\lambda)}\right\}$
%occur for infinitely many $n$.
%Hence, with positive probability, $\limsup M_n^{1/n}\geq \mathrm{e}^{\omega_\Gamma(\lambda)}$.
%By \cite[Lemma~4.7]{SWX}, $\limsup M_n^{1/n}$ is almost surely a constant and so by Proposition~\ref{upperboundvgrBRW},
%$$\limsup M_n^{1/n}=\mathrm{e}^{\omega_\Gamma(\lambda)}.$$

%We conclude as in \cite{SWX}.
%For every $\epsilon>0$, we can construct a probability measure $\nu'$ with mean $\lambda-\epsilon$ and with finite second moment so that $\nu$ stochastically dominates $\nu'$.
%Denoting by $M_n'$ the number of vertices in $S_n$ ever visited by a branching random walk driven by $\mu$ and $\nu'$, we see that $M_n$ stochastically dominates $M_n'$ and so
%$$\mathbb{P}\left(\limsup M_n^{1/n}\geq \mathrm{e}^{\omega_\Gamma(\lambda-\epsilon)}\right)\geq \mathbb{P}\left(\limsup (M_n')^{1/n}\geq \mathrm{e}^{\omega_\Gamma(\lambda-\epsilon)}\right)=1.$$
%Since $\omega_P(\lambda)$ is a continuous function of $\lambda$, we deduce that
%$\limsup M_n^{1/n}\geq \mathrm{e}^{\omega_\Gamma(\lambda)}$ almost surely.
%\end{proof}













%\section{Hausdorff dimension, first version. Failure, but there still are interesting comments so I left everything}


%{\color{blue}Setting $\partial \hat{\Gamma}$ the set of conical limit points, we want to obtain
%$$\Lambda_{\mathcal{F}}\cap \partial \hat{\Gamma}=\big \{X_\xi, \xi \in \partial\mathcal{T}\big \}$$
%and
%$$\Lambda_{\mathcal{B}}\cap \partial \hat{\Gamma}=\big \{X_\xi, \xi \in \partial\mathcal{T}\big \}.$$
%Here is a possible strategy. If the underlying random walk has finite support, we could try to adapt the arguments of Theorem~5.2 of \cite{SistoTaylor} to prove that
%$$\mathbb{P}(\exists P\ \mathrm{ parabolic}, d(e,\pi_P(X_u))\geq A \log|u|)$$
%tends to 0 as $|u|$ tends to infinity, for suitable $A$.
%}


%{\color{red}
%Let $u\in \mathcal{T}$.
%Define $\Lambda_u\subset \Lambda$ as the set of all $X_\gamma$, $\gamma\in \partial \mathcal{T}_u$.
%Also, for $\gamma\in \partial \mathcal{T}_u$, set $\xi=X_\gamma$ and define $\pi_{\gamma}(u)$ as
%the projection of $X_u$ on a geodesic from $e$ to $\xi$
%\begin{itemize}
%\item the projection of $X_u$ on a relative geodesic from $e$ to $\xi$ if $\xi$ is conical,
%\item or the projection of $X_u$ on $P_\xi$, if $\xi$ is parabolic, where $P_\xi$ is the stabilizer of $\xi$.
%\end{itemize}
%Finally, for $0\leq m<n$, set
%$$f_{m,n}=\sum_{k=m+1}^n\sup_{|u|=k}\sup_{\gamma\in\partial \mathcal{T}_u}d(X_u,\pi_\gamma(u)).$$

%By definition, $f_{0,n}\leq f_{0,m}+f_{m,n}$.

%We would need to prove the following two results. I'm skeptical about the first one, so we might need to change our process $f_{m,n}$.

%\begin{lem}
%The sets $\{f_{m,n},0\leq m<n\}$ and $\{f_{m+1,n+1},0\leq m<n\}$ have the same distribution and the induced shift transformation $f_{m,n}\mapsto f_{m+1,n+1}$ is ergodic.
%\end{lem}

%\lw{
%These two sequences cannot have the same distributions. For example, $f_{0,1}$ is bounded (since the random walk has finite support), but $f_{k,k+1}$ may be quite large for large $k$.
%}

%\begin{prop}
%Almost surely, we have
%$$\frac{f_{0,n}-f_{0,n-1}}{n}\underset{n\to \infty}{\longrightarrow}0.$$
%\end{prop}

%\begin{proof}
%This proposition might be adapted from \cite[Lemma~7]{Tiozzo}.
%In \cite{Tiozzo}, the proof is a bit more involved because the author doesn't know a priori that $f=f_{0,1}$ is an integrable function. But here, with the choice of $f_{0,1}$, this is true.
%Hence, once we know that $1/nf_{0,n}$ almost surely converges to something, we are done.
%I think we don't even need to use ergodicity at this point : in the classical Birkhoff's case where $f_{0,n}=f\circ T^n$, $1/nf_{0,n}$ converges to some function $\tilde{f}$ and ergodicity is used to prove that $\tilde{f}$ is constant. But even if $\tilde{f}$ is not a constant, then $1/nf_{0,n}-1/nf_{0,n-1}$ converges to 0.
%I'm not sure about our situation because I'm not a specialist on processes defined by two indices $m,n$.
%\end{proof}


%For $g\in \Gamma$ and $R\geq 0$, denote by $\Pi(g,R)$ the big shadow at $g$ of width $R$.
%\begin{prop}
%For any $\epsilon>0$ and for any $n\geq 0$,
%$$\Lambda\subset \bigcup_{m=n}^\infty \bigcup_{g\in \mathcal{P}_m}\Pi(g,\epsilon m).$$
%\end{prop}

%\begin{proof}
%Fix $\epsilon>0$ and $n\geq 0$.
%Consider $\xi\in \Lambda$.
%Then, for $n$ large enough, there exists $u\in \mathcal{T}_n$ such that $X_u$ is within $\epsilon |u|$ of a geodesic from $e$ to $\xi$.
%Also, if $n$ is large enough, then $d(e,X_u)\geq l_0|u|$.
%Hence, $X_u$ is within $C\epsilon d(e,X_u)$ of such a geodesic and so
%$\xi \in \Pi(X_u,C\epsilon d(e,X_u))$.
%This concludes the proof.
%\end{proof}

%\begin{lem}
%For any $m$, for any $g\in S_m$,
%$$\mathrm{diam}(\Pi(g,\epsilon m))\leq Ca^{(1-2\epsilon)m}.$$
%\end{lem}

%\begin{proof}
%Consider $\xi,\zeta\in \Pi(g,\epsilon m)$.
%Then, there exist a geodesic ray $[e,\xi]$ and a point $x$ on this geodesic and there exist a geodesic ray $[e,\zeta]$ and a point $y$ on this geodesic such that $x$ and $y$ are within a distance at most $\epsilon m$ of $g$.
%Consider the path $\alpha$ obtained by concatenating the sub-geodesic $\alpha_1$ from $\xi$ to $x$, a geodesic $\alpha_2$ from $x$ to $g$, a geodesic $\alpha_3$ from $g$ to $y$ and the sub-geodesic $\alpha_4$ from $y$ to $\zeta$.
%Then, the $f$-lengths of $\alpha_1$ and $\alpha_4$ is at most $C_1a^{m}$.
%Now, the geodesic from $x$ to $g$ stays at distance at least $(n-\epsilon n)$ from $e$, so the $f$-length of $\alpha_2$ is bounded by $\epsilon ma^{(1-\epsilon)m}$.
%Similarly, the $f$-length of $\alpha_3$ is bounded by $\epsilon m a^{(1-\epsilon)m}$.
%Thus, the $f$-length of $\alpha$ is at most $2(C_1+\epsilon m)a^{(1-\epsilon)m}$ which is bounded in turn by $Ca^{(1-2\epsilon)m}$.
%\end{proof}

%\begin{prop}
%The Hausdorff dimension of $\Lambda$ is at most $h_a(\lambda)$.
%\end{prop}

%\begin{proof}
%Fix $h>h_a(\lambda)$ and let $\epsilon>0$ be such that $h_a(\lambda)+h(1-2\epsilon)>0$.
%Then, for any $n$,
%\begin{align*}
%\sum_{m\geq n}\sum_{g\in \mathcal{P}_m}\mathrm{diam}(\Pi(g,\epsilon m))^h&\leq C^h \sum_{m\geq n}H(\lambda)^m a^{h(1-2\epsilon)m}\\
%& \leq C^h\sum_{m\geq n}a^{(h_a(\lambda)+h(1-2\epsilon)m}.
%\end{align*}
%This last quantity tends to 0 as $n$ goes to infinity.
%\end{proof}

%The proof of the lower bound can be reproduced from \cite[Section~5.3]{SWX}.
%}

%{\color{blue}
%Sketch of proof that the limit sets of two independent BRW are almost surely disjoint, whenever $\lambda<\rho^{-1}$.

%Define the partial shadow $\Omega(x,R)$ (maybe we need a new name, since I use the Floyd distance, rather than transition points in the formulation) as follows.
%$$\Omega_\delta(x=\big \{y\in \Gamma\cup \partial \Gamma, \delta_x(e,y)\geq \delta\big \}.$$

%\begin{prop}
%Let $\Omega_\delta(x)$ be the partial shadow at $x$ of width $R$.
%Then,
%$$\mathbb{P}\big ( \Lambda \cap \Omega_\delta(x) \neq \emptyset \big) \lesssim G_\lambda(e,x).$$
%\end{prop}

%\begin{proof}
%First,
%$$\mathbb{P}\big ( \Lambda \cap \Omega_\delta(x) \neq \emptyset \big)\leq \mathbb{P}\big ( \exists \ \mathrm{infinitely}\ \mathrm{ many } \ u, X_u\in \Omega_{\delta/2}(x) \big ).$$
%Now, write $\Omega=\Omega_{\delta/2}(x)$ for simplicity.

%\begin{claim}
%We have
%$$\mathbb{P}\big ( \exists \ \mathrm{infinitely}\ \mathrm{ many } \ u, X_u\in \Omega_{\delta/2}(x) \big )\asymp \sum_{y_1,...y_n,...\in \Omega}G_\lambda(e,y_1;\Omega^c)...G_\lambda(y_k,y_{k+1};\Omega^c)...$$
%\end{claim}

%\begin{proof}
%This would be proved conditioning on successive passages through $\Omega$ and using an elaborated version of the many-to-one formula.
%\end{proof}
%{\color{cyan}
%  This argument only works for random walk. For branching random walk, we cannot get successive passages for $\Omega$, since it is possible that many particles may first visit $\Omega$ at the same generation. Also the next passage may come from particles outside of $\Omega$, not the ones already visit $\Omega$. So I think the claim is incorrect.

%  A possible strategy to estimate the probability $\mathbb{P} \left( \Lambda \cap \Omega_{\delta}(x) \neq \emptyset \right)$ is as follows. We run a BRW starting at $e$ and freeze all particles when they first visit the set $\Omega_{\delta}(x)$. (there may be particles go to infinity and do not visit $\Omega_{\delta}$ throughout their whole lifetime.) Let $\mathcal{Z}$ be all the particles frozen in $\Omega_{\delta}(x)$. By the superposition property (something like the Markov property) of BRW, if we run an independent BRW starting at each particle in $\mathcal{Z}$, and put all the particles together, we will get the original BRW. Note that a BRW starting at $y \in \Omega_{\delta}(x)$ may visit $\Omega$ infinity times with positive probability. (I think it should be true, but do not have a proof.) So the probability $\mathbb{P} \left( \Lambda \cap \Omega_{\delta}(x) \neq \emptyset \right)$ is more or less comparable to $\mathbb{P} (\mathcal{Z} \geq 1)$, which is bounded by $\mathbb{E} [\mathcal{Z}]$.

%  By the many-to-one formula (or \cite[Lemma 4.2]{CandelleroGilchMuller}),
%  \[ \mathbb{E}[\mathcal{Z}] = F_{\lambda}(e, \Omega_{\delta}) = \sum_n \lambda^n \mathbb{P} \left( Y_0 = e, Y_n \in \Omega_{\delta}, Y_k \notin \Omega_{\delta} \text{ for } 1 \leq k \leq n-1 \right), \]
%  which is the first visit Green function to $\Omega_{\delta}$. But I do not have idea how to prove that $F_{\lambda}(e, \Omega_{\delta}(x)) \leq c G_{\lambda}(e, x)$. 
%}


%Now, $G_\lambda(e,y_1;\Omega^c)\leq CG_\lambda(e,x)G_\lambda(x,y_1;\Omega^c)$.
%Similarly to the claim, we have that
%$$\sum_{y_1,...y_n,...\in \Omega}G_\lambda(x,y_1;\Omega^c)...G_\lambda(y_k,y_{k+1};\Omega^c)...$$
%can be expressed as a probability, so this sum is uniformly bounded.
%Thus, 
%$$\mathbb{P}\big ( \Lambda \cap \Omega_\delta(x) \neq \emptyset \big)\leq C'G_\lambda(e,x).$$
%\end{proof}



%}



\section{A lower bound for the Hausdorff dimension of the limit set}\label{SectionlowerboundHausdorff}
\subsection{The limit set in the Floyd and Bowditch boundaries}
For $r\leq \rho^{-1}$, we let $\Lambda_\mathcal{F}(r)$ and $\Lambda_\mathcal{B}(r)$ be the limit sets of the branching random walk whose offspring distribution has mean $r$ in the Floyd and Bowditch boundary respectively.
\begin{prop}\label{lowerboundHdim}
Let $r\leq \rho^{-1}$.
Almost surely,
$$\Hdim(\Lambda_\mathcal{F}(r),\delta_e)\geq \frac{-1}{\log \lambda}\omega_{\Gamma}(r)$$
and
$$\Hdim(\Lambda_\mathcal{B}(r),\overline{\delta}_e)\geq \frac{-1}{\log \lambda}\omega_{\Gamma}(r).$$
\end{prop}

By~(\ref{shortcutsmallerFloyd}), it is enough to prove the lower bound for the shortcut distance on the Bowditch boundary.
For simplicity, we write $\Lambda=\Lambda_\mathcal{B}(r)$ and $X=\Gamma\cup \partial_{\mathcal{B}}\Gamma$.

Let $h<\frac{-1}{\log \lambda}\omega_\Gamma(r)$.
We will prove that with positive probability, there exists a positive finite measure $\chi$ on the limit set $\Lambda$ such that
$$\int_{\Lambda}\int_{\Lambda}\overline{\delta}_e(x,y)^{-h}d\chi(x)d\chi(y)<+\infty.$$
Recall that $\mathcal{P}_{n,L}$ is the set of $L$-transitional points in $S_n$ that are ever visited by the branching random walk and $M_{n,L}$ is the cardinality of $\mathcal{P}_{n,L}$.
Using Lemma~\ref{omegaL}, we fix $L$ such that
\begin{equation}\label{choiceofL}
h<\frac{-1}{\log \lambda}\omega_{\Gamma,L}(r).
\end{equation}

Let $A_n$ be the event
$$\left \{M_{n,L}\geq \frac{1}{2}\mathbf E[M_{n,L}]\right \}.$$
By the Paley-Zygmund inequality and Corollary~\ref{GreenLmoment}, there exists $p>0$ such that
$\mathbf P(A_n)\geq p$.
Also, for any $C\geq 0$,
$\mathbf P(M_{n,L}\geq C\mathbf E[M_{n,L}])\leq 1/C$.
Therefore, for large enough $C$, the event $B_n$ defined by
$$\left \{\frac{1}{2}\mathbf E[M_{n,L}]\leq M_{n,L}\leq C\mathbf E[M_{n,L}]\right \}$$
satisfies $\mathbf P(B_n)\geq p/2$.
We define for every $n$ a random measure $\chi_n$ by
\begin{equation}\label{defchin}
  \chi_n = \mathbf 1_{B_n}\frac{1}{\mathbf{E} \left[ M_{n, L} \right]} \sum_{x \in \mathcal{P}_{n, L}}\mathbf D(x)
\end{equation}
where $\mathbf D(x)$ is the Dirac measure at $x$.
% By Lemma~\ref{RelThinTriangleLem2}, there exists a point $w_n$ which is within a bounded distance of transition points on $[e,X_n]$ and on $[e,X_n']$.
% We denote by $d_n$ the supremum of $d(e,w_n)$ for such a point $w_n$.

% Recall that $[e,X_n]$ and $[e,X_n']$ are $L$-transitional.
% A small adaptation of the proof of \cite[Proposition~5.13]{PY} yields the following result.
% \begin{lem}
% There exists $c_L$ such that
% $$\hat{\delta}_e(X_n,X_n')\geq c_L\lambda^{d_n}.$$
% \end{lem}



% \begin{lem}
%   \label{boundWn}
% The expectation $\mathbf{E}[\mathbb{1}_{A_n}W_n]$ is uniformly bounded.
% \end{lem}

% \begin{proof}
% There exists $C$ such that $d_n\leq n+C$.
% Hence,
% \begin{equation}\label{boundAnWn1}
% \begin{split}
% \mathbf{E}[\mathbb{1}_{A_n}W_n]&\leq c_L\mathbf E \left [ \mathbb 1 _{A_n}\mathbf{E} \left [ \lambda^{-h d_n}\big | \mathcal{G}\right ]\right ]\\
% &\leq c_L\sum_{k=0}^{n+C} \lambda^{-hk}\mathbb E \left [ \mathbb 1 _{A_n}\mathbf{P} \left ( d_n=k\big | \mathcal{G}\right )\right ].
% \end{split}
% \end{equation}
% Furthermore, if $d_n=k$, then $2n-k-C\leq d(X_n,X_n')\leq 2n-k+C$.
% Thus,
% \begin{equation}\label{boundAnWn2}
% \mathbf P \left (d_n=k\big | \mathcal{G}\right )\leq \frac{1}{M_{n,L}^2}\sum_{\underset{2n-k-C\leq d(x,y)\leq 2n-k+C}{x,y\in S_{n,L}}}\mathbb 1_{\{X_n=x,X_n'=y\}}.
% \end{equation}
% By Lemma~\ref{4.3SWX},
% \begin{align*}
% &\mathbf E \left [\sum_{\underset{2n-k-C\leq d(x,y)\leq 2n-k+C}{x,y\in S_{n,L}}}\mathbb 1_{\{X_n=x,X_n'=y\}}\right ]\\
% &\hspace{.3cm}\leq C \sum_{z\in \Gamma}\sum_{\underset{2n-k-C\leq d(x,y)\leq 2n-k+C}{x,y\in S_{n,L}}}G_r(e,z)G_r(z,x)G_r(z,y).
% \end{align*}
% Recall that $r<\rho^{-1}$.
% Applying Proposition~\ref{GreenL} and Lemma~\ref{Correlationxy}, we get
% \begin{equation}\label{boundAnWn3}
% \mathbf E \left [\sum_{\underset{2n-k-C\leq d(x,y)\leq 2n-k+C}{x,y\in S_{n,L}}}\mathbb 1_{\{X_n=x,X_n'=y\}}\right ]\leq C_L \mathrm{e}^{\omega_{\Gamma,L}(r)(2n-k)}.
% \end{equation}
% Finally, by definition of $A_n$,
% $$\mathbb 1_{A_n}\frac{1}{M_{n,L}^2}\leq \mathbb 1_{A_n}\frac{4}{\mathbf{E}\left [M_{n,L}\right ]^2}$$
% and by Proposition~\ref{Coro4.4L},
% \begin{equation}\label{boundAnWn4}
% \mathbb 1_{A_n}\frac{1}{M_{n,L}^2}\leq C_L\mathbb{1}_{A_n}\mathrm{e}^{-2n\omega_{\Gamma,L}(r)}.
% \end{equation}
% Combining~(\ref{boundAnWn1}),~(\ref{boundAnWn2}),~(\ref{boundAnWn3}) and~(\ref{boundAnWn4}), we get
% $$\mathbf{E}[\mathbb{1}_{A_n}W_n]\leq C_L\mathrm{e}^{-2n\omega_{\Gamma,L}(r)}\sum_{k=0}^{n+C}\lambda^{-hk}\mathrm{e}^{\omega_{\Gamma,L}(r)(2n-k)}\leq C'_L\sum_{k=0}^{n+C}(\lambda^{-h}\mathrm{e}^{-\omega_{\Gamma,L}(r)})^k.$$
% Since $h<\frac{-1}{\log \lambda}\omega_{\Gamma,L}(r)$, this last quantity is uniformly bounded.
% \end{proof}


Our goal is to apply a compactness theorem to the sequence $\chi_n$ and find a limit random measure $\chi$.
In \cite{Crauel}, the author proves a compactness criterion for random probability measures.
Here the measure $\chi_n$ is not almost surely a probability measure, however some of the results of \cite{Crauel} still hold in our context.
References and proofs are postponed to the Appendix.

Note that
$$\mathbf E[\chi_n(X)]\leq C,$$
so $\chi_n$ is a random finite measure in the sense of Definition~\ref{defrandomfinitemeasure}.
We define the measure $\pi_\Omega(\chi_n)$ on $\Omega$ by setting, for every event $A$,
$$\pi_\Omega(\chi_n)(A)=\mathbf E\left [\mathbf 1_A\chi_n(X)\right]=\mathbf E\left [\mathbf 1_A \mathbf 1_{B_n}\frac{M_{n,L}}{\mathbf E[M_{n,L}]}\right ].$$
Then, 
$$\pi_\Omega(\chi_n)(A)\leq C\mathbf P(A).$$
Since $X$ is compact, any subset of $\mathcal{M}_\Omega(X)$ is tight in the sense of Definition~\ref{deftightrandom}.
Moreover, any point of $\Gamma$ is isolated in $X$.
By Corollary~\ref{coroProkhorovisolated}, the closure of $\left \{\chi_n\right \}$ is compact for the weak topology on random finite measures.
This is the smallest topology such that for every random bounded continuous function, the map $\mu\mapsto \mu(f)$ is continuous, where a random bounded continuous function is a function
$f:(x,\omega)\mapsto f(x,\omega)$ such that for every $\omega$, $f(\cdot,\omega)$ is bounded continuous, for every $x$, $f(x,\cdot)$ is measurable and the map $\omega \mapsto \|f(\cdot,\omega)\|_\infty$ is $\mathbf P$-essentially bounded.

Thus, there exists a sub-net $(\chi_\alpha)_{\alpha\in A}$ such that $\chi_{\alpha}$ converges to some random finite measure $\chi$.
This means that there exists a directed set $A$ and a monotone final function $h:A\to \mathbb N$ such that for every $\alpha\in A$,
$\chi_\alpha=\chi_{h(\alpha)}$ and such that $\chi_\alpha$ eventually lies in every neighborhood of $\chi$.
We refer to \cite[Definition~2.11, Definition~2.15, Theorem~2.31]{AliprantisBorder} for more details on nets and a characterization of compactness in terms of convergent sub-nets.
This implies that for every random bounded continuous function $f$, we have
\begin{equation}\label{equationconvergencechi}
\chi_{\alpha}(f)\underset{\alpha\to \infty}{\longrightarrow}\chi(f).
\end{equation}

Since
$$\frac{p}{4}\leq \mathbf E[\chi_{\alpha}(X)]\leq C,$$
the same holds for $\chi$, applying~(\ref{equationconvergencechi}) to the function $\mathbf 1_\Omega\mathbf 1_X$.
Therefore, with positive probability, $\chi$ is not the null-measure.

Note that $\chi_{\alpha}(\mathcal{P}\cup \Lambda)=\chi_{\alpha}(X)$ and that $\mathcal{P}\cup \Lambda$ is a random closed set in the sense of Definition~\ref{defrandomclosedset}.
Thus, by Proposition~\ref{randomPortmanteau}, which is analogous to the classical Portmanteau theorem, we have
$$\chi(\mathcal{P}\cup \Lambda)\geq \limsup_{\alpha}\chi_{\alpha}(\mathcal{P}\cup \Lambda)=\limsup_\alpha\chi_{\alpha}(X)=\chi(X).$$
Also, the topology on $\Gamma \cup \partial_{\mathcal{B}}\Gamma$ extends the discrete topology of $\Gamma$.
Thus, any compact $K\subset \Gamma$ is closed and open, so the function $\mathbf 1_K$ is continuous.
Applying convergence to this function and noting that for large enough $n$,we have $\chi_{n}(K)=0$, hence for large enough $\alpha$, $\chi_\alpha(K)=0$ we get that
$\chi(K)=0$.
Since $\Gamma$ can be written as a countable union of compact sets, we get that
$$\mathbf E[\chi(\mathcal{P})]\leq \mathbf E[\chi(\Gamma)]=0.$$
Thus, $\chi$ almost surely gives full measure to $\Lambda$.



We can now finish the proof of Proposition~\ref{lowerboundHdim}.
\begin{proof}[Proof of Proposition~\ref{lowerboundHdim}]
First, as in the proof of Proposition~\ref{P:vgrBRW}, we can construct a branching random walk $(\Gamma,\mu,\nu')$ such that $\nu'$ has mean $r-\epsilon$, has finite second momment and is stochastically dominated by $\nu$.
As a consequence, we can assume that $r<\rho^{-1}$ and that $\nu$ has finite second moment.


We slightly modify the distance $\overline{\delta}_e$ and set for $x,y\in \Gamma \cup \partial_{\mathcal{B}}\Gamma$
$$\hat{\delta}_e(x,y)=\begin{cases}
\overline{\delta}_e(x,y) &\text{ if }x\neq y \\
\lambda^{d(e,x)} &\text{ if }x=y
\end{cases}$$
where by definition $\lambda^{\infty}=0$.
Note that for $x,y\in \partial_{\mathcal{B}}\Gamma$, $\hat{\delta}_e(x,y)=\overline{\delta}_e(x,y)$.


\begin{claim}
The function $(x,y)\in (\Gamma\cup \partial_\mathcal{B}\Gamma)\times (\Gamma\cup \partial_\mathcal{B}\Gamma)\mapsto \hat{\delta}_e(x,y)$ is continuous.
\end{claim}
\begin{proof}[Proof of the claim]
Let $x_n,y_n$ converge to $x,y$.
If $x\neq y$, then $x_n\neq y_n$, so $\hat{\delta}_e(x_n,y_n)=\overline{\delta}_e(x_n,y_n)$ for large enough $n$, which converges to $\overline{\delta}_e(x,y)=\hat{\delta}_e(x,y)$.

Now if $x=y$, there are two cases.
First, if $x\in \Gamma$, then $x_n=y_n=x$ and so $\hat{\delta}_e(x_n,y_n)=\hat{\delta}_e(x,y)$ for large enough $n$.

Second, assume that $x=y\in \partial_\mathcal{B}\Gamma$.
We have to prove that $\hat{\delta}_e(x_n,y_n)$ converges to 0 as $n$ tends to infinity.
Up to extracting sub-sequences, we can assume that either $x_n=y_n$ for every $n$ or that $x_n\neq y_n$ for every $n$.
In the former sub-case, $\hat{\delta}_e(x_n,y_n)=\lambda^{d(e,x_n)}$ which tends to 0 since $x_n$ tends to infinity.
In the latter, $\hat{\delta}_e(x_n,y_n)=\overline{\delta}_e(x_n,y_n)$, which concludes the proof.
\end{proof}
  Let
\begin{equation}\label{defWn}
\begin{split}
  W_n &= \int\int \hat{\delta}_e(x, y)^{-h} d \chi_n(x) d \chi_n(y)\\
  &= \mathbf 1_{B_n}\frac{1}{\mathbf{E} \left[ M_{n, L} \right]^2} \sum_{x, y \in \mathcal{P}_{n, L}} \hat{\delta}_e(x, y)^{-h}. 
  \end{split}
\end{equation}


\begin{claim}
  \label{boundWn}
The expectation $\mathbf{E}[W_n]$ is uniformly bounded.
\end{claim}

\begin{proof}[Proof of the claim]
  By Lemma~\ref{RelThinTriangleLem2}, for every $x$, $y \in \mathcal{P}_{n, L}$ there exists a point $w_n$ which is within a bounded distance of transition point on $[e, x]$ and on $[e, y]$.  We denote by $d_n$ the supremum of $d(e,w_n)$ for such a point $w_n$.  Note that $[e, x]$ and $[e, y]$ are $L$-transitional.  A small adaptation of the proof of \cite[Proposition~5.13]{PY} yields that
  \[
    \hat{\delta}_e(x, \, y) \geq c_L \lambda^{d_n}.  
  \]
  There exists $C$ such that $d_n\leq n+C$.  Furthermore, if $d_n = k$, then
  $$2 n - k - C \leq d(x, y) \leq 2n - k + C.$$
  By Lemma~\ref{4.3SWX},
  \begin{align*}
    & \mathbf{E} \left[ \sharp \left\{ (x, y) \in \mathcal{P}_{n, L} \colon d_n = k \right\} \right] \\
    & \hspace{.3cm} \leq
          \mathbf{E} \left[ \sum_{\substack{x, y \in S_{n, L} \\ 2n - k - C \leq d(x, y) \leq 2n - k + C}} \mathbb{1}_{\left\{ x, y \text{ are visited by }\mathrm{BRW}(\Gamma,\nu,\mu) \right\}} \right] \\
    & \hspace{.3cm} \leq 
           C \sum_{z \in \Gamma} \sum_{\substack{x, y \in S_{n, L} \\ 2n - k - C \leq d(x, y) \leq 2n - k + C}} G_r(e, z) G_r(z, x) G_r(z, y). 
  \end{align*}
Recall that $r<\rho^{-1}$.
Applying Proposition~\ref{GreenL} and Lemma~\ref{Correlationxy}, we get
\begin{equation}\label{boundAnWn3}
  \mathbf{E} \left[ \sharp \left\{ (x, y) \in \mathcal{P}_{n, L} \colon d_n = k \right\} \right] 
  \leq
  C_L \mathrm{e}^{\omega_{\Gamma,L}(r)(2n-k)}.
\end{equation}
Combining~\eqref{defWn}, (\ref{boundAnWn3}) and Proposition~\ref{GreenL}, we get
$$\mathbf{E}[W_n]\leq C_L\mathrm{e}^{-2n\omega_{\Gamma,L}(r)}\sum_{k=0}^{n+C}\lambda^{-hk}\mathrm{e}^{\omega_{\Gamma,L}(r)(2n-k)}\leq C'_L\sum_{k=0}^{n+C}(\lambda^{-h}\mathrm{e}^{-\omega_{\Gamma,L}(r)})^k.$$
By our choice of $L$~(\ref{choiceofL}), this last quantity is uniformly bounded.
\end{proof}

For any $\kappa\geq 0$, the function
$\kappa\wedge \hat{\delta}_e(x,y)^{-h}$
is bounded continuous on $X\times X$.
Thus,
$$\mathbf E\left [\int \int \kappa \wedge \hat{\delta}_e(x,y)^{-h} d\chi_{\alpha}(x)d\chi_{\alpha}(y)\right ]\underset{\alpha\to \infty}{\longrightarrow}\mathbf E\left [\int \int \kappa \wedge \hat{\delta}_e(x,y)^{-h} d\chi(x)d\chi(y)\right ].$$
By what precedes, we have that for every $\alpha$,
$$\mathbf E\left [\int \int \hat{\delta}_e(x,y)^{-h} d\chi_{\alpha}(x)d\chi_{\alpha}(y)\right ]\leq C$$ for some uniform $C$.
By the Fatou Lemma applied to the measure $\pi_X(\chi)$ on $X$ defined for every Borelian subset $B$ of $X$ by
$$\pi_X(B)=\mathbf E[\chi(B)],$$ we have
$$\mathbf E\left [\int \int \hat{\delta}_e(x,y)^{-h} d\chi(x)d\chi(y)\right ]\leq \liminf_{\kappa \to \infty}\mathbf E\left [\int \int \kappa \wedge \hat{\delta}_e(x,y)^{-h} d\chi(x)d\chi(y)\right ].$$
For every $\kappa$, for every $\alpha$,
$$\mathbf E\left [\int \int \kappa \wedge \hat{\delta}_e(x,y)^{-h} d\chi_{\alpha}(x)d\chi_{\alpha}(y)\right ]\leq \mathbf E\left [\int \int \hat{\delta}_e(x,y)^{-h} d\chi_{\alpha}(x)d\chi_{\alpha}(y)\right ]\leq C$$
and letting $\alpha$ go to infinity, we get
$$\mathbf E\left [\int \int \kappa \wedge \hat{\delta}_e(x,y)^{-h} d\chi(x)d\chi(y)\right ]\leq C.$$
This bound being uniform in $\kappa$, we finally get that
$$\mathbf E\left [\int \int \hat{\delta}_e(x,y)^{-h} d\chi(x)d\chi(y)\right ]\leq C.$$
Consequently, $\mathbf P$-almost surely,
$$\int \int \hat{\delta}_e(x,y)^{-h} d\chi(x)d\chi(y)<+\infty.$$

Recall that $\hat{\delta}_e=\overline{\delta}_e$ on $\partial_\mathcal{B}\Gamma$.
Thus, with positive probability, there exists a finite measure $\chi$ on $\Lambda$ which is not the null-measure and such that
$$\int \int \overline{\delta}_e(x,y)^{-h} d\chi(x)d\chi(y)<+\infty.$$
By Frostman Lemma for metrizable spaces (see \cite[Theorem~2.6]{Shah}), this shows that with positive probability, $h<\Hdim(\Lambda,\overline{\delta}_e)$.
The same argument as in \cite[Lemma~4.7]{SWX} shows that $ \Hdim(\Lambda,\overline{\delta}_e)$ is almost surely a constant.
This concludes the proof.
\end{proof}

\subsection{The limit set in the ends boundary}\label{sectionends}
We prove here Theorem~\ref{coroaccessiblegroups}.
We briefly introduce infinitely ended groups and refer to \cite[Section~4]{DussauleYang} and references therein for more details on those groups and on the link between random walks and the end boundary.
Let $(V,E)$ be a locally finite graph and let $F$ be a finite set of $V$.
We denote by $\mathcal{C}(F)$ an infinite connected component of the complement of $F$ in $(V,E)$.
An end $\xi$ of $(V,E)$ is a collection of infinite connected components $\mathcal{C}(F)$, where $F$ is finite, such that for any two such $F,F'$, we have that the intersection of $\mathcal{C}(F)$ and $\mathcal{C}(F')$ is infinite.
We will also say for simplicity that $\xi$ lies in the connected component $\mathcal{C}(F)$ if $\mathcal{C}(F)$ is part of the collection defining $\xi$.
We denote by $\partial_{\mathcal{E}}(V,E)$ the set of ends.
We can endow the end compactification $V\cup \partial_{\mathcal{E}}(V,E)$ with a topology which extends the discrete topology on $V$ such that $V\cup \partial_{\mathcal{E}}(V,E)$ is compact and $V$ is dense in its ends compactification.
If $\Gamma$ is a finitely generated group, we define its end boundary as the set of ends of a Cayley graph with respect to a finite generating system.
Its topology does not depend on the choice of the finite generating system.
We denote by $\partial_{\mathcal{E}}\Gamma$ the end boundary of $\Gamma$.

Let $0<\lambda<1$.
We define the visual distance $\tilde{\delta}_e$ of parameter $\lambda$ on $\Gamma \cup \partial_{\mathcal{E}}\Gamma$ by setting
$\tilde{\delta}_e(x,y)=\lambda^n$, where $n$ is the minimal integer such that $x$ and $y$ lie in two distinct connected components of the complement of $B(e,n)$.
It is well known that the end boundary is covered by the Floyd boundary, see for instance \cite[Proposition~11.1]{GGPY} or \cite{Karlssonfreesubgroups}.
Moreover, we can be more precise and by \cite[Lemma~4.3]{DussauleYang}, the identity of $\Gamma$ extends to an 1-Lipschitz continuous and equivariant map $\psi$ from the Floyd compactification to the end compactification, so in particular
$$\delta_e(x,y)\geq \tilde{\delta}_e(\psi(x),\psi(y)).$$

If $\Gamma$ is a group with infinitely many ends, we denote by $\Lambda_{\mathcal{E}}(r)$ the limit set of a branching random walk $(\Gamma,\nu,\mu)$ with $\mathbf E[\nu]=r$.
We prove the following.
\begin{prop}\label{lowerboundends}
Let $r\leq \rho^{-1}$.
Almost surely,
$$\Hdim(\Lambda_\mathcal{E}(r),\tilde{\delta}_e)\geq \frac{-1}{\log \lambda}\omega_{\Gamma}(r).$$
\end{prop}


By a celebrated result of Stallings \cite{Stallings}, a group with infinitely many ends $\Gamma$ splits as an HNN extension $A_{*C}$ or an amalgamated product $A*_{C}B$, where $C$ is a finite group.
The action on the corresponding Bass-Serre tree satisfies the conditions of \cite[Definition~2]{Bowditch} and so $\Gamma$ is relatively hyperbolic.
However, there is no clear relation between the shortcut distance on the Bowditch boundary and the visual distance on the end boundary.
Indeed, the ends boundary is in general larger than the Bowditch boundary and even if they coincide, the shortcut distance is the largest distance on the Bowditch boundary satisfying~(\ref{shortcutsmallerFloyd}), so it is bounded from below by the visual distance.
Thus, we cannot deduce Proposition~\ref{lowerboundends} from Proposition~\ref{lowerboundHdim}.

\begin{proof}[Proof of Proposition~\ref{lowerboundends}]
We follow the same strategy as for the Bowditch boundary.
Since $\Gamma$ has infinitely many ends, it is relatively hyperbolic.
It is either an HNN extension $A*_{C}$ or an amalgamated product $A*_{C}B$, where $C$ is finite.
In the former case, the parabolic subgroups are the conjugates of $A$ and one can choose $\mathbb P_0=  \{A\}$, in the latter case, they are the conjugates of $A$ and $B$ and one can choose $\mathbb P_0=\{A,B\}$.
In both cases, every element of $\Gamma$ can be written uniquely in a normal form, see \cite[(9.2),(9.4)]{Woessends}.
For simplicity, we only give details of the proof when $\Gamma=A*_{C}B$.
The case of an HNN extension is treated similarly.
In this situation, the normal form is described as follows.
We choose a set of representatives $\mathfrak{A}$ of $A/C$ and a set of representatives $\mathfrak{B}$ of $B/C$.
Then, any element $x$ of $\Gamma$ can be uniquely written as
\begin{equation}\label{normalform}
x=a_1b_1...a_nb_nc,
\end{equation}
where $a_i\in \mathfrak{A}$, $b_i\in \mathfrak{B}$, $c\in C$.
Moreover, any path from $e$ to $x$ in the Cayley graph of $\Gamma$ has to pass within a bounded distance of every prefix of $x$ in the normal form~(\ref{normalform}).
This follows from the fact that $A*_CB$ is quasi-isometric to the space $X$ obtained by taking copies of $A$ and $B$ for each coset $gA$ and $hB$, $g,h\in \Gamma$ and connecting
\begin{itemize}
    \item $ga\in gA$ to $ga\in gaB$ by adding $\sharp C$ edges between the coset $gaC$ in $gA$ and the coset $gaC$ in $gaB$.
    \item $gb\in gB$ to $gb\in gbA$ by adding $\sharp C$ edges between the coset $gbC$ in $gB$ and the coset $gbC$ in $gbA$.
\end{itemize}
The construction of the space $X$ is performed in \cite{ScottWall}.
It is similar to the construction of a tree of spaces modeling the free product $A*B$ obtained by adding one single edge between every element $ga\in gA$ and $ga\in gaB$ and one single edge between every element $gb\in gB$ and $gb\in gbA$.
This space is also used in \cite{PapasogluWhyte} to prove that $A*B$ is quasi-isometric to $A*_CB$.

In particular, if a geodesic $[e,x]$ is $L$-transitional in the sense of Definition~\ref{deftransitional}, then $|e,x]$ cannot travel long in parabolic subgroups and thus every word $a_i$ and $b_i$ in the normal form of $x$~(\ref{normalform}) satisfies $|a_i|\leq D_L$ and $|b_i|\leq D_L$, where $D_L$ only depends on $L$.


Let $h<\frac{-1}{\log \lambda}\omega_\Gamma(r)$.
Using Lemma~\ref{omegaL}, we fix $L$ such that $h<\frac{-1}{\log \lambda}\omega_{\Gamma,L}(r)$.
The sequence of random finite measures $\chi_n$ on $\Gamma$ defined by~(\ref{defchin}) converges, up to a sub-net, to a random finite measure $\chi$ on the end boundary that gives full measure to $\Lambda_{\mathcal{E}}(r)$ and which is not the null measure with positive probability.
We slightly modify the distance $\tilde{\delta}_e$ by setting $\hat{\delta}_e(x,y)=\tilde{\delta}_e(x,y)$ if $x\neq y$ and $\hat{\delta}_e(x,y)=\lambda^{|x|}$ if $x=y$ and we define
\begin{align*}
  W_n &= \int\int \hat{\delta}_e(x, y)^{-h} d \chi_n(x) d \chi_n(y)\\
  &= \mathbf 1_{B_n}\frac{1}{\mathbf{E} \left[ M_{n, L} \right]^2} \sum_{x, y \in \mathcal{P}_{n, L}} \hat{\delta}_e(x, y)^{-h}. 
\end{align*}
For $x,y\in \mathcal{P}_{n,L}$, we set $d_n$ to be the maximal length of a common prefix of $x$ and $y$ in their normal form.
\begin{claim}
There exists $c_L$ only depending on $L$ such that $\hat{\delta}_e(x,y)\geq c_L\lambda^{d_n}$.
\end{claim}
\begin{proof}[Proof of the Claim]
Let $w_n$ be a common prefix of $x$ and $y$ of length $d_n$ and write
$$x=w_nx_1...x_mc,\ y=w_ny_1...y_lc'$$
where $x_i,y_i$ are either in $\mathfrak{A}$ or $\mathfrak{B}$ and $c,c'\in C$.
Then, any path from $x$ to $y$ has to pass within a bounded distance of $w_nx_1$.
Since $|x_1|\leq D_L$ and $|w_n|=d_n$, we see that $x$ and $y$ lie in distinct components of $B(e,d_n+D_L+C)$.
Thus, $\hat{\delta}_e(x,y)\geq \lambda^{d_n+D_L+C}$.
\end{proof}
Using this claim, we prove as above that $\mathbf E[W_n]$ is uniformly bounded, which allows us to prove that
$$\int \int \tilde{\delta}_e(x,y)^{-h}d\chi(x)d\chi(y)$$
is almost surely finite.
We then use the Frostman Lemma to conclude.
\end{proof}
\begin{rem}
Without involving \cite{PapasogluWhyte}, an alternative proof uses the bottleneck property  introduced in \cite{DussauleYang}.
Indeed, with $r, F$ be given in \cite[Lemma~5.4]{DussauleYang} replacing Lemma~\ref{ExtensionLem}, any two elements $g, h$ can be concatenated via some $f\in F$ such that any path from $e$ to $gfh$ intersects $B(g,r)$.
Such points are referred to as bottleneck points on $[e, gfh]$.
If a path contains a sequence of bottleneck points with consecutive distance at most $L$, then it is said to have the $L$-bottleneck property.
A triangle with sides having the $L$-bottleneck property satisfies the conclusion of Lemma~\ref{RelThinTriangleLem2} and Lemma~\ref{existenceofcenter}.
%\begin{comment}
 %\ywy{Since [e,x] has L-bottleneck property, the path [e,y][y,x] intersects the $r$-nbhd of every bottlenecks of [e,x]. Let p be the last bottleneck that [e,y] intersect, so [e,y] does not intersect the others after p. But the next bottleneck q after p on [e,x] has to intersect [x,y]. As $d(p,q)\le L$, the $L$-center for the triangle is found. 
%As the argument shows, we only need one side is L-bottlenecking, the other two sides can be any paths.Maybe I miss something?}
%\end{comment}.
We define similarly  $S_{n,L}$ to be the set of points $x \in S_n$ so that $[e, x]$ has the $L$-bottleneck property.
Then  Proposition~\ref{P:vgrBRW} follows verbatim the same argument  with transition point replaced with bottleneck points.
The remaining modification goes as explained above.
\end{rem}

%{{\color{blue}Okay I'm convinced. I was indeed missing something, sorry about the confusion and thank you for your explanation. I was thinking it was good to give a clarified situation for the relatively hyperbolic structure in terms of normal forms, but I agree this is a less suited argument and it is too bad we have to separate HNN extensions and amalgamated products. Anyway, I'm happy I clarified this for myself and I you want to replace my proof with your argument, I completely agree. } \ywy{It is great you also agree with it. We can of course have two proofs for the same result in a collaboration. I suggest keep yours which seems more independent to the readers, and also mine as it is just very short remark:)}

%{\color{blue}I chose to keep the same formulation as for the Bowditch boundary.
%Alternatively, we can introduce the bottleneck property defined in \cite{DussauleYang} : an element $y$ is an $L$-bottleneck between $e$ and $x$ if any path from $e$ to $x$ has to pass within $L$ of $y$.
%A geodesic is said to have the $L$-bottleneck property if every element on this geodesic is an $L$-bottleneck between the endpoints of the geodesic.
%Then, we introduce $\mathcal{P}_{n,L}$ as the set of $x\in S_n$ ever visited by the BRW such that $[e,x]$ has the $L$-bottleneck property.
%We use the extension lemma given by \cite[Lemma~5.4]{DussauleYang} and the fact that if $y$ is an $L$-bottleneck between $e$ and $x$, then the Floyd distance $\delta_y(e,x)$ is bounded from below by a constant that only depends on $L$ to reprove that the growth rate $\omega_{\Gamma,L}(r)$ converges to $\omega_\Gamma(r)$ as $L$ tends to infinity.
%We also need to prove that for every $x,y$, there exists a point which is an $L$-bottleneck for both $e,x$ and $e,y$ and an $L'$-bottleneck between $x$ and $y$ i.e.\ the existence of a center for the triangle $e,x,y$ mimicking Lemma~\ref{RelThinTriangleLem2} and Lemma~\ref{existenceofcenter}.\begin{comment}
%As recalled above, any path from $e$ to $x$ has to intersects a fixed neighborhood of a bottleneck point and the bottleneck points are uniformly $L$-spaced, I think all the statements are obviously true.   
%\end{comment}
%Then, we use directly the notion of bottleneck rather than transition points to prove Proposition~\ref{lowerboundends}.
%This option is longer, we have to introduce many definitions and I don't know how to prove the existence of a center without introducing the normal form of elements.  \begin{comment}
%I think all the defintions are already introduced as you described. It is not long in my opinion (6 lines), and does not involve the others quasi-isometry rigidity result and seperating free product and HNN. 
%\end{comment}I think it is more readable as I wrote it above, but if you find it unsatisfying, we can rewrite the whole proof.
%If both options seem bad to you, I vote we give up and don't speak about infinitely ended groups but only about free products.}
%\ywy{
%You argument is nice. Let us take it. I added a remark to point the other possibility I prefer.
%}

A particular class of groups with infinitely many ends are free products of the form $\Gamma=\Gamma_0*\Gamma_1$ where at least one the free factors $\Gamma_i$ is not $\mathbb Z/2\mathbb Z$.
For such groups, the authors of \cite{CandelleroGilchMuller} prove that the Hausdorff dimension of the limit set in the end boundary endowed with a visual distance is exactly $\frac{-1}{\log \lambda}\omega_{\Gamma}(r)$.
Their proof for the upper bound applies to any group with infinitely many ends.
Indeed, it consists in saying that for any end $\xi\in \Lambda_{\mathcal{E}}(r)$, for every $n$, the branching random walk has to visit some point $x\in \bigcup_{0\leq l\leq l_0}S_{n+l}$ such that $\xi$ and $x$ are in the same infinite connected component of the complement of $B(e,n)$.
The constant $l_0$ only depends on the support of the measure $\mu$.
Thus, $\Lambda_{\mathcal{E}}(r)$ can be covered with $\mathcal{P}_n$ sets of diameter bounded by $C\lambda^n$.

For the lower bound, they first show that
$\Hdim (\Lambda_{\mathcal{E}}(r),\tilde{\delta}_e)$ is bounded from below by some number $z^*$ and then prove that $z^*=\omega_\Gamma(r)$, see \cite[Lemma~4.7]{CandelleroGilchMuller}.
However, in order to prove that $z^*=\omega_\Gamma(r)$, they use the same invalid argument as the one described in Section~\ref{sectionexamples}, namely that the quantity $H_r(n)$ is sub-multiplicative.
Proposition~\ref{lowerboundends} fills their gap and
combining it with the proof for the upper bound described above, we get Theorem~\ref{coroaccessiblegroups}.

\begin{rem}\label{CGMCor}
The proof of the upper bound described above also works for the parabolic cosets.
Namely, for every $P\in \mathbb P$, the Hausdorff dimension of $\Lambda_{\mathcal{E}}(r)\cap \partial_{\mathcal{E}}P$ is bounded from above by $\omega_{P}(r)$.
Thus, assuming further that the Green function has a parabolic gap, we recover \cite[Corollary~3.7]{CandelleroGilchMuller}, i.e.\
$$\Hdim(\Lambda_{\mathcal{E}}(r)\cap \partial_{\mathcal{E}}P)<\Hdim (\Lambda_{\mathcal{E}}(r)).$$
\end{rem}

\section{An upper bound for the Hausdorff dimension of the limit set}\label{SectionupperboundHausdorff}
Let $\Gamma$ be a  relatively hyperbolic group and let $\Lambda_\mathcal{F}(r)$ be the limit set of the branching random walk in the Floyd boundary of $\Gamma$, endowed with the Floyd distance $\delta_e$.
Recall that by Theorem~\ref{mapGerasimov}, there exists a map $\phi$ from the Floyd boundary to the Bowditch boundary such that the preimage of a conical limit point is reduced to a single point.
We can thus see the set of conical limit points $\partial_{\mathcal{B}}^{con}\Gamma$ in the Bowditch boundary as a subset of the Floyd boundary.
We set
$$\Lambda^{con}_{\mathcal{F}}(r)=\Lambda_{\mathcal{F}}(r)\cap \phi^{-1}(\partial_{\mathcal{B}}^{con}\Gamma).$$
We prove here the following proposition. %\begin{comment}
    %A little issue is that conical points in Floyd boundary contains potentially  more points than those in Bowditch boundary. This extra conical points are contained in the limit set of parabolic subgroups, so it could not be argued by transition points. But I believe there is a work around to bypass this issue. 
%\end{comment}
%{\color{blue}I'm answering your comment. Yes you're right, sorry my terminology is confusing. I had forgotten that conical limit point already meant something for limit points in the Floyd boundary. I actually meant the preimage of conical limit points of the Bowditch boundary. Either we write the proposition for the preimage of conical limit points of the Bowditch boundary and there is nothing to change, or we adapt a bit the proof for all conical limit points in the Floyd boundary, but I fear this will be more difficult, because we make crucial use of the thin triangle property in the current proof.}

\begin{prop}\label{upperboundHausdorff}
Let $r\leq \rho^{-1}$.
Almost surely, %\begin{comment}
%    I would suggest to give a notation for conical points in the limit set: perhaps $\Lambda^c_{\mathcal F}$ (may be confuses with complement), $\Lambda^{Rad}_F$, or just state the result for limit set in Bowditch boundary. I do not think it deserves to make a distinction in Floyd boundary case.
%\end{comment}
$$\Hdim(\Lambda^{con}_{\mathcal{F}}(r),\delta_e)\leq \frac{-1}{\log \lambda}\omega_{\Gamma}(r).$$
\end{prop}

Under additional assumption on the volume growth of parabolic subgroup, we have the upper bound on the full limit set in Floyd boundary.
\begin{cor}\label{corowholeFloyd}
Let $r\leq \rho^{-1}$. Assume that $v_S(P)\le \omega_{\Gamma}(r)$ for every parabolic subgroup $P$. Then
almost surely,  
$$\Hdim(\Lambda_{\mathcal{F}}(r),\delta_e)\leq \frac{-1}{\log \lambda}\omega_{\Gamma}(r).$$
\end{cor}
\begin{proof}
The Bowditch boundary consists of conical points and countably many parabolic points and the preimage of each parabolic point is exactly the limit set of a parabolic subgroup (see \cite[Theorem A]{GePoJEMS}).
So the limit set $\Lambda_{\mathcal{F}}(r)$   is contained in the  union of $\Lambda^{con}_{\mathcal{F}}(r)$ with countably many limit sets of parabolic cosets $P\in\mathbb P$.
As $P$ is quasi-convex  in  the Cayley graph of $\Gamma$, any geodesic from $e$ to the limit points of $P$ is contained in a fixed neighborhood of $P$, see for instance \cite[Lemma~4.3]{DrutuSapir}.
Using a suited covering by shadows based at $S_n\cap P$, we can see that $\Hdim(\Lambda_{\mathcal{F}} P)\le \frac{-1}{\log \lambda}v_S(P)$ (e.g. \cite[Lemma 4.1]{PY}).
The conclusion now follows from Proposition \ref{upperboundHausdorff}.   
\end{proof}

To prepare the proof for Proposition \ref{upperboundHausdorff}, we first need a few geometric lemmas.
Denote as usual by $\mathbb P=\{gP: g\in \Gamma, P\in\mathbb P_0\}$ the collection of all parabolic cosets, and by $\ell(\gamma)$ the length of a path $\gamma$.
Fix a geodesic $[x,z]$ and $\epsilon_1,\epsilon_2\in [0,1]$. An \textit{$[\epsilon_1,\epsilon_2]$-percentage} of $[x,z]$ consists of points $w\in [x,z]$ such that $\epsilon_1 \le d(x,w)/d(x,z)\le \epsilon_2$.


\begin{lem}\label{DeepContainedLem}
Let $\gamma$ be a geodesic segment such that $[\epsilon, 1-\epsilon]$-percentage of $\gamma$ contains no $(\eta, L)$-transition point. Then there exists a unique $P\in\mathbb P$  such that the entry and exit points of $\gamma$ in $N_\eta(P)$ have distance at most $\epsilon \ell(\gamma)$ to the corresponding endpoints of $\gamma$.  
\end{lem}
\begin{proof}
By assumption, the middle point  $m\in \gamma$ is $(\eta, L)$-deep in a unique $P\in \mathbb P$. Let $u, v$ be the corresponding entry and exit point of $\gamma$ in $N_\eta(P)$. By Lemma \ref{EntryTransLem}, $u, v$ are $(\eta, L)$-transition points, so $d(x,u)\leq \epsilon d(x,y)$ and $d(v, y)\le \epsilon d(x,y)$, which concludes the proof.   
\end{proof}

 
%The following technical lemma will be used in Lemma \ref{LogTrackingLem}. It follows directly from the bounded intersection of $\mathbb P$ and we leave the proof to the interested reader. 
%\begin{lem}\label{SameParaCosetLem}
%For any $B>0$ there exists $L=L(B)>0$ with the following property.
%Let $\alpha, \beta$ be two geodesics so that $\alpha_-=\beta_-$ and $d(\alpha_+,\beta_+)\le B$. Assume that their   $[\epsilon, 1-\epsilon]$-percentage of $\alpha$ and $\beta$ are contained $N_\eta(P)$ and $N_\eta(P')$ respectively.  If $(1-2\epsilon)\ell(\alpha)>L$, then $P=P'$.      
%\end{lem}



Fix  $C>0$ and $x\in S_n$. The \textit{$C$-{\color{blue}}partial cone} $\Omega(x,C)$ consists  of points $z\in G$ such that $[e,z]$ contains an $(\eta, L)$-transition point $C$-close to $x$. %\begin{comment}
%The partial cone is a bit complicated here (already $U(x)$ is difficult to understand at the first glance). I prefer to use the standard cone which is easier to understand for the readers. 
%{\color{blue}I totally agree, but I think that we need to use the partial cone, see my remark below.}
%\end{comment}

%Similarly, for given a peripheral coset $P\in \mathbb P$, let $\Omega(P,r)$ be the $r$-cone at $P$   consisting  of points $z\in G$ such that $[e,z]\cap N_r(P)\ne\emptyset$. 
%Then$$
%\Omega(P,r)\subseteq \bigcup_{x\in P\cap S_n} \Omega(x, r)
%$$ where $n\ge d(e, P)$.

Let $B([x,z])$ be the ball centered at the middle point of $[x,z]$  of radius ${d(x,z)/2}$. Define $U(x)$ to be the union of the partial cone $\Omega(x, C)$ and the balls $B([x,z])$ for all geodesics $[x,z]$ between $x$ and $z\in \Omega(x,C)$. That is,
$$
U(x):= \Omega(x,C) \cup \left( \bigcup\left\{B([x,z]): \forall [x,z], \forall z\in  \Omega(x,C)\right\} \right). 
$$
Note that any ball of centered at $w\in[x,z]$ of radius $\min\{d(w,x), d(w,z)\}$ is contained in $B([x,z])$.

 In what follows, let $C>0$ be given by Lemma \ref{ThinTransitional2}. 
\begin{lem}\label{AvoidBallLem}
Let $\alpha$ be a path starting from $e$ and first entering  at a point $z\in U(x)$. Let $ w\in [x,z]$ be a transition point.    Set $S:=\min\{d(w,x), d( w, z)\}$. Then $B(w, S-2C)$ is contained in $U(x)$ so $\alpha$ has distance at least $S-2C$ to   the point $w$.     
\end{lem}
\begin{proof}
We first consider the case $z\in \Omega(x,C)$. By definition of $U(x)$,  any ball centered at $w\in [x,z]$ of radius $S=\min\{d(x,w), d(z,w)\}$ is contained in $ U(x)$. The statement follows immediately.  

Assume now that $z$  lies in a ball $B([x,\hat z])$ where $[x,\hat z]$ is a geodesic between $x$ and  some $\hat z$ in $\Omega(x,C)$. Let  $\hat w\in [x,\hat z]$    be the middle point.

Consider the triangle with vertices $x,\hat w,z$. As $w$ is a transition point on $[x,z]$, Lemma \ref{ThinTransitional2}  shows that $d(w, w')\le C$ for some $w'\in [\hat w,x]\cup[\hat w, z]$. Thus, $B(w, \kappa-2C)\subset B(w', \kappa-C)$. As $w'$ is on the radius $[\hat w,x]$ or $[\hat w, z]$ of the ball $B([x,\hat z])$,   we have the ball $B(w', S-C)$ is contained in $B(\hat w, \eta)\subseteq U(x)$. 
\end{proof}

 
\begin{lem}\label{ConvConicalLem}
Let $\xi\in \partial_{\mathcal B} \Gamma$ be a conical point and consider a sequence of points $z_n\to \xi$.
Let $x\in [e,\xi]$ be an $(\eta, L)$-transition point. Then   for for all  but finitely many $z_n$,   there exists an $(\eta,L)$-transition point $x_n$ on $[e, z_n]$ such that $d(x_n, x)\le C$.  In particular,  $z_n\in \Omega(x, C)$ for large enough $n$.

%Let $x\in [e, \xi]$ be a deep point in $P\in\mathbb P$. Then except for finitely many points, the sequence $z_n$ is   contained in $\Omega(P, r)$.
\end{lem}
\begin{proof}
If $z_n\to \xi$ then $\delta_x(z_n, \xi) < \lambda^C \delta$ for all large enough $n$, where $\delta=\delta(\eta,L)$ is given by Lemma~\ref{TransLargeLem}.
Applying Lemma~\ref{ThinTransitional2} to the triangle with vertices $e$, $\xi$, $z_n$, there is an $(\eta, L)$-transition point $x_n$ on $[z_n, \xi]$ or $[e, z_n]$ such that $d(x, x_n) \leq C$. 
It suffices to prove that $x_n \in [e, z_n]$. 
Let $y$ be any transition point on $[z_n, \xi]$. 
Then by Lemma~\ref{TransLargeLem}, $\delta_y(z_n, \xi) > \delta$.  It follows from~(\ref{changeofbasepointFloyd}) that $\lambda^C \delta > \delta_x(z_n, \xi) \geq \lambda^{d(x, y)} \delta_y(z_n, \xi)>\lambda^{d(x,y)}\delta$ and hence $d(x, y) > C$. The conclusion follows.
%so by Lemma~\ref{Karlssonlem}, we obtain $d(x, [e,z_n])\le C$. 
%The conclusion follows.
\end{proof}
%{\color{blue}Karlsson's lemma does not give a transition point. {\color{red}\begin{comment}
%Yes, you are right. But  Floyd large distance is indeed a transition point, see the second part of Lemma \ref{TransLargeLem} I just added.{\color{blue}Okay thanks I didn't know the converse was already proved in the cited papers. However, Longmin's proof is short and explicit, and we do not need to change $L$, so maybe it is better to keep it, what do you think ? \lw{I suggest to keep the original proof and it seems more natural. }}
%\end{comment}} But the lemma still holds with $x_n$ a transition point (which is important, see my remark below).
%Indeed, if $z_n$ converges to $\xi$, then $z_n$ converges to $\xi$ in the relative Cayley graph.
%Thus, the distance from $x$ to the relative geodesic $[z_n,\xi]$ goes to infinity, and so for large enough $n$, $x$ is far from a transition point on the Cayley graph geodesic $[z_n,\xi]$.
%Thus, by the relative thin triangle property, $x$ is within $C$ of a transition point on $[e,z_n]$.
%}

%\lw{Is the following proof correct?  If $z_n \to \xi$ then $\delta_x(z_n, \xi) < \lambda^C \delta$ for all large enough $n$, where $\delta=\delta(\eta,L)$ is given by Lemma~\ref{TransLargeLem}.
%Applying Lemma~\ref{ThinTransitional2} to the triangle with vertices $e$, $\xi$, $z_n$, there is an $(\eta, L)$-transition point $x_n$ on $[z_n, \xi]$ or $[e, z_n]$ such that $d(x, x_n) \leq C$. 
%It suffices to prove that $x_n \in [e, z_n]$. 
%Let $y$ be any transition point on $[z_n, \xi]$. 
%Then by Lemma~\ref{TransLargeLem}, $\delta_y(z_n, \xi) > \delta$.  It follows from~(\ref{changeofbasepointFloyd}) that $\lambda^C \delta > \delta_x(z_n, \xi) \geq \lambda^{d(x, y)} \delta_y(z_n, \xi)$ and hence $d(x, y) > C$. The conclusion follows. }
%{\color{blue}Yes this is great.}


Fix $\epsilon\in (0, 1/2)$.
Let $U_\epsilon(x)$ be the set of points $z\in U(x)$ such that $[x,z]$ contains a transition point $w$ being at distance at least $\epsilon d(x,z)$ to one of the endpoints:  $$ \max\{d( w, x), d(w,z)\}\ge \epsilon d(x,z).$$  

By  Lemma \ref{EntryTransLem}, the set $U(x)\setminus U_\epsilon(x)$
consists of points $z\in U(x)$ such that the $[\epsilon, 1-\epsilon]$-percentage of $[x,z]$ is contained in the $\eta$-neighborhood of a unique peripheral coset $P\in\mathbb P$. Explicitly, there exists a subsegment of $[x,z]$ with length at least $(1-2\epsilon)d(x,z)$ contained in $N_\eta(P)$.    

% By triangle inequality, we see that any (connected) subpath is also $\tau$-taut.
\begin{lem}\label{AvoidTransBall}
Let $\alpha$ be a path starting from $e$ and first entering $U(x)$ at a point $z\in U_\epsilon(x)$. 
Assume that $\epsilon d(x,z)>10C$. Then there exists a  transition point $y$   on $[e, z]$ such that $\alpha$ lies outside the ball around $y$ with radius $\epsilon d(x,z)-3C$.
\end{lem}
 
\begin{proof}
By definition of $z\in U_\epsilon(x)$, $[x,z]$ contains a transition point $w$ such that
$$\epsilon d(x,z)\le \max\{d(w,x),  d( w, z)\}.$$
Setting $S=\epsilon d(x,z)$,   Lemma \ref{AvoidBallLem} implies that $B(w,\kappa-2C)\subset U(x)$, so $\alpha$ does not intersect $B(w,S-2C)$. To conclude the
proof, it remains to find a $(\eta,L)$-transition point $y\in [e,z]$ such that $d(w,y)\le C$. In particular,  $\alpha$ does not intersect $B(y,S-3C)$, completing the proof.

Indeed, as in the proof of Lemma \ref{AvoidBallLem}, $z$ lies on the ball $B([x,\hat z])$ centered at  the middle point  $\hat w$  of $[x,\hat z]$ for some $\hat z\in \Omega(x,C)$. 
By assumption,   $w$ is  an $(\eta,L)$-transition point on $[x,z]$.   Lemma \ref{ThinTransitional2} applied for   the  triangle with vertices $x,\hat w,z$ shows that $d(w,w')\le C$ for some $w'\in [\hat w,x]\cup[\hat w,z]$. A schematic figure is shown below. 

Similarly  for   the triangle with vertices $e,x,z$,    Lemma \ref{ThinTransitional2}  implies that either $[e,z]$ or $[e,x]$ contains an $(\eta,L)$-transition point $y$ such that $d(w,y)\le C$ and then $d(y,w')\le 2C$.   As $\hat z\in \Omega(x,C)$, $d(x,[e,\hat z])\le C$ holds, so the triangle inequality shows
$$
d(e,x)+d(x,\hat z) -2C \le d(e,\hat z). 
$$
In a different term, this implies that the  path $p:=[e,x][x,\hat z]$ is   \textit{$(2C)$-taut}: $\ell(p)\le d(p_-,p_+)+2C$.
It follows from triangle inequality  that any subpath of a $(2C)$-taut path $p$ is $(2C)$-taut.

We first claim that either $w'\notin [x,\hat w]$ or $y\notin [e,x]$. Otherwise, we have  $w'\in [x,\hat w]$ and $y\in [e,x]$. Since $[y,x][x,w']$ is a $(2C)$-taut subpath of $p$,  we obtain that $d(y,x)+d(x,w')\le d(y,w')+2C\le 4C$, so $d(x,w)\le d(x,w')+d(w',w)\le 5C$. This is a contradiction because  $d(x,w)\ge \epsilon d(x,z)> 5C$.
\begin{figure}[htb] % You can use the htb to tell latex: place here if possible, if not try top of page, if not try bottom of page.
 
%\centering \scalebox{0.6}{
\includegraphics[width=0.8\linewidth]{Cone.pdf} % Do not put file extension, only name
%}
\label{figure1}
\end{figure}

Let us assume now  $y\in [e,x]$ to derive a contradiction.
The above claim shows   $w'\notin [x,\hat w]$ and then $w'\in [\hat w,z]$.  Using the $(2C)$-taut path   $[y,x][x,\hat w]$, we have $$d(y, \hat w) +2C \ge d(y,x)+d(x, \hat w).$$ which combined with          $d(x,\hat w)=d( z,\hat w)\ge d(w', \hat w)$ gives $$d(y, \hat w) +2C \ge d(y,x)+d(w', \hat w).$$
On the other hand,   $|d(y,\hat w)-d(w',\hat w)|\le d(y,w') \le 2C$.
These together show that $d(y, x)\le 4C$ and then $d(w',x)\le d(w',y)+d(y,x)\le 6C$.
Thus, $|d(x,\hat w)-d( w',\hat w)|\le 6C$.
Since $d(x,\hat w)=d( z,\hat w)$ and $w'$ is on $[\hat{w},z]$, we have $d(w',z)\le 6C$, so   $d(x,z)\le 12C$. This contradicts the assumption $\epsilon d(x,z)>10C$, hence we proved that $y\in [e,z]$ is impossible, so $y$ is  the desired transition point on $[e,z]$. 
\end{proof}

%{\color{blue}
%Okay I proofread this lemma and I agree.
%But I had trouble understanding why a sub-path of a taut path is also taut.
%I find the definition too vague.
%For example if $\alpha$ is any path, not a concatenation of geodesics and satisfies that for some point $x$ on $\alpha$, we have $d(\alpha_-,x)+d(x\alpha_+)\leq 2r+d(\alpha_-,\alpha_+)$, then there is no reason why a sub-path $\beta$ of $\alpha$ satisfies also this property for $\beta_-$ and $\beta_+$.
%And even if $\alpha$ is the concatenation of two geodesics $[a,b]$ and $[b,c]$ which is taut at the point $b$, I'm not sure that a sub-path of $\alpha$ is taut at another point than $b$.
%However, what is true is the following : if $[a,b][b,c]$ satisfies $d(a,b)+d(b,c)\leq d(a,c)+r$, and if $x$ is on $[a,b]$ and $y$ is on $[b,c]$, then $d(x,b)+d(b,y)\leq d(x,y)+r$.
%You just need to rewrite $d(a,b)=d(a,x)+d(x,b)$ and $d(b,c)=d(b,y)+d(y,c)$ and to bound $d(a,c)$ by $d(a,x)+d(x,y)+d(y,c)$ to conclude.
%Maybe we can either give a more precise definition and a more precise claim or we can omit to introduce this terminology and just prove the desired properties for the paths $[y,x]|x,w']$ and $[y,x][x,\hat w]$ above.
%As you wish.} \ywy{I think this follows just triangle inequality and definition of length of a curve. The term of taut path is due to Bowditch in his GGT notes. Indeed, by defintion  we have $\ell(\alpha)\le d(\alpha_-,\alpha_+)+\tau\le d(\alpha_-,\beta_-)+d(\beta_+,\alpha_+)+d(\beta_-,\beta_+)+\tau$. By defintion of curve length, $\ell(\alpha)\ge \ell(\beta)+d(\alpha_-,\beta_-)+d(\beta_+,\alpha_+).$ So $\ell(\beta)\le  d(\beta_-,\beta_+)+\tau$}
%{\color{blue}Okay I didn't know it was standard terminology. Thank you for the explanation.
%I still think it deserves either to be properly introduced (and then proved that the union of two geodesics $[e,x][x,\hat z]$ as above is taut) or not introduced at all (and we just prove the desired properties for the two paths we use). Indeed this is the only place where we use the notion of length of a curve, so it deserves some details.}\ywy{That is OKay, this is not standard terms, I think ( I only see Bowditch use it). Yes, I agree. We should avoid this term in the proof.}

For any $m\ge 1$, let $U_{\epsilon}(x,m)$ be the set of elements  $z\in U_{\epsilon}(x)$ such that $d(x,z)\ge m$.

\begin{lem}\label{TraceAvoidTransBall}
For any $\epsilon\in (0,1/2)$, there exists $\kappa>0$ with the following property.
Almost  surely,   there exists $n_0>0$ such that  for all $n>n_0$ and all $x\in S_n$: if $\mathrm{BRW}(\Gamma,\nu,\mu)$ first enters $U(x)$ at a point $z$, then   $z\in U(x) \setminus U_{\epsilon}(x,\kappa \log |x|)$.  %the   trace     of BRW intersects $U_\epsilon(x)$, then it does not intersect the subset $U_{\epsilon}(x,\kappa \log |x|)$. 



\end{lem}
 
\begin{proof}
Let us freeze all particles of the branching random walk %which      first enters $U(x)$ at a point $z\in U_{\epsilon}(x,m)$. Denote     $\mathcal Z({x,m})$  the   collection of  frozen particles $z$  so that  
when they enter $U(x)$ at the first time.  Denote by $\mathcal{Z}(x, m)$ the collection of frozen particles $z \in U_{\varepsilon}(x, m)$.  Then for $z \in \mathcal{Z}(x, m)$ we have 
\begin{enumerate}
    \item 
     $d(x,z)\ge m$,  
    \item 
     $\max\{d(y,x), d(y,z)\}>\epsilon d(x,z)$ where $y$ is an $(\eta,L)$-transition point on $[e,z]$ given by Lemma~\ref{AvoidTransBall}.  
\end{enumerate}    
Note that  the genealogy path from $e$ to $z$ does not intersect $B(y,\epsilon d(x,z))$.  

Let $\delta=\delta(\eta, L)$ be given by Lemma \ref{SupExpDecayLem}. Then     the expected number of  particles frozen%\begin{comment}
%It should be upper bounded by $F_{r}(e,z; [U_\epsilon(x)]^c)$, right?   Is it true that $F_{r}(e,z; [U_\epsilon(x)]^c)\le G_{r}(e,z; [U_\epsilon(x)]^c)$?   
% {\color{blue}Yes, because $F_r$ count all trajectories from $e$ to $z$ never visiting $z$ before the endpoint and never visiting $U_\epsilon$, while $G_r$ counts all trajectories from $e$ to $z$ never visiting $U_\epsilon$.}
% \ywy{In general, it only follows that $F_{r}(e,z; [U_\epsilon(x)]^c)G_{r}(z,z; [U_\epsilon(x)]^c)= G_{r}(e,z; [U_\epsilon(x)]^c)$. The quantity $G_{r}(z,z; [U_\epsilon(x)]^c)$ could be very small.   Here that is because $z$ is already contained in $U_\epsilon(x)$, so $F_r=G_r$, right? Now I'm understood.}
% {\color{blue}Well yes here both quantities agree, you're right, but in general $G_r(z,z;U_\epsilon^c)$ cannot be very small. You have the trivial path of length 0 from $z$ to $z$, so $G_r(z,z;U_\epsilon^c)\geq 1$.
% And here, since $z\in U_\epsilon$, $G_r(z,z;U_\epsilon^c)=1$, so both quntities agree indeed.}
%\end{comment} 
at $z\in U_\epsilon(x)$   is upper bounded by
$$
G_{r}(e,z; [U_\epsilon(x)]^c) \le G_{r}(e,z; [B(y,\epsilon d(x,z))]^c) \le e^{-   e^{\delta \epsilon d(x,z)}}. 
$$  

By Lemma~\ref{purelyexponentialgrowth} there exists $c>0$ such that $\sharp S_n\le c e^{vn}$ for any $n\ge 1$. Thus,  there exist $\epsilon_1=\epsilon_1(\epsilon,\delta,v)$ and $m_0>0$ such that for any $m>m_0$, we have
$$
\mathbf E[\sharp \mathcal Z({x,m})]\le \sum_{z\in U_\epsilon(x,m)} G_{r}\bigl (e,z; [B(y,\epsilon d(x,z)]^c)\bigr) \le \sum_{k=m}^{\infty} ce^{kv} e^{-   e^{\delta\epsilon k}} \le e^{-e^{\epsilon_1 m}}.
$$
 
%Consider the events   
%$$
%A_n=\{ \text{BRW visits } z\in U_{\epsilon}(x, \kappa \log n) \text{ of first type for some } x\in S_n \}
%$$
%and 
%$$
%B_n=\{ \text{BRW visits } z\in U_{\epsilon}(x, \kappa \log n) \text{ of second type for some } x\in S_n \}
%$$
% Let $A_n$  be the   event that if  BRW visits $U_x$ for some $x\in S_n$, then the first killed particle is in $U_{\epsilon}(x,m)$. 

 
% Choose $\kappa$ so that $ \kappa \epsilon_1 >1$. Setting   $m=\kappa \log n$, we have
Choose $\kappa$ so that $ \kappa \epsilon_1 >1$ and let $m=\kappa \log n$.  Consider the event
$$A_n = \left\{ \mathcal{Z}(x, m) \ge 1 \, \text{for some } x \in S_n \right\},$$ i.e. \ $A_n$ is the   event such that if  the branching random walk visits $U_x$ for some $x\in S_n$, then the first frozen particle is in $U_{\epsilon}(x,m)$.  Then  
$$
\mathbf P(A_n) \le \sum_{x\in S_n}\mathbf P(\sharp  \mathcal Z({x,m})\ge 1) \le\sum_{x\in S_n} \mathbf E[\sharp \mathcal Z({x,m})] \le c e^{vn-e^{\epsilon_1 \kappa\log n}} \le c e^{vn-n^{\epsilon_1\kappa}}. 
$$ 
Therefore,  $\sum_{n=1}^{\infty} \mathbf P(A_n)  <\infty$ and so
the conclusion follows from the Borel-Cantelli Lemma.      
\end{proof}





Similarly, we prove the following.
\begin{lem}\label{TraceAvoidTransBall2}
For every $K\geq 0$ and $C\geq 0$, there exists $\kappa>0$ such that the following holds.
Almost surely, for all but finitely many $x$, if the branching random walk ever visits a point $y$ with $d(e,y)\leq Kd(e,x)$ and such that $x$ is within $C$ of a transition point on $[e,y]$, then it first enters $B(x,\kappa \log |x|)$.
\end{lem}

\begin{proof}
Given $x\in S_n$, we consider the set $U(x)$ of $y\in\Gamma$ so that $|y|\le K|x|$ and $x$ is within $C$ of a transition point on $[e,y]$. We freeze particles when they first reach a point $y\in U(x)$, without entering $B(x,\kappa \log |x|)$.
We denote by $\mathcal{Z}_n$ the set of frozen particles for all $x\in S_n$.
By Lemma \ref{SupExpDecayLem}, for $n$ large enough,    the expected number of  particles frozen at $y\in U(x)$   is upper bounded by
$$
G_{r}(e,y; [U(x)]^c) \le G_{r}(e,y; [B(x,\kappa \log |x|)]^c) \le e^{-   n^{\delta \kappa }}. 
$$  

By Lemma~\ref{purelyexponentialgrowth}, there exist $c>0$ such that $\sharp S_n\le c e^{vn}$ and $\sharp U(x)\le c e^{vKn}$ for any $n\ge 1$ and $x\in \Gamma$. Thus,   we have
$$
\mathbf E[\sharp \mathcal Z_n]\le \sum_{x\in S_n}  \sum_{y\in U(x)} G_{r}(e,y; [U(x)]^c) \le  c e^{v(K+1)n} e^{-   n^{\delta\kappa}}.   
$$
If $\kappa$ is chosen large enough, then $\sum_{n=1}^\infty \mathbf E[\sharp \mathcal Z_n]<\infty$, so the proof follows from the  Borel-Cantelli Lemma.
\end{proof}
%\begin{comment}
%{\color{red}Regarding my above remark before Prop 6.1, to get control of conical points in parabolic subgroup, we need something replacing this lemma.}
%Yes you're right, but see my answer to the above comment.
%\end{comment}
Recall that if $P\in \mathbb P$ is a parabolic coset, $\eta\geq 0$ and $x\in \Gamma$, we denote by $$\pi_{N_\eta(P)}(x) :=\{y\in N_\eta(P): d(x,y)=d(x, N_\eta(P))\}$$ the set of its shortest projections on the $\eta$-neighborhood $N_\eta(P)$ of $P$. 
Also, for $x,y\in \Gamma$, we denote by $$d_{N_\eta(P)}(x,y):=\diam{\pi_{N_\eta(P)}(x)\cup\pi_{N_\eta(P)}(y)}$$
and \cite[Corollary~8.2]{Hruska} the shortest projection is coarsely Lipschitz,
$$d_{N_\eta(P)}(x,y)\le kd(x,y)+k$$ for a fixed $k\ge 1$ depending only on $\eta$. Thus, $\pi_{N_\eta(P)}(x)$ has  bounded diameter. 





\begin{lem}\label{ProjectionLem}\cite[Lemma~1.15]{Sistoprojections}
For every large enough $\eta$, there exists $C=C(\eta)>0$ such that
for every $x\in \Gamma$, $P\in \mathbb P$ and for every geodesic $\gamma$ starting at $x$ and entering $N_\eta(P)$, we have
$$
\pi_{N_\eta(P)}(x) \subset B(y,C) 
$$
where $y$ is the entrance point of $\gamma$ in $N_\eta(P)$.
\end{lem}

%Since projections on parabolic cosets are well defined up to a bounded distance, projections on their $\eta$-neighborhood also are well defined up to a bounded distance depending on $\eta$.
%We arbitrarily choose one projection $\pi_{N_\eta(P)}(x)$ for every $x$.
%Also, for two points $x,y\in \Gamma$, we denote by $d_{P}(x,y)$ the diameter of the union   $\pi_{P}(x)\cup \pi_{P}(y)$.

\begin{lem}\label{exponentialmoments}
There exists $\eta_0\geq 0$ such that for $\eta\geq \eta_0$, the following holds.
Almost surely, for all but finitely many parabolic cosets $P\in \mathbb P$,  if the branching random walk ever visits a point $z$ satisfying 
both that $d_{N_\eta(P)}(e,z)\geq d(e,N_\eta(P))$ and $d(z,N_\eta(P))\leq  d_{N_\eta(P)}(e,z)$, then it first needs to enter $N_\eta(P)$ at a point $w$ such that $d_{N_\eta(P)}(e,w)\leq d(e,N_\eta(P))$.

In particular, if the branching random walk ever enters $N_\eta(P)$, then the first entrance point must be within $d(e,N_\eta(P))$ of the projection of $e$ on $N_\eta(P)$.
\end{lem}


%\begin{comment}
% {\color{red}Here, it seems you claimed the statement for any $\eta$ and $\eta$ is independent of $k$? }
% Yes I think so, $k$ is just the notation for $d(e,gP)$ and I'm saying that since the sum on $k$ converges, there are finitely many $k$ such that the event $A_k$ occurs which implies there are finitely many cosets $gP$ such that the statement in the lemma occurs. I clarified a bit the writing.
% \ywy{I see your point. Here I would suggest to isolate your estimates from \cite{DG21} as independent lemma which is crucial here and seems of interests elsewhere.}
%\end{comment}

\begin{proof}
Given $P\in\mathbb P$, we freeze particles when they first visit some point $z\in \Gamma$ with
$$d_{N_\eta(P)}(e,z)\geq d(e,N_\eta(P))$$ and
$$d(z,N_\eta(P))\leq  d_{N_\eta(P)}(e,z),$$ without having entered $\{w\in N_\eta(P): d_{N_\eta(P)}(e,w)\leq d(e,N_\eta(P))\}$. Denote  by $\mathbb P_k$  the collection  of   parabolic cosets $P$  with $d(e,N_\eta(P))=k$.
We denote by $\mathcal{Z}_k$ the set of such frozen particles for those $P\in \mathbb P_k$.
Then,
$$\mathbf E[\sharp \mathcal{Z}_k]\leq \sum_{P\in \mathbb P_k}\sum_{\underset{d_{N_\eta(P)}(x,e)\geq k}{x\in N_\eta(P)} }\sum_{\underset{d(x,y)\leq k}{y \colon \pi_{N_\eta(P)}(y)=x}}G_R(e,y;N_\eta(P)^c).$$%\begin{comment}
 %   there is no $K$, right?
%    {\color{blue} No there is no $K$ anymore.
%    I first had stated the lemma with a $K$, but this was because I hadn't noticed that in the below proof, up to choosing $\epsilon$ small enough, we could show that $d(z,N_\eta(P))\leq  d_{N_\eta(P)}(e,z)$.
%    I thought we could only prove that $d(z,N_\eta(P))\leq  K d_{N_\eta(P)}(e,z)$ for some $K$.
%    I find it simpler to state the lemma without $K$.} OK, so I removed the $K$ in the formual.
%\end{comment}
%\ywy{I would rewrite as follows: Then,
%$$\mathbf E[\sharp \mathcal{Z}_k]\leq \sum_{P\in \mathbb P_k}\sum_{z\in Z(P)}G_R(e,z;N_\eta(P)^c),$$
%where  $Z(P)$ denotes the set of $z\in \Gamma$ whose projection to   $ {N_\eta(P)}$  contains a point $x\in {N_\eta(P)}$ so that $d(x,z)\le k$ and $d_{N_\eta(P)}(e,z)\geq k$. }
As $x\in \pi_{N_\eta(P)}(z)$,  there exists a trajectory for the $\mu$-random walk from $z$ to $x$ that stays outside $N_\eta(P)$ of linear length in $d(x,z)$, since the support of $\mu$ is finite.
In particular,
$$G_R(z,x;N_\eta(P)^c)\geq \mathrm{e}^{-\alpha d(x,z)}\ge \mathrm{e}^{-\alpha k}$$ for some positive $\alpha$.
Note that by Lemma~\ref{purelyexponentialgrowth}, for a fixed $x$, the set of such elements $z$ being contained in a ball of radius $k$ grows as an exponential function in $k$.  Thus,
\begin{equation}\label{proofexponentialmoments}
    \mathbf E[\sharp \mathcal{Z}_k]\leq \sum_{P\in\mathbb P_k}\sum_{\underset{d(x,\pi_{N_\eta(P)}(e))\geq k}{x\in N_\eta(P)} }G_R(e,x;N_\eta(P)^c)\mathrm{e}^{\beta k},
\end{equation}
where $\beta$ depends both on $\alpha$ and on the growth rate $v$ of the word distance.
%The term $G_R(e,x;N_\eta(P)^c)$ is exactly the first entrance transition kernel in $N_\eta(P)$, starting at $e$.
By  Lemma~\ref{lemmaexponentialmoments}, for every $M\geq 0$, there exists $\eta_0$ such that for all $\eta\geq \eta_0$,
$$G_R(e,x;N_\eta(P)^c)\leq C \mathrm{e}^{-Md(x,\pi_{\eta,P}(e))}.$$
%\begin{comment}{\color{red}
%It would be helpful to the reader if you can precise the equality as a additive constant or multiplicative constant. (I think it can be additive, right?)  Again, I would suggest to move the nice result as a lemma. It would be a bit techical in a proof, in my opinion.}
%{\color{blue} No actually, the constant is multiplicative, because you project points on the trajectory on $P$ and you want the projection between two successive particles to be between $a/2$ and $2a$ of each other.
%About stating the result in \cite{DG21} as a separate lemma, I don't know, since it is only used here.} Okay, you decided it.
%\end{comment}
Choosing $M>v+\beta$, where $v$ is the growth rate of the word distance, we have by Lemma~\ref{purelyexponentialgrowth} and by~(\ref{proofexponentialmoments})
that for $\eta\geq \eta_0$,
$$\mathbf E[\sharp \mathcal{Z}_k]\leq C\mathrm{e}^{-k(M-v-\beta)}.$$
%We denote by $A_k$ the event $\{\sharp \mathcal{Z}_k\geq 1\}$.
%Then,
%$$\sum_{k\geq 0}\mathbf P(A_k)\leq  \sum_{k\geq 0}\mathbf E[\mathcal{Z}_k]\leq \sum_{k\geq 0}\mathrm{e}^{-(M-v-\beta)k}.$$
By the choice of $M$, the sum $\sum_{k\geq 0}\mathbf E[\sharp \mathcal{Z}_k]$is finite.
The result again follows from the Borel-Cantelli lemma.
%that almost surely, there are only finitely many $k$ such that $A_k$ occurs.
\end{proof}


%We now consider a tight augmentation of the trace of the branching random walk.
%\begin{defn}\label{TightAugDef2}
%Let $X\subset \Gamma$ and $X_n=X\cap S_n$. Denote by ${\rm Tran}_\epsilon(X_n)$ the union of transition points $x$ on  $[e,z]$ for some $z\in X_n$ and $|x|\le \epsilon n$. Let $\widetilde{X}_\epsilon$ be the union of all $X_n$ with $ {\rm Tran}_\epsilon(X_n)$ over $n\ge 1$. 
%\end{defn}

\begin{lem}\label{LogTrackingLem}
There exists $\kappa$ such that almost surely,   for any conical limit  point $\xi$ in $\Lambda$
and {for all but finitely many transition point $x$ on $[e,\xi]$, $\mathcal{P}$ intersects $B(x,\kappa \log |x|)$.}
\end{lem}
\begin{proof}
Let  $\xi\in \Lambda$  be a conical point. Then $[e,\xi]$ contains infinitely many $(\eta, L)$-transition points (see  \cite[Lemma 2.20]{YangPS}). Consider any  $(\eta, L)$-transition point $x\in [e,\xi]$ so that $|x|>n_0$. According to  Lemma \ref{ConvConicalLem}, for $\Omega(x,C)\subset U(x)$, the branching random walk must enter $  U(x)$.  

Set $\epsilon\in (0,1/2)$ so that $K=\epsilon^{-1}\ge 4$.  Let $\kappa$ be given by Lemmas \ref{TraceAvoidTransBall} and \ref{TraceAvoidTransBall2}.  Up to enlarging $\eta$, we may assume that it is big enough to apply Lemma~\ref{exponentialmoments}.

Let $z\in U_\epsilon(x)\setminus U_{\epsilon}(x,\kappa \log |x|)$ be the first entrance point by Lemma \ref{TraceAvoidTransBall}, so one of  the following statements is true:
\begin{enumerate}
    \item 
      $d(z, x)\le  \kappa \log |x|$,  
    \item
    the $[\epsilon, 1-\epsilon]$-percentage  of $[x,z]$ does not contain any $(\eta,L)$-transition point.  
\end{enumerate}

 

If the case (1) happens, then  we are done. We now assume     $d(x,z)>\kappa \log |x|$. 
%This will concludes the proof, as  $x\in \widetilde {\mathcal P}_\epsilon$ and $\xi\in \Omega(x,C)$.  

By Lemma \ref{DeepContainedLem}, there exist a unique   coset $P\in \mathbb P$ such that if $y_1,y_2$ are the entrance and exit points of  $[x,z]$ in $N_\eta(P)$, then
$$\max\{d(x,y_1), d(y_2,z)\}\le \epsilon d(x,z),$$
so
$$d(y_1,y_2)\geq (1-2\epsilon)d(x,z)\ge (1-2\epsilon) \kappa \log n_0.$$
By definition,   $z$ is contained in a ball $B([x,\hat z])$ centered at $\hat w$ for some $\hat z\in \Omega(x,C)$.
%Consider the geodesic triangle $\Delta(x\hat w\hat z)$ where $d(x,\hat w)=d(\hat z,\hat w)$. By Lemma \ref{ThinTransitional2}, we have $$\diam{[x,\hat w]\cap N_\eta(P)}>(1-2\epsilon)d(x,z)/2-4\eta.$$



By definition of $\Omega(x,C)$, there exists a transition point $\hat{x}$ on $[e,\hat{z}]$ such that $d(x,\hat x)\le C$.
By Lemma~\ref{ThinTransitional2} for the triangle with vertices $e,z,\hat{z}$, we see that $\hat{x}$ is within $C$ and so $x$ is within $2C$ of a transition point on $[e,z]$.
According to Lemma~\ref{TraceAvoidTransBall2}, if the branching random walk does not enter $B(x, \kappa\log |x|)$, then we have $ |z|> K |x|=\epsilon^{-1}|x|$.
Noting as above that $x$ is within $2C$ of a transition point on $[e,z]$,     Lemma~\ref{ProjectionLem} implies that $\pi_{N_\eta(P)}(e)$ is within a bounded   distance   of the entry point $y_1$ of $[x,z]$ into $N_\eta(P)$,  which implies that
$d(y_1,y_2) \leq d_{N_\eta(P)}(e,z)+C'$ for some constant $C'$ depending on $C$.
Moreover, if $n_0$ is large enough, then $d(y_1,y_2)\geq C'/2$, so we get
$d(y_1,y_2)\leq 2d_{N_\eta(P)}(e,z)$, hence
$$d(z,N_\eta(P))\leq \frac{2\epsilon}{1-2\epsilon}d_{N_\eta(P)}(e,z).$$
Furthermore, $d(e,N_\eta(P))\leq |x|+d(x,y_1)$.
Since $|x|\leq \epsilon|z|$ and $d(x,y_1)\leq \epsilon d(x,z)$, we get
$$d(e,N_\eta(P))\leq \left (\frac{\epsilon}{1-\epsilon}+\epsilon\right )d(x,z)\leq \left (\frac{\epsilon}{1-\epsilon}+\epsilon\right )(1-2\epsilon)2d_{N_\eta(P)}(e,z).$$

Thus, if $\epsilon$ is small enough, the conditions of Lemma~\ref{exponentialmoments} holds, hence the branching random walk first enters $N_\eta(P)$ at a point $w$ such that $d(\pi_{N_\eta(P)}(e),w)\leq d(e,N_\eta(P))$.
%\begin{comment}
%    Here, again I think you need specify the choice of $\eta$ as somewhere which should be large.
%
%    And $gP$ should be $P$, right?
%\end{comment}
Since $x$ is within $2C$ of a transition point on $[e,w]$, $d(e,N_\eta(P))\geq |x|+d(x,N_\eta(P))-2C$
and so
$$d(x,w)\leq d(x, N_\eta(P))+d_{N_\eta(P)}(e,w)\leq 2 d(e, {N_\eta(P)})   +2C\leq 2d(e,w)+2C.$$Thus, for $n_0$ large enough, we have $d(e,w)\ge d(e,x)-2C\ge  C$ so $d(x,w)\le 4d(e,w)$. Applying   Lemma~\ref{TraceAvoidTransBall2} with $K\ge 4$ again, we see that the branching random walk necessarily enters $B(x,\kappa \log |x|)$.
This concludes the proof.
\end{proof}  
%\ywy{
%    I agree with your nice argument. Here is a suggestion of the way to conclude: your   Lemma~\ref{TraceAvoidTransBall2} corresponds to the event $B$ in   Lemma~\ref{TraceAvoidTransBall}. So I would suggest two ways to conclude: 
    
%1. Keep Lemma~\ref{TraceAvoidTransBall} as it is. And state   Lemma~\ref{exponentialmoments} in another way as follows:  "Almost surely, for all but finitely many parabolic cosets $P\in \mathbb P$,  if the branching random walk ever visits a point $z$ satisfying $d_P(e,z)> K \eta$, then it first needs to enter the $\eta$-neighborhood $N_\eta(P)$ at a point $w$ satisfying $d_P(e, w)\le K\eta$.". Then continue the above arguments.

%2. Split the event B from Lemma~\ref{TraceAvoidTransBall} (i.e. restore to the old version) and merge into Lemma~\ref{exponentialmoments}. Now I think this is a better option. 
%}
%{\color{blue}I agree. Feel free to edit as you wish !
%I think I also cleaned the use of $\eta$ in the statements and in the proofs;} \ywy{I'm done with the edits of these proofs. Feel free to do any changes.}
%Let $y$ be the entry point of $[x, z]$ in $N_\eta(P)$. Let $\hat y$ be the entry point of $[x, \hat z]$ in $N_\eta(P)$. Then $d(y,\hat y)\le 4\eta$.

%As $|S_n|\prec e^{vn}$, the following lemma is straightforward.\begin{comment}
%{\color{blue}We do not need to state this as a lemma anymore. It is used inside the proof of Proposition~\ref{upperboundHausdorff}.} Okay, I see. But I think    Lemma \ref{ComputeHdim} gives a more transparent formulation for the computation of Hdim. It is much more simple and  easy to follow. \ywy{I think here the same growth rate of tight augmentation contains a gap: even we added the elements vert short which takes very small proportion, but their length is also very small, so their shadows may very large. So I think the above fix still cannot give a desired covering of the limit sets. }   
%\end{comment}
%\begin{lem}\label{TightAugmentLem}
%There exists $\epsilon>0$ such that the following holds.  Then almost surely,      the tight augmentation $\widetilde{\mathcal P}_\epsilon$ of the trace of BRW has the same growth rate as $\mathcal P$.   
%\end{lem}
%\begin{proof}
%Almost surely, there exists $n_0>0$ such that for any $n>n_0$ and any $z\in \mathcal P_n$,  any transition point on    $[e,z]$ is $\epsilon n$-close to $\mathcal P$. Indeed, by Lemma \ref{SupExpDecayLem}, this relative version follows the same   proof for the inequality (5.2) in \cite{SWX}. For the convienence of the reader, we re-produce the proof.

%Let $A_n$ be the event such that for some $z\in \mathcal P_n$, the trace   $\mathcal P$  is  $\epsilon n$-close to some transition point in  $[e,z]$. Let $\delta$ be given by Lemma \ref{SupExpDecayLem}. For any $n\gg 0$, we have 
%\begin{align*}
%\mathbf P(A_n) &\le \sum_{z\in S_n} \sum_{x\in {\rm Tran}([e,z])} \mathbf{P}(z\in \mathcal P_n \text{ and } d(y,\mathcal P)\le \epsilon n)\\  
%&\le  \sum_{z\in S_n}  \sum_{x\in {\rm Tran}([e,z])}  G_r(e, z; [B(x,\epsilon n)]^c)\\
%& \le c n e^{vn} e^{-e^{\delta \epsilon n}}
%\end{align*}
%where $c$ is a positive constant. The above claim follows by  Borel-Cantelli Lemma.

%Now by definition of ${\rm Hull}_\epsilon(\mathcal P_n)$, any transition point $x\in [e,z]$ for $z\in \mathcal P_n$ is $\epsilon n$-close to a point in $\mathcal P$. 
%By Proposition \ref{P:vgrBRW}, the growth rate of $\mathcal P$ is almost surely a positive constant say $v_0>0$.   Thus, $|{\rm Hull}_\epsilon(\mathcal P_n)| \le \sum_{0\le m\le 2\epsilon n}|\mathcal P_{m}| \le c e^{v_0 3\epsilon n}$.
%It suffices to recall that $v(X)$ is the convergence radius of the series $$
%\Theta_s(\mathcal P) = \sum_{x\in X} e^{-s|x|} = \sum_{n\ge 0} |\mathcal P_n| e^{-sn}
%$$    
%\end{proof}

%As $\widetilde{\mathcal P}\setminus \mathcal P$ is growth tight, the following corollary is proved in the above proof.
%\begin{cor}\label{LogTrackingCor}
%Let $(\epsilon,\kappa)$ satisfy      Lemmas \ref{TraceAvoidTransBall} and \ref{TightAugmentLem}. Then almost surely,   for any conical limit  point $\xi\in\Lambda$ with an exception  subset of strictly smaller Hausdorff dimension\begin{comment}{\color{blue}
%I don't think that we can say that the exceptional subset has smaller Hausdorff dimension, since we do not know that the Hausdorff dimension of this set is bounded by the residual growth rate} 
%\end{comment},    the trace of BRW intersects the ball $B(x,2\kappa \log |x|)$  for all but finitely many transition points $x$ on $[e,\xi]$.
%\end{cor}
%\begin{proof}
%By Lemma \ref{LogTrackingLem},      
%\end{proof}


We can now end the proof of the upper-bound. 
 
 \begin{proof}[Proof of Proposition~\ref{upperboundHausdorff}]
Let $r\leq \rho^{-1}$ and fix $h$ such that
 $$h>\frac{\omega_\Gamma(r)}{-\log \lambda}.$$
 Let $0<\epsilon<1/2$.
 Let $\xi\in \Lambda^{con}_{\mathcal{F}}(r)$ and let $x$ be a transition point on $[e,\xi]$.
 By Lemma~\ref{LogTrackingLem}, almost surely, there exists $n_0$ such that if $|x|\geq n_0$, we can find $z\in \mathcal{P}$ such that $z\in B(x,\kappa \log |x|)$.
 In particular, $|x|\leq |z|+\kappa \log |x|$ and if $|x|$ is large enough, then $|z|\geq (1-\epsilon)|x|$, hence
 $$x\in B\left (z,\kappa \log \frac{|z|}{1-\epsilon}\right ).$$
 Consequently, for every $m$,
 $$\Lambda^{con}_{\mathcal{F}}(r)\subset \bigcup_{n\geq m}\bigcup_{z\in \mathcal{P}_n}\Pi\left (z,\kappa \log \frac{|z|}{1-\epsilon}\right ).$$
By Lemma~\ref{diametershadow}, the diameter of $\Pi\left (z,\kappa \log \frac{|z|}{1-\epsilon}\right )$ is bounded by $C\lambda^{|z|}|z|^\alpha \log |z|$.
Thus,
$$\sum_{n\geq m}\sum_{z\in \mathcal{P}_n}\mathbf{diam}\left (\Pi\left (z,\kappa \log \frac{|z|}{1-\epsilon}\right )\right )^h\leq C \sum_{n\geq m}M_n\lambda^{hn}n^{h\alpha} (\log n)^h.$$
Since $\limsup \frac{1}{n}\log M_n\leq \omega_{\Gamma}(r)$, by the choice of $h$, this last quantity converges to 0 as $m$ tends to infinity.
%Since $[e,\xi]$ contains infinitely many transition points, $[x|$ can be chosen arbitrarily large, hence $n$ can be chosen arbitrarily large.
%This proves that for every $m$, there almost surely exists $n\geq m$ such that $\xi\in \Pi\left (z,2\kappa \log \frac{|z|}{1-\epsilon}\right )$ for some $z\in \widetilde{(\mathcal{P}_n)}_\epsilon$ satisfying further $|z|\geq m(1-\epsilon)$.
%We set $U_n(z)=\Pi\left (z,2\kappa \log \frac{|z|}{1-\epsilon}\right )$ to simplify notations.
%In other words,
%$$\Lambda \cap \{\text{ conical limit points}\}\subset \bigcup_{n\geq m}\bigcup_{z\in \widetilde{(\mathcal{P}_n)}_\epsilon}U_n(z).$$
%By Lemma~\ref{diametershadow}, the diameter of $U_n(z)$ is bounded by $C_1\lambda^{|z|}|z|^{\alpha}\log |z|$
%and so
%$$\mathbf{diam} (U_n(z))^h\leq C_2\lambda^{|z|h}|z|^{\alpha h}(\log |z|)^h.$$
%Summing over $z\in \mathcal{P}_n$ and over $z$ such that $(1-\epsilon)m\leq |z|\leq \epsilon n$, we get
%\begin{align*}
%\sum_{n\geq m}\sum_{z\in \widetilde{(\mathcal{P}_n)}_\epsilon}\mathbf{diam} (U_n(z))^h & \leq  C_3\sum_{n\geq m}M_n \lambda^{nh}n^{\alpha h}(\log n)^h\\
%&+C_3\sum_{n\geq m}\sum_{k=m(1-\epsilon)}^{\epsilon n}\mathrm{e}^{(v+h\log \lambda) k}k^{\alpha h}(\log k)^h.
%\end{align*}
%By Proposition~\ref{upperboundvgrBRW},
%%$\limsup \frac{1}{n}\log M_n\leq \omega_\Gamma(r)$, so for large enough $n$, we have
%$$M_n\leq \mathrm{e}^{n(\omega_\Gamma(r)+\epsilon)}.$$
%Therefore,
%$$\sum_{n\geq m}\sum_{z\in \widetilde{(\mathcal{P}_n)}_\epsilon}\mathbf{diam}(U_n(z))^h\leq C_3\sum_{n\geq m}(\mathrm{e}^{n(\omega_\Gamma(r)+\epsilon)}+C\mathrm{e}^{v\epsilon n})\lambda^{nh}n^{\alpha h}(\log n)^h.$$
%If $\epsilon$ is small enough, this quantity converges to 0 as $m$ tends to infinity.
%Finally the complement of conical limit points in $\Lambda$ is at most countable.
This concludes the proof.
 \end{proof}


We deduce the following result.
We denote by $\Lambda_\mathcal{B}(r)$ the limit set of the branching random walk inside the Bowditch boundary, endowed with the shortcut distance $\overline{\delta}_e$.
\begin{cor}\label{coroupperboundHausdorff}
Let $r\leq \rho^{-1}$.
Almost surely,
$$\Hdim(\Lambda_\mathcal{B}(r),\overline{\delta}_e)\leq \frac{-1}{\log \lambda}\omega_{\Gamma}(r).$$
\end{cor}

\begin{proof}
Combining Proposition~\ref{upperboundHausdorff} and~(\ref{shortcutsmallerFloyd}), we get that
$$\Hdim(\Lambda_\mathcal{B}(r)\cap \partial^{con}_{\mathcal{B}}\Gamma,\overline{\delta}_e)\leq \frac{-1}{\log \lambda}\omega_{\Gamma}(r).$$
The complement of the set of conical limit points in the Bowditch boundary is the set of parabolic limit points, which is countable.
This yields the desired upper-bound.
\end{proof}
Theorem~\ref{ThmHdimconst} is a consequence of Proposition~\ref{lowerboundHdim} and Corollary~\ref{coroupperboundHausdorff}.
Recall that if $\Gamma$ is hyperbolic, then it is also relatively hyperbolic and its Bowditch, Floyd and Gromov boundaries coincide.
Moreover, the shortcut distance and the visual distance are bi-Lipschitz by \cite[Proposition~6.1]{PY}.
Thus, Corollary~\ref{corohausdorffdimensionhyperbolic} follows from Theorem~\ref{ThmHdimconst}.

%We also deduce Corollary~\ref{coroaccessiblegroups} from Proposition~\ref{upperboundHausdorff}.
%Indeed, an infinitely-ended groups $\Gamma$ is relatively hyperbolic.
%Moreover, if $\Gamma$ is accessible, we can define a relatively hyperbolic structure on $\Gamma$ such that the complement of conical limit points is countable, see the discussion in \cite[Section~4.3]{DussauleYang}.
%Also, by \cite{Karlssonfreesubgroups}, the Floyd boundary covers the end boundary and by \cite[Lemma~4.3]{DussauleYang}, the Floyd distance is bounded from below by the visual distance on the end boundary.
%This proves the upper-bound
%$$\Hdim(\Lambda(r))\leq \frac{-1}{\log \lambda}\omega_{\Gamma}(r).$$
%{\color{blue}We need an argument to prove the lower bound without rewriting the whole proof.
%Maybe it is enough to say that like in Section 5, we can prove the lower bound, but I don't find this very satisfying.
%Also, in \cite{CandelleroGilchMuller}, the authors cover the case of free products by amalgamation, see Corollary~3.18, so our Corollary~\ref{coroaccessiblegroups} is not so new.
%And the proof for the upper-bound seems correct to me there.
%To avoid complication, I think the following would be more reasonable.
%We do not state Corollary~\ref{coroaccessiblegroups}, but
%in the end of Section 5, we state as a remark that the same proof would apply to the end boundary of an infinitely ended group, and so we would recover the results in \cite{CandelleroGilchMuller}.
%But if you hold on to this Corollary, I agree with keeping it, as long as we have a clean argument.
%Anyway, as I said, their proof for the upper-bound is correct and they use sub-multiplicativity of the Green function along spheres only for the lower bound.
%Thus, our proof of the lower bound fills their gap.
%So I think it is worth saying this somewhere.

%\ywy{I feel like their proof for the lower bound is correct: they use a similar idea to restrict BRW not go deep in parabolic cosets. Maybe I'm wrong?

%I agree Cor 1.3 is not that new. But I think it deserves to be recorded in the introduction, as it first complements the discussion about the situation of various boundaries, and secondly our methods are different and the result is more general.  The upper bound holds automatically as you explained, the lower bound only touches the definition of $\hat \delta_e$  and the second claim about $\mathbf E[W_n]$ in  Proposition \ref{lowerboundHdim}. 
%}

%\lw{Their proof has a gap. In fact, the upper and lower bound in their paper are respectively $1/R(\mathcal{F})$ and $1/z^{*}$.  They proved the two bounds are the same in Lemma 4.7, using the unproved inequality $H_{m + n} \leq H_m H_n$.  To fix this gap I think  something like our Lemma \ref{omegaL} might be necessary.   
%}\begin{comment}
%   {\color{red} Thanks for correcting me. Now I doubt the claim that critical exponent of $\mathcal L^{(d)}(\lambda|z)$ is the solution $z^{\star}_{d,L}$ in the proof of Lemma 4.11. It seems they overlooked the singularity pf $\mathcal L_i^{(d)}(\lambda|z)$??}
%   {\color{blue}The problem is that they use the same argument that proves the parabolic gap property to deduce that the singularity of the two power series are the same, so that $z^*=\omega_\Gamma(r)$.}
%\end{comment}

%{\color{blue}Well, you are absolutely right. And I realized that their proof for the upper bound also works for any infinitely ended group. So we do not need to reprove it here. I edited Corollary~\ref{coroaccessiblegroups} groups accordingly and I started writing a proof for the lower bound in the previous section.
%It is rather short and I think it definitely deserves to be written, since we both fill their gap and extend their result to all infinitely ended groups.}

%\medskip




%\section{Possible strategy for the upper bound of the whole limit set in the Floyd boundary}
%Assuming that $r<\rho^{-1}$, we have a continuous (actually H\"older) map from the compactification of the GW tree to the Floyd compactification of the group
%$$\phi : \mathcal{T}\cup \partial \mathcal{T}\to \Gamma \cup \partial_{\mathcal{F}}\Gamma.$$
%Define a Patterson-Sullivan like measure on $\partial \mathcal{T}$ as follows.
%Set
%$$\Upsilon_n=\frac{1}{\sharp B_{\mathcal{T}}(u_0,n)}\sum_{u\in B_{\mathcal{T}}(u_0,n)}\mathbf{D}(u),$$
%where $u_0$ is the root of $\mathcal{T}$ and $B_{\mathcal{T}}(u_0,n)$ is the ball of radius $n$ centered at the root $u_0$ in $\mathcal{T}$.
%Then, $\Upsilon_n$ converges weakly to a limit measure $\Upsilon$ supported on $\partial \mathcal{T}$.
%Next, we work on the level of the group and define
%$$\chi_n=\frac{1}{\widetilde{M}_n}\sum_{x\in \widetilde{P}_n}\mathbf{D}(x),$$
%where $\widetilde{P}_n$ is the intersection of $\mathcal{P}$ with the ball $B(e,n)$ in $\Gamma$ and $\widetilde M_n$ is the cardinality of $\widetilde P_n$.
%Then, $\chi_n$ is a random probability measure on the Floyd compactification and thus weakly converges to some random measure $\chi$ up to a sub-sequence for the weak topology on random probability measures.

%The following two claims seem reasonable.
%\begin{claim}
%We have almost surely $\chi=\phi_*\Upsilon$.
%\end{claim}

%\begin{claim}
%We have for every random measurable set $A$ on $\partial \mathcal{T}$
%$$\mathbf P(\exists u_n \in \mathcal{T}, u_n\to u_\infty \in A)\leq \mathbf E[\Upsilon(A)].$$
%\end{claim}

%{\color{blue}After second thought, this claim is completely wrong, since as long as $A$ is non empty, the probability on the left hand side is 1. Still, I think it could give information to look at the Patterson-Sullivan type measure on $\partial \mathcal{T}$.
%But maybe this would deserve a whole new research program and should be postponed to future work.
%I think I'm ready now to give up the Hausdorff dimension of the whole limite set in the Floyd boundary and to publish the paper as it is.}

%\lw{Yes, I agree.  The measure $\Upsilon$ is the uniform measure on $\partial \mathcal{T}$ and so its image $\chi = \phi_{*} \Upsilon$ may be also viewed as the uniform measure on the limit set $\Lambda$. I guess the intensity measure $\mathbf{E}[\chi]$ may have some relation to the Patterson-Sullivan measure on $\partial_{\mathcal{F}} \Gamma$ defined by the Green function. I think $\chi$ is an interesting object even in the hyperbolic case. }

%{\color{blue}My guess is that almost surely, the Hausdorff dimension of the measure $\chi$ is $\omega_\Gamma(r)$ and that $\chi$ is equivalent to the $\omega_\Gamma(r)$-dimensional Hausdorff measure.
%Then, if there is a parabolic gap, the discussion below still shows that $\chi(\partial_{\mathcal{F}}(gP)=0$, even though the claims do not hold.
%Maybe there is a way to deduce from this that the Hausodrff dimension of $\partial_{\mathcal{F}}(gP)$ is smaller than $\omega_\Gamma(r)$.}

%Fix then a parabolic coset $gP$.
%Then, for all $k$, the set $U_k$ defined by
%$$U_k=\{x\in \Gamma, d_{gP}(e,x)\geq k\}$$
%contains an open neighborhood $V_k$ of $\partial_{\mathcal{F}}(gP)$.
%Now, for every $\epsilon>0$, we can define an event $A_\epsilon$ by
%$$A_\epsilon=\{M_n\geq c_\epsilon \mathbf{E}[M_n]\},$$
%for suited $c_\epsilon$ such that $\mathbf P(A_n^c)\leq \epsilon$.
%Under the event $A_\epsilon$, we have
%$\chi_{n}(U_k)\leq c_\epsilon^{-1}\mathrm{e}^{-k\alpha}$
%for suited $\alpha$, assuming a parabolic gap.
%Moreover, by definition of the weak topology on random probability measures, the measure $\mathbf 1_{A_\epsilon}\chi_n$ weakly converges to $\mathbf 1_{A_\epsilon}\chi$, which proves by the random Portmanteau Theorem that almost surely,
%$$\mathbf 1_{A_\epsilon}\chi(\partial_{\mathcal{F}}(gP))=0.$$
%Thus, $\mathbf E[\chi(\partial_{\mathcal{F}}(gP))]\leq \epsilon$ and since this is true for every $\epsilon>0$, we have almost surely
%$$\chi(\partial_{\mathcal{F}}(gP))=0.$$
%Finally,
%$$\mathbf P(\Lambda \cap \partial_{\mathcal{F}}(gP)\neq \emptyset)\leq \mathbf P\bigg(\exists u_n\in \mathcal{T}, u_n\to u_\infty \in \phi^{-1}(\partial_{\mathcal{F}}(gP))\bigg).$$
%By the second above claim,
%$$\mathbf P(\Lambda \cap \partial_{\mathcal{F}}(gP)\neq \emptyset )\leq \mathbf E[\Upsilon(\phi^{-1}(\partial_{\mathcal{F}}(gP))]$$
%and by the first above claim,
%$$\mathbf P(\Lambda\cap \partial_{\mathcal{F}}(gP)\neq \emptyset)\leq \mathbf E[\chi(\partial_{\mathcal{F}}(gP))]=0.$$
%Thus, almost surely, $\Lambda$ is contained in the set of conical limit points, and thus its Hausdorff dimension is upper bounded by $\omega_\Gamma(r)$.


%\section{Another upper bound for the Hausdorff dimension}
%\ywy{The upper bound in  section is not any new, so will be removed}
%{\color{blue}Yes, but maybe we can improve the proof and only assume a parabolic gap.} \ywy{This time, I think this section could be removed: I still do not find a way to control the BRW in a parabolic subgroup.}
%{\color{blue}Okay let us comment this out. But I like very much your strategy, I'm still trying to adapt your arguments.}
%In this section, we prove the upper bound on Hausdorff dimension of the limit set of the branching random walk in the Floyd boundary.
%We fix $r<\rho^{-}$ and
%we assume that $\Gamma$ has a parabolic gap for the Green function at $r$, i.e. $\omega_P(r)<\omega_\Gamma(r)$ for every parabolic subgroup $P$.

%The branching random walk $\mathrm {BRW}(\Gamma, \nu,\mu)$ can be viewed as tree $\mathcal T$ indexed random walks, where the tree is randomly sampled by a Galton-Watson process.   

 
%\begin{lem}\label{5.2SWX} 
%Let $D=\diam{\supp(\mu)}$.  Then almost surely, we have $|X_u|\leq D|u|$ for any $u\in \mathcal T$.
%\end{lem}

%The following lemmae are proved in \cite{SWX}. Their proof does not use hyperbolicity.



%\begin{lem}\label{5.3SWX}\cite[Lemma~5.3]{SWX}
%Suppose $r\in [1,\rho^{-1})$ and $0<l<-v^{-1}\log (r\rho)$.
%Then almost surely, there are finitely many $u\in \mathcal T$ such that $|X_u|<l|u|$.
%In particular,
%$$\liminf_{n\to \infty} \frac{1}{n} \inf_{|u|=n}|X_u|\geq -v^{-1}\log (r\rho)$$ almost surely.
%\end{lem}


%\begin{lem}\label{5.4SWX}\cite[Lemma~5.4]{SWX}
%Assume that $r\in [1,\rho^{-1}]$.
%Then, for any $\epsilon>0$, almost surely there are finitely many $u$ such that $d(X_{u^-},X_u)>\epsilon |X_{u^-}|$.
%\end{lem}



%We deduce the following.
%\begin{prop}\label{EmbedGWTProp}
%Let $r\in [1, \rho^{-1})$, $0<l_0<-\frac{\log rp}{v}$  and  $\alpha=l_0\log_2 a$. Then almost surely, $u\in (\mathcal T, d_{\mathcal T})\mapsto X_u\in (\Gamma, \rho_\lambda)$ defines a Hölder continuous map of index $\alpha$, and this map can be extended to Hölder continuous map of index $\alpha$ from $(\mathcal T\cup \partial \mathcal T, d_{\mathcal T})$ to $(\Gamma\cup \partial_{\mathcal{F}} \Gamma, \delta_{e})$.
%\end{prop}

%Denote by $\Lambda_{\mathcal{F}}$, respectively $\Lambda_{\mathcal{B}}$, the limit set of the BRW in the Floyd compactification, respectively the Bowditch compactification.

%\begin{cor}
%For $1\leq \lambda<\rho^{-1}$, almost surely
%$$\Lambda_{\mathcal{F}}=\big \{X_\xi, \xi \in \partial\mathcal{T}\big \}$$
%and
%$$\Lambda_{\mathcal{B}}=\big \{X_\xi, \xi \in \partial\mathcal{T}\big \}.$$
%\end{cor}

 
 




%Fix $\lambda\in (0,1)$.
%Let $\mathcal P$ be the (random) trace   of BRW on $\Gamma$. Set $\mathcal P_n=\mathcal P\cap S_n$ for $n\ge 1$. 



%\begin{lem}\label{SublinearVisitTransLem}
%Let $\epsilon>0$.  Almost surely, there exists $n_0>0$ such that for any $n>n_0$ and any $z\in \mathcal P_n$,  any transition point on    $[e,z]$ is $\epsilon n$-close to $\mathcal P$.  
%\end{lem}
%\begin{proof}
%Let $A_m$ be the event such that for some $z\in \mathcal P_m$, the trace   $\mathcal P$  is  $\epsilon m$-close to some transition point in  $[1,z]$. By Lemma \ref{SupExpDecayLem}, 
%the probability of the trajectories from $1$ to $z$ without visiting $\epsilon n$-ball of transition points in $[1,z]$ decays super-exponentially in $\epsilon n$. The proof proceeds exactly as the proof for (5.2) in \cite{SWX}.
%\end{proof}


%Fix $\epsilon>0$. We define the tight augmentation by $\epsilon$-saturated parabolic cosets as follows.  
%\begin{defn}\label{TightAugDef1}
%Let $x\in S_n$ for $n\ge 1$.
%We say a parabolic coset $P\in \mathbb P$ is \textit{$\epsilon$-saturated} in some $[e,x]$  if  $$\diam{[e,x]\cap N_\eta(P)}\ge 2\epsilon |x|.$$ Let $P_\epsilon({x})$ be the set of points $y\in P\cap B(e,n)$  such that $$
% \diam{[e,y]\cap N_\eta(P)}\ge 2\epsilon |x|.
%$$

%Define  $\widetilde X_\epsilon$ to be the union of $X$ with all $P_{\epsilon}(x)$ for $\epsilon$-saturated parabolic cosets $P\in\mathbb P$ and $x\in X$. 

%\end{defn}





%\begin{lem}\label{ }
%Assume that $\Gamma$ has parabolic gap for Green function. Then almost surely, that the growth rate $v(P\cap \mathcal P)$ of each $P\in \mathbb P$ is strictly less than the growth rate $v(\mathcal P)$ of $\mathcal P$.
%\end{lem}

%begin{proof}
%We set $M_{n,P}=\sharp \mathcal{P}_n\cap P$.
%By definition, the growth rate of $P\cap \mathcal{P}$ is $\limsup \frac{1}{n}\log M_{n,P}$ and the growth rate of $\mathcal{P}$ is $\limsup \frac{1}{n}\log M_n$.
%Similarly to Proposition~\ref{upperboundvgrBRW}, it follows from the Borel-Cantelli Lemma that almost surely,
%$$\limsup \frac{1}{n}\log M_{n,P}\leq \omega_P(r).$$
%Moreover, by Proposition~\ref{P:vgrBRW}, almost surely,
%$$\limsup \frac{1}{n}\log M_{n}\geq \omega_\Gamma(r).$$
%Thus, if $\omega_P(r)<\omega_\Gamma(r)$, we have almost surely
%$$\limsup \frac{1}{n}\log M_{n,P}<\limsup \frac{1}{n}\log M_{n},$$
%which concludes the proof.
%\end{proof}

%In the following lemma, we consider $X$ to be the random trace graph $\mathcal P$ of BRW provided by Lemma \ref{SublinearVisitTransLem}.
%\begin{lem}\label{SameGrowthLem}
%Fix $1>\epsilon>0$. Assume that $\Gamma$ has parabolic gap for Green function.\begin{comment}
%By the current proof, this assumption is not enough: we need to assume that the growth rate of $P$ (not $P\cap \mathcal P)$ is strictly less than that $\mathcal P$.       
%\end{comment} Then almost surely,  we have  $v(\mathcal P)=v(\widetilde {\mathcal P}_\epsilon)$.
%\end{lem}
%\begin{proof}
%According to the above discussion, it suffices to prove that the growth rate of $\widetilde{X}\setminus X$ is strictly less than $v(X)$. To be precise, we need to find $v'<v$ so that 
%$$
%|\cup_{x\in X_n} P_\epsilon(x)| \le e^{v' n}
%$$
% for all $n\gg 0$ and $X_n:=S_n\cap X$.

%Let $x\in X_n$ and $P\in\mathbb P$ be $\epsilon$-saturated in $[e,x]$. The   entry point $z$ of $[e,x]$ in $N_\eta(P)$   is transitional in $[e,x]$ by Lemma \ref{EntryTransLem}. For any $y\in S_n\cap P_{\epsilon}(x)$, we have   $d(z,[e,y])\le C$, which yields    $$|d(e,z)+d(z,y)-n|\le 2C$$ Since $P$ is $\epsilon$-saturated for $[e,x]$,   we have $|z|\le (1-\epsilon)|x|\le n(1-\epsilon)$ and $d(z,y)\ge \epsilon n-2C$.

%By Lemma \ref{SublinearVisitTransLem} for $\theta= \frac{v-v_1}{2(v+v_1)}$, there exists $x'\in X$ such that $d(z,x')\le \theta n$. Thus,  we have $$|P_{\epsilon}(x)|\prec e^{\omega_1|x'^{-1}y|} \prec e^{\omega_1 (|z^{-1}y|+\theta n)}.$$ 
%As $X$ has purely exponential growth, we have $|X_{n(1-\epsilon)}|\prec e^{v (|z|+\theta n)}$. The cardinality of the set $\cup_{x\in X_n} P_\epsilon(x)$  is upper bounded by
%\begin{align*}
%& C  e^{v (|z|+\theta n)} e^{v_1 (|z^{-1}y|+\theta n)}\\
%\le & e^{v n} e^{(v_1-v)|z^{-1}y| +\theta n(v+v_1)} \\
%\le & e^{v n} e^{(v_1-v)\epsilon n +\theta n (v+v_1)}  \le e^{v n} e^{(v_1-v)\epsilon n/2}   
%\end{align*}
%for some uniform constant $C$.  This concludes the proof.
%\end{proof}


%The following fact will be used in the next proof.
%\begin{lem}\label{FloydgeodLem}
%There exists a constant $C=C(\eta,L)$ with the following property.
%Let $\gamma$ be a Floyd geodesic between $x$ and $\xi$. If $y$ is a $(\eta,L)$-transition point on a word geodesic $[x,\xi]$, then $d(y, \gamma)\le C$.    
%\end{lem}

%\begin{proof}
%This follows from the results in \cite[Section 5]{PY}. We first recall the facts proven there. By definition, a path is called \textit{$l$-tight} for $l>0$ if any subpath of diameter at most $l$ is a word geodesic.  By   \cite[Lemma 5.3]{PY}, the Floyd length between $x_n$ and $\xi$ is the limit of Floyd lengths of   a sequence of tight paths in $\Gamma$ between them. A  truncation  of tight paths introduced there is to replace the maximal connected components contained in parabolic cosets by word geodesics, and the resulting path is a uniform quasi-geodesic whose transitional points are contained in the original tight path (\cite[Lemma 5.8]{PY}).  

%By Lemma \ref{ThinTransitional2},  the transition point $y$ in $[x,\xi]$ is $C$-close to a transitional point in the truncated tight paths. Assuming the above facts, we obtain that  $y$ is $C$-close to the Floyd geodesic $\gamma$. 
%\end{proof}

%\begin{thm}\begin{comment}
%{\color{blue}I have the impression this works for the Floyd boundary too. This would prove that assuming a parabolic gap and assuming $r<\rho^{-1}$, we also have the upper bound for the Floyd boundary.}\ywy{Yes, but the assumption \ref{SameGrowthLem} requires a stronger assumption, which seems known to be true only for amenable parabolic subgroups. The good point is that here the proof uses a different argument.}
%{\color{blue} Why would this be true for amenable groups ? We prove that the growth rate of $P\cap \mathcal{P}$ is non-positive for amenable $P$, but in general we have no comparison between the growth rate of $P$ and that of $\mathcal{P}$. But this is true for parabolic subgroups of sub-epoxnential growth, which is already very good.} Yes, You are right. Still, if we know that parabolic subgroups have sub-exp growth, then their limit sets in Floyd boundary is zero. The spectral radius case is also true in this case. How to deal with this result? Maybe we should keep it for further reference, though it does not give any real new.
%\end{comment}
%Assume $r \in[1,\rho^{-1})$. Then  almost surely, the Hausdorff dimension of $(\Lambda, \rho_\lambda)$ is less than $h(r)=-\log_\lambda H(r).$ 
%\end{thm}

%\begin{proof}
%For any $\xi\in \Lambda$, there exists a geodesic ray $[u_n]$ in $\mathcal T$ where $n=d_{\mathcal T}(o, u_n)$, so that $\xi=\lim_{n\to\infty} x_n$ where   $x_n:=X_{u_n}$ is an element in $\Gamma$. By Proposition \ref{EmbedGWTProp}, there exist $l_0>0$ and $c>0$ such that  
%$$\delta_{e}(x_n, \xi)\le c\lambda^{l_0n}$$
%and 
%$$
%l_0 n \le |x_n| \le D n
%$$
%where   $D$ is the diameter of the finite support of the    $\mu$-random walk. 

%Fix an arbitrary small number $\epsilon>0$.  Consider the    sub-segment $\alpha$ of $[e,x_n]$ with length $2\epsilon nl_0$  consisting of points $y$ so that 
%\begin{equation}\label{alphalengthEQ}
%(1-3\epsilon)nl_0\le d(y, x_n)\le (1-\epsilon)nl_0.    
%\end{equation}

 
%Let $\delta$ be the   Floyd constant in Lemma \ref{TransLargeLem}.  For any $y\in \alpha$ and $n\gg 0$, we have   $$c\lambda^{l_0n}\lambda^{-|y|}\le c \lambda^{3\epsilon l_0 n } <\delta/2,$$
%which in turn implies $$\delta_{y}(x_n,\xi) \le \lambda^{|y|} \delta_{e}(x_n,\xi)< \delta/2.$$ 



%We consider two cases according to whether $\alpha$ contains a transition point in $[1,x_n]$.

%\textbf{Case 1.}
%Suppose that $\alpha$ contains a transition point $y$ of $[e,x_n]$. Thus, $\delta_{y}(e,x_n)\ge \delta$ and then $$\delta_{y}(e,\xi)>\delta-\delta/2\ge \delta/2.$$ By Lemma \ref{Karlssonlem}, we obtain that $d(y, [e,\xi])\le \kappa=\kappa(\delta/2)$.



%Assume that $|x_n|\ge n_0$. By Lemma \ref{SublinearVisitTransLem}, since $y$ is transition  in $[e,x_n]$, there exists $z\in \mathcal P$ such that $d(z, y)\le \epsilon|x_n|\le  D \epsilon n$, so $d(z, [e,\xi])\le \kappa+ D\epsilon n.$ That is to say, $\xi\in \Pi(z, \kappa+D\epsilon n)$.

%\textbf{Case 2.} Suppose that $\alpha$ contains no transition point. By Lemma \ref{EntryTransLem},  $\alpha$ is contained in the $\eta$-neighborhood of some parabolic coset $P\in\mathbb P$. If   $u, v$ are the corresponding entry and exit points of $[e,x_n]$ in $N_\eta(P)$,  we have $d(u, v)\ge |\alpha|\ge 2\epsilon l_0 n$. It is possible    that $u=e, v=x_n$. 

%Since $u$  is a transition point on $[e,x_n]$, it follows by   Lemma \ref{ThinTransitional2} that  $u$ is $C$-close to a transition point either on $[x_n,\xi]$ or on $[e,\xi]$.  Let $\gamma$ be a Floyd geodesic between $x_n$ and $\xi$. If $u$ is $C$-close  to a transition point on $[x_n,\xi]$, then by Lemma \ref{FloydgeodLem}, we obtain   $d(u, \gamma)\le 2C$, and thus it follows from   (\ref{alphalengthEQ}) that    $$d(e, \gamma)\le d(e,u)+2C\le  d(e,\alpha_-)+2C\le l_0 n-3\epsilon l_0 n+2C$$
%Recalling   $\delta_{e}(x_n, \xi)\le c\lambda^{l_0n}$ as above,  there exists   $R=R(c)$ by  Lemma \ref{FloydRayLem} such  that $$d(e, \gamma)>l_0n- R.$$ For any given $\epsilon>0$, this gives a  contradiction for any sufficiently large $n\gg 0$. Thus, we proved that $u$ is $C$-close to $[e,\xi]$. 

%Let $z$ be the exit point of $[e,\xi]$ at $N_\eta(P)$, so $z$ is a transitional point.  Applying  Lemma \ref{ThinTransitional2} to the triangle $\Delta(e x_n\xi )$,  we have $d(z, [x_n,\xi])\le C$ or $d(z, [e,x_n])\le C$. In the former case,  we have $d(z, \gamma)\le 2C$ by Lemma \ref{FloydRayLem}, so     $$d(e,z)> d(e,\gamma)-d(z,\gamma) \ge l_0n-R-2C$$
 
%In the latter case, we have   $d(z,z')\le C$ for some $z'\in [e,x_n]$. By Lemma \ref{ThinTransitional2} applied for triangle $\Delta(ezz')$, as   $v$ is a transition  point on   $[e,z']$, we have $d(v,[z,z'])\le C$ and so $d(v, z)\le 2C$. Using (\ref{alphalengthEQ}), we obtain $$d(e,z)\ge d(e,v) - 2C\ge d(e, \alpha^+)-2C\ge l_0 n(1-\epsilon)-2C.$$
 
%Summarizing the two cases, we proved that $d(e,z)\ge l_0 n(1-\epsilon)-R-2C$. We choose $w\in [e,z]$ so that 
%\begin{equation}\label{wLenEQ}
%|w|=l_0 n(1-\epsilon)-R-2C    
%\end{equation} We are going to prove that $w$ is contained in the $P_{\epsilon'}(x_n)$ for some $\epsilon'>0$ decided below. 

%First of all, $P$ is $\epsilon'$-saturated in $[e,x_n]$: for   $|x_n|\le Dn$ and $\epsilon':=2\epsilon l_0/D$, we have  \begin{align*}
%diam([e,x_n]\cap N_\eta(P))\ge |\alpha|\ge 2\epsilon l_0n\ge 2\epsilon l_0 |x_n|/D\ge \epsilon'|x_n|.    
%\end{align*}  


%Since $u$ is a transition point, we have $d(u, [e,z])\le C$. Noting that $d(e,u)\le l_0n(1-3\epsilon)$, we see that $w\in N_\eta(P)$. On account of (\ref{wLenEQ}), up to decreasing $\epsilon'$, we obtain       $$d(u,w)\ge d(e,w)-d(e,u)\ge 2\epsilon l_0 n-R-2C\ge   \epsilon'|w|$$ As $u, w\in N_\eta(P)$, we have $\diam{N_\eta(P)\cap [e,w]}\ge d(u,w)\ge \epsilon'|w|$, so   $w\in P_{\epsilon'}(x_n)$. 

%For $w\in [e,\xi]$, we have  $\xi\in \Pi(w, C)$. In conclusion, for given $\epsilon$,  we have  the following covering of $\Lambda$:
%\begin{align*}
%\Lambda &\subset \cup_{n\ge n_0}^{\infty} ( \{ \Pi(x, \epsilon n): x\in \mathcal P_n\} \bigcup \cup \{\Pi(w, C): w\in P_{x,\epsilon'}\})\\
%&\subset \cup_{n\ge n_0}^{\infty}  \{\Pi(x, \epsilon n): x\in \widetilde{\mathcal P}_n\}     
%\end{align*}
 
%By Lemma \ref{SameGrowthLem}, the subset $\widetilde X_{\epsilon'}$ has the same growth rate as $X$. By Lemma \ref{ComputeHdim}, the Hausdorff dimension of $\Lambda$ is upper bounded by $v(X)/-\log\lambda$. 
%\end{proof}

\appendix
\section{Convergence of random finite measures}

\subsection{Random finite measures}
Let $(X,\mathcal{B})$ be a Polish space endowed with its Borelian $\sigma$-algebra.
Let $(\Omega,\mathcal{F},\mathbf P)$ be a probability space.
\begin{defn}\label{defrandomfinitemeasure}
A random finite measure on $X$ is a map
$$\mu : (\omega,B)\in \Omega\times \mathcal{B}\mapsto \mu_\omega(B)\in \mathbb R$$
such that
\begin{enumerate}[(a)]
    \item for every $B\in \mathcal{B}$, the map $\omega\mapsto \mu_\omega(B)$ is measurable,
    \item for $\mathbf P$-almost every $\omega$, $\omega\mapsto \mu_\omega$ is a finite Borelian measure on $X$,
    \item the expectation $\mathbf E[\mu_\omega(X)]$ is finite,
\end{enumerate}
We identify $\mu$ with the family of maps $\mu_\omega:\mathcal{B}\to \mathbb R$ and we write $\mu=(\mu_\omega)_\omega$.
\end{defn}

We denote by $\mathcal{M}_\Omega(X)$ the set of random finite measure on $X$ and by $\mathcal{M}(X)$ the set of finite measures on $X$.
Given a random finite measure $\mu=(\mu_\omega)_\omega$, we define the measure $\pi_X(\mu)$ on $X$ by
$$\pi_X(\mu)(B)=\mathbf E\left [\mu_\omega(B) \right]$$
for every Borelian set $B\in \mathcal B$.
We also define the measure $\pi_\Omega(\mu)$ on $\Omega$ by
$$\pi_\Omega(\mu)(A)=\mathbf E[\mathbf 1_A\mu_\omega(X)].$$
We call $\pi_X(\mu)$, respectively $\pi_\Omega(\mu)$, the $X$-marginal, respectively the $\Omega$-marginal of $\mu$.

If $\mu=(\mu_\omega)_\omega$ is a random finite measure, then one can define a finite measure $\tilde{\mu}$ on $X\times \Omega$  by setting for every measurable set $A$ of $X\times \Omega$
$$\mu(A)=\mathbf E\left [\int \mathbf 1_{A}(x,\omega)d\mu_\omega(x)\right ].$$
Then, $\pi_X(\mu)$ and $\pi_\Omega(\mu)$ are the push-forward measures of $\mu$ by the canonical projections $\pi_X:X\times \Omega \to X$ and $\pi_\Omega:X\times \Omega\to \Omega$.

\begin{defn}
A random finite measure $\mu$ is called $\mathbf P$-regular if $\mathbf P$ and $\pi_\Omega(\mu)$ are absolutely continuous with respect to each other.
\end{defn}


\begin{lem}
Let $\mu=(\mu_\omega)_\omega$ be a random finite measure.
Then $\pi_\Omega(\mu)$ is absolutely continuous with respect to $\mathbf P$.
Moreover, $\mu$ is $\mathbf P$-regular if and only if for $\mathbf P$-almost every $\omega$, $\mu_\omega$ is not the null measure.
\end{lem}

\begin{proof}
If $A$ is  such that $\mathbf P(A)=0$, then $\mathbf P$-almost surely, $\mathbf 1_A \mu_\omega(X)=0$, so $\pi_\Omega(\mu)(A)=\mathbf E[\mathbf 1_A\mu_\omega(X)]=0$.
This concludes the first part of the lemma.

For the second part, assume that the event $A=\{\omega,\mu_\omega(X)=0\}$ has positive probability.
Then, $\mathbf E[\mathbf 1_A \mu_\omega(X)]=0$.
So $\pi_\Omega(\mu)(A)=0$, but $\mathbf P(A)>0$.
Conversely, assume that $\mathbf{P}$-almost surely, $\mu_\omega(X)>0$ and let $A$ be such that
$\pi_\Omega(\mu)(A)=0$, hence $\mathbf P$-almost surely, $\mathbf 1_A \mu_\omega(X)=0$.
Then, we necessarily have $\mathbf 1_A=0$ $\mathbf P$-almost surely, i.e. $\mathbf P(A)=0$. 
\end{proof}

\begin{rem}
In fact, by definition of $\pi_\Omega(\mu)$, we have that the Radon-Nikodym derivative of $\pi_\Omega(\mu)$ with respect to $\mathbf P$ is given by
$$\frac{d\pi_\Omega(\mu)}{d\mathbf P}(\omega)=\mu_\omega(X).$$
Thus, $\mathbf P$ is absolutely continuous with respect to $\pi_\Omega(\mu)$ if and only if this Radon-Nikodym derivative is almost surely positive and then,
$$\frac{d\mathbf P}{d\pi_\Omega(\mu)}(\omega)=\frac{1}{\mu_\omega(X)}$$
$\mathbf P$-almost surely.
\end{rem}

If $\mu=(\mu_\omega)_\omega$ is a random probability measure, i.e.\ $\mathbf P$-almost surely, $\mu_\omega$ is a probability measure on $X$, then the $\Omega$-marginal of $\mu$ is $\mathbf P$.
Thus, $\mathbf P$-regularity is automatic in this context.
However, it is easy to construct an example where $\mathbf P$-regularity fails, since one only needs that $\mathbf P(\mu_\omega(X)=0)>0$.
Indeed, let $\mu$ be any random finite measure and let $A$ be an event such that $\mathbf P(A)\leq 1/2$, then $\mathbf 1_A\mu$ is also a random finite measure and
$\mathbf P(\mathbf 1_A\mu(X)=0)\geq 1/2$.
Restricting our attention to $\mathbf P$-regular measures will be important in the following, mainly because of the following result.



\begin{prop}\label{disintegrationrandom}
Every finite measure $\mu$ on $X\times \Omega$ such that $\pi_\Omega(\mu)$ and $\mathbf P$ are absolutely continuous with respect to each other is defined by a random finite measure.
\end{prop}

\begin{proof}
Let $\mu$ be a finite measure on $X\times \Omega$ and denote by $\|\mu\|=\mu(X\times \Omega)$ its mass.
Then, $\nu=\mu/\|\mu\|$ is a probability measure on $X\times \Omega$ with marginal $\mathbf Q=\pi_\Omega(\mu)/\|\mu\|$.
By the disintegration theorem \cite[Proposition~3.6]{Crauel}, there exists a map
$$(B,\omega)\in \mathcal{B}\times \Omega\mapsto \nu_\omega(B)$$
such that
$\omega\mapsto \nu_\omega$ is $\mathbf Q$-almost surely a probability measure on $X$ and for every Borelian $B\in \mathcal{B}$, $\omega \mapsto \nu_\omega(B)$ is measurable.

By assumption, the probability measure $\mathbf Q$ is absolutely continuous with respect to $\mathbf P$.
Let $f=d\mathbf Q/d\mathbf P$ be the corresponding Radon-Nikodym derivative.
Define then $\mu_\omega=\|\mu\|f(\omega)\nu_\omega$.
Then, $\mu_\omega$ is $\mathbf Q$-almost surely, hence $\mathbf P$-almost surely, a finite measure on $X$, for every Borelian set $B$, the map $\omega\mapsto \mu_\omega(B)$ is measurable and for every non-negative measurable function $\phi$ on $X\times \Omega$, we have
\begin{align*}
\mu(\phi)&=\|\mu\|\nu(\phi)\\&=\|\mu\|\int \int \phi(x,\omega) d\nu_\omega(x) d\mathbf Q(\omega)\\&=\int \int \phi(x,\omega)\|\mu\|d\nu_\omega(x) f(\omega)d\mathbf P(\omega)\\
&=\mathbf E \left [\int \phi(x,\omega)d\mu_\omega(x)\right ].
\end{align*}
Thus, $\mu$ is defined by the random finite measure $(\mu_\omega)_\omega$.
\end{proof}

In general, the disintegration theorem allows one to decompose a finite measure on $X\times \Omega$ along its $\Omega$-marginal, so without assuming $\mathbf P$-regularity, it cannot be seen as a random finite measure in the sense of Definition~\ref{defrandomfinitemeasure}.

We now introduce some definitions based on \cite{Crauel}.

\begin{defn}\label{defrandomclosedset}
A random closed set is a map $\omega \in \Omega \mapsto C(\omega)\in 2^X$ taking values in closed subsets of $X$ and such that the map $\omega\mapsto d(x,C(\omega))$ is measurable for every $x\in X$.
A random open set is a map $\omega \mapsto U(\omega)$ such that the complement map $\omega\mapsto U^c(\omega)$ is a random closed set.
\end{defn}

Let $C=\omega\mapsto C(\omega)$ be a random closed set, $U=\omega\mapsto U(\omega)$ a random open set and $\mu=(\mu_\omega)_\omega$ a random finite measure.
We set
$$\mu(C)=\mathbf E\left [\mu_\omega(C(\omega))\right]$$
and
$$\mu(U)=\mathbf E\left [\mu_\omega(U(\omega))\right].$$

\begin{defn}\label{defrandomcontinuousfunction}
A random bounded continuous function is a map
$f:X\times \Omega\to \mathbb R$ such that
\begin{enumerate}[(a)]
    \item for every $x\in X$, $\omega\mapsto f(x,\omega)$ is measurable,
    \item for every $\omega\in \Omega$, $x\mapsto f(x,\omega)$ is continuous and bounded,
    \item There exists $C\geq 0$ such that for $\mathbf P$-almost every $\omega$, $\|f(\cdot,\omega)\|_\infty\leq C$.
\end{enumerate}
\end{defn}

\begin{rem}
The third condition can be reformulated as $\|f(\cdot,\omega)\|_\infty$ is in $L_\infty(\Omega,\mathbf P)$.
In \cite{Crauel}, the author introduces several spaces of random functions, replacing the third condition by $\|f(\cdot,\omega)\|_\infty\in L_p(\Omega,\mathbf P)$.
For $p=1$ the corresponding space of function is called the space of random continuous functions there.
\end{rem}

We denote by $C_{\Omega,b}(X)$ the space of random bounded continuous functions and endow $C_{\Omega,b}(X)$ with the $L_\infty\times L_\infty$-norm $\|\cdot\|_\infty$,
defined by
$$\|f\|_\infty=\inf \left\{C\geq 0, \mathbf P\big (\omega,\|f(\cdot,\omega)\|_\infty>C\big)=0\right\},$$
for every $f\in C_{\Omega,b}(X)$.
If $f$ is a random bounded continuous function and $\mu$ is a random finite measure, then the integral
$$\mu(f)=\mathbf E\left [\int f(x,\omega)d\mu_\omega(x)\right]$$
is well defined.

Recall that a Lipschitz function on $X$ is a function $f:X\to \mathbb R$ such that
$$\|f\|_L=\sup_{x,y\in X}\frac{|f(x)-f(y)|}{d(x,y)}$$
is finite.
We then set
$$\|f\|_{BL}=\sup \{\|f\|_\infty,\|f\|_L\}$$
and say that $f$ is bounded Lipschitz if $\|f\|_{BL}$ is finite.
We denote by $BL(X)$ the set of bounded Lipschitz functions on $X$.

\begin{defn}\label{defrandomlipschitzfunction}
A random bounded Lipschitz function is a random bounded continuous function $f$ such that there exists $C\geq0$ such that for $\mathbf P$-almost every $\omega$, the map 
$x\mapsto f(x,\omega)$ is bounded Lipschitz and $\|f(\cdot,\omega)\|_{BL}\leq C$.
\end{defn}

\begin{rem}
In \cite{Crauel}, random bounded Lipschitz functions are called random Lipschitz functions.
We changed the terminology to insist on the fact that random functions we are considering here are $\mathbf P$-essentially bounded.
\end{rem}
We denote by $BL_\Omega(X)$ the set of random bounded Lipschitz functions on $X$.

\begin{lem}\label{measuredeterminedbylipschitz}
If two random finite measures $(\mu_\omega)_\omega$ and $(\nu_\omega)_\omega$ coincide on random bounded Lipschitz functions, i.e. for every $f\in BL_\Omega(X)$, $\mu(f)=\nu(f)$, then for $\mathbf P$-almost every $\omega$, $\mu_\omega=\nu_\omega$.
\end{lem}

\begin{proof}
For every closed set $C\in X$, the sequence of functions
$$f_n:x\mapsto 1-(1\wedge nd(x,C))$$ is non-increasing and converges to $\mathbf 1_C$.
Moreover, for every $n$, $f_n$ is bounded Lipschitz.
Therefore, for every event $A$, the maps $(x,\omega)\mapsto \mathbf{1}_A(\omega)f_n(x)$ are random bounded Lipschitz functions.
Thus by monotone convergence, it suffices to prove that if for all closed set $C$, for all event $A$, $\mu(C\times A)=\nu(C\times A)$, then $(\mu_\omega)_\omega$ and $(\nu_\omega)_\omega$ coincide $\mathbf P$-almost surely.
Since closed sets are closed under finite intersection and the measures we are considering are finite, this follows from the monotone class theorem \cite[Theorem~3.4]{Bill86}.
\end{proof}


\subsection{The weak topology on random finite measures}

Recall that the weak topology on $\mathcal{M}(X)$ is the smallest topology such that for all bounded continuous function $f:X\to \mathbb R$, the map $\mu\mapsto \mu(f)$ is continuous.

\begin{defn}\label{defweaktopology}
The weak topology on $\mathcal{M}_\Omega(X)$ is the topology generated by the maps $\mu \mapsto \mu(f)$, for every $f\in C_{\Omega,b}(X)$, i.e. it is the smallest topology on $\mathcal{M}_\Omega(X)$ such that for every $f\in C_{\Omega,b}(X)$, the map $\mu\in \mathcal{M}_\Omega(X)\mapsto \mu(f)\in \mathbb R$ is continuous. 
\end{defn}

\begin{rem}
In \cite{Crauel}, the author defines the narrow topology on the space of random probability measures as the topology generated by the maps $\mu\mapsto \mu(f)$ for every random continuous function $f$.
    Recall that a random continuous function as defined there is a function $f$ such that for all $\omega$, $f(\cdot,\omega)$ is bounded continuous and $\|f(\cdot,\omega)\|_\infty$ is in $L_1(\Omega,\mathbf P)$.
\begin{itemize}
    \item First, we preferred to use the terminology weak topology which is more common, although both exist in literature.
    \item Second, by \cite[Lemma~3.16]{Crauel}, for random probability measures, the induced weak topology on the set of measures is the same when choosing either random bounded continuous functions or random continuous functions.
    However, in our context, the proof of this lemma does not apply and it seems that choosing different spaces of functions can yield different notions of weak topologies.
\end{itemize}
\end{rem}

\begin{lem}\label{lipschitzgenerateweaktopology}
The weak topology is generated by the maps $\mu \mapsto \mu(f)$ for every $f\in BL_\Omega(X)$, i.e. it is the smallest topology on $\mathcal{M}_\Omega(X)$ such that for every random bounded Lipschitz function $f$, $\mu\mapsto \mu(f)$ is continuous.
\end{lem}

\begin{proof}
This follows from the fact that bounded continuous functions can be approximated by bounded Lipschitz functions, see \cite[Proposition~4.9]{Crauel} for more details.
\end{proof}

We now prove the following generalization of the classical Portmanteau theorem in terms of convergent nets.
We refer to \cite[Definition~2.11, Definition~2.15]{AliprantisBorder} for more details on nets.

\begin{prop}\label{randomPortmanteau}
Let $\mu_{\alpha}=((\mu_\alpha)_\omega)_\omega$ be a net of random finite measures and let $\mu=(\mu_\omega)_\omega$ be a random finite measure.
Then, the following assertions are equivalent.
\begin{enumerate}
    \item The net $\mu_\alpha$ converges to $\mu$ for the weak topology.
    \item For all random closed set $C=\omega\mapsto C(\omega)$, $\limsup_\alpha\mu_\alpha(C)\leq \mu(C)$ and $\mu_\alpha(X\times \Omega)$ converges to $\mu(X\times \Omega)$.
    \item For all random open set $U=\omega\mapsto U(\omega)$, $\liminf_\alpha\mu_n(U)\geq \mu(U)$ and $\mu_\alpha(X\times \Omega)$ converges to $\mu(X\times \Omega)$.
\end{enumerate}
\end{prop}

\begin{proof}
Taking complements, the second and third assertions are equivalent.
Also, if $\mu_\alpha$ converges to $\mu$, applying the definition of the weak topology to the constant function $\mathbf 1$, we see that $\mu_\alpha(X\times \Omega)$ converges to $\mu(X\times \Omega)$.

Let us assume that $\mu_\alpha$ converges to $\mu$ and let $C$ be a closed random set.
For every $k\in \mathbb N$, set
$$f_k(x,\omega)=1-(1\wedge kd(x,C(\omega))).$$
Then, for every $k$, $f_k$ is a random bounded Lipschitz function and the sequence $f_k$ is non-increasing and converges to $\mathbf 1_{C(\omega)}(x)$.
Thus, for every $k$,
$$\limsup_\alpha \mu_\alpha(C)\leq \lim_\alpha \mu_\alpha(f_k)=\mu(f_k),$$
so
$$\limsup_\alpha \mu_\alpha(C)\leq \inf_k(\mu(f))=\mu(C).$$
Consequently, the first assertion implies the second one.

Assume now that for all closed random set $C$, $\limsup \mu_\alpha(C)\leq \mu(C)$ and that $\mu_\alpha(X\times \Omega)$ converges to $\mu(X\times \Omega)$.
We show that for every non-negative random bounded continuous function $f$,
\begin{equation}\label{equationPortmanteau}
\limsup \mu_\alpha(f)\leq \mu(f).
\end{equation}
Fix $m\in \mathbb N$ and set for every $0\leq k\leq m$
$$C_k(\omega)=\left \{x\in X,f(x,\omega)\geq \frac{k}{m}\|f(\cdot,\omega)\|_\infty\right \}.$$
Then, $C_k=\omega\mapsto C_k(\omega)$ is a random closed set.
By \cite[Lemma~1.4]{Crauel}, for every random finite measure $(\nu_\omega)_\omega$,
$$\frac{1}{m}\sum_{k=1}^m\nu(C_k)\leq \nu(f)\leq \frac{1}{m}\sum_{k=0}^m\nu(C_k).$$
Applying this both to $\mu_\alpha$ and $\mu$, we get

\begin{align*}
\mu(f)\geq \frac{1}{m}\sum_{k=1}^m\mu(C_k)&\geq \limsup_\alpha\frac{1}{m}\sum_{k=1}^m\mu_\alpha(C_k)\\
&=\limsup_\alpha\left (\frac{1}{m}\sum_{k=0}^m\mu_\alpha(C_k)-\frac{\mu_\alpha(X\times \Omega)}{m}\right )\\
&\geq \limsup_\alpha\mu_\alpha(f)-\frac{\mu(X\times\Omega)}{m}.
\end{align*}
Since $m$ is arbitrary, this proves~(\ref{equationPortmanteau}).
Using that $\mu_\alpha(\mathbf 1)=\mu_\alpha(X\times \Omega)$ converges to $\mu(\mathbf 1)=\mu(X\times \Omega)$ and applying~(\ref{equationPortmanteau}) to the function $\Vert f\Vert_{\infty}-f$, we get that
$\mu_\alpha(f)$ converges to $\mu(f)$.
This is true for all non-negative random bounded continuous function, so $\mu_\alpha$ converges to $\mu$ for the weak topology.
\end{proof}

\begin{rem}
By Lemma~\ref{measuredeterminedbylipschitz}, the weak topology is Hausdorff.
We do not attempt to study metrizability of the weak topology in here to avoid lengthily arguments,
but in \cite[Theorem~4.16]{Crauel}, the author proves that the weak topology on the space of random probability measures is metrizable, provided that the probability space $(\Omega,\mathcal{F},\mathbf P)$ is countably generated (mod. $\mathbf P$).
This might also holds in our situation.
\end{rem}

\subsection{A compactness criterion for random finite measures}
In all this section, $\mathcal{M}(X)$ and $\mathcal{M}_\Omega(X)$ are endowed with the weak topology.
Let us recall the following definition.
\begin{defn}\label{deftightnonrandom}
A subset $M$ of $\mathcal{M}(X)$ is tight if for every $\epsilon>0$, there exists a compact subset of $X$ such that for every $\mu\in M$, we have
$$\mu(K^c)\leq \epsilon.$$
\end{defn}

Following \cite{Crauel}, we define tightness for random finite measures as tightness under the projection map $\pi_X$.
\begin{defn}\label{deftightrandom}
A subset $M_\Omega$ of $\mathcal{M}_\Omega(X)$ is tight if $\pi_X(M_\Omega)$ is tight, i.e. for every $\epsilon>0$, there exists a compact subset of $X$ such that for all $(\mu_\omega)_\omega\in M_\Omega$, we have
$$\mathbf E[\mu_\omega(K^c)]\leq \epsilon.$$
\end{defn}

The classical Prokhorov theorem states that a set $M$ of probability measures on $X$ is relatively compact if and only if it is tight.
The following generalizes this result to finite measures.

\begin{thm}\label{classicalProkhorov}\textbf{Prokhorov Theorem for finite measures} \cite[Lemma~4.4]{K17}.
A subset $M$ of $\mathcal{M}(X)$ is relatively compact if and only if $M$ is tight and uniformly bounded, in the sense that $\sup_{\mu\in M} \mu(X)$ is finite.
\end{thm}

Our goal is to generalize this to random finite measures.
Unfortunately, we will not get a necessary and sufficient condition for compactness as in the Prokhorov Theorem but only a sufficient condition.

We follow the strategy of \cite[Theorem~4.4]{Crauel} and first prove the following representation result.
Recall that $BL_\Omega(X)$ denotes the set of random bounded Lipschitz functions and $BL(X)$ denotes the set of bounded Lipschitz functions on $X$.
Then, if $f\in BL(X)$, $f$ can be viewed as an element of $BL_\Omega(X)$ by setting $f(x,\omega)=f(x)$.
Similarly, if $f$ is $\mathbf P$ -essentially bounded, i.e.\ $f\in L_\infty(\Omega,\mathbf P)$, then $f$ can be viewed as an element of $BL_\Omega(X)$ by setting $f(x,\omega)=f(\omega)$.
If $L:BL_\Omega(X)\to \mathbb R$ is a function, we denote by $\pi_X(L)$, respectively $\pi_\Omega(L)$ its restriction to bounded Lipschitz functions, respectively to $\mathbf P$-essentially bounded functions.

\begin{lem}\label{StoneDaniell}
Let $L:BL_\Omega(X)\to \mathbb R$. be a function.
Assume that the following conditions hold
\begin{enumerate}[(a)]
    \item $L$ is linear,
    \item $L$ is non-negative, i.e.\ for every non-negative function $f$ in $BL_\Omega(X)$, we have $L(f)\geq 0$,
    \item there exists $\kappa\in \mathcal{M}(X)$ such that for every $f\in BL(X)$,
    we have
    $$\pi_X(L)(f)=\kappa(f)=\int fd\kappa,$$
    \item there exists a constant $C>0$ such that for every $f\in L_\infty(\Omega,\mathbf P)$, we have
    $$\frac{1}{C}\pi_\Omega(L)(f)\leq \mathbf E [f]\leq C\pi_\Omega(L)(f).$$
\end{enumerate}
Then, there exists a random finite measure $\mu=(\mu_\omega)_\omega$ such that $L(f)=\mu(f)$ for every $f\in BL_\Omega(X)$.
\end{lem}

\begin{proof}
The proof relies on the general Stone-Daniell representation theorem.
We claim that if $f_n$ is a non-increasing sequence of functions of $BL_\Omega(X)$ converging to 0, then $L(f_n)$ is non-increasing and converges to 0.
By \cite[Theorem~4.11]{Crauel}, we deduce that there exists a measure $\mu$ on $X\times \Omega$ such that $L(f)=\mu(f)$ for every $f\in BL_\Omega(X)$.
Conditions~(c) and~(d) ensure that $\mu$ is finite.
Moreover, the $\Omega$-marginal $\pi_\Omega(\mu)$ of $\mu$ is $\pi_\Omega(L)$ and so the fourth condition shows $\pi_\Omega(\mu)$ and $\mathbf P$ are absolutely continuous with respect to each other.
Therefore, we can apply Proposition~\ref{disintegrationrandom}, so that $\mu$ is defined by a random finite measure.

We just need to prove the claim to conclude the proof.
Let $f_n$ be a non-increasing sequence of random bounded Lipschitz functions converging to 0.
Since $L$ is linear and non-negative, we get that $L(f_n)$ is non-inscreasing.
We need to prove that $L(f_n)$ converges to 0.
Fix $\epsilon>0$.
Let $K$ be a compact such that $\kappa(K^c)\leq \epsilon$.
For every $n$, there exists $C_n$ such that $\|f_n(\cdot,\omega)\|_{BL}\leq C_n$, for $\mathbf P$-almost every $\omega$.
Set $\delta_n=\epsilon/C_n$ and consider the function
$$\chi_n:x\mapsto \chi_n(x)=1-\big(1\wedge \delta_n^{-1}d(x,K)\big).$$
Then, $\chi_n$ vanishes outside the $\delta_n$-neighborhood of $K$.
Moreover, the function
$g_n:(x,\omega)\mapsto f_n(x,\omega)\chi_n(x)$ is in $BL_\Omega(X)$.
Now, $f_n=g_n+f_n(1-\chi_n)$ and so
\begin{equation}\label{equationStoneDaniell1}
    L(f_n)=L(g_n)+L(f_n(1-\chi_n)).
\end{equation}

We first deal with the second term in the right-hand side of~(\ref{equationStoneDaniell1}).
Since $f_n$ is non-increasing, we have that $f_n(x,\omega)\leq \|f_1(\cdot,\omega)\|_\infty$ Moreover, there exists $M\geq 0$ such that the event $A=\{\|f_1(\cdot,\omega)\|_\infty\leq M\}$ satisfies $\mathbf P(A^c)=0$.
Note that we can write $f_n(1-\chi_n)\leq M(1-\chi_n)\mathbf 1_A+\|f_1(\cdot,\omega)\|\mathbf 1_{A^c}$.
By Condition~(d), $L(\|f_1(\cdot,\omega)\|\mathbf 1_{A^c})=0$.
Thus,
$$L(f_n(1-\chi_n))\leq ML(1-\chi_n)=M\int (1-\chi_n)d\kappa.$$
Note that $1-\chi_n$ vanishes in $K$, so
\begin{equation}\label{equationStoneDaniell2}
L(f_n(1-\chi_n))\leq M\kappa (K^c)\leq M\epsilon.
\end{equation}

We now deal with the first term in the right-hand side of~(\ref{equationStoneDaniell1}).
For every $\omega$ such that $\|f_n(\cdot,\omega)\|_{BL}\leq C_n$ and for every $x,y\in X$, we have
$$|g_n(x,\omega)-g_n(y,\omega)|\leq C_nd(x,y)+2f_n(y,\omega).$$
If $x$ is in the $\delta_n$-neighborhood of $K$, choose $y\in K$ such that
$d(x,y)\leq \delta_n$.
Then, by what precedes,
$$g_n(x,\omega)\leq 3\sup_{y\in K}f_n(y,\omega)+C_n\delta_n\leq 3\sup_{y\in K}f_n(y,\omega)+\epsilon.$$
Since $\chi_n$ vanishes outside the $\delta_n$-neighborhood of $K$, this yields
$$\sup_{x\in X}g_n(x,\omega)\leq 3\sup_{y\in K}f_n(y,\omega)+\epsilon.$$
We write $h_n(\omega)=\sup_{y\in K}f_n(y,\omega)$.
The event $A_n= \{\omega$, $\|f_n(\cdot,\omega)\|_{BL}>C_n\}$ satisfies $\mathbf P(A_n)=0$, hence the same manipulation as above shows that
\begin{equation}\label{equationStoneDaniell3}
L(g_n)\leq 3L(h_n)+\epsilon L(\mathbf 1)\leq 3C\mathbf E(h_n)+\kappa(X)\epsilon,
\end{equation}
using Conditions~(c) and~(d).

Combining~(\ref{equationStoneDaniell1}),~(\ref{equationStoneDaniell2}) and~(\ref{equationStoneDaniell3}), we get
$$L(f_n)\leq 3C\mathbf E(h_n)+(\kappa(X)+M)\epsilon.$$
Using that $f_n$ is non-increasing and converges to 0, that for all $\omega$, $f_n(\cdot,\omega)$ is continuous and that $K$ is compact, we get that $h_n$ is non-increasing and converges to 0.
By monotone convergence, $\mathbf E(h_n)$ converges to 0,
so for large enough $n$, $\mathbf E(h_n)\leq \epsilon$.
Thus, for large enough $n$,
$$L(f_n)\leq (3C+\kappa(X)+M)\epsilon.$$
Since $\epsilon$ is arbitrary, this concludes the proof of the claim.
\end{proof}

The main ingredient in proving a compactness criterion for random finite measures is the following proposition.

\begin{prop}
Let $C>0$.
If $K$ is compact in $\mathcal{M}(X)$, then
$$\pi_X^{-1}(K)\cap \left \{\mu \in \mathcal{M}_\Omega(X), \frac{1}{C}\pi_\Omega(\mu)\leq \mathbf P\leq C\pi_\Omega(\mu)\right \}$$ is compact in $\mathcal{M}_\Omega(X)$.
\end{prop}

\begin{proof}
Let $f\in BL_\Omega(X)$ and set
$$M_f^+=\sup_{\mu\in \pi_X^{-1}(K)}\mu(f),$$
$$M_f^-=\inf_{\mu\in \pi_X^{-1}(K)}\mu(f).$$
Recall that if $f\in BL_\Omega(X)$, then in particular, $f\in C_{\Omega,b}(X)$ and so for $\mathbf P$-almost every $\omega$, $|f(\cdot,\omega)|$ is bounded and
$\mathbf E[\|f(\cdot,\omega)\|_\infty]$ is bounded.
Thus, $\pi_X(\mu)(f)$ is finite.
Since $K$ is compact, this proves that $M_f^+$ and $M_f^-$ are finite.

We set
$$\mathbf C=\prod_{f\in BL_\Omega(X)}[M_f^-,M_f^+]$$
and we endow $\mathbf C$ with the product topology.
Then $\mathbf C$ is compact by the Tychonoff Theorem \cite[Theorem~2.61]{AliprantisBorder}.
We can identify $\mathbf C$ with the set of all functions $L:BL_\Omega(X)\to \mathbb R$ such that $L(f)\in [M_f^-,M_f^+]$.
A neighborhood basis of an element $L_0$ of $\mathbf C$ is given by
$$U_{\delta}(f_1,...f_n)(L_0)=\{L\in \mathbf C,|L(f_k)-L_0(f_k)|<\delta,1\leq k\leq n\},$$
where $\delta>0$ and $f_1,...,f_n\in BL_\Omega(X)$.
We define the map
$$I:\mu\in \pi_X^{-1}(K)\mapsto (\mu(f))_{f\in BL_\Omega(X)}\in \mathbf C.$$
Then, $I$ is one-to-one and continuous.
Moreover, by Lemma~\ref{lipschitzgenerateweaktopology}, a neighborhood basis of an element $\mu_0\in \pi_X^{-1}(K)$ for the weak topology is given by
$$V_\delta(f_1,...,f_n)(\mu_0)=\{\mu \in \pi_X^{-1}(K), |\mu(f_k)-\mu_0(f_k)|<\delta,1\leq k\leq n\},$$
where $\delta>0$ and $f_1,...,f_n\in BL_\Omega(X)$.
Since
$$I\big (V_\delta(f_1,...,f_n)(\mu_0)\big)=U_\delta(f_1,...,f_n)(I(\mu_0)),$$
we see that $I$ is an open map and thus a homeomorphism onto its image in $\mathbf C$.

We make a similar construction for non-random finite measures.
Recall that $BL(X)$ is the set of (non-random) bounded Lipschitz functions on $X$ and that $BL(X)$ can be seen as a subset of $BL_\Omega(X)$ by setting $f(x,\omega)=f(x)$ for any $f\in BL(X)$.
We define
$$\mathbf c=\prod_{f\in BL(X)}[M_f^-,M_f^+]$$
and
$$i(\kappa)=(\kappa(f))_{f\in BL(X)}.$$
Then, $i$ is a homeomorphism onto its image in $\mathbf c$.
If $L\in \mathbf C$, denote by $\pi(L)\in \mathbf c$ its restriction to non-random bounded Lipschitz functions.

For every $f\in BL(X)$, seen as an element of $BL_\Omega(x)$ and for every random finite measure $\mu$, we have
$$\mu(f)=\mathbf E\left [\int f(x)d\mu_\omega(x)\right ]=\pi_X(\mu)(f).$$
Thus, $\pi \circ I=i\circ \pi_X$.
Now, fix $C>0$.
By Lemma~\ref{StoneDaniell}, we have
\begin{align*}
&I\left (\pi_X^{-1}(K)\cap \left \{\mu, \frac{1}{C}\leq \pi_\Omega(\mu)\leq \mathbf P\leq C\pi_\Omega(\mu)\right\}\right)\\
=&\hspace{.2cm}\pi^{-1}(i(K))\cap \{L \text{ linear}\}\cap \{ L\text{ non-negative}\}\\
&\hspace{.3cm}\cap \left \{L, \frac{1}{C}\pi_\Omega(L)\leq \mathbf E[\cdot]\leq C\pi_\Omega(L)\right \}.
\end{align*}
The four sets on the right-hand side are closed.
Since $\mathbf C$ is compact and $I$ is a homeomorphism onto its image, we get that
$$ \pi_X^{-1}(K)\cap \left \{\mu, \frac{1}{C}\leq \pi_\Omega(\mu)\leq \mathbf P\leq C\pi_\Omega(\mu)\right\}$$
is compact.
\end{proof}

%\begin{definition}
%A subset $M_\Omega$ of $\mathcal{M}_\Omega(X)$ is called uniformly bounded in expectation if there exists $C\geq 0$ such that
%for every $\mu\in M_\Omega$,
%$$\sup_{(\mu_\omega)_\omega\in M_\Omega}\mathbf E[\mu_\omega(X)]\leq C.$$
%\end{definition}

\begin{cor}\label{coroProkhorov}
Consider a subset $M_\Omega$ of $\mathcal{M}_\Omega(X)$.
Assume that $M_\Omega$ is tight.
Also assume that there exists $C\geq 0$ such that for every $\mu\in M_\Omega$,
$$\frac{1}{C}\pi_\Omega(\mu)\leq \mathbf P \leq C \pi_\Omega(\mu).$$
Then, $M_\Omega$ is relatively compact.
\end{cor}

\begin{proof}
If $M_\Omega$ satisfies the assumptions of the corollary, then $\pi_X(M_\Omega)\subset \mathcal{M}(X)$ is tight.
Also, the condition $\pi_\Omega(\mu)\leq C \mathbf P$ can be reformulated by
$$\frac{d\pi_\omega(\mu)}{d\mathbf P}(\omega)=\mu_\omega(X)\leq C$$
$\mathbf P$-almost surely.
In particular,
$$\pi_X\mu(X)=\mathbf E[\mu_\omega(X)]\leq C,$$
so $\pi_X(\mu)(X)$ is uniformly bounded.
By Theorem~\ref{classicalProkhorov}, $\pi_X(M_\Omega)$ is relatively compact, i.e.\ $\mathbf{cl}(\pi_X(M_\Omega))$ is compact.
Moreover,
$$M_\Omega\subset \pi_X^{-1}\big(\mathbf{cl}(\pi_X(M_\Omega))\big )\cap \left \{\mu,\frac{1}{C}\pi_\Omega(\mu)\leq \mathbf P\leq C\pi_\Omega(\mu)\right \}.$$
Thus, $\mathbf{cl}(M_\Omega)$ is compact.
\end{proof}





This corollary does not give a necessary and sufficient condition for compactness, but only a sufficient one.
On the contrary, \cite[Theorem~4.4]{Crauel} gives a necessary and sufficient condition for compactness in the context of random probability measures : a set is relatively compact if and only if it is tight.
However, assuming that for $\mathbf P$-almost every $\omega$, $\mu_\omega$ is a probability measure, we get that $\pi_X(\mu)$ is a probability measures and that the $\Omega$-marginal of $\mu$ is exactly $\mathbf P$.
This ensures that the others conditions of Corollary~\ref{coroProkhorov} are automatically satisfied and so tightness is sufficient to get compactness.

Recall that the Radon-Nikodym derivative of $\pi_\Omega(\mu)$ with respect to $\mathbf P$ is given by $\mu_\omega(X)$.
In our context, on the one hand, the assumption
$$\frac{1}{C}\pi_\Omega(\mu)\leq \mathbf P\leq C\pi_\Omega(\mu)$$
which is equivalent to
$$\frac{1}{C}\leq \mu_\omega(X)\leq C$$
for $\mathbf P$-almost every $\omega$
seems strong.
Assuming only $\mathbf P$-regularity, there is no reason for this Radon-Nikodym derivative to be $\mathbf P$-essentially bounded from above and below.
Even when restricting our attention to random finite measures satisfying this property, the maps
$$\mu\mapsto \mathbf P-\textbf{ess}\sup \mu_\omega(X)$$
and
$$\mu\mapsto \mathbf P-\textbf{ess}\inf \mu_\omega(X)$$
have no reason to be continuous.
Thus, compactness for the weak topology does not seem to imply the existence of a uniform $C$ such that the condition 
$$\frac{1}{C}\pi_\Omega(\mu)\leq \mathbf P\leq C\pi_\Omega(\mu)$$
holds.
On the other hand, it is not realistic to expect a sufficient condition for compactness without assuming anything on the $\Omega$-marginals.

\medskip
In the particular case where $X$ has an isolated point $x_0$, then the following trick allows us to weaken a bit this assumption.
For any subset of $M_\Omega$ of $\mathcal{M}_\Omega(X)$, define
$$\widetilde{M}_\Omega=\left \{\nu\in \mathcal{M}_\Omega(X), \nu=\mathbf D(x_0)+\mu,\mu\in M_\Omega\right \},$$
where $\mathbf D$ is the Dirac measure on $x_0$.
Assume that for every $\mu\in M_\Omega$, we have $\mu_\omega(X)\leq C$ for $\mathbf P$-almost every $\omega$, for some constant $C$.
Then for every $\nu\in \widetilde{M}_\Omega$,
\begin{equation}\label{coboundedMtilde}
1\leq \nu_\omega(X)\leq 1+C
\end{equation}
for $\mathbf P$-almost every $\omega$.

\begin{cor}\label{coroProkhorovisolated}
Assume that $X$ has an isolated point $x_0$.
Consider a subset $M_\Omega$ of $\mathcal{M}_\Omega(X)$.
Assume that $M_\Omega$ is tight and that there exists $C\geq 0$ such that for every $\mu\in M_\Omega$,
$\pi_\Omega(\mu)\leq C \mathbf P$.
Then, $M_\Omega$ is relatively compact.
\end{cor}

\begin{proof}
The set $M_\Omega$ is tight, so there is a compact $K_0$ such that for every $\mu\in M_\Omega$,
$\pi_X(\mu)(K_0^c)\leq \epsilon$.
Then, the set $K=K_0\cup \{x_0\}$ is also compact and for every $\nu\in \widetilde{M}_\Omega$,
$\pi_X(\nu)(K^c)\leq \epsilon$.
Thus, $\widetilde{M}_\Omega$ is also tight.
By~(\ref{coboundedMtilde}), $\widetilde{M}_\Omega$ satisfies the assumptions of Corollary~\ref{coroProkhorov},
so it is relatively compact.

Let $\nu\in \mathbf{cl}(\widetilde{M}_\Omega)$.
By \cite[Theorem~2.14]{AliprantisBorder}, there exists a net $(\nu_\alpha)_{\alpha\in A}$ converging to $\nu$.
Assume by contradiction that $\mathbf P(\nu_\omega(\{x_0\})<1)>0$.
Then, there exists $c<1$ such that
the event $A=\{\nu_\omega(\{x_0\})\leq c\}$ satisfies $\mathbf P(A)>0$.
Since $x_0$ is isolated and $X$ is Hausdorff, $\{x_0\}$ is both closed and open, hence the function $\mathbf 1_{x_0}$ is bounded continuous
and $(x,\omega)\mapsto \mathbf{1}_A(\omega)\mathbf{1}_{x_0}(x)$ is a random bounded continuous function.
Since for every $\alpha$, $\nu_\alpha(\{x_0\})\geq 1$, applying convergence to this function, we get that
$$\mathbf P(A)\leq \mathbf E[\mathbf 1_A\nu_\alpha(\{x_0\})]\underset{\alpha\to \infty}{\longrightarrow}\mathbf E[\mathbf 1_A\nu(\{x_0\})]\leq c\mathbf P(A),$$
which is a contradiction.
Thus,
$$\mu_\omega=\nu_\omega-\mathbf D(x_0)$$
is a well-defined random finite measure on $X$.
In other words, the map
$$F:\nu\in \mathbf{cl}(\widetilde{M}_\Omega)\mapsto \nu-\mathbf{D}(x_0)\in \mathcal{M}_\Omega(X)$$
is well defined.
Moreover, it is continuous, therefore $F(\mathbf{cl}(\widetilde{M}_\Omega))$ is compact and by Lemma~\ref{measuredeterminedbylipschitz}, the weak topology is Hausdorff, so $F(\mathbf{cl}(\widetilde{M}_\Omega))$ is closed.
Now, $M_\Omega\subset F(\mathbf{cl}(\widetilde{M}_\Omega))$ and so  $\mathbf{cl}(M_\Omega)\subset F(\mathbf{cl}(\widetilde{M}_\Omega))$.
Thus, $\mathbf{cl}(M_\Omega)$ is compact, i.e.\ $M_\Omega$ is relatively compact.
\end{proof}




\bibliographystyle{alpha}
\bibliography{main}
\end{document}

%------------------------------------------------------------------------------
% End of journal.tex
%------------------------------------------------------------------------------
%%% Local Variables:
%%% TeX-engine: default
%%% End:
\begin{thebibliography}{GGPY21}

\bibitem[AB06]{AliprantisBorder}
Charalambos Aliprantis and Kim Border.
\newblock {\em Infinite Dimensional Analysis: a hitchhiker's guide}.
\newblock Springer, 2006.

\bibitem[AN04]{AN04}
Krishna Athreya and Peter Ney.
\newblock {\em Branching processes}.
\newblock Dover Publications, Inc., Mineola, NY, 2004.
\newblock Reprint of the 1972 original [Springer, New York; MR0373040].

\bibitem[Anc88]{Ancona}
Alano Ancona.
\newblock Positive harmonic functions and hyperbolicity.
\newblock In {\em Potential theory-surveys and problems}, pages 1--23. Lecture
  notes in mathematics, Springer, 1988.

\bibitem[BB07]{BlachereBrofferio}
S\'ebastien Blach\`ere and Sara Brofferio.
\newblock Internal diffusion limited aggregation on discrete groups having
  exponential growth.
\newblock {\em Probab. Theory Relat. Fields}, 137:323–343, 2007.

\bibitem[Bil86]{Bill86}
Patrick Billingsley.
\newblock {\em Probability and Measure}.
\newblock John Wiley and Sons, second edition, 1986.

\bibitem[BM12]{BM12}
Itai Benjamini and Sebastian M\"uller.
\newblock On the trace of branching random walks.
\newblock {\em Groups Geom. Dyn.}, 6:231--247, 2012.

\bibitem[Bow12]{Bowditch}
Brian Bowditch.
\newblock Relatively hyperbolic group.
\newblock {\em Int. J. Algebra Comput.}, 22:1--66, 2012.

\bibitem[BP94]{BP94}
Itai Benjamini and Yuval Peres.
\newblock Markov chains indexed by trees.
\newblock {\em Ann. Probab.}, 22(1):219--243, 1994.

\bibitem[Car88]{Cartwright1}
Donald Cartwright.
\newblock Some examples of random walks on free products of discrete groups.
\newblock {\em Ann. di Mat. Pura ed Appl.}, 151:1--15, 1988.

\bibitem[Car89]{Cartwright2}
Donald Cartwright.
\newblock On the asymptotic behaviour of convolution powers of probabilities on
  discrete groups.
\newblock {\em Monatsh. Math.}, 107:287--290, 1989.

\bibitem[CG12]{CandelleroGilch}
Elisabetta Candellero and Lorenz Gilch.
\newblock Phase transitions for random walk asymptotics on free products of
  groups.
\newblock {\em Random Struct. Algorithms}, 40:150--181, 2012.

\bibitem[CGM12]{CandelleroGilchMuller}
Elisabetta Candellero, Lorenz Gilch, and Sebastian M\"uller.
\newblock Branching random walks on free products of groups.
\newblock {\em Proc. London Math. Soc.}, 104:1085--1120, 2012.

\bibitem[Cor17]{Cordes}
Matthew Cordes.
\newblock Morse boundaries of proper geodesic metric spaces.
\newblock {\em Groups Geom. Dyn.}, 11:1281–1306, 2017.

\bibitem[Cra02]{Crauel}
Hans Crauel.
\newblock {\em Random probability measures on Polish spaces}.
\newblock CRC Press, 2002.

\bibitem[DFW19]{DFW}
Fran\c{c}ois Dahmani, David Futer, and Daniel~T. Wise.
\newblock Growth of quasiconvex subgroups.
\newblock {\em Math. Proc. Cambridge Philos. Soc.}, 167(3):505--530, 2019.

\bibitem[DG20]{DG20}
Matthieu Dussaule and Ilya Gekhtman.
\newblock Entropy and drift for word metrics on relatively hyperbolic groups.
\newblock {\em Groups Geom. Dyn.}, 14:1455--1509, 2020.

\bibitem[DG21]{DG21}
Matthieu Dussaule and Ilya Gekhtman.
\newblock Stability phenomena for {M}artin boundaries of relatively hyperbolic
  groups.
\newblock {\em Probab. Theory Relat. Fields}, 179:201--259, 2021.

\bibitem[DOP00]{DOP}
Fran\c{c}oise Dal'bo, Jean-Pierre Otal, and Marc Peign\'e.
\newblock S\'eries de {P}oincar\'e des groupes g\'eom\'etriquement finis.
\newblock {\em Isr. J. Math.}, 118:109--124, 2000.

\bibitem[DPPS11]{DPPS}
Fran\c{c}oise Dal'bo, Marc Peign\'e, Jean-Claude Picaud, and Andrea Sambusetti.
\newblock On the growth of quotients of {K}leinian groups.
\newblock {\em Ergod. Theory Dyn. Syst.}, 31:835--851, 2011.

\bibitem[DPT22]{DPT}
Matthieu Dussaule, Marc Peign\'e, and Samuel Tapie.
\newblock A local limit theorem for convergent random walks on relatively
  hyperbolic groups.
\newblock arXiv:2202.11339, 2022.

\bibitem[DS05]{DrutuSapir}
Cornelia Dru\c{t}u and Mark Sapir.
\newblock Tree graded spaces and asymptotic cones of groups.
\newblock {\em Topology}, 44:959–1058, 2005.
\newblock with an {A}ppendix by {D}enis {O}sin and {M}ark {S}apir.

\bibitem[Dus22a]{DussauleLLT1}
Matthieu Dussaule.
\newblock Local limit theorems in relatively hyperbolic groups {I} : rough
  estimates.
\newblock {\em Ergod. Theory Dyn. Syst.}, 42:1926 --1966, 2022.

\bibitem[Dus22b]{DussauleLLT2}
Matthieu Dussaule.
\newblock Local limit theorems in relatively hyperbolic groups {I}{I} : the non
  spectrally degenerate case.
\newblock {\em Compos. Math.}, 158:764--830, 2022.

\bibitem[DY20]{DussauleYang}
Matthieu Dussaule and Wenyuan Yang.
\newblock The {H}ausdorff dimension of the harmonic measure for relatively
  hyperbolic groups.
\newblock arXiv:2010.07671, 2020.

\bibitem[Far98]{Farb}
Benson Farb.
\newblock Relatively hyperbolic groups.
\newblock {\em GAFA}, 8:810--840, 1998.

\bibitem[FK60]{FurstenbergKesten}
Harry Furstenberg and Harry Kesten.
\newblock Products of random matrices.
\newblock {\em Ann. Math. Statist.}, 31:457--469, 1960.

\bibitem[Flo80]{Floyd}
William Floyd.
\newblock Group completions and limit sets of {K}leinian groups.
\newblock {\em Invent. Math.}, 57:205--2018, 1980.

\bibitem[Fol99]{Folland}
Gerald Folland.
\newblock {\em Real analysis : Modern techniques and their applications}.
\newblock John Wiley and sons, 1999.

\bibitem[Fre45]{Freudenthal}
Hans Freudenthal.
\newblock \"{U}ber die enden diskreter ra\"ume und gruppen.
\newblock {\em Comment. Math. Helvetici}, 17:1--38, 1945.

\bibitem[Ger12]{Gerasimov}
Victor Gerasimov.
\newblock Floyd maps for relatively hyperbolic groups.
\newblock {\em GAFA}, 22:1361--1399, 2012.

\bibitem[GGPY21]{GGPY}
Ilya Gekhtman, Victor Gerasimov, Leonid Potyagailo, and Wenyuan Yang.
\newblock Martin boundary covers {F}loyd boundary.
\newblock {\em Invent. Math.}, 223:759–809, 2021.

\bibitem[GL13]{GouezelLalley}
S\'ebastien Gou\"ezel and Steven Lalley.
\newblock Random walks on co-compact {F}uchsian groups.
\newblock {\em Ann. Sci. \'Ec. Norm. Sup\'er.}, 46:129--173, 2013.

\bibitem[GM06]{GM06}
Nina Gantert and Sebastian M\"uller.
\newblock The critical branching {M}arkov chain is transient.
\newblock {\em Markov Process. Related Fields}, 12:805--814, 2006.

\bibitem[Gou14]{Gouezel}
S\'ebastien Gou\"ezel.
\newblock Local limit theorem for symmetric random walks in {G}romov-hyperbolic
  groups.
\newblock {\em JAMS}, 27:893--928, 2014.

\bibitem[GP13]{GePoJEMS}
Victor Gerasimov and Leonid Potyagailo.
\newblock Quasi-isometric maps and {F}loyd boundaries of relatively hyperbolic
  groups.
\newblock {\em JEMS}, 15:2115--2137, 2013.

\bibitem[GP15]{GePoGGD}
Victor Gerasimov and Leonid Potyagailo.
\newblock Non-finitely generated relatively hyperbolic groups and {F}loyd
  quasiconvexity.
\newblock {\em Groups, Geom. Dyn.}, 9:369--434, 2015.

\bibitem[Gui80]{Guivarch}
Yves Guivarc'h.
\newblock Sur la loi des grands nombres et le rayon spectral d’une marche
  al\'eatoire.
\newblock In {\em Conference on Random Walks}, pages 47--98. Astérisque, vol.
  74, Soc. Math. France, 1980.

\bibitem[HL00]{HL00}
Irene Hueter and Steven Lalley.
\newblock Anisotropic branching random walks on homogeneoustrees.
\newblock {\em Probab. Theory Relat. Fields}, 116:57--88, 2000.

\bibitem[Hru10]{Hruska}
Geoffrey Hruska.
\newblock Relative hyperbolicity and relative quasiconvexity for countable
  groups.
\newblock {\em Algebraic Geom. Topol.}, 10:1807–1856, 2010.

\bibitem[Hut20]{Hutchcroft20}
Tom Hutchcroft.
\newblock Non-intersection of transient branching random walks.
\newblock {\em Probab. Theory Relat. Fields}, 178:1--23, 2020.

\bibitem[Kai87]{Kaimanovichsymmetricspaces}
Vadim Kaimanovich.
\newblock Lyapunov exponents, symmetric spaces and a multiplicative ergodic
  theorem for semisimple {L}ie groups.
\newblock {\em Zap. Nauchn. Sem. Leningrad. Otdel. Mat. Inst. Steklov (LOMI)},
  184:29--46, 1987.

\bibitem[Kai00]{Kaimanovich}
Vadim Kaimanovich.
\newblock The {P}oisson formula for groups with hyperbolic properties.
\newblock {\em Ann. Math.}, 152:659--692, 2000.

\bibitem[Kal17]{K17}
Olav Kallenberg.
\newblock {\em Random measures, theory and applications}, volume~77 of {\em
  Probability Theory and Stochastic Modelling}.
\newblock Springer, Cham, 2017.

\bibitem[Kar03]{Karlssonfreesubgroups}
Anders Karlsson.
\newblock Free subgroups of groups with nontrivial {F}loyd boundary.
\newblock {\em Commun. Algebra}, 31:5361--5376, 2003.

\bibitem[KM99]{KarlssonMargulis}
Anders Karlsson and Gregory Margulis.
\newblock A multiplicative ergodic theorem and nonpositively curved spaces.
\newblock {\em Commun. Math. Phys.}, 208:107--123, 1999.

\bibitem[KW22]{KaimanovichWoess}
Vadim Kaimanovich and Wolfgang Woess.
\newblock Limit distributions of branching random walks.
\newblock arXiv:2205.13615, 2022.

\bibitem[Led01]{Ledrappier}
Fran\c{c}ois Ledrappier.
\newblock Some asymptotic properties of random walks in free groups.
\newblock In {\em Topics in Probability and Lie Groups: Boundary Theory}, pages
  117--152. CRM Proceedings and Lecture Notes, 2001.

\bibitem[Lig96]{Liggett96}
Thomas~M. Liggett.
\newblock Branching random walks and contact processes on homogeneous trees.
\newblock {\em Probab. Theory Related Fields}, 106(4):495--519, 1996.

\bibitem[MT18]{MaherTiozzo}
Joseph Maher and Giulio Tiozzo.
\newblock Random walks on weakly hyperbolic groups.
\newblock {\em J. für die Reine und Angew. Math.}, 742:187--239, 2018.

\bibitem[Osi06]{Osin}
Denis Osin.
\newblock Relatively hyperbolic groups: Intrinsic geometry, algebraic
  properties, and algorithmic problems.
\newblock {\em Mem. Am. Math. Soc.}, 20, 2006.

\bibitem[Pei11]{Peigne}
M.~Peign\'e.
\newblock On some exotic {S}chottky groups.
\newblock {\em Discrete Contin. Dyn. Syst.}, 118(31):559 -- 579, 2011.

\bibitem[PS18]{PS18}
Vincent {Pit} and Barbara {Schapira}.
\newblock {Finitude des mesures de {G}ibbs sur les vari\'et\'es non compactes
  \`a courbure n\'egative pinc\'ee}.
\newblock {\em {Ann. Institut Fourier}}, 68:457--510, 2018.

\bibitem[PTV20]{PTV}
Marc Peign\'e, Samuel Tapie, and Pierre Vidotto.
\newblock Orbital counting for some convergent groups.
\newblock {\em Ann. de l'Institut Fourier}, 70:1307--1340, 2020.

\bibitem[PW02]{PapasogluWhyte}
Panos Papasoglu and Kevin Whyte.
\newblock Quasi-isometries between groups with infinitely many ends.
\newblock {\em Comment. Math. Helvetici}, 77:133–144, 2002.

\bibitem[PY19]{PY}
Leonid Potyagailo and Wenyuan Yang.
\newblock Hausdorff dimension of boundaries of relatively hyperbolic groups.
\newblock {\em Geom. Topol.}, 23:1779--1840, 2019.

\bibitem[QRT20]{QRT}
Yulan Qing, Kasra Rafi, and Giulio Tiozzo.
\newblock Sublinearly {M}orse boundary {I}{I}: {P}roper geodesic spaces.
\newblock arXiv:2011.03481, 2020.

\bibitem[{Rob}03]{Rob}
Thomas {Roblin}.
\newblock {Ergodicit\'e et \'equidistribution en courbure n\'egative}.
\newblock {\em {M\'em. Soc. Math. Fr., Nouv. S\'er.}}, 95:96, 2003.

\bibitem[Sha09]{Shah}
Jay Shah.
\newblock Hausdorff dimension and its applications.
\newblock available at http://www.math.uchicago.edu/\\
  ~may/VIGRE/VIGRE2009/REUPapers/Shah.pdf, 2009.

\bibitem[Shi12]{Shi}
Zhan Shi.
\newblock Branching random walks.
\newblock In {\em \'Ecole d'\'et\'e de Probabilit\'es de Saint-Flour XLII},
  volume 2151. Lecture notes in mathematics, Springer, 2012.

\bibitem[Sis12]{Sistometricrelhyp}
Alessandro Sisto.
\newblock On metric relative hyperbolicity.
\newblock arXiv:1210.8081, 2012.

\bibitem[Sis13]{Sistoprojections}
Alessandro Sisto.
\newblock Projections and relative hyperbolicity.
\newblock {\em Enseign. Math.}, 59:165–181, 2013.

\bibitem[ST21]{ST21}
Barbara {Schapira} and Samuel {Tapie}.
\newblock {Regularity of entropy, geodesic currents and entropy at infinity}.
\newblock {\em {Ann. Sci. \'Ec. Norm. Sup\'er.}}, 54:1--68, 2021.

\bibitem[Sta71]{Stallings}
John Stallings.
\newblock {\em Group theory and three-dimensional manifolds}.
\newblock Yale University Press, 1971.

\bibitem[SW79]{ScottWall}
Peter Scott and Terry Wall.
\newblock Topological methods in group theory.
\newblock In {\em London Math. Soc. Lecture Notes Ser. 36}, pages 137--203.
  Cambridge University Press, 1979.

\bibitem[SWX20]{SWX}
Vladas Sidoravicius, Longmin Wang, and Kainan Xiang.
\newblock Limit set of branching random walks on hyperbolic groups.
\newblock arXiv:2007.13267, 2020.

\bibitem[Tio15]{Tiozzo}
Giulio Tiozzo.
\newblock Sublinear deviation between geodesics and samples paths.
\newblock {\em Duke Math. J.}, 164:511--539, 2015.

\bibitem[Tra18]{Tran}
Hung~Cong Tran.
\newblock Divergence spectra and {M}orse boundaries of relatively hyperbolic
  groups.
\newblock {\em Geom. Dedicata}, 194:99--129, 2018.

\bibitem[Var86]{Varapoulos}
{Nicholas Theodoros} Varapoulos.
\newblock Th\'eorie du potentiel sur des groupes et des vari\'et\'es.
\newblock {\em C. R. Acad. Sci. Paris}, 302:203–205, 1986.

\bibitem[{Vid}19]{Vidotto}
Pierre {Vidotto}.
\newblock Ergodic properties of some negatively curved manifolds with infinite
  measure.
\newblock {\em {M\'em. Soc. Math. Fr., Nouv. S\'er.}}, 160:1--32, 2019.

\bibitem[Woe86a]{WoessMartin}
Wolfgang Woess.
\newblock A description of the {M}artin boundary for nearest neighbour random
  walks on free products.
\newblock In {\em Probability measures on groups VIII}, pages 203--215. Lecture
  notes in mathematics, Springer, 1986.

\bibitem[Woe86b]{WoessLLT}
Wolfgang Woess.
\newblock Nearest neighbour random walks on free products of discrete groups.
\newblock {\em Bolletino dell Unione Mat. Ital.}, 5-B:961--982, 1986.

\bibitem[Woe89]{Woessends}
Wolfgang Woess.
\newblock Boundaries of random walks on graphs and groups with infinitely many
  ends.
\newblock {\em Isr. J. Math.}, 68:271--301, 1989.

\bibitem[Woe00]{Woessbook}
Wolfgang Woess.
\newblock {\em Random Walks on Infinite Graphs and Groups}.
\newblock Cambridge Press University, 2000.

\bibitem[Yan14]{Yanggrowthtightness}
Wenyuan Yang.
\newblock Growth tightness for groups with contracting elements.
\newblock {\em Math. Proc. Camb. Philos. Soc.}, 157:297--319, 2014.

\bibitem[Yan19]{YangSCC}
Wenyuan Yang.
\newblock Statistically convex-cocompact actions of groups with contracting
  elements.
\newblock {\em Int. Math. Res.}, 23:7259--7323, 2019.

\bibitem[Yan22]{YangPS}
Wenyuan Yang.
\newblock Patterson-{S}ullivan measures and growth of relatively hyperbolic
  groups.
\newblock {\em Peking Math. J.}, 5:153--212, 2022.

\end{thebibliography}
